% Generado por revision/generar_matriz_pruebas.py; no editar a mano.
\begin{tablasimple}{>{\hsize=0.82\hsize\linewidth=\hsize}Y >{\hsize=1.04\hsize\linewidth=\hsize}Y >{\hsize=0.65\hsize\linewidth=\hsize}Y >{\hsize=1.49\hsize\linewidth=\hsize}Y}
{\multicolumn{4}{l}{\textbf{Matriz individual: bloque 1 de 26.}} \\
\cab{Requisito y origen} & \cab{Componente y EDT} & \cab{Prueba prevista y estado} & \cab{Criterio y fuente}}
\textbf{RF-\allowbreak{}01.01}\newline Caso, §§ 4.1 y 9.4. & M1 Recepción (monolito modular en Laravel 13 sobre PHP 8.5 y PostgreSQL 16 por sitio; aplicación nativa Android 15 en Kotlin 2.4 con SQLite); terminal Zebra MC9400 estándar; compara SKU y cantidades recibidas contra la orden de compra y emite una alerta por cada discrepancia sin impedir el cierre; el comparativo queda disponible para Calidad\newline \textbf{EDT:} 3.4.1 & /P01 a /P03\newline Cumple; Alta; E1; H5 & El comparativo pedido vs. recibido queda registrado al cerrar la recepción y toda discrepancia queda disponible en el módulo de calidad.\newline \emph{SD3, Tabla 3.A.1, fila RF-\allowbreak{}01.01} \\
\textbf{RF-\allowbreak{}01.02}\newline Caso, §§ 4.1 y 7.3, § 16.1 (dec. 2) y cap. 18 (crit. 2); SD2, F-02 y S-02. & M1 Recepción, Kotlin 2.4/Room y Zebra MC9400; captura el lote del proveedor mediante AI 10 o ingreso manual auditado. Si el SKU exige trazabilidad y el lote está vacío, bloquea la línea hasta completarlo.\newline \textbf{EDT:} 3.4.1 & /P01 a /P05\newline Cumple; Crítica; E1; H5 & El 100 \% de los ítems con trazabilidad obligatoria registran lote; campo vacío muestra «Lote requerido para este producto» e impide continuar.\newline \emph{SD3, Tabla 3.A.1, fila RF-\allowbreak{}01.02} \\
\textbf{RF-\allowbreak{}01.03}\newline Caso, §§ 4.1, 4.8 y 9.10; SD2, F-02 y F-03. & M1 Recepción (monolito modular en Laravel 13 sobre PHP 8.5 y PostgreSQL 16 por sitio; aplicación nativa Android 15 en Kotlin 2.4 con SQLite); terminal Zebra MC9400 estándar; vincula vencimiento, lote y ubicación para alimentar FEFO; la fecha se captura y se muestra en formato DD/MM/AAAA\newline \textbf{EDT:} 3.4.1 & /P01 a /P05\newline Cumple; Crítica; E1; H5 & El 100 \% de los SKU con control de vencimiento habilitado tienen fecha registrada y asociada al lote al cerrar la recepción; campo vacío bloquea el avance.\newline \emph{SD3, Tabla 3.A.1, fila RF-\allowbreak{}01.03} \\
\textbf{RF-\allowbreak{}01.04}\newline Caso, § 4.9, cap. 8 (Calidad) y cap. 15 (RT-17.06). & M1 Recepción (monolito modular en Laravel 13 sobre PHP 8.5 y PostgreSQL 16 por sitio; aplicación nativa Android 15 en Kotlin 2.4 con SQLite); terminal Zebra MC9400 estándar; M1 registra la no conformidad con su tipo elegido de una lista parametrizada, SKU, lote, cantidad afectada y al menos una fotografía capturada con el terminal\newline \textbf{EDT:} 3.4.1 & /P01 a /P03\newline Cumple; Alta; E1; H5 & La no conformidad queda registrada con todos los campos requeridos y la fotografía es descargable desde el módulo de calidad.\newline \emph{SD3, Tabla 3.A.1, fila RF-\allowbreak{}01.04} \\
\textbf{RF-\allowbreak{}01.05}\newline Caso, §§ 4.2 y 4.9, § 16.1 (dec. 4); SD2, S-04. & M1 Recepción (monolito modular en Laravel 13 sobre PHP 8.5 y PostgreSQL 16 por sitio; aplicación nativa Android 15 en Kotlin 2.4 con SQLite); terminal Zebra MC9400 estándar; al confirmar la no conformidad, M2 pasa a cuarentena la cantidad afectada, bloqueada para la preparación, y M1 alerta en la aplicación al perfil de Calidad dentro de cinco minutos\newline \textbf{EDT:} 3.4.1 & /P01 a /P05\newline Cumple; Crítica; E1; H5 & La cantidad del lote queda bloqueado para picking de forma inmediata y el perfil de calidad recibe la notificación en $\le$ 5 min\newline \emph{SD3, Tabla 3.A.1, fila RF-\allowbreak{}01.05} \\
\textbf{RF-\allowbreak{}01.06}\newline Caso, cap. 15 (RT-05.23 y RT-17.06) y § 16.2 (GS1); BTT, RT-05.23. & M1 Recepción (monolito modular en Laravel 13 sobre PHP 8.5 y PostgreSQL 16 por sitio; aplicación nativa Android 15 en Kotlin 2.4 con SQLite); terminal Zebra MC9400 estándar; M1 decodifica los AI 01, 10, 17 y 00, autocompleta GTIN, lote, vencimiento y SSCC y, ante un código inválido o no registrado, habilita el ingreso manual auditado\newline \textbf{EDT:} 3.4.1 & /P01 a /P03\newline Cumple; Alta; E1; H5 & Código GS1 válido autocompleta los campos en pantalla; código no encontrado muestra «Producto no encontrado» y habilita ingreso manual.\newline \emph{SD3, Tabla 3.A.1, fila RF-\allowbreak{}01.06} \\
\textbf{RF-\allowbreak{}01.07}\newline Caso, § 16.1 (dec. 2) y § 16.2 (GS1); BTT, RT-05.23; SD2, S-02. & M1 genera un SSCC único al cerrar el pallet, lo vincula con GTIN, lote, vencimiento y fecha de recepción y ordena su impresión GS1-128 en Zebra ZT411. AI 00 identifica el SSCC; la fecha de recepción se conserva como dato interno y texto legible, sin codificarla como AI 11.\newline \textbf{EDT:} 3.4.1 & /P01 a /P02\newline Cumple; Media; E1; H5 & SSCC único con AI 00 impreso en GS1-128 y asociado a GTIN, lote, vencimiento y fecha de recepción. La fecha de recepción es dato interno y texto legible; AI 11 identifica producción y no se reutiliza como recepción.\newline \emph{SD3, Tabla 3.A.1, fila RF-\allowbreak{}01.07} \\
\textbf{RF-\allowbreak{}01.08}\newline Caso, §§ 4.2 y 4.9 y § 16.2 (GS1); SD2, F-03 y S-02. & M1 Recepción (monolito modular en Laravel 13 sobre PHP 8.5 y PostgreSQL 16 por sitio; aplicación nativa Android 15 en Kotlin 2.4 con SQLite); terminal Zebra MC9400 estándar; el operario escanea el SSCC y la etiqueta de la ubicación, identificada por pasillo, columna, nivel y posición con coordenada única en mae\_ubicacion; antes de confirmar se valida la compatibilidad térmica de la zona y la capacidad, y se registra el traslado con auditoría\newline \textbf{EDT:} 3.4.1 & /P01 a /P02\newline Cumple; Media; E1; H5 & La asociación SSCC–ubicación queda registrada; intento de ubicación incompatible es bloqueado con mensaje descriptivo.\newline \emph{SD3, Tabla 3.A.1, fila RF-\allowbreak{}01.08} \\
\textbf{RF-\allowbreak{}01.09}\newline Caso, §§ 4.1 y 4.2, caps. 5 y 11. & M1 Recepción (monolito modular en Laravel 13 sobre PHP 8.5 y PostgreSQL 16 por sitio); ACL INT-06; M1 publica RecepcionConfirmada por la ACL al ERP y conserva reintentos de la transacción\newline \textbf{EDT:} 3.4.1 & /P01 a /P05\newline Cumple; Crítica; E1; H5 & La conciliación diaria no muestra diferencias entre el inventario del sistema y el del ERP; ante ERP no disponible, el sistema informa al usuario y reintenta sin intervención.\newline \emph{SD3, Tabla 3.A.1, fila RF-\allowbreak{}01.09} \\
\textbf{RF-\allowbreak{}01.10}\newline Caso, § 14.1 y § 16.1 (dec. 11); SD2, S-11. & M1 Recepción (Laravel 13 sobre PHP 8.5 y PostgreSQL 16); mae\_producto exige descripción, formato, GTIN, SKU interno único, tipo de almacenamiento, temperatura mínima y máxima, peso neto y bruto, alto, ancho y largo con volumen derivado, vida útil, indicadores de lote y de vencimiento, clasificación sanitaria aprobada por Calidad y proveedor con su RUT; una ficha con campos obligatorios vacíos se rechaza (FICHA\_PRODUCTO\_INCOMPLETA) y los cambios de GTIN o SKU quedan auditados\newline \textbf{EDT:} 3.4.1, 3.7.2 & /P01 a /P03\newline Cumple; Alta; E1; H5 & El sistema no permite guardar una ficha con campos obligatorios vacíos y registra historial de cambios de GTIN o SKU con fecha y usuario.\newline \emph{SD3, Tabla 3.A.1, fila RF-\allowbreak{}01.10} \\
\textbf{RF-\allowbreak{}01.11}\newline Caso, §§ 4.2 y 4.9 y cap. 12; SD2, F-03. & M2 Inventario (monolito modular en Laravel 13 sobre PHP 8.5 y PostgreSQL 16 por sitio; aplicación nativa Android 15 en Kotlin 2.4 con SQLite); terminal Zebra MC9400 estándar; el putaway compara el tipo y el rango térmico de la ficha del producto con el tipo y el rango de mae\_zona de la ubicación destino (AMBIENTE equivale a SECO) y bloquea la operación incompatible con INCOMPATIBILIDAD\_TERMICA\newline \textbf{EDT:} 3.4.2 & /P01 a /P05\newline Cumple; Crítica; E1; H5 & Intento de ubicar un SKU congelado fuera de zona congelada es bloqueado con mensaje «Incompatibilidad térmica: requiere zona congelado».\newline \emph{SD3, Tabla 3.A.1, fila RF-\allowbreak{}01.11} \\
\textbf{RF-\allowbreak{}02.01}\newline Caso, §§ 2.3, 4.2 y 4.9, caps. 3 y 6 y cap. 10 (restr. 2); SD2, F-03 y F-05. & M2 Inventario (monolito modular en Laravel 13 sobre PHP 8.5 y PostgreSQL 16 por sitio; Aurora PostgreSQL); cada instalación se configura con zonas SECO, REFRIGERADO y CONGELADO con su rango de temperatura (mae\_zona) y posiciones por pasillo, rack/columna, nivel y posición (mae\_ubicacion); no se crea una posición sin zona; la topología versionada se publica de forma central y se distribuye como copia local de solo lectura para consulta offline\newline \textbf{EDT:} 3.4.2 & /P01 a /P03\newline Cumple; Alta; E1; H5 & (a) Al configurar un nuevo CD, el sistema permite crear y mapear zonas, pasillos y posiciones con coordenada única. (b) Cada posición tiene zona térmica asociada con rango de temperatura. (c) No se permite crear posiciones sin zona asignada. (d) Cambios de topología quedan registrados con usuario, fecha y detalle del cambio.\newline \emph{SD3, Tabla 3.A.1, fila RF-\allowbreak{}02.01} \\
\textbf{RF-\allowbreak{}02.02}\newline Caso, § 4.2, cap. 5 y § 9.2; SD2, F-03 y S-14. & M2 Inventario (monolito modular en Laravel 13 sobre PHP 8.5 y PostgreSQL 16 por sitio; Aurora PostgreSQL); M2 consolida el stock físico de los sitios operativos para consulta central\newline \textbf{EDT:} 3.4.2 & /P01 a /P03\newline Cumple; Alta; E1; H5 & El sistema expone un panel donde el usuario visualiza: (a) el stock físico totalizado por SKU, (b) el desglose por instalación, y (c) el desglose por zona operativa dentro de cada instalación.\newline \emph{SD3, Tabla 3.A.1, fila RF-\allowbreak{}02.02} \\
\textbf{RF-\allowbreak{}02.03}\newline Caso, § 4.2 y § 16.2 (ubicación por rotación); SD2, F-03. & M2 Inventario (monolito modular en Laravel 13 sobre PHP 8.5 y PostgreSQL 16 por sitio; Aurora PostgreSQL); la propuesta de slotting (inv\_reubicacion\_propuesta) conserva período ABC configurable, clase, motivo (ABC, peso, volumen o zona), ubicación actual y propuesta; la clase A prioriza ubicaciones de menor distancia al andén de salida sin exceder peso, volumen ni zona; el movimiento solo se instruye después de la aprobación del jefe de bodega, que queda registrada con aprobador y fecha\newline \textbf{EDT:} 3.4.2 & /P01 a /P02\newline Cumple; Media; E1; H5 & (a) Al ejecutar el cálculo, el sistema presenta una propuesta con: producto, ubicación actual, ubicación propuesta, motivo (clasificación ABC, peso, volumen o zona). (b) El jefe de bodega puede aprobar, rechazar o modificar cada movimiento individualmente. (c) Los movimientos aprobados generan instrucciones de reubicación asignables a dispositivos de operarios, con indicación de ubicación origen y destino. (d) Cada reubicación ejecutada queda registrada con operario, fecha/hora, y ubicaciones origen/destino.\newline \emph{SD3, Tabla 3.A.1, fila RF-\allowbreak{}02.03} \\
\textbf{RF-\allowbreak{}02.04}\newline Caso, §§ 4.2, 7.2 y 14.1; SD2, F-03 y Tabla 2.1. & M2 Inventario (monolito modular en Laravel 13 sobre PHP 8.5 y PostgreSQL 16 por sitio); aplicación nativa Android 15 en Kotlin 2.4 con SQLite; terminal Zebra MC9400 estándar y Zebra MC9400 Cold Storage Freezer; plan mensual de conteo cíclico por rotación, valor y discrepancia histórica, en conteo ciego sin cantidad teórica; la diferencia se registra con su causa y, sobre el umbral del recinto, queda pendiente de autorización de la jefatura de bodega\newline \textbf{EDT:} 3.4.2 & /P01 a /P03\newline Cumple; Alta; E1; H5 & (a) El operario recibe en su dispositivo la lista de posiciones a contar, sin cantidad teórica visible. (b) Al ingresar la cantidad física, si la diferencia valorizada es menor o igual al umbral configurado por el jefe de bodega para ese recinto, el sistema registra el ajuste automáticamente con la causa seleccionada (sobrante, falta, error de ubicación, daño, vencimiento). (c) Si la diferencia supera el umbral, el sistema marca la tarea como 'Pendiente de Autorización' y notifica al jefe de bodega.\newline \emph{SD3, Tabla 3.A.1, fila RF-\allowbreak{}02.04} \\
\textbf{RF-\allowbreak{}02.05}\newline Caso, § 9.2, cap. 10 (restr. 3), § 16.1 (dec. 8) y cap. 18 (crit. 5); SD2, S-08; SD3, RNG-02. & M2 Inventario (monolito modular en Laravel 13 sobre PHP 8.5 y PostgreSQL 16 por sitio; Aurora PostgreSQL); M2 consulta disponibilidad online y muestra la copia offline con su fecha de actualización; la reserva considera pedidos confirmados\newline \textbf{EDT:} 3.4.2 & /P01 a /P03\newline Cumple; Alta; E1; H5 & (a) Al consultar stock, el preventista visualiza la cantidad disponible por SKU para su CD asignado. (b) La consulta refleja el stock actualizado tras cada transacción confirmada. (c) En modo offline, la consulta muestra «Dato de las HH:MM del DD/MM». (d) Al sincronizar un pedido tomado offline, si el stock disponible es insuficiente para alguna línea, el sistema marca esa línea como «sin stock disponible» y notifica al preventista. (e) Los conflictos entre pedidos simultáneos se resuelven por orden de confirmación en el servidor central (S-08, RNG-01); un pedido capturado sin señal toma su lugar al sincronizar.\newline \emph{SD3, Tabla 3.A.1, fila RF-\allowbreak{}02.05} \\
\textbf{RF-\allowbreak{}02.06}\newline Caso, cap. 3, §§ 4.2 y 4.5 y cap. 5; SD2, F-05. & M2 Inventario deriva el tránsito desde carga confirmada hasta recepción física confirmada mediante eventos de custodia; lo separa del stock físico, comprometido y disponible, con llegada estimada. Recepción offline y sincronización idempotente; transferencia inter-bodegas al ERP por INT-06 con cola y reintento automático.\newline \textbf{EDT:} 3.4.2 & /P01 a /P02\newline Cumple; Media; E1; H5 & (a) La consulta distingue stock físico en sitio, comprometido, disponible y en tránsito al sitio, con fecha/hora estimada. (b) El tránsito no suma a disponible hasta la recepción confirmada. (c) Al confirmarla, pasa automáticamente de tránsito a físico en destino. (d) Sin enlace, la recepción se registra localmente y sincroniza al reconectar. (e) Cada transferencia genera el movimiento interbodegas en el ERP; si no está disponible, se encola para reintento automático. (f) El tránsito se muestra como disponibilidad futura separada del stock físico del sitio.\newline \emph{SD3, Tabla 3.A.1, fila RF-\allowbreak{}02.06} \\
\textbf{RF-\allowbreak{}02.07}\newline Caso, §§ 2.3 y 4.2 y cap. 5 (WMS); SD2, F-03. & M2 Inventario (monolito modular en Laravel 13 sobre PHP 8.5 y PostgreSQL 16 por sitio; Aurora PostgreSQL); hech\_ocupacion\_diaria calcula por sitio, zona y día el porcentaje de posiciones con stock físico positivo sobre las posiciones habilitadas, contando una vez cada posición aunque tenga varios lotes o SSCC, con al menos doce meses de historia\newline \textbf{EDT:} 3.4.2 & /P01 a /P02\newline Cumple; Media; E1; H5 & (a) Al consultar la ocupación de un CD, el sistema muestra el porcentaje de ocupación total y por zona (b) El reporte incluye un histórico diario de al menos 12 meses para visualizar tendencias. (c) El reporte es exportable a formato de hoja de cálculo.\newline \emph{SD3, Tabla 3.A.1, fila RF-\allowbreak{}02.07} \\
\textbf{RF-\allowbreak{}02.08}\newline Caso, § 4.5, cap. 8 (Jefa de Bodega), cap. 15 (RT-09.01) y cap. 18 (crit. 15); SD2, F-04. & M2 Inventario (monolito modular en Laravel 13 sobre PHP 8.5 y PostgreSQL 16 por sitio); aplicación nativa Android 15 en Kotlin 2.4 con SQLite; terminal Zebra MC9400 estándar y Zebra MC9400 Cold Storage Freezer; M5 dirige misiones GS1: cada línea indica ubicación y secuencia de recorrido, y la asignación aporta producto, lote y cantidad con FEFO; el preparador confirma cada línea y un faltante exige causa estructurada (sin stock, producto dañado, producto vencido, error de ubicación u otro con detalle) antes de confirmar\newline \textbf{EDT:} 3.4.2, 3.4.3 & /P01 a /P03\newline Cumple; Alta; E1; H5 & (a) El preparador recibe en su dispositivo la lista de líneas a recolectar, ordenadas por secuencia de recorrido óptima. (b) La confirmación de cada línea no supera 1 segundo de respuesta en el dispositivo. (c) Al registrar un faltante, el sistema exige la selección de una causa antes de continuar. (d) Al finalizar la preparación, el sistema actualiza el stock automáticamente y envía los faltantes al área comercial para gestión con el cliente.\newline \emph{SD3, Tabla 3.A.1, fila RF-\allowbreak{}02.08} \\
\textbf{RF-\allowbreak{}02.09}\newline Caso, §§ 4.2, 4.8 y 9.10; SD2, F-03. & M2 Inventario (monolito modular en Laravel 13 sobre PHP 8.5 y PostgreSQL 16 por sitio; Aurora PostgreSQL); M2 aplica FEFO, prioriza el vencimiento más próximo y alerta a la jefatura de bodega cuando un lote queda bajo el umbral configurable de vida útil remanente\newline \textbf{EDT:} 3.4.2 & /P01 a /P03\newline Cumple; Alta; E1; H5 & (a) Al generar las tareas de picking, el sistema prioriza las ubicaciones con menor vida útil remanente para cada SKU. (b) Si un producto en stock tiene una vida útil remanente inferior al umbral configurado (ej. 30 días), el sistema genera una alerta al jefe de bodega con la acción recomendada (descuento, devolución al proveedor, destrucción). (c) El sistema genera un reporte semanal de productos próximos a vencer por zona y por SKU.\newline \emph{SD3, Tabla 3.A.1, fila RF-\allowbreak{}02.09} \\
\textbf{RF-\allowbreak{}02.10}\newline Caso, §§ 4.9 y 9.5, § 16.1 (dec. 2) y cap. 18 (crit. 1); SD2, F-11 y S-02. & M9 consulta el ciclo del lote por SKU, fecha o proveedor: estado, vencimiento, cantidad por ubicación, recepción y proveedor/OC, salidas por cliente y guía, mermas, devoluciones y ajustes. Los eventos EPCIS y las referencias de M1/M2/M5/M6 permiten exportar el detalle para retiro sanitario.\newline \textbf{EDT:} 3.4.4 & /P01 a /P03\newline Cumple; Alta; E1; H5 & (a) El sistema provee un buscador de lotes por SKU, fecha o proveedor, devolviendo una lista con el estado (activo, agotado, en cuarentena). (b) Al seleccionar un lote, se despliega su ficha detallando: SKU, fecha vencimiento y cantidad física remanente por ubicación (pasillo/posición). (c) Se detalla el origen: fecha de recepción, proveedor, cantidad recibida inicial y documento. (d) Se detalla la salida de unidades: guías de despacho, cliente, fecha y cantidad entregada. (e) Se detallan excepciones: mermas, devoluciones o ajustes asociados. (f) Todo el detalle es exportable a hoja de cálculo para retiro sanitario inmediato.\newline \emph{SD3, Tabla 3.A.1, fila RF-\allowbreak{}02.10} \\
\textbf{RF-\allowbreak{}02.11}\newline Caso, sección 2.3 y cap. 3; SD2, Anexo 2.1, F-05. & M2 Inventario: registro de lectura de unidades en cada cross-docking al ingreso del camión de línea y carga del camión de reparto, con hora y excepción/causa; Laravel 13.x, PHP 8.5.x, PostgreSQL 16.x y consolidación central Aurora PostgreSQL\newline \textbf{EDT:} 3.4.2 & /P01 a /P03\newline Cumple; Alta; E1; H5 & Por cada plataforma y día se obtiene el tiempo real entre consolidación, llegada, desconsolidación y despacho, y el detalle de excepciones, para medir la brecha de 16 horas (Subdocumento 2, Anexo 2.2, SP-02).\newline \emph{SD3, Tabla 3.A.3b, fila RF-\allowbreak{}02.11} \\
\textbf{RF-\allowbreak{}03.01}\newline Caso, § 9.2 y cap. 18 (crit. 5); SD2, F-01. & M3 Preventa (aplicación nativa Android 15 en Kotlin 2.4 con SQLite; terminal Zebra TC58e; monolito modular en Laravel 13 sobre PHP 8.5; ElastiCache con Valkey 8.1); M3 consulta por API el stock actual por SKU y T-11 declara respuesta menor a dos segundos\newline \textbf{EDT:} 3.4.6 & /P01 a /P05\newline Cumple; Crítica; E1; H5 & El preventista ve la disponibilidad vigente antes de ofrecer un producto.\newline \emph{SD3, Tabla 3.A.1, fila RF-\allowbreak{}03.01} \\
\textbf{RF-\allowbreak{}03.02}\newline Caso, § 9.2 y cap. 10 (restr. 3); SD2, F-01 y S-08. & M3 Preventa (aplicación nativa Android 15 en Kotlin 2.4 con SQLite; terminal Zebra TC58e; monolito modular en Laravel 13 sobre PHP 8.5; ElastiCache con Valkey 8.1); M3 mantiene copia local con marca de sincronización para consulta sin conexión\newline \textbf{EDT:} 3.4.6 & /P01 a /P05\newline Cumple; Crítica; E1; H5 & El preventista ve el dato cacheado y su antigüedad.\newline \emph{SD3, Tabla 3.A.1, fila RF-\allowbreak{}03.02} \\
\textbf{RF-\allowbreak{}03.03}\newline Caso, §§ 9.2 y 16.1 (dec. 8); SD2, S-08; SD3, RNG-01. & M3 Preventa (aplicación nativa Android 15 en Kotlin 2.4 con SQLite; terminal Zebra TC58e; monolito modular en Laravel 13 sobre PHP 8.5; ElastiCache con Valkey 8.1); M3 confirma reservas en el servidor central, ordenadas por recepción en ese servidor y sin depender del ERP\newline \textbf{EDT:} 3.4.6 & /P01 a /P03\newline Cumple; Alta; E1; H5 & Si dos preventistas comprometen el mismo producto, el segundo en confirmar queda con quiebre parcial registrado.\newline \emph{SD3, Tabla 3.A.1, fila RF-\allowbreak{}03.03} \\
\end{tablasimple}

\begin{tablasimple}{>{\hsize=0.82\hsize\linewidth=\hsize}Y >{\hsize=1.04\hsize\linewidth=\hsize}Y >{\hsize=0.65\hsize\linewidth=\hsize}Y >{\hsize=1.49\hsize\linewidth=\hsize}Y}
{\multicolumn{4}{l}{\textbf{Matriz individual: bloque 2 de 26.}} \\
\cab{Requisito y origen} & \cab{Componente y EDT} & \cab{Prueba prevista y estado} & \cab{Criterio y fuente}}
\textbf{RF-\allowbreak{}03.04}\newline Caso, §§ 4.3 y 9.2 y cap. 18 (crit. 5); SD2, F-01. & M3 Preventa (aplicación nativa Android 15 en Kotlin 2.4 con SQLite; terminal Zebra TC58e; monolito modular en Laravel 13 sobre PHP 8.5; ElastiCache con Valkey 8.1); con conectividad, M3 muestra el crédito disponible y la deuda vencida del cliente al tomar el pedido\newline \textbf{EDT:} 3.4.6 & /P01 a /P05\newline Cumple; Crítica; E1; H5 & El preventista ve crédito y deuda antes de finalizar el pedido.\newline \emph{SD3, Tabla 3.A.1, fila RF-\allowbreak{}03.04} \\
\textbf{RF-\allowbreak{}03.05}\newline Caso, §§ 4.3 y 9.2, cap. 10 (restr. 3); SD2, F-01. & M3 Preventa (aplicación nativa Android 15 en Kotlin 2.4 con SQLite; terminal Zebra TC58e; monolito modular en Laravel 13 sobre PHP 8.5; ElastiCache con Valkey 8.1); sin conectividad, M3 muestra el crédito y la deuda vencida de la última sincronización con su antigüedad\newline \textbf{EDT:} 3.4.6 & /P01 a /P05\newline Cumple; Crítica; E1; H5 & El preventista ve el dato cacheado y su antigüedad.\newline \emph{SD3, Tabla 3.A.1, fila RF-\allowbreak{}03.05} \\
\textbf{RF-\allowbreak{}03.06}\newline Caso, § 9.2; SD2, F-01. & M3 Preventa (aplicación nativa Android 15 en Kotlin 2.4 con SQLite; terminal Zebra TC58e; monolito modular en Laravel 13 sobre PHP 8.5; ElastiCache con Valkey 8.1); M3 muestra el historial de compra del cliente con productos y frecuencia\newline \textbf{EDT:} 3.4.6 & /P01 a /P02\newline Cumple; Media; E1; H5 & El preventista consulta el historial e identifica productos habituales o quiebres.\newline \emph{SD3, Tabla 3.A.1, fila RF-\allowbreak{}03.06} \\
\textbf{RF-\allowbreak{}03.07}\newline Caso, §§ 4.3 y 9.2. & M3 Preventa (aplicación nativa Android 15 en Kotlin 2.4 con SQLite; terminal Zebra TC58e; monolito modular en Laravel 13 sobre PHP 8.5; ElastiCache con Valkey 8.1); M3 mantiene en la copia local las promociones vigentes y calcula el descuento de la promoción aplicada a la línea\newline \textbf{EDT:} 3.4.6 & /P01 a /P02\newline Cumple; Media; E1; H5 & El preventista aplica la promoción y el sistema calcula el descuento.\newline \emph{SD3, Tabla 3.A.1, fila RF-\allowbreak{}03.07} \\
\textbf{RF-\allowbreak{}03.08}\newline Caso, §§ 4.3 y 9.2 y cap. 8 (Comercial y Preventa). & M3 Preventa (aplicación nativa Android 15 en Kotlin 2.4 con SQLite; terminal Zebra TC58e; monolito modular en Laravel 13 sobre PHP 8.5; ElastiCache con Valkey 8.1); M3 alerta al preventista antes del envío cuando el pedido más la deuda vencida supera el límite de crédito\newline \textbf{EDT:} 3.4.6 & /P01 a /P03\newline Cumple; Alta; E1; H5 & El preventista es alertado antes de enviar el pedido.\newline \emph{SD3, Tabla 3.A.1, fila RF-\allowbreak{}03.08} \\
\textbf{RF-\allowbreak{}03.09}\newline Caso, §§ 4.3 y 9.2; SD3, RNG-03. & M3 implementa bloqueo y excepción auditada del supervisor. RNG-03 somete la fórmula y el umbral del 110 \% a la validación de Finanzas antes de configurarlos.\newline \textbf{EDT:} 3.4.6 & /P01 a /P03\newline Parcial; Alta; E1; H5 & El pedido no puede enviarse mientras el bloqueo esté activo.\newline \emph{SD3, Tabla 3.A.1, fila RF-\allowbreak{}03.09} \\
\textbf{RF-\allowbreak{}03.10}\newline Caso, §§ 4.3 y 9.2; SD3, RNG-03. & M3 Preventa (aplicación nativa Android 15 en Kotlin 2.4 con SQLite; terminal Zebra TC58e; monolito modular en Laravel 13 sobre PHP 8.5; ElastiCache con Valkey 8.1); un supervisor autorizado aprueba el envío del pedido bloqueado, con su identidad, fecha y hora en la auditoría\newline \textbf{EDT:} 3.4.6 & /P01 a /P03\newline Cumple; Alta; E1; H5 & Queda registrado supervisor, fecha y hora, y el pedido puede enviarse.\newline \emph{SD3, Tabla 3.A.1, fila RF-\allowbreak{}03.10} \\
\textbf{RF-\allowbreak{}03.11}\newline Caso, § 16.1 (dec. 9); SD2, S-09; SD3, RNG-08. & M3 Preventa (aplicación nativa Android 15 en Kotlin 2.4 con SQLite; terminal Zebra TC58e; monolito modular en Laravel 13 sobre PHP 8.5; ElastiCache con Valkey 8.1); M3 conserva el precio del pedido y lo transmite al ERP por el contrato de integración\newline \textbf{EDT:} 3.4.6 & /P01 a /P02\newline Cumple; Media; E1; H5 & El pedido conserva el precio acordado y lo transmite al ERP para la facturación.\newline \emph{SD3, Tabla 3.A.1, fila RF-\allowbreak{}03.11} \\
\textbf{RF-\allowbreak{}03.12}\newline Caso, § 16.1 (dec. 9); SD2, S-09; SD3, RNG-08. & M3 Preventa (aplicación nativa Android 15 en Kotlin 2.4 con SQLite; terminal Zebra TC58e; monolito modular en Laravel 13 sobre PHP 8.5; ElastiCache con Valkey 8.1); el precio pactado es inmutable y prv\_pedido\_linea conserva el precio recibido del ERP para facturación y su diferencia; una discrepancia deja la línea BLOQUEADA, impide la facturación y genera una incidencia a Comercial, que la concilia con actor y fecha sin modificar lo acordado\newline \textbf{EDT:} 3.4.6 & /P01 a /P02\newline Cumple; Media; E1; H5 & La diferencia queda alertada al responsable comercial y se corrige antes de facturar, conservando el precio acordado.\newline \emph{SD3, Tabla 3.A.1, fila RF-\allowbreak{}03.12} \\
\textbf{RF-\allowbreak{}03.13}\newline Caso, § 9.1; SD2, § 2.2.3 (tensión de entrega); SD3, S-21 y RNG-15. & M3 Preventa (aplicación nativa Android 15 en Kotlin 2.4 con SQLite; terminal Zebra TC58e; monolito modular en Laravel 13 sobre PHP 8.5; ElastiCache con Valkey 8.1); antes de confirmar, M3 informa la fecha y la ventana comprometibles según RNG-15, el stock reservable, la capacidad de despacho de la ruta y las franjas del cliente, y conserva la promesa\newline \textbf{EDT:} 3.4.6 & /P01 a /P05\newline Cumple; Crítica; E1; H5 & El preventista visualiza la ventana disponible antes de confirmar; la promesa queda registrada.\newline \emph{SD3, Tabla 3.A.1, fila RF-\allowbreak{}03.13} \\
\textbf{RF-\allowbreak{}03.14}\newline Caso, § 9.2. & M3 Preventa (aplicación nativa Android 15 en Kotlin 2.4 con SQLite; terminal Zebra TC58e; monolito modular en Laravel 13 sobre PHP 8.5; ElastiCache con Valkey 8.1); M3 sugiere productos según historial, frecuencia, productos habituales y promociones vigentes, con los criterios que configura Comercial\newline \textbf{EDT:} 3.4.6 & /P01 a /P02\newline Cumple; Media; E1; H5 & Al seleccionar un cliente, se muestran sugerencias según criterios configurados por comercial.\newline \emph{SD3, Tabla 3.A.1, fila RF-\allowbreak{}03.14} \\
\textbf{RF-\allowbreak{}03.15}\newline Caso, § 9.2, cap. 10 (restr. 3) y cap. 18 (crit. 6). & M3 Preventa (aplicación nativa Android 15 en Kotlin 2.4 con SQLite; terminal Zebra TC58e; monolito modular en Laravel 13 sobre PHP 8.5; ElastiCache con Valkey 8.1); M3 genera UUID en el dispositivo para cada pedido capturado sin conexión\newline \textbf{EDT:} 3.4.6 & /P01 a /P05\newline Cumple; Crítica; E1; H5 & Cada pedido creado sin conectividad tiene un identificador único generado localmente.\newline \emph{SD3, Tabla 3.A.1, fila RF-\allowbreak{}03.15} \\
\textbf{RF-\allowbreak{}03.16}\newline Caso, § 9.2, cap. 10 (restr. 3) y cap. 18 (crit. 6); SD3, RNG-12. & M3 Preventa (aplicación nativa Android 15 en Kotlin 2.4 con SQLite; terminal Zebra TC58e; monolito modular en Laravel 13 sobre PHP 8.5; ElastiCache con Valkey 8.1); M3 deduplica reenvíos con el mismo UUID al sincronizar\newline \textbf{EDT:} 3.4.6 & /P01 a /P05\newline Cumple; Crítica; E1; H5 & Ningún pedido se pierde ni se duplica al sincronizar; todo reintento queda auditado.\newline \emph{SD3, Tabla 3.A.1, fila RF-\allowbreak{}03.16} \\
\textbf{RF-\allowbreak{}04.01}\newline Caso, §§ 4.4 y 9.3, § 16.2 (ruteo) y cap. 18 (crit. 7); SD2, F-06. & M4 Rutas (monolito modular en Laravel 13 sobre PHP 8.5; motor propio de optimización en contenedor ECS Fargate; INT-10, mapas y geocodificación); mae\_vehiculo declara capacidad nominal en kg y m³ y el punto del cliente su zona de reparto; M4 asigna pedidos a la flota calculando peso bruto y volumen desde las líneas, las conversiones de unidad y la ficha del producto, y secuencia la ruta por zona logística\newline \textbf{EDT:} 3.4.7, 3.6.3 & /P01 a /P03\newline Cumple; Alta; E1; H5 & Se despliega una ruta propuesta con la carga asignada a los vehículos de la flota.\newline \emph{SD3, Tabla 3.A.1, fila RF-\allowbreak{}04.01} \\
\textbf{RF-\allowbreak{}04.02}\newline Caso, § 9.3, § 16.1 (dec. 16) y cap. 18 (crit. 7); SD2, F-06 y S-16. & M4 Rutas (monolito modular en Laravel 13 sobre PHP 8.5; motor propio de optimización en contenedor ECS Fargate; INT-10, mapas y geocodificación); el planificador modifica la secuencia y reasigna pedidos entre vehículos, y M4 recalcula de inmediato la carga y las horas estimadas\newline \textbf{EDT:} 3.4.7, 3.6.3 & /P01 a /P03\newline Cumple; Alta; E1; H5 & La tabla de ruteo y el ETA se recalculan inmediatamente tras realizar alguna reasignación\newline \emph{SD3, Tabla 3.A.1, fila RF-\allowbreak{}04.02} \\
\textbf{RF-\allowbreak{}04.03}\newline Caso, § 4.4 y § 16.2 (ruteo con capacidad); SD2, F-06. & M4 Rutas (monolito modular en Laravel 13 sobre PHP 8.5; motor propio de optimización en contenedor ECS Fargate; INT-10, mapas y geocodificación); mae\_vehiculo contiene los límites nominales de peso y volumen y mae\_producto el peso bruto y las dimensiones; la adición de una línea que excede cualquiera de los dos límites se rechaza y bloquea el cierre de la ruta hasta redistribuir la carga\newline \textbf{EDT:} 3.4.7, 3.6.3 & /P01 a /P03\newline Cumple; Alta; E1; H5 & Se despliega un error bloqueante impidiendo el cierre del camión hasta reubicar el excedente.\newline \emph{SD3, Tabla 3.A.1, fila RF-\allowbreak{}04.03} \\
\textbf{RF-\allowbreak{}04.04}\newline Caso, §§ 9.1 y 9.9 y § 16.2 (ruteo con ventanas); SD2, F-06. & M4 Rutas (monolito modular en Laravel 13 sobre PHP 8.5; motor propio de optimización en contenedor ECS Fargate; INT-10, mapas y geocodificación); M4 considera ventanas de clientes en la secuencia de ruta\newline \textbf{EDT:} 3.4.7, 3.6.3 & /P01 a /P03\newline Cumple; Alta; E1; H5 & El planificador secuencia todas las entregas en ventanas donde el cliente está disponible\newline \emph{SD3, Tabla 3.A.1, fila RF-\allowbreak{}04.04} \\
\textbf{RF-\allowbreak{}04.05}\newline Caso, §§ 2.3 y 4.9 y § 9.6; SD2, F-06 y F-07. & M4 Rutas (monolito modular en Laravel 13 sobre PHP 8.5; motor propio de optimización en contenedor ECS Fargate; INT-10, mapas y geocodificación); M4 asigna pedidos refrigerados o congelados solo a vehículos con equipo de frío y rechaza su asignación a un camión seco\newline \textbf{EDT:} 3.4.7, 3.6.3 & /P01 a /P03\newline Cumple; Alta; E1; H5 & Se impide la adición de cadena de frío a un camión seco.\newline \emph{SD3, Tabla 3.A.1, fila RF-\allowbreak{}04.05} \\
\textbf{RF-\allowbreak{}04.06}\newline Caso, §§ 9.1 y 9.9 y cap. 15 (RT-16.21). & M4 genera por INT-11, al confirmar la ruta, el aviso al teléfono registrado en la ficha del cliente con ETA y ventana de 30 minutos, antes de la partida; aceptación por recepción efectiva\newline \textbf{EDT:} 3.4.7 & /P01 a /P02\newline Cumple; Media; E1; H5 & El cliente recibe de forma exitosa el mensaje antes de la partida del camión.\newline \emph{SD3, Tabla 3.A.1, fila RF-\allowbreak{}04.06} \\
\textbf{RF-\allowbreak{}04.07}\newline Caso, §§ 9.1 y 9.9 y § 16.2 (ruteo con ventanas). & Franjas de recepción por cliente en parámetros versionados INTERVALO mantenidos por el planificador; M4 no asigna ETA fuera de franja vigente\newline \textbf{EDT:} 3.4.7 & /P01 a /P03\newline Cumple; Alta; E1; H5 & Ningún horario estimado de llegada es asignado en un horario fuera de la ventana horaria establecida por el cliente.\newline \emph{SD3, Tabla 3.A.1, fila RF-\allowbreak{}04.07} \\
\textbf{RF-\allowbreak{}04.08}\newline Caso, § 9.8, § 16.1 (dec. 7) y cap. 18 (crit. 11); SD2, F-14. & M4 Rutas (monolito modular en Laravel 13 sobre PHP 8.5; motor propio de optimización en contenedor ECS Fargate; INT-10, mapas y geocodificación); M4 alimenta el costo de servir y el modelo calcula costo por entrega y cliente\newline \textbf{EDT:} 3.5.4, 3.4.7 & /P01 a /P02\newline Cumple; Media; E2; H10 & El sistema calcula el costo de cada entrega basado en kilómetros, tiempo de descarga, tamaño del pedido y reentregas.\newline \emph{SD3, Tabla 3.A.1, fila RF-\allowbreak{}04.08} \\
\textbf{RF-\allowbreak{}05.01}\newline Caso, § 4.5 y § 16.2 (preparación de pedidos); SD2, F-04. & M5 Preparación (monolito modular en Laravel 13 sobre PHP 8.5 y PostgreSQL 16 por sitio; aplicación nativa Android 15 en Kotlin 2.4 con SQLite; terminal Zebra MC9400 estándar y Zebra MC9400 Cold Storage Freezer); la jefatura de bodega genera las misiones manualmente o las programa al cierre de pedidos del turno, agrupadas por pedido o por ruta\newline \textbf{EDT:} 3.4.3 & /P01 a /P03\newline Cumple; Alta; E1; H5 & (a) Al generar misiones, la interfaz lista los ítems agrupados por pasillo lógico en secuencia de proximidad. (b) El sistema exige y permite confirmar cada línea digitalmente, exhibiendo: cantidad solicitada, SKU, lote asignado (FEFO) y coordenada exacta de ubicación.\newline \emph{SD3, Tabla 3.A.1, fila RF-\allowbreak{}05.01} \\
\textbf{RF-\allowbreak{}05.02}\newline Caso, §§ 4.5 y 9.5, cap. 15 (RT-17.06) y § 16.2 (GS1); SD2, F-04 y S-02. & M5 Preparación (monolito modular en Laravel 13 sobre PHP 8.5 y PostgreSQL 16 por sitio; aplicación nativa Android 15 en Kotlin 2.4 con SQLite; terminal Zebra MC9400 estándar y Zebra MC9400 Cold Storage Freezer); M5 exige lectura GS1 en la preparación para validar el SKU y registrar el lote exacto; el ingreso manual del código es una contingencia registrada en la auditoría\newline \textbf{EDT:} 3.4.3 & /P01 a /P03\newline Cumple; Alta; E1; H5 & Si el código de barras escaneado difiere del SKU o Lote instruido, el sistema bloquea la confirmación. Como mecanismo de contingencia, el sistema debe permitir el ingreso manual del código numérico, pero debe registrar obligatoriamente el evento en la bitácora como 'Ingreso Manual' para futura auditoría de calidad.\newline \emph{SD3, Tabla 3.A.1, fila RF-\allowbreak{}05.02} \\
\textbf{RF-\allowbreak{}05.03}\newline Caso, § 4.5 y cap. 18 (crit. 15); SD2, Figura 2.3. & M5 exige causa de la lista parametrizada cuando la cantidad es inferior y bloquea Confirmar Línea; causa en lectura local y auditoría\newline \textbf{EDT:} 3.4.3 & /P01 a /P03\newline Cumple; Alta; E1; H5 & Si la cantidad recolectada es inferior a la solicitada, el sistema despliega el campo obligatorio 'Causa de Faltante' con lista predefinida y deshabilita la acción 'Confirmar Línea' hasta su selección.\newline \emph{SD3, Tabla 3.A.1, fila RF-\allowbreak{}05.03} \\
\textbf{RF-\allowbreak{}05.04}\newline Caso, §§ 4.5 y 4.9; SD2, Figura 2.3 y F-04. & M5 Preparación aplica RNG-13: picking en orden seco, refrigerado y congelado; preparación de frío próxima a la carga y conservación en su zona durante esperas e interrupciones.\newline \textbf{EDT:} 3.4.3 & /P01 a /P03\newline Cumple; Alta; E1; H5 & Ante un pedido mixto, el sistema exige recolectar el 100\% de los SKUs de ambiente seco antes de liberar las ubicaciones correspondientes a las cámaras de frío en la interfaz del operario.\newline \emph{SD3, Tabla 3.A.1, fila RF-\allowbreak{}05.04} \\
\end{tablasimple}

\begin{tablasimple}{>{\hsize=0.82\hsize\linewidth=\hsize}Y >{\hsize=1.04\hsize\linewidth=\hsize}Y >{\hsize=0.65\hsize\linewidth=\hsize}Y >{\hsize=1.49\hsize\linewidth=\hsize}Y}
{\multicolumn{4}{l}{\textbf{Matriz individual: bloque 3 de 26.}} \\
\cab{Requisito y origen} & \cab{Componente y EDT} & \cab{Prueba prevista y estado} & \cab{Criterio y fuente}}
\textbf{RF-\allowbreak{}05.05}\newline Caso, §§ 4.2, 4.8 y 9.10; SD2, F-03. & M5 Preparación (monolito modular en Laravel 13 sobre PHP 8.5 y PostgreSQL 16 por sitio; aplicación nativa Android 15 en Kotlin 2.4 con SQLite; terminal Zebra MC9400 estándar y Zebra MC9400 Cold Storage Freezer); M5 asigna inventario aplicando FEFO por vencimiento, bloquea la confirmación de un lote posterior si el lote FEFO está disponible y registra el motivo de toda excepción\newline \textbf{EDT:} 3.4.3 & /P01 a /P03\newline Cumple; Alta; E1; H5 & (a) La misión dirige al operario a la coordenada del lote con fecha de vencimiento más próxima. (b) Al escanear un lote más nuevo cuando hay inventario del lote FEFO disponible, el sistema bloquea la confirmación con mensaje «Lote (código de lote) con vencimiento anterior disponible en [ubicación]». (c) Si el lote FEFO no está disponible (ej. dañado), el preparador puede seleccionar el siguiente lote registrando el motivo (lista estructurada: no encontrado, dañado, cantidad insuficiente). (d) La excepción FEFO queda registrada en bitácora de auditoría\newline \emph{SD3, Tabla 3.A.1, fila RF-\allowbreak{}05.05} \\
\textbf{RF-\allowbreak{}05.06}\newline Caso, §§ 4.5 y 9.5 y caps. 5 y 11; SD2, F-04 y S-02. & M5, al cerrar la misión, envía por erp-sync/ACL INT-06 el ID del pedido y las líneas con SKU, cantidad efectivamente preparada y lote. El ERP recibe el conjunto sin digitación para emitir la guía; la salida exige después carga y guía vigentes.\newline \textbf{EDT:} 3.4.3, 3.3.2 & /P01 a /P03\newline Cumple; Alta; E1; H5 & Al ejecutar 'Cerrar Misión', el sistema transmite al ERP el objeto completo de la recolección asociado al ID del pedido. El ERP procesa este conjunto de datos ajustando las cantidades finales y registrando los lotes exactos para la emisión de la Guía de Despacho, sin permitir digitación manual.\newline \emph{SD3, Tabla 3.A.1, fila RF-\allowbreak{}05.06} \\
\textbf{RF-\allowbreak{}05.07}\newline Caso, § 4.5 y § 16.2 (GS1); SD2, Figura 2.3 y S-02. & M5 Preparación (monolito modular en Laravel 13 sobre PHP 8.5 y PostgreSQL 16 por sitio; aplicación nativa Android 15 en Kotlin 2.4 con SQLite; terminal Zebra MC9400 estándar y Zebra MC9400 Cold Storage Freezer); la carga sigue el orden inverso de la ruta con lectura del SSCC de cada bulto, alerta fuera de secuencia y bloqueo de la carga completa si falta un bulto, con aviso a la jefatura de bodega\newline \textbf{EDT:} 3.4.3 & /P01 a /P03\newline Cumple; Alta; E1; H5 & (a) El dispositivo del cargador indica el siguiente bulto/pallet a subir, mostrando: SSCC, cliente destino y posición en la secuencia. (b) El cargador confirma cada bulto escaneando el SSCC; si escanea un bulto fuera de secuencia, el sistema alerta y solicita confirmación del motivo. (c) Al completar la carga, el sistema verifica que todos los bultos asignados al camión han sido escaneados. (d) Si falta algún bulto, el sistema bloquea la confirmación de carga completa y genera una alerta al jefe de bodega.\newline \emph{SD3, Tabla 3.A.1, fila RF-\allowbreak{}05.07} \\
\textbf{RF-\allowbreak{}06.01}\newline Caso, §§ 4.6 y 9.4 y cap. 18 (crit. 8); SD2, F-08. & M6 Reparto (aplicación nativa Android 15 en Kotlin 2.4 con SQLite; terminal Zebra TC58e); M6 registra cantidad recibida, rechazos y evidencia asociada a la guía; el conductor confirma por línea o marca en bloque toda la guía\newline \textbf{EDT:} 3.4.8 & /P01 a /P03\newline Cumple; Alta; E1; H5 & Los elementos marcados como entregados quedan registrados de forma exitosa en la guía.\newline \emph{SD3, Tabla 3.A.1, fila RF-\allowbreak{}06.01} \\
\textbf{RF-\allowbreak{}06.02}\newline Caso, § 9.4, cap. 15 (RT-16.14 y RT-17.06) y § 16.1 (dec. 15); SD2, F-08 y S-15. & M6 Reparto (aplicación nativa Android 15 en Kotlin 2.4 con SQLite; terminal Zebra TC58e); M6 captura la firma del receptor y la vincula al POD y a la guía de despacho\newline \textbf{EDT:} 3.4.8 & /P01 a /P03\newline Cumple; Alta; E1; H5 & La firma queda correctamente registrada en la guía de despacho.\newline \emph{SD3, Tabla 3.A.1, fila RF-\allowbreak{}06.02} \\
\textbf{RF-\allowbreak{}06.03}\newline Caso, §§ 4.6 y 9.4 y § 16.1 (dec. 12); SD2, S-12. & M6 Reparto (aplicación nativa Android 15 en Kotlin 2.4 con SQLite; terminal Zebra TC58e); M6 registra rechazos parciales o totales, cantidades y causal; el modelo exige causal identificada y tipifica el resultado\newline \textbf{EDT:} 3.4.8 & /P01 a /P02\newline Cumple; Media; E1; H5 & Los productos rechazados quedan registrados de forma correcta en la solicitud de rechazo.\newline \emph{SD3, Tabla 3.A.1, fila RF-\allowbreak{}06.03} \\
\textbf{RF-\allowbreak{}06.04}\newline Caso, § 4.6 y § 16.1 (dec. 3); SD2, S-03. & M6 Reparto (aplicación nativa Android 15 en Kotlin 2.4 con SQLite; terminal Zebra TC58e); M6 registra intentos de entrega y distingue entregas no atendidas; el modelo conserva la dirección del punto cliente\newline \textbf{EDT:} 3.4.8 & /P01 a /P02\newline Cumple; Media; E1; H5 & El envío se registra como fallido y se agenda para ser entregado en la ventana horaria siguiente mas cercana del cliente.\newline \emph{SD3, Tabla 3.A.1, fila RF-\allowbreak{}06.04} \\
\textbf{RF-\allowbreak{}06.05}\newline Caso, §§ 4.7 y 9.7 y § 16.1 (dec. 6); SD2, F-10 y S-06. & M7 Cobranza y rendición (monolito modular en Laravel 13 sobre PHP 8.5; modelo de cobro con monto y medio EFECTIVO); M7 permite capturar cobros en ruta; el modelo registra monto exacto, medio EFECTIVO y vínculo con la entrega\newline \textbf{EDT:} 3.4.10 & /P01 a /P05\newline Cumple; Crítica; E1; H5 & El monto recibido en efectivo por el conductor queda registrado en el sistema.\newline \emph{SD3, Tabla 3.A.1, fila RF-\allowbreak{}06.05} \\
\textbf{RF-\allowbreak{}06.06}\newline Caso, §§ 9.1 y 9.4; SD3, RNG-09. & M6 Reparto: aplicación Kotlin 2.4/Room en Zebra TC58e registra hora real de llegada y término de entrega, con evidencia en Amazon S3\newline \textbf{EDT:} 3.4.8 & /P01 a /P03\newline Cumple; Alta; E1; H5 & Cada entrega tiene hora real de llegada y término, que alimenta el cálculo de OTIF (RF-11.03).\newline \emph{SD3, Tabla 3.A.3b, fila RF-\allowbreak{}06.06} \\
\textbf{RF-\allowbreak{}06.07}\newline Caso, cap. 10 (restr. 3), cap. 15 (RT-03.13 y RT-17.01) y § 9.4; BTT, RT-17.01. & M6 conserva toda la jornada en Room cifrado y sincroniza automáticamente por UUID al recuperar señal o regresar al CD, sin borrar antes del acuse durable. M7 consolida los registros aceptados en la planilla de cierre por camión y turno; el dispositivo de reparto completa la sincronización en no más de 10 min.\newline \textbf{EDT:} 3.4.8, 3.3.6, 3.4.10 & /P01 a /P05\newline Cumple; Crítica; E1; H5 & Al llegar al CD o recuperar señal, la app sube la data y el sistema genera automáticamente el reporte de rendición.\newline \emph{SD3, Tabla 3.A.1, fila RF-\allowbreak{}06.07} \\
\textbf{RF-\allowbreak{}06.08}\newline Caso, cap. 10 (restr. 6), cap. 15 (RT-12.11) y § 16.1 (dec. 5); SD2, S-05. & El jefe de flota solicita el OTP de un solo uso a la autoridad de identidad; despacho aprueba conductor, vehículo y turno. El alta de una persona nueva requiere conexión y no exige correo corporativo. Sin enlace solo inicia un conductor previamente enrolado, con asignación firmada y PIN personal; un sustituto no enrolado no obtiene identidad local.\newline \textbf{EDT:} 3.4.8, 3.3.1, 7.1.3 & /P01 a /P03\newline Cumple; Alta; E1; H5 & Un conductor nuevo de una empresa externa se autentica en 2 minutos con un código que le da su jefe de flota, y puede iniciar su ruta.\newline \emph{SD3, Tabla 3.A.1, fila RF-\allowbreak{}06.08} \\
\textbf{RF-\allowbreak{}06.09}\newline Caso, §§ 4.8 y 9.7, § 16.1 (dec. 10) y cap. 18 (crit. 12); SD2, S-10. & M8 captura cantidad, tipo y estado físico de canastillos y pallets; estado en la evidencia vinculada por env\_movimiento.evidencia\_obj\_id\newline \textbf{EDT:} 3.4.9 & /P01 a /P03\newline Cumple; Alta; E1; H5 & El conductor registra que el cliente devolvió 5 canastillos; el sistema descuenta del saldo del cliente y actualiza el inventario de envases.\newline \emph{SD3, Tabla 3.A.1, fila RF-\allowbreak{}06.09} \\
\textbf{RF-\allowbreak{}06.10}\newline Caso, § 9.7, cap. 10 (restr. 3), § 16.1 (dec. 10) y cap. 18 (crit. 12); SD2, S-10. & M8 Devoluciones y envases (aplicación nativa Android 15 en Kotlin 2.4 con SQLite; terminal Zebra TC58e; cuenta de envases por cliente y tipo); M8 actualiza en línea o de forma diferida los saldos de envases por cliente y parque, asociados al movimiento de entrega o guía\newline \textbf{EDT:} 3.4.9 & /P01 a /P03\newline Cumple; Alta; E1; H5 & El saldo del cliente indica que se devolvieron 5 canastillos y el sistema lo actualiza en el inventario de envases\newline \emph{SD3, Tabla 3.A.1, fila RF-\allowbreak{}06.10} \\
\textbf{RF-\allowbreak{}06.11}\newline Caso, § 9.4, cap. 10 (restr. 5) y § 16.1 (dec. 15); SD2, S-15. & M6 genera en la aplicación del conductor un QR único por entrega para que el cliente con teléfono confirme la recepción. El resultado queda vinculado a la guía y al POD; esa confirmación operacional no sustituye por sí sola el acuse tributario del ERP.\newline \textbf{EDT:} 3.4.8 & /P01 a /P02\newline Cumple; Media; E1; H5 & Se resuelve alternativamente la Decisión \#15; los clientes con smartphone pueden validar su recepción de forma autónoma.\newline \emph{SD3, Tabla 3.A.1, fila RF-\allowbreak{}06.11} \\
\textbf{RF-\allowbreak{}06.12}\newline Caso, § 9.4, cap. 10 (restr. 5), cap. 15 (RT-16.14 y RT-17.06) y § 16.1 (dec. 15); SD2, S-15. & M6 Reparto (aplicación nativa Android 15 en Kotlin 2.4 con SQLite; terminal Zebra TC58e); M6 captura firma en pantalla o evidencia fotográfica sin conexión y vincula el POD a la guía\newline \textbf{EDT:} 3.4.8 & /P01 a /P03\newline Cumple; Alta; E1; H5 & Los clientes sin conectividad y sin smartphone pueden validar su recepción de forma autónoma.\newline \emph{SD3, Tabla 3.A.1, fila RF-\allowbreak{}06.12} \\
\textbf{RF-\allowbreak{}06.13}\newline Caso, cap. 15 (RT-17.06); BTT, RT-17.06. & M6 en TC58e envía por Bluetooth el comprobante de la entrega a la Zebra ZQ620 Plus de cabina cuando el cliente requiere respaldo físico; conserva el vínculo a entrega y guía.\newline \textbf{EDT:} 3.4.8, 6.5.5 & /P01 a /P02\newline Cumple; Media; E1; H5 & El cliente del canal tradicional recibe un comprobante físico pequeño, igual que hoy con la guía, pero generado digitalmente.\newline \emph{SD3, Tabla 3.A.1, fila RF-\allowbreak{}06.13} \\
\textbf{RF-\allowbreak{}06.14}\newline SD2, Anexo 2.1, F-09. & M6 Reparto: aplicación Kotlin 2.4/Room en Zebra TC58e, impresora Zebra ZQ620 Plus y evidencia en Amazon S3; envío de prueba de entrega al ERP y conciliación del acuse por guía\newline \textbf{EDT:} 3.4.8 & /P01 a /P03\newline Cumple; Alta; E1; H5 & Cada guía queda con su estado de acuse (recibido conforme, con observaciones o sin acuse) visible el mismo día; las guías sin acuse generan una alerta para gestión. La solución no emite documentos tributarios.\newline \emph{SD3, Tabla 3.A.3b, fila RF-\allowbreak{}06.14} \\
\textbf{RF-\allowbreak{}07.01}\newline Caso, cap. 18, criterio 10. & M7 Cobranza y rendición (monolito modular en Laravel 13 sobre PHP 8.5); aplicación nativa Android 15 en Kotlin 2.4 con SQLite; terminal Zebra TC58e; M7 suma por camión y turno los cobros de cob\_cobro asociados a la rendición (entregas al contado y abonos en ruta) y guarda el resultado como monto esperado en cob\_cuadratura.monto\_sistema; la planilla de rendición (RF-07.02) lo contrasta con cob\_rendicion.monto\_declarado y registra la diferencia a justificar (RF-07.03)\newline \textbf{EDT:} 3.4.10 & /P01 a /P03\newline Cumple; Alta; E1; H5 & La planilla de rendición (RF-07.02) contrasta lo declarado con el monto esperado y muestra la diferencia a justificar (RF-07.03).\newline \emph{SD3, Tabla 3.A.3b, fila RF-\allowbreak{}07.01} \\
\textbf{RF-\allowbreak{}07.02}\newline Caso, §§ 4.7 y 9.7 y cap. 18 (crit. 10); SD2, F-10. & M7 Cobranza y rendición (monolito modular en Laravel 13 sobre PHP 8.5; ACL INT-06); M7 genera rendición individual por conductor y contrasta cobros con entregas y documentos\newline \textbf{EDT:} 3.4.10 & /P01 a /P03\newline Cumple; Alta; E1; H5 & Despliegue automático de la cuadratura en caja mostrando monto esperado vs recaudado y la diferencia resultante a justificar.\newline \emph{SD3, Tabla 3.A.1, fila RF-\allowbreak{}07.02} \\
\textbf{RF-\allowbreak{}07.03}\newline Caso, §§ 4.7 y 9.7 y cap. 18 (crit. 10); SD2, F-10. & M7 Cobranza y rendición (monolito modular en Laravel 13 sobre PHP 8.5; ACL INT-06); M7 clasifica y conserva las causales de descuadre de la rendición y la auditoría de su aprobación; la rendición no se cierra con saldo sin causal tipificada\newline \textbf{EDT:} 3.4.10 & /P01 a /P03\newline Cumple; Alta; E1; H5 & Cero pesos de descuadre sin motivo formal tipificado, firma digital y asignación de responsabilidad en bitácora.\newline \emph{SD3, Tabla 3.A.1, fila RF-\allowbreak{}07.03} \\
\textbf{RF-\allowbreak{}07.04}\newline Caso, §§ 4.3 y 4.7 y cap. 5 (cuentas por cobrar); SD2, F-10. & M7 Cobranza y rendición (monolito modular en Laravel 13 sobre PHP 8.5; ACL INT-06); la cartera a 30 días se sincroniza con el ERP por documento (cob\_saldo\_documento), cob\_cartera\_asignacion asigna cada cliente a un ejecutivo por intervalo sin solapamiento y cob\_gestion registra contactos, compromisos de pago con monto y fecha, y fechas de recontacto\newline \textbf{EDT:} 3.4.10 & /P01 a /P03\newline Cumple; Alta; E1; H5 & Listado diario de morosidad al día con registro auditable de promesas de pago y tareas automáticas de seguimiento.\newline \emph{SD3, Tabla 3.A.1, fila RF-\allowbreak{}07.04} \\
\textbf{RF-\allowbreak{}07.05}\newline Caso, §§ 4.3 y 9.2 y cap. 18 (crit. 5). & M3 Preventa (monolito modular en Laravel 13 sobre PHP 8.5; aplicación nativa Android 15 en Kotlin 2.4 con SQLite; terminal Zebra TC58e); cob\_saldo\_documento conserva el vencimiento ERP de cada factura y los días de mora a la fecha de corte; con mora superior a 30 días o cupo excedido la aplicación muestra el bloqueo, impide pedidos a crédito y mantiene la venta al contado y el cobro de deuda\newline \textbf{EDT:} 3.4.6 & /P01 a /P03\newline Cumple; Alta; E1; H5 & Alerta en terminal de preventa en menos de 2 s al seleccionar al cliente, impidiendo el ingreso de órdenes a crédito no autorizadas.\newline \emph{SD3, Tabla 3.A.1, fila RF-\allowbreak{}07.05} \\
\textbf{RF-\allowbreak{}07.06}\newline Caso, § 9.7 y cap. 15 (RT-17.06). & M7 Cobranza y rendición (monolito modular en Laravel 13 sobre PHP 8.5; ACL INT-06); aplicación nativa Android 15 en Kotlin 2.4 con SQLite; terminal Zebra TC58e; terminal de pago PAX A920 Pro; la aplicación recibe por Bluetooth del PAX A920 Pro el código de autorización y M7 lo concilia con la entrega en la rendición diaria, sin digitación\newline \textbf{EDT:} 3.4.10, 3.6.2 & /P01 a /P02\newline Cumple; Media; E1; H5 & Captura automática del código de autorización POS y conciliación contra la entrega en la rendición diaria de caja sin digitación manual.\newline \emph{SD3, Tabla 3.A.1, fila RF-\allowbreak{}07.06} \\
\textbf{RF-\allowbreak{}07.07}\newline Caso, cap. 15 (RT-16.30); SD2, F-12. & Portal autenticado de clientes (Angular 22; Keycloak 26; S3 privado, CloudFront, AWS WAF y API REST); el estado de cuenta y la cartola derivan los tramos corriente, 1 a 30, 31 a 60, 61 a 90 y más de 90 días desde el vencimiento de cada documento y la fecha de corte (hech\_cartera\_antiguedad), descargables en PDF y XLSX por el cliente autenticado desde E2\newline \textbf{EDT:} 3.5.2 & /P01 a /P02\newline Cumple; Media; E2; H10 & Acceso seguro del cliente autenticado a su saldo vivo y detalle de facturas pendientes con opción de descarga de cartola.\newline \emph{SD3, Tabla 3.A.1, fila RF-\allowbreak{}07.07} \\
\textbf{RF-\allowbreak{}07.08}\newline Caso, §§ 4.7 y 9.7 y cap. 15 (RT-16.30). & M7 Cobranza y rendición (monolito modular en Laravel 13 sobre PHP 8.5; ACL INT-06); Portal de Clientes; el pago informado en el portal queda «En revisión» en una bandeja donde el cajero de Tesorería verifica el comprobante bancario y autoriza su traspaso estructurado al ERP\newline \textbf{EDT:} 3.5.2, 3.4.10 & /P01 a /P02\newline Cumple; Media; E2; H10 & Pago web queda en estado «En Revisión» hasta que el cajero lo aprueba, momento en que se envía la rebaja contable al ERP.\newline \emph{SD3, Tabla 3.A.1, fila RF-\allowbreak{}07.08} \\
\end{tablasimple}

\begin{tablasimple}{>{\hsize=0.82\hsize\linewidth=\hsize}Y >{\hsize=1.04\hsize\linewidth=\hsize}Y >{\hsize=0.65\hsize\linewidth=\hsize}Y >{\hsize=1.49\hsize\linewidth=\hsize}Y}
{\multicolumn{4}{l}{\textbf{Matriz individual: bloque 4 de 26.}} \\
\cab{Requisito y origen} & \cab{Componente y EDT} & \cab{Prueba prevista y estado} & \cab{Criterio y fuente}}
\textbf{RF-\allowbreak{}07.09}\newline Caso, §§ 4.7 y 9.7, caps. 5 y 11 y cap. 18 (crit. 10). & M7 Cobranza y rendición (monolito modular en Laravel 13 sobre PHP 8.5; ACL INT-06); M7 publica los comprobantes y cierres comerciales aprobados al ERP de forma asíncrona e idempotente\newline \textbf{EDT:} 3.4.10 & /P01 a /P03\newline Cumple; Alta; E1; H5 & Abonos y cancelaciones contables reflejados en el ERP el mismo día sin intervención manual de digitación (Restricción R4).\newline \emph{SD3, Tabla 3.A.1, fila RF-\allowbreak{}07.09} \\
\textbf{RF-\allowbreak{}07.10}\newline Caso, §§ 4.7 y 9.7 y § 16.1 (dec. 6). & M7 Cobranza y rendición (monolito modular en Laravel 13 sobre PHP 8.5; ACL INT-06); aplicación nativa Android 15 en Kotlin 2.4 con SQLite; terminal Zebra TC58e; el conductor registra el abono parcial en una entrega al contado, la aplicación emite un comprobante digital y el saldo insoluto pasa a la cuenta corriente del cliente en el ERP por la ACL\newline \textbf{EDT:} 3.4.10 & /P01 a /P02\newline Cumple; Media; E1; H5 & Registro de abono parcial en app móvil, comprobante digital emitido al cliente y saldo pendiente registrado en cuenta corriente.\newline \emph{SD3, Tabla 3.A.1, fila RF-\allowbreak{}07.10} \\
\textbf{RF-\allowbreak{}07.11}\newline Caso, cap. 8 (Preventa), §§ 4.7 y 9.7 y cap. 18 (crit. 10). & M3 permite al preventista registrar efectivo de facturas vencidas; M7 conserva el cobro PREVENTA sin exigir vehículo ni entrega, aplica el abono a las facturas, emite recibo provisional digital con folio único e incorpora el monto a la rendición diaria obligatoria en caja.\newline \textbf{EDT:} 3.4.6, 3.4.10 & /P01 a /P03\newline Cumple; Alta; E1; H5 & Emisión de recibo digital de cobro en preventa y consolidación automática en la bolsa de dinero a rendir por el vendedor al volver a base.\newline \emph{SD3, Tabla 3.A.1, fila RF-\allowbreak{}07.11} \\
\textbf{RF-\allowbreak{}08.01}\newline Caso, §§ 4.6, 4.8 y 9.4, cap. 10 (restr. 3) y § 16.1 (dec. 12); SD2, S-12. & M8 Devoluciones y envases (aplicación nativa Android 15 en Kotlin 2.4 con SQLite; terminal Zebra TC58e); M8 permite registrar devoluciones durante la entrega desde la aplicación de reparto, incluso sin conexión, con SKU, cantidad y motivo (daño, vencimiento, producto no pedido o error)\newline \textbf{EDT:} 3.4.9 & /P01 a /P05\newline Cumple; Crítica; E1; H5 & La devolución queda registrada en el dispositivo al momento de la entrega, sin papeleo, y se sincroniza con el sistema central al reconectarse.\newline \emph{SD3, Tabla 3.A.1, fila RF-\allowbreak{}08.01} \\
\textbf{RF-\allowbreak{}08.02}\newline Caso, § 4.8, caps. 5 y 11 y § 16.1 (dec. 12); SD2, S-12. & M8 recibe la devolución en andén y registra destino, usuario y fecha; el ERP recibe el ajuste por INT-06. Si hay retención sanitaria, solo una resolución de Calidad en M9 habilita el destino: reingreso o descuento exigen aptitud; devolución al proveedor o destrucción respetan su autorización. Bodega no puede levantar el bloqueo.\newline \textbf{EDT:} 3.4.9, 3.4.4, 3.3.2 & /P01 a /P05\newline Cumple; Crítica; E1; H5 & El destino queda registrado y la transacción llega al ERP; toda decisión queda trazada con usuario, fecha y destino asignado.\newline \emph{SD3, Tabla 3.A.1, fila RF-\allowbreak{}08.02} \\
\textbf{RF-\allowbreak{}08.03}\newline Caso, §§ 4.8 y 9.7, § 16.1 (dec. 10) y cap. 18 (crit. 12); SD2, S-10. & M8 Devoluciones y envases (aplicación nativa Android 15 en Kotlin 2.4 con SQLite; cuenta de envases por cliente y tipo); M8 mantiene saldo de envases retornables por cliente y tipo, y actualiza movimientos de entrega y devolución\newline \textbf{EDT:} 3.4.9 & /P01 a /P03\newline Cumple; Alta; E1; H5 & El saldo de envases por cliente queda actualizado tras cada movimiento y es consultable desde el dispositivo del conductor antes de cada entrega.\newline \emph{SD3, Tabla 3.A.1, fila RF-\allowbreak{}08.03} \\
\textbf{RF-\allowbreak{}08.04}\newline Caso, § 9.7, § 16.1 (dec. 10) y cap. 18 (crit. 12); SD2, S-10. & M8 Devoluciones y envases (aplicación nativa Android 15 en Kotlin 2.4 con SQLite; cuenta de envases por cliente y tipo); M8 compara el saldo de envases de cada cliente con el umbral del administrador y alerta al preventista dentro de las 24 horas siguientes a superarlo\newline \textbf{EDT:} 3.4.9 & /P01 a /P02\newline Cumple; Media; E1; H5 & El preventista recibe la notificación en $\le$ 24 horas de superado el umbral.\newline \emph{SD3, Tabla 3.A.1, fila RF-\allowbreak{}08.04} \\
\textbf{RF-\allowbreak{}08.05}\newline Caso, §§ 4.8, 7.2 y 9.10 y § 16.2 (logística inversa); SD2, Tabla 2.1. & M2/M5 registran merma por vencimiento en conteo y preparación; M3 registra la detectada por el preventista en la góndola, y M8 tramita el retorno si ocurre. Cada registro identifica SKU, lote, cantidad, causa y origen. Un aviso sobre mercadería ya entregada no vuelve a descontar stock físico de la bodega.\newline \textbf{EDT:} 3.4.2, 3.4.3, 3.4.6, 3.4.9 & /P01 a /P03\newline Cumple; Alta; E1; H5 & La merma identifica SKU, lote, cantidad y origen (conteo, preparación o góndola). El ajuste afecta las existencias bajo custodia correspondientes. Un aviso en góndola de producto ya entregado no vuelve a descontar existencias físicas de la bodega; si retorna, M8 registra el evento.\newline \emph{SD3, Tabla 3.A.1, fila RF-\allowbreak{}08.05} \\
\textbf{RF-\allowbreak{}08.06}\newline Caso, §§ 4.8, 7.2 y 9.7 y cap. 18 (crit. 12); SD2, S-10. & M8 Devoluciones y envases (movimientos por cliente y tipo); M10 Analítica (S3, Glue, Redshift Serverless y QuickSight); M8 conserva movimientos de envases por cliente y tipo, y M10 informa a diario la pérdida por cliente, tipo y zona\newline \textbf{EDT:} 3.4.9 & /P01 a /P03\newline Cumple; Alta; E1; H5 & El reporte identifica los clientes y zonas con mayor pérdida de envases, permitiendo priorizar acciones de recuperación.\newline \emph{SD3, Tabla 3.A.1, fila RF-\allowbreak{}08.06} \\
\textbf{RF-\allowbreak{}08.07}\newline Caso, §§ 4.8 y 9.8, § 16.1 (dec. 7) y cap. 18 (crit. 11); SD2, F-14. & M7 y M10 (monolito modular en Laravel 13 sobre PHP 8.5; almacén analítico S3, Glue, Redshift Serverless y QuickSight); inv\_merma registra la merma por conteo, preparación, góndola o devolución con su cliente, y hech\_actividad\_costo suma al costo de servir de cada cliente el manejo de devoluciones y la merma valorizada al costo ERP, una sola vez por origen, desde E2\newline \textbf{EDT:} 3.5.4, 3.4.9 & /P01 a /P02\newline Cumple; Media; E2; H10 & El costo de servir de un cliente con alta tasa de devolución refleja ese costo adicional.\newline \emph{SD3, Tabla 3.A.1, fila RF-\allowbreak{}08.07} \\
\textbf{RF-\allowbreak{}09.01}\newline Caso, § 9.5 y cap. 18 (crit. 1); SD2, F-11 y S-02. & M9 Calidad y trazabilidad (EPCIS 2.0/CBV 2.0; AWS IoT Core; Amazon DynamoDB); M9 permite consultar por lote los clientes receptores y enlaza sus entregas, guías, fechas y cantidades\newline \textbf{EDT:} 3.4.4 & /P01 a /P05\newline Cumple; Crítica; E1; H5 & En un simulacro de retiro sanitario, se obtiene en menos de 2 horas la lista de clientes afectados, con trazabilidad desde el proveedor hasta los clientes y archivo exportable (RNF-09.01).\newline \emph{SD3, Tabla 3.A.1, fila RF-\allowbreak{}09.01} \\
\textbf{RF-\allowbreak{}09.02}\newline Caso, § 9.5 y § 16.1 (dec. 2); SD2, F-11 y S-02. & M9 Calidad y trazabilidad (EPCIS 2.0/CBV 2.0; AWS IoT Core; Amazon DynamoDB); M9 permite rastrear el lote de producto hasta la recepción y su proveedor\newline \textbf{EDT:} 3.4.4 & /P01 a /P05\newline Cumple; Crítica; E1; H5 & El equipo de calidad puede rastrear el origen de cualquier unidad devuelta, realizando la trazabilidad desde los clientes afectados hasta los productores\newline \emph{SD3, Tabla 3.A.1, fila RF-\allowbreak{}09.02} \\
\textbf{RF-\allowbreak{}09.03}\newline Caso, §§ 4.9, 7.3 y 9.6, cap. 15 (RT-17.06) y cap. 18 (crit. 3); SD2, F-07. & M9 Calidad y trazabilidad; 21 puntos de cámara B-01 (Ebyte ME31-XDXX0400 con PT100) y tres gateways B-02 Moxa UC-8200 con Greengrass V2. Registro inicial cada 5 min, 288 lecturas por punto y día; evaluación y bloqueo locales. El CLIENTE adquiere los sensores B-01 según la especificación de LafroX (\textit{Caso: Logística}, 2026, cap. 11, p. 20); LafroX los instala y mantiene, y provee los gateways B-02 como infraestructura on-premise.\newline \textbf{EDT:} 3.4.4, 6.5.3 & /P01 a /P05\newline Cumple; Crítica; E1; H5 & El sistema registrara las temperaturas. Los registros son auditable y se puede exportar para la autoridad sanitaria.\newline \emph{SD3, Tabla 3.A.1, fila RF-\allowbreak{}09.03} \\
\textbf{RF-\allowbreak{}09.04}\newline Caso, §§ 2.3, 4.9 y 9.6 y cap. 18 (crit. 3); SD2, F-07; SD3, S-28. & M9; kit en los 28 camiones: gateway B-04 con Greengrass V2, ioLogik B-05, tres PT100 B-06 y simuladores B-07, sensor independiente B-08 y termógrafo de aire B-03 EYE SENSOR EN12830. Sondas cada 5 s y serie por minuto; móvil e Iridium B-09 siempre activos. El CLIENTE adquiere el kit según la especificación de LafroX (\textit{Caso: Logística}, 2026, cap. 11, p. 20), y LafroX lo instala y mantiene.\newline \textbf{EDT:} 3.4.4, 3.4.5, 6.5.4 & /P01 a /P05\newline Cumple; Crítica; E1; H5 & El sistema registrara las temperaturas de los camiones. El cliente o Puelche pueden consultar el registro histórico de temperatura de su envío.\newline \emph{SD3, Tabla 3.A.1, fila RF-\allowbreak{}09.04} \\
\textbf{RF-\allowbreak{}09.05}\newline Caso, §§ 4.9 y 9.6 y § 16.1 (dec. 4); SD2, S-04; SD3, RNG-04. & M9 evalúa la regla por familia en Greengrass V2; la desviación de 2 °C durante 15 min se parametriza conforme a la ficha y los límites aplicables. Alarma cableada B-12 sin red y segunda evaluación en TC58e con B-08. Aviso por móvil e Iridium en paralelo, metas $\leq$ 2 min y $\leq$ 5 min respectivamente.\newline \textbf{EDT:} 3.4.4, 3.4.5, 6.5.4 & /P01 a /P05\newline Cumple; Crítica; E1; H5 & El sistema al registrar una temperatura fuera del rango aceptable, generara una alerta tanto al conductor como al equipo de calidad\newline \emph{SD3, Tabla 3.A.1, fila RF-\allowbreak{}09.05} \\
\textbf{RF-\allowbreak{}09.06}\newline Caso, §§ 4.9 y 9.6 y § 16.1 (dec. 4); SD2, S-04; SD3, RNG-04 y RNG-17. & M9 y TC58e retienen por SSCC e incidente, según su familia; M2/M5 impiden preparar o despachar lo retenido y M6 impide su entrega. La advertencia no libera una crítica. Calidad decide la disposición y, sin decisión en plazo, el SSCC vuelve retenido al CD.\newline \textbf{EDT:} 3.4.4, 3.4.2, 3.4.3, 3.4.8 & /P01 a /P05\newline Cumple; Crítica; E1; H5 & El lote queda retenido automáticamente ante excursión crítica; el conductor ve la alerta y registra observaciones; solo Calidad libera, bloquea o rechaza, y su decisión queda registrada.\newline \emph{SD3, Tabla 3.A.1, fila RF-\allowbreak{}09.06} \\
\textbf{RF-\allowbreak{}09.07}\newline Caso, §§ 4.6, 4.9 y 9.6 y § 16.1 (dec. 12). & M6 registra el rechazo por sospecha de ruptura con causal, SSCC, sensor, compartimento y tramo temporal de la ruta; M9 lo vincula a la historia térmica entre carga y bajada y al tramo de entrega. La temperatura no incorpora GPS.\newline \textbf{EDT:} 3.4.8, 3.4.4 & /P01 a /P03\newline Cumple; Alta; E1; H5 & Si se rechaza el pedido, el conductor registra el rechazo ingresando evidencia\newline \emph{SD3, Tabla 3.A.1, fila RF-\allowbreak{}09.07} \\
\textbf{RF-\allowbreak{}09.08}\newline Caso, §§ 4.5 y 9.5, § 16.1 (dec. 2) y cap. 18 (crit. 1); SD2, S-02. & M5 Preparación, aplicación Kotlin 2.4/Room y terminal Zebra MC9400, con variante Cold Storage Freezer para congelado, registra el lote elegido por FEFO en cada línea de picking. Los eventos de preparación y carga alimentan la trazabilidad M9 por lote y SSCC.\newline \textbf{EDT:} 3.4.3, 3.4.4 & /P01 a /P05\newline Cumple; Crítica; E1; H5 & El picking digital registra qué lote concreto sale en cada canastillo/pallet hacia el cliente.\newline \emph{SD3, Tabla 3.A.1, fila RF-\allowbreak{}09.08} \\
\textbf{RF-\allowbreak{}09.09}\newline Caso, §§ 9.6 y 12 (inocuidad) y cap. 18 (crit. 3); SD2, F-07. & M9 construye el expediente por envío y SSCC con la curva del compartimento durante su permanencia, exposición de andén, tres series conciliadas, eventos y decisiones. TC58e descarga ambos registradores, verifica integridad y firma hash; S3 Object Lock conserva el original y la temperatura cinco años. Consulta y exportación para cliente y autoridad.\newline \textbf{EDT:} 3.4.4, 3.4.5, 3.2.4 & /P01 a /P05\newline Cumple; Crítica; E1; H5 & La evidencia de cadena de frío de un envío se obtiene sin reconstrucción manual.\newline \emph{SD3, Tabla 3.A.1, fila RF-\allowbreak{}09.09} \\
\textbf{RF-\allowbreak{}09.10}\newline Caso, sección 16.1, decisión 4; SD2, Anexo 2.2, S-04. & M9 parametriza y versiona por familia límites, exposición, máximos de andén, plazos y tolerancias; Calidad aprueba con segundo aprobador. La regla se distribuye a Greengrass V2 de cámaras y camiones y a la implementación independiente del TC58e. Sin parámetros aprobados rige la regla preventiva de ficha; una versión nueva no libera retenciones históricas.\newline \textbf{EDT:} 3.4.4, 3.4.5 & /P01 a /P02\newline Cumple; Media; E1; H5 & La alerta y la retención de RF-09.05 y RF-09.06 aplican el parámetro vigente, y cada cambio queda registrado.\newline \emph{SD3, Tabla 3.A.3b, fila RF-\allowbreak{}09.10} \\
\textbf{RF-\allowbreak{}09.11}\newline Propio LafroX. & M5 registra los escaneos de salida de cámara y carga; M9 calcula exposición de andén por SSCC, la compara con el máximo de su familia aprobado por Calidad y bloquea antes de salir cuando lo supera. El vínculo de custodia y ventana conserva los instantes y el compartimento.\newline \textbf{EDT:} 3.4.3, 3.4.4, 4.1.4 & /P01 a /P05\newline Cumple; Crítica; E1; H5 & Exposición de andén trazable y bloqueo antes de salir si excede el máximo.\newline \emph{SD3, Tabla 3.A.3b, fila RF-\allowbreak{}09.11} \\
\textbf{RF-\allowbreak{}09.12}\newline Propio LafroX. & M6 en TC58e escanea la bajada y exige manifiesto vigente, latido firmado de menos de 30 s y ausencia de retención y contradicción. Cierra el intervalo térmico a bordo y habilita el POD posterior dentro del local sin Wi-Fi de cabina; una recarga abre otra ventana.\newline \textbf{EDT:} 3.4.8, 3.4.4, 4.1.4 & /P01 a /P05\newline Cumple; Crítica; E1; H5 & Fin de la ventana a bordo; POD posterior permitido dentro del local aun fuera del Wi-Fi de cabina.\newline \emph{SD3, Tabla 3.A.3b, fila RF-\allowbreak{}09.12} \\
\textbf{RF-\allowbreak{}09.13}\newline Propio LafroX. & TC58e, servicio en primer plano, evalúa el sensor independiente B-08 por Bluetooth y comprueba el latido firmado del gateway B-04 por Wi-Fi local. Crítica, contradicción o pérdida de latido sin bajada en curso impiden una nueva bajada; una entrega previamente habilitada puede terminar fuera del alcance. No depende de móvil ni satélite.\newline \textbf{EDT:} 3.4.8, 3.4.4, 4.1.4 & /P01 a /P05\newline Cumple; Crítica; E1; H5 & Bloqueo de una nueva bajada con estado no confiable; lo retenido vuelve al CD si no existe decisión en plazo.\newline \emph{SD3, Tabla 3.A.3b, fila RF-\allowbreak{}09.13} \\
\textbf{RF-\allowbreak{}09.14}\newline Propio LafroX. & M9 registra la resolución de Calidad: liberar, venta con descuento solo si apto, devolver al proveedor, uso no alimentario autorizado o destrucción autorizada. M8 tramita retorno/disposición y M2 mantiene el bloqueo en el área de no aptos marcada con X roja. La autorización de entrega exige además estado térmico vigente; los ajustes tributarios los emite el ERP.\newline \textbf{EDT:} 3.4.4, 3.4.9, 3.4.8, 7.2.4, 4.1.4 & /P01 a /P05\newline Cumple; Crítica; E1; H5 & Decisión trazable; área de no aptos separada y marcada con X roja u otro signo inequívoco (RSA, art. 105).\newline \emph{SD3, Tabla 3.A.3b, fila RF-\allowbreak{}09.14} \\
\textbf{RF-\allowbreak{}10.01}\newline Caso, sección 4.1 y cap. 5. & M2 Inventario genera sugerencias revisables de reposición por SKU y proveedor a partir de ventas y condiciones de abastecimiento\newline \textbf{EDT:} 3.4.2 & /P01 a /P02\newline Cumple; Media; E1; H5 & La sugerencia reemplaza la planilla de reposición y queda disponible para revisión (D-05).\newline \emph{SD3, Tabla 3.A.3b, fila RF-\allowbreak{}10.01} \\
\end{tablasimple}

\begin{tablasimple}{>{\hsize=0.82\hsize\linewidth=\hsize}Y >{\hsize=1.04\hsize\linewidth=\hsize}Y >{\hsize=0.65\hsize\linewidth=\hsize}Y >{\hsize=1.49\hsize\linewidth=\hsize}Y}
{\multicolumn{4}{l}{\textbf{Matriz individual: bloque 5 de 26.}} \\
\cab{Requisito y origen} & \cab{Componente y EDT} & \cab{Prueba prevista y estado} & \cab{Criterio y fuente}}
\textbf{RF-\allowbreak{}10.02}\newline Caso, sección 4.1. & M2 Inventario incorpora promociones comprometidas, estacionalidad y condiciones de proveedor al cálculo de reposición\newline \textbf{EDT:} 3.4.2 & /P01 a /P02\newline Cumple; Media; E1; H5 & La sugerencia incluye las promociones sin agregarlas a mano.\newline \emph{SD3, Tabla 3.A.3b, fila RF-\allowbreak{}10.02} \\
\textbf{RF-\allowbreak{}10.03}\newline Caso, sección 4.1; cap. 10. & M2 Inventario traspasa al ERP la sugerencia de reposición aprobada mediante INT-06\newline \textbf{EDT:} 3.4.2, 3.3.2 & /P01 a /P02\newline Cumple; Media; E1; H5 & El ERP recibe la sugerencia aprobada y conserva la emisión de la orden de compra.\newline \emph{SD3, Tabla 3.A.3b, fila RF-\allowbreak{}10.03} \\
\textbf{RF-\allowbreak{}11.01}\newline Caso, §§ 7.1, 7.2, 9.1, 9.8 y 9.10. & M10 Analítica (S3, Glue, Redshift Serverless y QuickSight embebido); M10 envía por correo los informes ejecutivos en PDF y XLSX cada lunes a las 06:00 y el día 2 de cada mes, con OTIF, fill rate, merma por vencimiento y costo de servir por zona y canal\newline \textbf{EDT:} 3.5.4, 3.4.11 & /P01 a /P02\newline Cumple; Media; E2; H10 & Reportes emitidos automáticamente a las casillas de los gerentes el día 2 de cada mes y cada lunes a las 06:00 AM sin intervención manual.\newline \emph{SD3, Tabla 3.A.1, fila RF-\allowbreak{}11.01} \\
\textbf{RF-\allowbreak{}11.02}\newline Caso, § 9.8, § 16.1 (dec. 7) y cap. 18 (crit. 11); SD2, F-14. & M10 Analítica (S3, Glue, Redshift Serverless y QuickSight embebido); M10 calcula el costo por entrega y cliente con combustible según kilometraje GPS, peajes de la ruta y tiempos de conducción y descarga, y prorratea los costos indirectos de bodega y administración\newline \textbf{EDT:} 3.5.4 & /P01 a /P02\newline Cumple; Media; E2; H10 & Ficha de rentabilidad por cliente actualizada tras cada liquidación, desglosando costos logísticos directos e indirectos y margen neto resultante.\newline \emph{SD3, Tabla 3.A.1, fila RF-\allowbreak{}11.02} \\
\textbf{RF-\allowbreak{}11.03}\newline Caso, §§ 9.1 y 16.1 (dec. 1) y cap. 18 (crit. 4); SD2, S-01; SD3, RNG-09 y R18-04. & M10 Analítica (S3, Glue, Redshift Serverless y QuickSight embebido); M10 calcula a diario OTIF por cliente, ruta, zona y consolidado, con definición de On-Time/In-Full y registro de causal ante desvíos\newline \textbf{EDT:} 3.4.11 & /P01 a /P03\newline Cumple; Alta; E1; H5 & Todas las gerencias visualizan el mismo valor de OTIF diario; ante entregas parciales o tardías se audita el motivo en el sistema.\newline \emph{SD3, Tabla 3.A.1, fila RF-\allowbreak{}11.03} \\
\textbf{RF-\allowbreak{}11.04}\newline Caso, § 7.1, § 9.1, cap. 10 (restr. 5) y § 16.2 (indicadores); SD2, S-01. & M10 Analítica (S3, Glue, Redshift Serverless y QuickSight embebido); M10 calcula fill rate e índice de pedido perfecto contra el pedido original con las entregas confirmadas en el terminal de reparto, desglosados por causa\newline \textbf{EDT:} 3.4.11 & /P01 a /P02\newline Cumple; Media; E1; H5 & Indicador mensual de Perfect Order desglosado por causa raíz (daño, faltante, atraso, documento) contrastado contra la orden de venta original.\newline \emph{SD3, Tabla 3.A.1, fila RF-\allowbreak{}11.04} \\
\textbf{RF-\allowbreak{}11.05}\newline Caso, §§ 4.4 y 7.2 y cap. 18 (crit. 14); SD3, S-27 y R18-14. & M10 Analítica (S3, Glue, Redshift Serverless y QuickSight embebido); hech\_viaje conserva por viaje despachado las capacidades de peso y volumen del vehículo y la carga de salida, calcula ambos porcentajes y alerta cuando la mayor utilización es inferior al umbral configurable y versionado, con referencia inicial de 65 \%\newline \textbf{EDT:} 3.4.11 & /P01 a /P03\newline Cumple; Alta; E1; H5 & Generación de bitácora de viajes subutilizados en la ventana de 05:30 a 07:00 para auditoría y balanceo de carga de transporte.\newline \emph{SD3, Tabla 3.A.1, fila RF-\allowbreak{}11.05} \\
\textbf{RF-\allowbreak{}11.06}\newline Caso, §§ 9.1, 9.6 y 9.7 y cap. 15 (RT-05.29); BTT, RT-05.29. & La consola operacional consume las API autorizadas de M4/M6/M8/M9/M7 para avance de rutas, entregas cumplidas, devoluciones en tránsito, alertas térmicas y efectivo recaudado, con antigüedad visible y latencia no mayor de 5 min. M10 conserva los históricos en Redshift/QuickSight sin consultar directamente la base operacional.\newline \textbf{EDT:} 3.4.11 & /P01 a /P03\newline Cumple; Alta; E1; H5 & Despliegue en vivo en pantallas gerenciales con actualización continua de la distribución y alertas críticas de ruta.\newline \emph{SD3, Tabla 3.A.1, fila RF-\allowbreak{}11.06} \\
\textbf{RF-\allowbreak{}11.07}\newline Caso, §§ 4.7, 4.8 y 9.7, cap. 15 (RT-05.29) y cap. 18 (crit. 10 y 12); SD2, F-10 y S-10. & M7 genera la liquidación individual al retorno y recepción de los registros del turno: documentos tributarios entregados, devoluciones recibidas de M8, efectivo y balance de envases. Cuadra monto esperado/declarado y exige causal de diferencias; M10 solo consulta el cierre. El cierre del día queda disponible en no más de 2 h tras el retorno del último camión.\newline \textbf{EDT:} 3.4.10, 3.4.9, 3.4.11 & /P01 a /P03\newline Cumple; Alta; E1; H5 & Planilla de cuadratura emitida inmediatamente al arribo del camión para tesorería y bodega, sin rezagos para el día siguiente.\newline \emph{SD3, Tabla 3.A.1, fila RF-\allowbreak{}11.07} \\
\textbf{RF-\allowbreak{}11.08}\newline Caso, § 9.8 y § 16.1 (dec. 7); SD2, F-14 y S-07. & M10 Analítica (S3, Glue, Redshift Serverless y QuickSight embebido); M10 clasifica cada mes a los clientes en rentables, no rentables y estratégicos según margen comercial y costo de servir, y sugiere ajustes de frecuencia de visita o de compra mínima\newline \textbf{EDT:} 3.5.4 & /P01 a /P02\newline Cumple; Media; E2; H10 & Clasificación automática de clientes en rentables, no rentables y estratégicos, con simulación de impacto de cambio de frecuencias.\newline \emph{SD3, Tabla 3.A.1, fila RF-\allowbreak{}11.08} \\
\textbf{RF-\allowbreak{}12.01}\newline Caso, § 9.9, §§ 13.2 y 16.2 (EDI) y cap. 18 (crit. 13); BTT, RT-05.23; SD2, F-13. & M11 Canal moderno (monolito modular en Laravel 13 sobre PHP 8.5; transporte OpenAS2 4.9 BSD; perfiles EDI por cadena; INT-08); M11 recibe mensajes comerciales EANCOM y GS1 XML mediante contratos configurados por cadena.\newline \textbf{EDT:} 3.5.1, 3.6.5, 2.1.4 & /P01 a /P03\newline Cumple; Alta; E2; H10 & El mensaje recibido queda disponible para su validación (RF-12.02) en un plazo no superior a 15 segundos, identificando la cadena de origen.\newline \emph{SD3, Tabla 3.A.1, fila RF-\allowbreak{}12.01} \\
\textbf{RF-\allowbreak{}12.02}\newline Caso, § 9.9 y cap. 18 (crit. 13); BTT, RT-05.23; SD2, F-13. & M11 Canal moderno (reglas del monolito Laravel 13 sobre PHP 8.5; transporte OpenAS2 4.9 BSD; perfiles EDI por cadena; INT-08); M11 valida estructura de intercambios y datos de pedido antes de pasarlos a M3\newline \textbf{EDT:} 3.5.1, 3.6.5, 2.1.4 & /P01 a /P03\newline Cumple; Alta; E2; H10 & El pedido queda marcado como válido o inválido, con el detalle de cada validación fallida.\newline \emph{SD3, Tabla 3.A.1, fila RF-\allowbreak{}12.02} \\
\textbf{RF-\allowbreak{}12.03}\newline Caso, § 9.9 y § 16.1 (dec. 11); SD2, S-11. & M11 Canal moderno (monolito modular en Laravel 13 sobre PHP 8.5; transporte OpenAS2 4.9 BSD; perfiles EDI por cadena; INT-08); M11 mantiene equivalencias de códigos por cadena y traduce cada código al producto canónico vigente.\newline \textbf{EDT:} 3.5.1, 3.6.5, 2.1.4 & /P01 a /P03\newline Cumple; Alta; E2; H10 & Cada línea queda asociada a un SKU válido del maestro interno, o marcada como discrepancia no resuelta.\newline \emph{SD3, Tabla 3.A.1, fila RF-\allowbreak{}12.03} \\
\textbf{RF-\allowbreak{}12.04}\newline Caso, § 9.9, caps. 5 y 11 y cap. 18 (crit. 13); SD2, F-13. & M11 Canal moderno (reglas del monolito Laravel 13 sobre PHP 8.5; transporte OpenAS2 4.9 BSD; perfiles EDI por cadena; INT-08); M11 transforma el mensaje válido y la ACL INT-06 lo deriva al ERP\newline \textbf{EDT:} 3.5.1, 3.6.5, 2.1.4 & /P01 a /P03\newline Cumple; Alta; E2; H10 & El pedido aparece como pedido único en el sistema de gestión dentro de 2 minutos desde su validación.\newline \emph{SD3, Tabla 3.A.1, fila RF-\allowbreak{}12.04} \\
\textbf{RF-\allowbreak{}12.05}\newline Caso, §§ 4.4, 9.3 y 9.9; SD2, F-06 y F-13. & M11 Canal moderno (reglas del monolito Laravel 13 sobre PHP 8.5; transporte OpenAS2 4.9 BSD; perfiles EDI por cadena; INT-08); M11 publica por ACL el pedido normalizado para su consumo por planificación y M3\newline \textbf{EDT:} 3.5.1, 3.6.5, 2.1.4 & /P01 a /P03\newline Cumple; Alta; E2; H10 & El pedido aparece en el listado de pedidos confirmados del módulo de planificación, sin digitación manual.\newline \emph{SD3, Tabla 3.A.1, fila RF-\allowbreak{}12.05} \\
\textbf{RF-\allowbreak{}12.06}\newline Caso, § 9.9 y § 16.1 (dec. 11); SD2, S-11. & M11 mantiene una bandeja por pedido con causa de validación o equivalencia, estado, fecha y responsable. El perfil comercial autorizado corrige, rechaza o reprocesa con auditoría; solo un pedido válido se entrega a M3.\newline \textbf{EDT:} 3.5.1, 3.6.5, 3.6.6, 2.1.4 & /P01 a /P03\newline Cumple; Alta; E2; H10 & Los pedidos con error aparecen en la bandeja con causa, estado, fecha y responsable.\newline \emph{SD3, Tabla 3.A.1, fila RF-\allowbreak{}12.06} \\
\textbf{RF-\allowbreak{}12.07}\newline Caso, § 9.9 y § 16.2 (mensajes EDI); SD2, F-13. & M11 Canal moderno (reglas del monolito Laravel 13 sobre PHP 8.5; transporte OpenAS2 4.9 BSD; perfiles EDI por cadena; INT-08); M11 responde a la cadena con la confirmación o el rechazo del pedido y su causa; el MDN de transporte no sustituye esa respuesta\newline \textbf{EDT:} 3.5.1, 3.6.5, 2.1.4 & /P01 a /P03\newline Cumple; Alta; E2; H10 & La cadena recibe automáticamente el mensaje correspondiente.\newline \emph{SD3, Tabla 3.A.1, fila RF-\allowbreak{}12.07} \\
\textbf{RF-\allowbreak{}12.08}\newline Caso, § 9.9 y cap. 18 (crit. 8 y 13); SD2, F-13. & M11 Canal moderno (reglas del monolito Laravel 13 sobre PHP 8.5; transporte OpenAS2 4.9 BSD; perfiles EDI por cadena; INT-08); M11 conserva el historial de estados de cada pedido electrónico: recepción, validación, confirmación, preparación, despacho, ASN y prueba de entrega\newline \textbf{EDT:} 3.5.1, 3.6.5, 2.1.4 & /P01 a /P03\newline Cumple; Alta; E2; H10 & Se puede consultar el estado actual y el historial de eventos con fecha, hora y responsable/sistema\newline \emph{SD3, Tabla 3.A.1, fila RF-\allowbreak{}12.08} \\
\textbf{RF-\allowbreak{}12.09}\newline Caso, §§ 9.9 y 13.2, cap. 8 (Comercial) y cap. 18 (crit. 13); SD2, F-13. & M11 Canal moderno (reglas del monolito Laravel 13 sobre PHP 8.5; transporte OpenAS2 4.9 BSD; perfiles EDI por cadena; INT-08); M11 envía por INT-08 el ASN al quedar el pedido cargado, con referencia a la guía emitida, productos, cantidades y hora estimada de llegada\newline \textbf{EDT:} 3.5.1, 3.6.5, 2.1.4 & /P01 a /P03\newline Cumple; Alta; E2; H10 & El cliente recibe el ASN dentro de los 2 minutos posteriores a la carga.\newline \emph{SD3, Tabla 3.A.1, fila RF-\allowbreak{}12.09} \\
\textbf{RF-\allowbreak{}12.10}\newline Caso, §§ 9.9 y 13.2 y cap. 15 (RT-10.05); SD2, F-13. & M4 Rutas (monolito Laravel 13 sobre PHP 8.5; motor propio de optimización en contenedor ECS Fargate; INT-10); M4 planifica la llegada al cliente del canal moderno dentro de la ventana de 30 minutos comprometida\newline \textbf{EDT:} 3.5.1, 3.4.7 & /P01 a /P03\newline Cumple; Alta; E2; H10 & El camión llega dentro de la ventana comprometida en el 80 \% de los casos.\newline \emph{SD3, Tabla 3.A.1, fila RF-\allowbreak{}12.10} \\
\textbf{RF-\allowbreak{}12.11}\newline Caso, § 13.2, cap. 15 (RT-10.05) y cap. 10 (restr. 1). & M4 Rutas (monolito Laravel 13 sobre PHP 8.5; motor propio de optimización en contenedor ECS Fargate; INT-10); M6 registra el cumplimiento o incumplimiento de la ventana de 30 minutos de cada entrega del canal moderno, separado del control de la ventana de despacho de 05:30 a 07:00\newline \textbf{EDT:} 3.5.1, 3.4.7, 3.4.8, 3.4.11 & /P01 a /P03\newline Cumple; Alta; E2; H10 & Toda desviación queda registrada con hora real de llegada.\newline \emph{SD3, Tabla 3.A.1, fila RF-\allowbreak{}12.11} \\
\textbf{RF-\allowbreak{}12.12}\newline Caso, § 9.4, cap. 15 (RT-16.14 y RT-17.06) y § 16.1 (dec. 15); SD2, S-15. & M6 Reparto (aplicación nativa Android 15 en Kotlin 2.4 con SQLite; terminal Zebra TC58e; POD vinculado a guía); M6 produce POD operativo con receptor, evidencia y asociación a guía; el Anexo 4-U separa acuse operativo y tributario\newline \textbf{EDT:} 3.5.1, 3.4.8 & /P01 a /P03\newline Cumple; Alta; E2; H10 & Queda un único registro de evidencia por entrega, con el mecanismo utilizado identificado.\newline \emph{SD3, Tabla 3.A.1, fila RF-\allowbreak{}12.12} \\
\textbf{RF-\allowbreak{}12.13}\newline Caso, §§ 9.4 y 9.9, cap. 15 (RT-16.14) y § 16.1 (dec. 15); SD2, F-09. & M6 Reparto (aplicación nativa Android 15 en Kotlin 2.4 con SQLite; terminal Zebra TC58e; POD vinculado a guía); M11 envía por INT-08 el acuse comercial digital dentro de 30 minutos desde la descarga, referido a la guía y a su ASN; el Anexo 4-U distingue el POD operativo del acuse tributario\newline \textbf{EDT:} 3.5.1, 3.4.8 & /P01 a /P03\newline Cumple; Alta; E2; H10 & El cliente recibe el comprobante dentro de los 30 minutos posteriores a la descarga.\newline \emph{SD3, Tabla 3.A.1, fila RF-\allowbreak{}12.13} \\
\textbf{RF-\allowbreak{}12.14}\newline Caso, cap. 15 (RT-16.30). & Portales N-01 a N-03 (Angular 22; S3 privado, CloudFront, AWS WAF y API REST pública); el portal público muestra catálogo sin precios ni stock; SD4 describe contenido institucional\newline \textbf{EDT:} 3.5.2 & /P01 a /P02\newline Cumple; Media; E2; H10 & Cualquier visitante ve el catálogo sin poder generar pedidos ni ver precios.\newline \emph{SD3, Tabla 3.A.1, fila RF-\allowbreak{}12.14} \\
\textbf{RF-\allowbreak{}12.15}\newline Caso, § 9.2 y cap. 15 (RT-16.30); SD2, F-12. & Portales N-01 a N-03 (Angular 22; S3 privado, CloudFront, AWS WAF y API REST pública); Keycloak 26; el portal autenticado permite crear pedidos en autoservicio y valida crédito y stock\newline \textbf{EDT:} 3.5.2 & /P01 a /P02\newline Cumple; Media; E2; H10 & El cliente genera un pedido válido sin intervención de un preventista.\newline \emph{SD3, Tabla 3.A.1, fila RF-\allowbreak{}12.15} \\
\end{tablasimple}

\begin{tablasimple}{>{\hsize=0.82\hsize\linewidth=\hsize}Y >{\hsize=1.04\hsize\linewidth=\hsize}Y >{\hsize=0.65\hsize\linewidth=\hsize}Y >{\hsize=1.49\hsize\linewidth=\hsize}Y}
{\multicolumn{4}{l}{\textbf{Matriz individual: bloque 6 de 26.}} \\
\cab{Requisito y origen} & \cab{Componente y EDT} & \cab{Prueba prevista y estado} & \cab{Criterio y fuente}}
\textbf{RF-\allowbreak{}12.16}\newline Caso, § 9.1 y cap. 15 (RT-16.30); SD2, F-12. & Portales N-01 a N-03 (Angular 22; S3 privado, CloudFront, AWS WAF y API REST pública); Keycloak 26; el portal de clientes permite consultar el estado de entregas y documentación disponible\newline \textbf{EDT:} 3.5.2 & /P01 a /P02\newline Cumple; Media; E2; H10 & El cliente visualiza el estado y compromiso de cada pedido.\newline \emph{SD3, Tabla 3.A.1, fila RF-\allowbreak{}12.16} \\
\textbf{RF-\allowbreak{}12.17}\newline Caso, cap. 15 (RT-16.30); SD2, F-12. & Portales N-01 a N-03 (Angular 22; S3 privado, CloudFront, AWS WAF y API REST pública); Keycloak 26; el portal de clientes permite consultar y descargar documentos disponibles de su cuenta\newline \textbf{EDT:} 3.5.2 & /P01 a /P02\newline Cumple; Media; E2; H10 & El cliente descarga los documentos correspondientes a sus operaciones.\newline \emph{SD3, Tabla 3.A.1, fila RF-\allowbreak{}12.17} \\
\textbf{RF-\allowbreak{}12.18}\newline Caso, cap. 5 y cap. 15 (RT-16.30); SD2, F-12. & Portales N-01 a N-03 (Angular 22; S3 privado, CloudFront, AWS WAF y API REST pública); Keycloak 26; el portal de clientes consulta saldo y estado de cuenta desde datos del sistema de gestión\newline \textbf{EDT:} 3.5.2 & /P01 a /P02\newline Cumple; Media; E2; H10 & El cliente consulta su saldo según la información del sistema de gestión.\newline \emph{SD3, Tabla 3.A.1, fila RF-\allowbreak{}12.18} \\
\textbf{RF-\allowbreak{}12.19}\newline Caso, §§ 4.4 y 9.3 y cap. 15 (RT-16.30); SD2, F-12. & Portales N-01 a N-03 (Angular 22; S3 privado, CloudFront, AWS WAF y API REST pública); Keycloak 26; el portal de transportistas consulta las rutas asignadas para el día siguiente\newline \textbf{EDT:} 3.5.3 & /P01 a /P02\newline Cumple; Media; E2; H10 & La empresa visualiza el día anterior las rutas asignadas.\newline \emph{SD3, Tabla 3.A.1, fila RF-\allowbreak{}12.19} \\
\textbf{RF-\allowbreak{}12.20}\newline Caso, §§ 4.5 y 9.4 y cap. 15 (RT-16.30); SD2, F-12. & Portal de transportistas N-01 a N-03 (Angular 22; S3 privado, CloudFront, AWS WAF y API REST); M6 Reparto entrega la ruta, la guía y los documentos de despacho vinculados\newline \textbf{EDT:} 3.5.3 & /P01 a /P02\newline Cumple; Media; E2; H10 & La empresa descarga la documentación correspondiente.\newline \emph{SD3, Tabla 3.A.1, fila RF-\allowbreak{}12.20} \\
\textbf{RF-\allowbreak{}12.21}\newline Caso, cap. 10 (restr. 1 y 6), cap. 15 (RT-16.30) y § 16.1 (dec. 5); SD2, S-05. & Portal de transportistas N-01 a N-03 (Angular 22; S3 privado, CloudFront, AWS WAF y API REST; Keycloak 26); el representante confirma conductor y vehículo antes de entregar la ruta y despacho aprueba la asignación conductor-vehículo-turno.\newline \textbf{EDT:} 3.5.3 & /P01 a /P02\newline Cumple; Media; E2; H10 & Queda registrado el conductor/vehículo real; si no se confirma, el sistema mantiene la asignación planificada como supuesto no verificado.\newline \emph{SD3, Tabla 3.A.1, fila RF-\allowbreak{}12.21} \\
\textbf{RF-\allowbreak{}12.22}\newline Caso, § 4.1 y cap. 15 (RT-16.30); SD2, F-12. & Portales N-01 a N-03 (Angular 22; S3 privado, CloudFront, AWS WAF y API REST pública); Keycloak 26; el portal de proveedores permite consultar en solo lectura el estado de sus órdenes de compra\newline \textbf{EDT:} 3.5.1 & /P01 a /P02\newline Cumple; Media; E2; H10 & El proveedor consulta el estado actualizado de sus órdenes.\newline \emph{SD3, Tabla 3.A.1, fila RF-\allowbreak{}12.22} \\
\textbf{RF-\allowbreak{}12.23}\newline Caso, § 4.1 y cap. 15 (RT-16.30); SD2, F-12. & Portales N-01 a N-03 (Angular 22; S3 privado, CloudFront, AWS WAF y API REST pública); Keycloak 26; el portal de proveedores permite consultar en solo lectura las recepciones asociadas a sus órdenes\newline \textbf{EDT:} 3.5.1 & /P01 a /P02\newline Cumple; Media; E2; H10 & El proveedor consulta las recepciones actualizadas.\newline \emph{SD3, Tabla 3.A.1, fila RF-\allowbreak{}12.23} \\
\textbf{RF-\allowbreak{}12.24}\newline Caso, § 4.8, cap. 15 (RT-16.30) y § 16.1 (dec. 12); SD2, F-12. & Portal de proveedores N-01 a N-03 (Angular 22; S3 privado, CloudFront, AWS WAF y API REST); M8 Devoluciones y envases registra devoluciones y enlaza notas de crédito/documentos del ERP por INT-06/07\newline \textbf{EDT:} 3.5.1 & /P01 a /P02\newline Cumple; Media; E2; H10 & El proveedor consulta las devoluciones/notas de crédito actualizadas.\newline \emph{SD3, Tabla 3.A.1, fila RF-\allowbreak{}12.24} \\
\textbf{RF-\allowbreak{}12.25}\newline Caso, cap. 15 (RT-12.12 y RT-16.30); BTT, RT-12.12. & Autorregistro en el Portal de Clientes con RUT, correo corporativo y contraseña mediante Keycloak 26, validación de cliente activo contra el ERP y correo de confirmación\newline \textbf{EDT:} 3.5.2, 3.3.1 & /P01 a /P02\newline Cumple; Media; E2; H10 & El cliente queda registrado y recibe un correo de confirmación de cuenta. Si el RUT no está habilitado, el sistema informa que debe contactar a su ejecutivo comercial.\newline \emph{SD3, Tabla 3.A.1, fila RF-\allowbreak{}12.25} \\
\textbf{RF-\allowbreak{}12.26}\newline Caso, cap. 15 (RT-12.12 y RT-16.30); BTT, RT-12.03 y RT-12.12. & Portales N-01 a N-03 (Angular 22; S3 privado, CloudFront, AWS WAF y API REST pública); Keycloak 26; Keycloak autentica usuarios del portal con credenciales y el diseño exige MFA a perfiles sensibles\newline \textbf{EDT:} 3.5.2, 3.3.1 & /P01 a /P02\newline Cumple; Media; E2; H10 & El cliente accede a su área privada. Tras tres intentos fallidos consecutivos, la cuenta queda bloqueada temporalmente.\newline \emph{SD3, Tabla 3.A.1, fila RF-\allowbreak{}12.26} \\
\textbf{RF-\allowbreak{}12.27}\newline Caso, cap. 15 (RT-12.12); BTT, RT-12.12. & Keycloak 26 entrega al correo registrado el enlace de recuperación solicitado por el cliente, de un solo uso y con vigencia máxima de 15 min. Al restablecer se invalidan las credenciales anteriores y se registra la recuperación, conforme a la política de identidad y al resultado de RF-12.27.\newline \textbf{EDT:} 3.5.2, 3.3.1 & /P01 a /P02\newline Cumple; Media; E2; H10 & El cliente recibe un enlace válido por tiempo limitado que le permite establecer una nueva contraseña. El enlace expira tras su uso o transcurridos 15 minutos.\newline \emph{SD3, Tabla 3.A.1, fila RF-\allowbreak{}12.27} \\
\textbf{RF-\allowbreak{}12.28}\newline Caso, cap. 15 (RT-16.30); BTT, RT-12.02 y RT-12.07. & Portales N-01 a N-03 (Angular 22; S3 privado, CloudFront, AWS WAF y API REST pública); Keycloak 26; el portal invalida el token al cerrar la sesión y aplica cierre por inactividad\newline \textbf{EDT:} 3.5.2, 3.3.1 & /P01 a /P02\newline Cumple; Media; E2; H10 & La sesión queda invalidada y el cliente es redirigido a la pantalla de inicio de sesión.\newline \emph{SD3, Tabla 3.A.1, fila RF-\allowbreak{}12.28} \\
\textbf{RF-\allowbreak{}12.29}\newline Caso, cap. 15 (RT-16.30); BTT, RT-12.07 (caducidad por inactividad). & Portales N-01 a N-03 (Angular 22; S3 privado, CloudFront, AWS WAF y API REST pública); Keycloak 26; el portal cierra automáticamente las sesiones por inactividad, con límites declarados en el diseño\newline \textbf{EDT:} 3.5.2, 3.3.1 & /P01 a /P02\newline Cumple; Media; E2; H10 & La sesión expira y el cliente debe iniciar sesión nuevamente para acceder a su área privada.\newline \emph{SD3, Tabla 3.A.1, fila RF-\allowbreak{}12.29} \\
\textbf{RF-\allowbreak{}12.30}\newline Caso, cap. 10 (restr. 5) y cap. 15 (RT-16.30); SD2, F-12; T-11, dependencias de salida. & El Portal de Clientes solicita por INT-09 Transbank el pago del pedido de autoservicio. M7 verifica el resultado con la pasarela, registra autorización y vínculo al pedido, deduplica respuestas y concilia; un resultado incierto se consulta antes de repetir el cobro.\newline \textbf{EDT:} 3.5.2, 3.6.2, 3.4.10 & /P01 a /P02\newline Cumple; Media; E2; H10 & El pago es procesado por Transbank y el portal registra la confirmación o rechazo de la transacción.\newline \emph{SD3, Tabla 3.A.1, fila RF-\allowbreak{}12.30} \\
\textbf{RF-\allowbreak{}13.01}\newline BTT, RT-02.01, RT-02.03 & Arquitectura lógica de ocho capas, módulos M1–M12 y servicios transversales; la arquitectura lógica organiza la solución en las ocho capas requeridas y asigna los componentes y módulos correspondientes\newline \textbf{EDT:} 2.1.1 & /P01 a /P03\newline Cumple; Alta; E1 y E2; H5 y H10 & El diagrama de arquitectura identifica cada capa, sus componentes y las interfaces entre ellas, siguiendo ISO/IEC/IEEE 42010.\newline \emph{SD3, Tabla 3.A.3, fila RF-\allowbreak{}13.01} \\
\textbf{RF-\allowbreak{}13.02}\newline BTT, Cap. 3.1, RT-03.01 & Tabla de emplazamiento nube/on-premise con criterios de latencia, criticidad, volumen, acoplamiento físico, elasticidad y regulación; la tabla de emplazamiento asigna los componentes a nube o a instalaciones y expone los criterios de latencia, criticidad, volumen, acoplamiento físico, elasticidad y regulación\newline \textbf{EDT:} 2.1.2 & /P01 a /P03\newline Cumple; Alta; E1 y E2; H5 y H10 & La tabla de emplazamiento cubre el 100\% de los componentes de la solución, con justificación documentada para cada uno.\newline \emph{SD3, Tabla 3.A.3, fila RF-\allowbreak{}13.02} \\
\textbf{RF-\allowbreak{}13.03}\newline BTT, RT-02.04 & Registro único ADR-01 a ADR-24, fechado, con alternativa escogida, descartadas, criterio, consecuencias y evidencia. El Comité de Arquitectura conserva su actualización durante toda la ejecución.\newline \textbf{EDT:} 2.1.1, 1.8.4 & /P01 a /P03\newline Cumple; Alta; E1 y E2; H5 y H10 & El registro ADR es entregable contractual y se actualiza con cada decisión arquitectónica relevante.\newline \emph{SD3, Tabla 3.A.3, fila RF-\allowbreak{}13.03} \\
\textbf{RF-\allowbreak{}13.04}\newline BTT, RT-02.13 & Modelo de dominio con entidades, relaciones y eventos canónicos\newline \textbf{EDT:} 2.1.3 & /P01 a /P03\newline Cumple; Alta; E1 y E2; H5 y H10 & El modelo de dominio es entregado como parte de la oferta técnica, con notación UML.\newline \emph{SD3, Tabla 3.A.3, fila RF-\allowbreak{}13.04} \\
\textbf{RF-\allowbreak{}14.01}\newline Caso, §§ 2.3 y 4.4 y caps. 5 y 15 (RT-17.06); SD3, S-19. & INT-15 consume de la plataforma existente posición GPS, velocidad y kilometraje de los 42 camiones propios; M12 los presenta con fecha y entrega kilometraje a M10\newline \textbf{EDT:} 3.3.4 & /P01 a /P03\newline Cumple; Alta; E1; H5 & El sistema muestra en el mapa la ubicación de los 42 camiones propios, igual que hoy, pero integrado en la nueva solución.\newline \emph{SD3, Tabla 3.A.1, fila RF-\allowbreak{}14.01} \\
\textbf{RF-\allowbreak{}14.02}\newline Caso, §§ 4.4 y 9.3, cap. 5 y § 16.1 (dec. 13); SD3, S-13 y S-19. & M12 compara plan y recorrido de los 42 propios desde INT-15 y de los 54 externos desde el GPS del terminal C-02 de reparto; cada registro conserva fuente, fecha y calidad, con envío diferido sin señal\newline \textbf{EDT:} 3.4.5, 3.6.3 & /P01 a /P03\newline Cumple; Alta; E1; H5 & El gerente de operaciones ve que el camión 5 se desvía 5 km en su ruta; investiga si es por tráfico, error del conductor o una parada no registrada.\newline \emph{SD3, Tabla 3.A.1, fila RF-\allowbreak{}14.02} \\
\textbf{RF-\allowbreak{}14.03}\newline Caso, §§ 4.4 y 9.3 y cap. 5; SD3, S-19; SD4, M12. & M12 evalúa desvíos mayores de 3 km para ambas flotas, con posiciones de INT-15 (42 propios) o del terminal C-02 (54 externos); sin conectividad el registro se conserva y la alerta central se entrega al sincronizar, mostrando su antigüedad\newline \textbf{EDT:} 3.4.5 & /P01 a /P02\newline Cumple; Media; E1; H5 & El centro de control recibe una alerta si un camión se desvía para investigar si hay un problema (accidente, robo, etc.).\newline \emph{SD3, Tabla 3.A.1, fila RF-\allowbreak{}14.03} \\
\textbf{RF-\allowbreak{}14.04}\newline Caso, § 9.8, cap. 5, § 16.1 (dec. 7) y cap. 18 (crit. 11); SD2, F-14. & M12 entrega posición, kilometraje y tiempos de los camiones propios desde INT-15; M4/M6 aportan ruta, viaje y entregas identificadas. M10 imputa kilómetros y tiempo de viaje a cada entrega y cliente desde E2, manteniendo la restricción de uso de geolocalización.; para los 54 externos, M4/M6 vinculan posiciones y tiempos del terminal C-02 al viaje y entrega, sin atribuirles acceso al proveedor de INT-15\newline \textbf{EDT:} 3.5.4, 3.4.5, 3.3.4 & /P01 a /P02\newline Cumple; Media; E2; H10 & El costo de servir por cliente incluye el costo real de los kilómetros recorridos para llegar a su local.\newline \emph{SD3, Tabla 3.A.1, fila RF-\allowbreak{}14.04} \\
\textbf{RF-\allowbreak{}14.05}\newline Caso, § 2.3, cap. 10 (restr. 6), cap. 15 (RT-12.11) y § 16.1 (dec. 5); SD2, S-05. & M12 Telemetría (adaptador de solo lectura INT-15; Amazon DynamoDB; fuente de posicionamiento existente); M12 recibe identificadores de conductor validados por el proceso OTP de despacho\newline \textbf{EDT:} 3.4.5, 3.3.1 & /P01 a /P03\newline Cumple; Alta; E1; H5 & Se resuelve la Decisión \#5; el sistema exige que el conductor se identifique (PIN o QR) antes de iniciar la ruta, incluso si es externo.\newline \emph{SD3, Tabla 3.A.1, fila RF-\allowbreak{}14.05} \\
\textbf{RF-\allowbreak{}14.06}\newline Caso, §§ 4.4, 4.6 y 9.1 y cap. 5 (GPS de flota). & M12 Telemetría (adaptador de solo lectura INT-15; Amazon DynamoDB; fuente de posicionamiento existente); el planificador administra geocercas de sitios y zonas de clientes, y M12 registra la entrada y la salida de cada camión para medir tiempos de llegada y salida\newline \textbf{EDT:} 3.4.5 & /P01 a /P02\newline Cumple; Media; E1; H5 & El sistema registra automáticamente la hora de llegada del camión al área del cliente, para validar el cumplimiento de la ventana horaria.\newline \emph{SD3, Tabla 3.A.1, fila RF-\allowbreak{}14.06} \\
\end{tablasimple}

\begin{tablasimple}{>{\hsize=0.82\hsize\linewidth=\hsize}Y >{\hsize=1.04\hsize\linewidth=\hsize}Y >{\hsize=0.65\hsize\linewidth=\hsize}Y >{\hsize=1.49\hsize\linewidth=\hsize}Y}
{\multicolumn{4}{l}{\textbf{Matriz individual: bloque 7 de 26.}} \\
\cab{Requisito y origen} & \cab{Componente y EDT} & \cab{Prueba prevista y estado} & \cab{Criterio y fuente}}
\textbf{RF-\allowbreak{}14.01-\allowbreak{}BTT}\newline BTT, RT-11.03 & Clasificación Restringido / Confidencial / Interno / Público con controles por nivel\newline \textbf{EDT:} 2.2.4 & /P01 a /P03\newline Cumple; Alta; E1 y E2; H5 y H10 & La matriz de clasificación es entregada y aplicada a todos los datos de la solución.\newline \emph{SD3, Tabla 3.A.3, fila RF-\allowbreak{}14.01-\allowbreak{}BTT} \\
\textbf{RF-\allowbreak{}14.02-\allowbreak{}BTT}\newline BTT, RT-11.05 & Matriz de controles ISO/IEC 27001/27002 con implementación y evidencia\newline \textbf{EDT:} 2.2.4 & /P01 a /P03\newline Cumple; Alta; E1 y E2; H5 y H10 & La matriz es entregada y actualizada durante toda la ejecución.\newline \emph{SD3, Tabla 3.A.3, fila RF-\allowbreak{}14.02-\allowbreak{}BTT} \\
\textbf{RF-\allowbreak{}14.03-\allowbreak{}BTT}\newline BTT, RT-11.13 & La sección 4.2.5.2 del Subdocumento 4 declara cada entrada externa con su nombre de dominio bajo el dominio del CLIENTE (\texttt{clientes.}, \texttt{transportistas.}, \texttt{proveedores.}, \texttt{api.}, \texttt{id.}, \texttt{consola.}, \texttt{admin-id.} y \texttt{as2.<dominio>}), su servicio, puerto y control; los túneles IPsec de los sitios usan direcciones IP públicas de AWS, sin nombre de dominio.\newline \textbf{EDT:} 2.1.2, 2.2.2 & /P01 a /P03\newline Cumple; Alta; E1 y E2; H5 y H10 & La declaración es entregada y mantenida actualizada.\newline \emph{SD3, Tabla 3.A.3, fila RF-\allowbreak{}14.03-\allowbreak{}BTT} \\
\textbf{RF-\allowbreak{}14.04-\allowbreak{}BTT}\newline BTT, RT-11.18 & Plan de respuesta con clasificación, escalamiento, plazos y responsables; aviso al CLIENTE dentro de 2 h para incidentes críticos\newline \textbf{EDT:} 2.2.3 & /P01 a /P03\newline Cumple; Alta; E1 y E2; H5 y H10 & El plan es entregado y aprobado por el CLIENTE.\newline \emph{SD3, Tabla 3.A.3, fila RF-\allowbreak{}14.04-\allowbreak{}BTT} \\
\textbf{RF-\allowbreak{}14.05-\allowbreak{}BTT}\newline BTT, RT-11.19 & SOC 24$\times$7 de 8.1.5: toda brecha de seguridad o de datos personales se notifica al CLIENTE con informe preliminar dentro de 24 horas de su detección y análisis de causa raíz dentro de los cinco días hábiles siguientes a la detección.\newline \textbf{EDT:} 8.1.5, 2.2.3 & /P01 a /P03\newline Cumple; Alta; E1 y E2; H5 y H10 & El CLIENTE recibe la notificación dentro de las 24 horas y el informe de causa raíz dentro de los 5 días hábiles.\newline \emph{SD3, Tabla 3.A.3, fila RF-\allowbreak{}14.05-\allowbreak{}BTT} \\
\textbf{RF-\allowbreak{}15.01}\newline BTT, RT-12.01 & Keycloak 26.x centraliza identidad con OIDC y OAuth 2.1, o SAML 2.0 cuando la integración lo requiera, e integración con el directorio del CLIENTE por LDAP o equivalente en nube\newline \textbf{EDT:} 3.3.1 & /P01 a /P03\newline Cumple; Alta; E1 y E2; H5 y H10 & La identidad de todas las personas usuarias internas es gestionada desde un único punto.\newline \emph{SD3, Tabla 3.A.3, fila RF-\allowbreak{}15.01} \\
\textbf{RF-\allowbreak{}15.02}\newline BTT, RT-12.02 & Keycloak 26 centraliza la autenticación OIDC y propaga el cierre de sesión a los servicios conectados.\newline \textbf{EDT:} 3.3.1 & /P01 a /P03\newline Cumple; Alta; E1 y E2; H5 y H10 & El usuario se autentica una vez y accede a todos los módulos sin reautenticación.\newline \emph{SD3, Tabla 3.A.3, fila RF-\allowbreak{}15.02} \\
\textbf{RF-\allowbreak{}15.03}\newline BTT, RT-12.03 & Keycloak 26 exige MFA/FIDO2 a administradores y Verified Access exige MFA y postura en consolas externas. Clientes, proveedores y transportistas usan en los portales código TOTP o clave de acceso de Keycloak como segundo factor; terreno usa el terminal enrolado por MDM más el PIN personal.\newline \textbf{EDT:} 3.3.1 & /P01 a /P03\newline Cumple; Alta; E1 y E2; H5 y H10 & El acceso a perfiles administradores o desde fuera de la red requiere segundo factor.\newline \emph{SD3, Tabla 3.A.3, fila RF-\allowbreak{}15.03} \\
\textbf{RF-\allowbreak{}15.04}\newline BTT, RT-12.05 & Control RBAC complementado con atributos de instalación, turno, ruta, dispositivo y empresa transportista, con matriz de segregación de funciones.\newline \textbf{EDT:} 3.3.1 & /P01 a /P03\newline Cumple; Alta; E1 y E2; H5 y H10 & La matriz de roles y permisos es entregada y aprobada por el CLIENTE.\newline \emph{SD3, Tabla 3.A.3, fila RF-\allowbreak{}15.04} \\
\textbf{RF-\allowbreak{}15.05}\newline BTT, RT-12.10 & Altas, cambios y bajas del ciclo laboral activan trabajos Laravel 13 que crean, actualizan o deshabilitan cuentas Keycloak 26 y revocan sesiones; baja efectiva en menos de 24 h\newline \textbf{EDT:} 3.3.1 & /P01 a /P03\newline Cumple; Alta; E1 y E2; H5 y H10 & Las cuentas son creadas/eliminadas automáticamente según el estado laboral del usuario.\newline \emph{SD3, Tabla 3.A.3, fila RF-\allowbreak{}15.05} \\
\textbf{RF-\allowbreak{}16.01}\newline BTT, RT-14.01 & OpenTelemetry (PHP 8.5/Laravel 13, Kotlin 2.4) y AWS Distro for OpenTelemetry (colector 0.50) hacia Amazon CloudWatch, con transaction\_id\newline \textbf{EDT:} 3.1.5 & /P01 a /P03\newline Cumple; Alta; E1 y E2; H5 y H10 & El equipo de operaciones del CLIENTE puede consultar métricas, logs y trazas desde una única interfaz.\newline \emph{SD3, Tabla 3.A.3, fila RF-\allowbreak{}16.01} \\
\textbf{RF-\allowbreak{}16.02}\newline BTT, RT-14.02 & Tablero del mandante de solo lectura con auditoría de consultas: estado del servicio con métricas de CloudWatch de resolución de un minuto e indicadores de negocio de QuickSight con las latencias de RT-05.29 (operación del día $\leq$ 5 minutos, cierre comercial $\leq$ 2 horas, gestión $\leq$ 4 horas); exportación en CSV y Excel.\newline \textbf{EDT:} 3.1.5 & /P01 a /P03\newline Cumple; Alta; E1 y E2; H5 y H10 & El CLIENTE visualiza los tableros en tiempo real y exporta datos.\newline \emph{SD3, Tabla 3.A.3, fila RF-\allowbreak{}16.02} \\
\textbf{RF-\allowbreak{}16.03}\newline BTT, RT-14.04 & Alertas por síntomas con agrupación, supresión de duplicados, escalamiento automático por criticidad y turnos NOC/SOC 24$\times$7\newline \textbf{EDT:} 3.1.5 & /P01 a /P03\newline Cumple; Alta; E1 y E2; H5 y H10 & Las alertas se generan por síntomas de negocio (ej. OTIF bajo, pedidos no preparados), no solo por CPU alta.\newline \emph{SD3, Tabla 3.A.3, fila RF-\allowbreak{}16.03} \\
\textbf{RF-\allowbreak{}16.04}\newline BTT, RT-14.06 & Gestión de incidentes y problemas de 8.1.2: análisis obligatorio de causa raíz de cada incidente crítico, entregado al CLIENTE dentro de cinco días hábiles desde el incidente, y seguimiento por SRE de las acciones correctivas con responsable y plazo hasta su cierre verificado.\newline \textbf{EDT:} 8.1.2 & /P01 a /P03\newline Cumple; Alta; E1 y E2; H5 y H10 & El informe de causa raíz es entregado al CLIENTE dentro de los 5 días hábiles.\newline \emph{SD3, Tabla 3.A.3, fila RF-\allowbreak{}16.04} \\
\textbf{RF-\allowbreak{}17.01}\newline BTT, RT-16.01 & Consola Angular 22: el CLIENTE gestiona personas, roles, permisos y unidades organizacionales con jerarquías mediante grupos y atributos de Keycloak 26, con auditoría\newline \textbf{EDT:} 3.3.1 & /P01 a /P03\newline Cumple; Alta; E1 y E2; H5 y H10 & El administrador del CLIENTE crea, modifica o elimina usuarios/roles desde la interfaz sin desarrollo.\newline \emph{SD3, Tabla 3.A.3, fila RF-\allowbreak{}17.01} \\
\textbf{RF-\allowbreak{}17.02}\newline BTT, RT-16.02 & Consola Angular 22/API Laravel 13 administra umbrales, plazos, montos, tolerancias, catálogos, listas y textos de reglas implementadas; mae\_parametro/mae\_parametro\_version y auditoría conservan versión, vigencia, autor y fecha.\newline \textbf{EDT:} 3.3.5 & /P01 a /P03\newline Cumple; Alta; E1 y E2; H5 y H10 & El administrador del CLIENTE modifica parámetros de negocio sin desarrollo, con auditoría de cambios.\newline \emph{SD3, Tabla 3.A.3, fila RF-\allowbreak{}17.02} \\
\textbf{RF-\allowbreak{}17.03}\newline BTT, RT-16.03 & Cambio de parámetro con impacto operacional sin aplicar hasta aprobación de un segundo perfil; versión y auditoría con proponente, aprobador, valores y justificación\newline \textbf{EDT:} 3.3.5 & /P01 a /P03\newline Cumple; Alta; E1 y E2; H5 y H10 & El cambio se aplica solo después de la aprobación de un segundo perfil autorizado.\newline \emph{SD3, Tabla 3.A.3, fila RF-\allowbreak{}17.03} \\
\textbf{RF-\allowbreak{}17.04}\newline BTT, RT-16.06, RT-16.07 & gob\_auditoria registra toda escritura con actor/sistema, fecha y hora con zona, origen, dispositivo y valores anteriores/posteriores. La evidencia inalterable se conserva con S3 Object Lock durante siete años y no puede alterarla ni eliminarla un perfil administrador.\newline \textbf{EDT:} 3.3.5, 3.2.4 & /P01 a /P03\newline Cumple; Alta; E1 y E2; H5 y H10 & La auditoría es consultable desde la interfaz y no puede ser alterada ni eliminada.\newline \emph{SD3, Tabla 3.A.3, fila RF-\allowbreak{}17.04} \\
\textbf{RF-\allowbreak{}17.05}\newline BTT, RT-16.08 & Consola Angular 22/API Laravel 13 consulta y exporta gob\_auditoria con filtros por persona, período, entidad y operación; resultado filtrado CSV UTF-8/JSON, con permisos y sin acceso a base de datos.\newline \textbf{EDT:} 3.3.5 & /P01 a /P03\newline Cumple; Alta; E1 y E2; H5 y H10 & El auditor consulta y exporta la auditoría con los filtros deseados.\newline \emph{SD3, Tabla 3.A.3, fila RF-\allowbreak{}17.05} \\
\textbf{RF-\allowbreak{}17.06}\newline BTT, RT-16.11, RT-16.12 & Flujos con estados y transiciones; el CLIENTE configura responsables, plazos y niveles; escalamiento automático por vencimiento y delegación temporal a suplente en gob\_asignacion\newline \textbf{EDT:} 3.3.1 & /P01 a /P03\newline Cumple; Alta; E1 y E2; H5 y H10 & El administrador del CLIENTE configura flujos de trabajo (ej. aprobación de pedidos, gestión de devoluciones) sin desarrollo.\newline \emph{SD3, Tabla 3.A.3, fila RF-\allowbreak{}17.06} \\
\textbf{RF-\allowbreak{}17.07}\newline BTT, RT-16.13 & Bandeja unificada por persona en la pantalla de inicio con aprobaciones, revisiones y excepciones, ordenadas por criticidad y plazo y con alertas de vencimiento\newline \textbf{EDT:} 3.3.1 & /P01 a /P03\newline Cumple; Alta; E1 y E2; H5 y H10 & El usuario visualiza y gestiona todas sus tareas desde una única bandeja.\newline \emph{SD3, Tabla 3.A.3, fila RF-\allowbreak{}17.07} \\
\textbf{RF-\allowbreak{}17.08}\newline BTT, RT-16.15, RT-16.16 & Repositorio S3 versionado, cifrado y con hash; consulta documental por contenido y metadatos y previsualización en Angular 22 sin descarga\newline \textbf{EDT:} 3.3.5 & /P01 a /P03\newline Cumple; Alta; E1 y E2; H5 y H10 & Los documentos son almacenados, versionados y accesibles según permisos del usuario.\newline \emph{SD3, Tabla 3.A.3, fila RF-\allowbreak{}17.08} \\
\textbf{RF-\allowbreak{}17.09}\newline BTT, RT-16.19 & El CLIENTE administra plantillas y Laravel 13 combina datos de la transacción para producir HTML; la vista se guarda en PDF desde el navegador.\newline \textbf{EDT:} 3.6.4 & /P01 a /P02\newline Parcial; Media; E1 y E2; H5 y H10 & El usuario genera el documento con los datos de la transacción.\newline \emph{SD3, Tabla 3.A.3, fila RF-\allowbreak{}17.09} \\
\textbf{RF-\allowbreak{}17.10}\newline BTT, RT-16.20, RT-16.21 & INT-11 ofrece CORREO, SMS, WHATSAPP y AVISO\_EN\_PORTAL, con catálogo común en plantillas, preferencias e intentos; el CLIENTE administra las plantillas y variables versionadas en consola Angular 22.\newline \textbf{EDT:} 3.6.4 & /P01 a /P03\newline Cumple; Alta; E1 y E2; H5 y H10 & El usuario recibe la notificación por el canal configurado.\newline \emph{SD3, Tabla 3.A.3, fila RF-\allowbreak{}17.10} \\
\textbf{RF-\allowbreak{}17.11}\newline BTT, RT-16.27 & Búsqueda global con texto completo en PostgreSQL 16, tolerancia a errores, facetas y resultados limitados a recursos autorizados; respuesta de hasta 2 s\newline \textbf{EDT:} 3.3.1 & /P01 a /P03\newline Cumple; Alta; E1 y E2; H5 y H10 & El usuario encuentra información relevante mediante una búsqueda que no sobrepase más de 2 segundos en la respuesta y ademas de poder estar filtrada.\newline \emph{SD3, Tabla 3.A.3, fila RF-\allowbreak{}17.11} \\
\end{tablasimple}

\begin{tablasimple}{>{\hsize=0.82\hsize\linewidth=\hsize}Y >{\hsize=1.04\hsize\linewidth=\hsize}Y >{\hsize=0.65\hsize\linewidth=\hsize}Y >{\hsize=1.49\hsize\linewidth=\hsize}Y}
{\multicolumn{4}{l}{\textbf{Matriz individual: bloque 8 de 26.}} \\
\cab{Requisito y origen} & \cab{Componente y EDT} & \cab{Prueba prevista y estado} & \cab{Criterio y fuente}}
\textbf{RF-\allowbreak{}17.12}\newline BTT, RT-16.28 & Listados ordenables, filtrables y paginados; exportación CSV UTF-8 o JSON del conjunto filtrado completo, conservando filtros y orden\newline \textbf{EDT:} 3.3.5 & /P01 a /P03\newline Cumple; Alta; E1 y E2; H5 y H10 & El usuario ordena, filtra y exporta los datos del listado.\newline \emph{SD3, Tabla 3.A.3, fila RF-\allowbreak{}17.12} \\
\textbf{RF-\allowbreak{}17.13}\newline BTT, RT-16.31 & Portal público estático servido desde Amazon S3 privado y Amazon CloudFront, sin autenticación; T-14 fija WCAG 2.2 AA para la aceptación de accesibilidad de los portales.\newline \textbf{EDT:} 3.5.2, 3.8.2 & /P01 a /P02\newline Cumple; Media; E1 y E2; H5 y H10 & Cualquier visitante accede al portal público y visualiza la información. En E1 se publican la información institucional y el contacto; en E2 se agrega el catálogo de RF-12.14, que concreta este requisito para el caso.\newline \emph{SD3, Tabla 3.A.3, fila RF-\allowbreak{}17.13} \\
\textbf{RF-\allowbreak{}18.01}\newline BTT, RT-16.17, RT-16.18 & Firma por acto: guía de despacho electrónica con firma electrónica avanzada del CLIENTE aplicada por el ERP ante el SII, con verificación de la firma del XML y de la vigencia del certificado a la fecha de firma antes de vincularla; acuse de receptor electrónico habilitado con su firma verificada del mismo modo; firma electrónica simple del receptor persona natural (Ley N.º 19.799, art. 2, letra f) como evidencia vinculada, con el acuse según el mecanismo que acuerde el CLIENTE (V-18); el caso no exige firma avanzada para otros actos. XML, acuses y POD con hash, marca de tiempo y cadena de verificación en S3 Object Lock por seis años.\newline \textbf{EDT:} 3.3.2, 3.4.8 & /P01 a /P03\newline Cumple; Alta; E1 y E2; H5 y H10 & El documento queda firmado mediante la modalidad aplicable y se conserva su evidencia para la verificación posterior.\newline \emph{SD3, Tabla 3.A.3, fila RF-\allowbreak{}18.01} \\
\textbf{RF-\allowbreak{}18.02}\newline BTT, RT-16.22 & Preferencias personales de canal y frecuencia como atributos de identidad en Keycloak 26, aplicadas antes de cada envío; los avisos que el CLIENTE defina como obligatorios usan su canal establecido; not\_preferencia conserva la elección y las bajas por cliente\newline \textbf{EDT:} 3.6.4 & /P01 a /P02\newline Cumple; Media; E1 y E2; H5 y H10 & El usuario recibe notificaciones por el canal y frecuencia configurados.\newline \emph{SD3, Tabla 3.A.3, fila RF-\allowbreak{}18.02} \\
\textbf{RF-\allowbreak{}18.03}\newline BTT, RT-16.23 & INT-11 envía asíncronamente, reintenta y deduplica, y registra entrega y error. La apertura se registra cuando el proveedor o el portal la informa; en SMS, con la primera apertura del enlace corto al aviso en el portal, y sin ella queda NO\_SOPORTADA.\newline \textbf{EDT:} 3.6.4 & /P01 a /P03\newline Cumple; Alta; E1 y E2; H5 y H10 & La notificación es enviada y registrada, si falla, se reintenta según política configurada.\newline \emph{SD3, Tabla 3.A.3, fila RF-\allowbreak{}18.03} \\
\textbf{RNF-\allowbreak{}01.01}\newline Cap. 15, RT-17.01 (recepción móvil desconectada). & Perfil wms\_only local en el CD (PostgreSQL 16.x); registro de cada ítem en no más de 2 s desde la confirmación del escaneo\newline \textbf{EDT:} 3.4.1, 3.3.6 & /P01 a /P03\newline Cumple; Alta; E1; H5 & El tiempo de procesamiento para el registro de cada ítem en recepción, desde la confirmación del escaneo hasta la respuesta en pantalla, no debe superar los 2 segundos.\newline \emph{SD3, Tabla 3.A.4, fila RNF-\allowbreak{}01.01} \\
\textbf{RNF-\allowbreak{}01.02}\newline Caso, cap. 15, RT-05.10. & Retención por dominio con AWS Backup y S3 Object Lock\newline \textbf{EDT:} 3.2.4 & /P01 a /P05\newline Cumple; Crítica; E1; H5 & El sistema debe conservar los registros de recepción según el plazo por tipo de dato: trazabilidad sanitaria y no conformidades por mínimo 5 años; documentos de referencia tributaria por mínimo 6 años.\newline \emph{SD3, Tabla 3.A.4, fila RNF-\allowbreak{}01.02} \\
\textbf{RNF-\allowbreak{}01.03}\newline Cap. 15, RT-03.10 y RT-03.13 (24 horas del CD y sincronización en 2 horas). & Perfil wms\_only con PostgreSQL 16.x y RabbitMQ 4.x locales (24 h)\newline \textbf{EDT:} 3.3.6 & /P01 a /P05\newline Cumple; Crítica; E1; H5 & El sistema debe operar en modo desconectado durante mínimo 24 horas continuas sin pérdida de datos y sincronizar todas las transacciones pendientes con el ERP en $\le$ 2 horas tras restablecer la conectividad.\newline \emph{SD3, Tabla 3.A.4, fila RNF-\allowbreak{}01.03} \\
\textbf{RNF-\allowbreak{}02.01}\newline Caso, R2 (Cap. 10), RT-03.10, RT-03.13 & Perfil wms\_only con PostgreSQL 16.x y RabbitMQ 4.x locales (24 h)\newline \textbf{EDT:} 3.3.6 & /P01 a /P05\newline Cumple; Crítica; E1; H5 & El sistema de gestión de almacén (inventario, movimientos, conteos, consultas de stock) debe operar en modo desconectado durante un mínimo de 24 horas continuas sin pérdida de datos. Al restablecer la conectividad, las transacciones pendientes deben sincronizarse con el sistema central en $\le$ 2 horas, con resolución determinista de conflictos de stock (ej. si dos operarios movieron el mismo pallet).\newline \emph{SD3, Tabla 3.A.4, fila RNF-\allowbreak{}02.01} \\
\textbf{RNF-\allowbreak{}03.01}\newline Caso, cap. 15, RT-09.01; cap. 10, restricción 3. & Aplicación Kotlin 2.4/Room con registro local de la línea\newline \textbf{EDT:} 3.4.6 & /P01 a /P03\newline Cumple; Alta; E1; H5 & El tiempo para registrar una línea de pedido no debe superar 1,5 segundos, con o sin señal parcial.\newline \emph{SD3, Tabla 3.A.4, fila RNF-\allowbreak{}03.01} \\
\textbf{RNF-\allowbreak{}03.02}\newline Caso, cap. 15, RT-09.01. & Caché ElastiCache con Valkey 8.1 para stock y crédito (N-07)\newline \textbf{EDT:} 3.4.6, 3.2.2 & /P01 a /P03\newline Cumple; Alta; E1; H5 & El tiempo para consultar stock y crédito con conectividad no debe superar 2 segundos.\newline \emph{SD3, Tabla 3.A.4, fila RNF-\allowbreak{}03.02} \\
\textbf{RNF-\allowbreak{}03.03}\newline Cap. 9.2 (consulta de historial). & API Laravel 13.x con copia local de historial en Room; consulta en no más de 2 s con conectividad\newline \textbf{EDT:} 3.4.6 & /P01 a /P03\newline Cumple; Alta; E1; H5 & El tiempo para consultar el historial de compra con conectividad no debe superar 2 segundos.\newline \emph{SD3, Tabla 3.A.4, fila RNF-\allowbreak{}03.03} \\
\textbf{RNF-\allowbreak{}04.01}\newline Cap. 18, criterio 7 (menos de 20 minutos). & Motor de optimización en contenedor propio (ECS Fargate), < 20 min por corrida\newline \textbf{EDT:} 3.4.7 & /P01 a /P02\newline Cumple; Media; E1; H5 & El motor de ruteo debe procesar la asignación y secuencia completa de la matriz diaria de entregas en menos de 20 minutos por ejecución completa, medidos desde el inicio del cálculo hasta disponer de la asignación y secuencia de todas las entregas de la matriz diaria.\newline \emph{SD3, Tabla 3.A.4, fila RNF-\allowbreak{}04.01} \\
\textbf{RNF-\allowbreak{}05.01}\newline Caso, RT-09.01 & Kotlin 2.4/Room confirma escaneo, validación y almacenamiento durable local en no más de 1 s, online u offline; el servicio WMS recibe después la sincronización. En congelado se usa MC9400 Cold Storage Freezer a $-$22 °C y se mide con el perfil completo de aceptación.\newline \textbf{EDT:} 3.4.3 & /P01 a /P03\newline Cumple; Alta; E1; H5 & La confirmación de una línea de picking (escaneo + validación + almacenamiento local) no debe superar 1 segundo de tiempo de respuesta en el dispositivo del preparador, tanto en modo online como en modo offline, incluyendo la operación dentro de la cámara de congelado (-22 °C).\newline \emph{SD3, Tabla 3.A.4, fila RNF-\allowbreak{}05.01} \\
\textbf{RNF-\allowbreak{}05.02}\newline Caso, cap. 10, restricción 7; cap. 15, RT-13.08 y RT-12.11. & Interfaz de picking en MC9400 Cold Storage Freezer con controles de al menos 15 $\times$ 15 mm, sin gestos complejos; aceptación de 30 min continuos a $-$22 °C\newline \textbf{EDT:} 3.4.3, 3.8.3 & /P01 a /P05\newline Cumple; Crítica; E1; H5 & La interfaz de picking debe ser operable sin precisión de toque fino, con elementos interactivos de al menos 15$\times$15 mm y sin requerir gestos complejos (deslizar, pellizcar, doble toque). El dispositivo debe mantener funcionalidad operativa durante al menos 30 minutos continuos en cámara de congelado a $-$22\,°C.\newline \emph{SD3, Tabla 3.A.4, fila RNF-\allowbreak{}05.02} \\
\textbf{RNF-\allowbreak{}05.03}\newline Caso, cap. 8 (entrevista de Ximena Bravo); cap. 10, restricción 11; SD3, Tabla 3.A.4a. & Entrenamiento asistido en el puesto dentro de la bodega: preparador nuevo sin experiencia previa completa al menos 20 líneas con error máximo de 5 \% tras un máximo de dos horas; sin sesiones fuera del recinto ni materiales impresos de más de dos páginas.\newline \textbf{EDT:} 7.1.2 & /P01 a /P05\newline Cumple; Crítica; E1; H5 & Un nuevo operario de preparación, sin experiencia previa con el sistema, debe completar una misión de picking de al menos 20 líneas con una tasa de error $\le$ 5\% tras un máximo de 2 horas de entrenamiento asistido en el puesto de trabajo. La capacitación no debe requerir materiales impresos de más de 2 páginas ni sesiones fuera del recinto de bodega.\newline \emph{SD3, Tabla 3.A.4, fila RNF-\allowbreak{}05.03} \\
\textbf{RNF-\allowbreak{}06.01}\newline Cap. 10, restricción 3; Cap. 15, RT-03.10 (turno de 14 horas sin cobertura). & Aplicación Kotlin 2.4/Room con cola cifrada para 14 h sin señal\newline \textbf{EDT:} 3.3.6 & /P01 a /P02\newline Cumple; Media; E1; H5 & El sistema debe persistir el 100\% de los datos necesarios en el almacenamiento cifrado local para ejecutar la ruta sin cobertura.\newline \emph{SD3, Tabla 3.A.4, fila RNF-\allowbreak{}06.01} \\
\textbf{RNF-\allowbreak{}07.01}\newline Cap. 15, RT-03.13; Cap. 18, criterio 10. & Sincronización de cobros por API Laravel 13 y SQS FIFO\newline \textbf{EDT:} 3.4.10 & /P01 a /P05\newline Cumple; Crítica; E1; H5 & El sistema debe sincronizar la información de preconciliación de cobranzas capturada en la aplicación móvil en un tiempo no superior a 10 minutos tras conectar a red, permitiendo la liquidación de la flota el mismo día.\newline \emph{SD3, Tabla 3.A.4, fila RNF-\allowbreak{}07.01} \\
\textbf{RNF-\allowbreak{}07.02}\newline Cap. 15, RT-11.10 (cifrado a nivel de campo). & Cifrado por campo con AES-256 y claves de AWS KMS; TLS 1.3 en tránsito; PAN tokenizado por la pasarela\newline \textbf{EDT:} 2.2.4, 3.6.2 & /P01 a /P05\newline Cumple; Crítica; E1; H5 & El sistema debe cifrar a nivel de campo (AES-256 en reposo, TLS 1.3 en tránsito) los datos de pagos en efectivo, transacciones con tarjeta POS, comportamiento comercial de clientes y geolocalización de personas (conductores y preventistas) (Ley N° 21.719).\newline \emph{SD3, Tabla 3.A.4, fila RNF-\allowbreak{}07.02} \\
\textbf{RNF-\allowbreak{}08.01}\newline Caso, § 4.8; SD3, Tabla 3.A.4 (umbral ofertado). & Registro local en andén (perfil wms\_only), en no más de 30 segundos por ítem con escaneo y destino\newline \textbf{EDT:} 3.4.9 & /P01 a /P03\newline Cumple; Alta; E1; H5 & El registro de una devolución en el andén del CD (escaneo + selección de destino) no debe superar 30 segundos por ítem.\newline \emph{SD3, Tabla 3.A.4, fila RNF-\allowbreak{}08.01} \\
\textbf{RNF-\allowbreak{}08.02}\newline Cap. 15, RT-05.10 (retención). & Retención por dominio\newline \textbf{EDT:} 3.2.4 & /P01 a /P03\newline Cumple; Alta; E1; H5 & Los registros de devoluciones y mermas deben conservarse según el plazo por tipo de dato: trazabilidad sanitaria por mínimo 5 años; documentos con efectos tributarios (notas de crédito asociadas) por mínimo 6 años.\newline \emph{SD3, Tabla 3.A.4, fila RNF-\allowbreak{}08.02} \\
\textbf{RNF-\allowbreak{}09.01}\newline Cap. 18, criterio 1 (menos de dos horas con evidencia). & M9 con trazabilidad EPCIS evento a evento\newline \textbf{EDT:} 3.4.4 & /P01 a /P05\newline Cumple; Crítica; E1; H5 & El sistema tiene un tiempo de respuesta para generar el listado de clientes afectados por un lote (RF-09.01) inferior a 2 horas (incluyendo la generación del archivo exportable).\newline \emph{SD3, Tabla 3.A.4, fila RNF-\allowbreak{}09.01} \\
\textbf{RNF-\allowbreak{}09.02}\newline Caso, cap. 15, RT-05.10. & Retención sanitaria: vida útil + 6 meses, mínimo 5 años\newline \textbf{EDT:} 3.2.4 & /P01 a /P05\newline Cumple; Crítica; E1; H5 & El sistema debe conservar los registros sanitarios y de trazabilidad de lote durante el período mínimo legal: 'vida útil del producto + 6 meses', con un piso mínimo de 5 años.\newline \emph{SD3, Tabla 3.A.4, fila RNF-\allowbreak{}09.02} \\
\textbf{RNF-\allowbreak{}09.03}\newline Cap. 15, RT-16.09 (consultas sensibles y modificaciones). & Registro de consultas y modificaciones en gob\_auditoria\newline \textbf{EDT:} 3.3.5 & /P01 a /P03\newline Cumple; Alta; E1; H5 & El sistema debe registrar todas las consultas y modificaciones de información sensible —trazabilidad (lotes, clientes afectados), antecedentes comerciales, comportamiento de pago y geolocalización de personas— en un log de auditoría, con usuario, fecha/hora, objeto consultado y finalidad cuando aplique. Toda exportación de esa información registra además quién exportó qué y cuándo (RT-16.30).\newline \emph{SD3, Tabla 3.A.4, fila RNF-\allowbreak{}09.03} \\
\textbf{RNF-\allowbreak{}09.04}\newline Cap. 10, restricción 7; Cap. 15, RT-13.08 (operación sin señal en congelado). & Registro local en terminal Zebra MC9400 y sincronización al salir de la cámara\newline \textbf{EDT:} 3.3.6, 3.4.3 & /P01 a /P05\newline Cumple; Crítica; E1; H5 & El sistema debe funcionar Sin cobertura de señal y sincronizando al recuperar conectividad fuera de la cámara de frío\newline \emph{SD3, Tabla 3.A.4, fila RNF-\allowbreak{}09.04} \\
\end{tablasimple}

\begin{tablasimple}{>{\hsize=0.82\hsize\linewidth=\hsize}Y >{\hsize=1.04\hsize\linewidth=\hsize}Y >{\hsize=0.65\hsize\linewidth=\hsize}Y >{\hsize=1.49\hsize\linewidth=\hsize}Y}
{\multicolumn{4}{l}{\textbf{Matriz individual: bloque 9 de 26.}} \\
\cab{Requisito y origen} & \cab{Componente y EDT} & \cab{Prueba prevista y estado} & \cab{Criterio y fuente}}
\textbf{RNF-\allowbreak{}09.05}\newline Propio LafroX. Fuente: RF-09.05. & M9 configura frecuencias distintas por origen: cámaras cada 5 min; sondas de camión cada 5 s con serie por minuto; termógrafo certificado de aire cada 1 min y sensor independiente en la segunda cadena. Gateway y TC58e detectan la regla aplicable y registran intervalos sin lectura; se concilian las tres series al regreso.\newline \textbf{EDT:} 3.4.4, 3.4.5, 6.5.3, 6.5.4 & /P01 a /P05\newline Cumple; Crítica; E1; H5 & Cada cámara y cada vehículo con frío registra temperatura con una frecuencia, parametrizable por tipo de producto, suficiente para detectar una desviación de 2\,°C sostenida durante 15 minutos (RF-09.05). La serie no admite brechas: un intervalo sin lectura se registra como evento.\newline \emph{SD3, Tabla 3.A.5b, fila RNF-\allowbreak{}09.05} \\
\textbf{RNF-\allowbreak{}09.06}\newline Propio LafroX. & Sondas cableadas al ioLogik/gateway, luz y zumbador B-12 por relé, Wi-Fi local gateway–TC58e y Bluetooth del sensor independiente sostienen la protección sin ninguna red externa. TC58e retiene ante latido inválido en no más de 30 s. Verificación V6–V9.\newline \textbf{EDT:} 3.4.4, 3.4.5, 2.4.3, 4.1.4 & /P01 a /P05\newline Cumple; Crítica; E1; H5 & La detección, alarma al conductor y bloqueo de la entrega operan sin red móvil ni satélite; ante pérdida de latido, el TC58e retiene en no más de 30 s.\newline \emph{SD3, Tabla 3.A.5b, fila RNF-\allowbreak{}09.06} \\
\textbf{RNF-\allowbreak{}09.07}\newline Propio LafroX. & Gateway B-04 con dos operadores móviles e Iridium Edge B-09; IoT Core y CloudConnect/SQS o receptor SBD alternativo. Ambos caminos bidireccionales permanecen activos: aviso a Calidad $\leq$ 2 min móvil y $\leq$ 5 min satélite en al menos 95 \% de ensayos por camino. V2/V18.\newline \textbf{EDT:} 3.4.5, 2.4.3, 4.1.4, 5.3.4 & /P01 a /P05\newline Cumple; Crítica; E1; H5 & Red móvil de dos operadores e Iridium SBD transmiten en paralelo. Aviso a Calidad $\leq$ 2 min por red móvil y $\leq$ 5 min por satélite en al menos 95 \% de los ensayos por camino.\newline \emph{SD3, Tabla 3.A.5b, fila RNF-\allowbreak{}09.07} \\
\textbf{RNF-\allowbreak{}09.08}\newline Propio LafroX. & Simuladores B-06/B-07 en el punto crítico de cada tipo de carrocería fijado por V4; sensor independiente B-08 vigila la segunda evaluación. El termógrafo B-03 EYE SENSOR EN12830 registra el aire cada minuto para evidencia certificada (V3). Los rangos y tolerancias se acreditan antes de seleccionar y aceptar los equipos.\newline \textbf{EDT:} 3.4.4, 2.4.3, 2.4.4, 4.1.4 & /P01 a /P05\newline Cumple; Crítica; E1; H5 & La decisión usa simuladores de producto en el punto más caliente fijado por mapeo térmico; el termógrafo certificado EN 12830 registra el aire cada 1 min.\newline \emph{SD3, Tabla 3.A.5b, fila RNF-\allowbreak{}09.08} \\
\textbf{RNF-\allowbreak{}09.09}\newline Propio LafroX. & Cola persistente Greengrass V2 hasta A2 más archivo de 72 h; M9 procesa una vez productor/arranque/secuencia, detecta huecos y solicita faltantes. La serie fina reemplaza el resumen satelital. Una crítica tardía retiene o abre evaluación de retiro; una recuperación tardía no libera. V11.\newline \textbf{EDT:} 3.4.4, 3.4.5, 2.4.3, 4.1.4 & /P01 a /P05\newline Cumple; Crítica; E1; H5 & Sin pérdidas ni efectos duplicados; identidad por productor, arranque y secuencia, huecos detectados y crítica tardía siempre procesada. Una recuperación tardía no libera retenciones.\newline \emph{SD3, Tabla 3.A.5b, fila RNF-\allowbreak{}09.09} \\
\textbf{RNF-\allowbreak{}11.01}\newline Caso, cap. 15, RT-05.29. & Latencia $\leq$ 5 min de los indicadores del día\newline \textbf{EDT:} 3.4.11, 3.2.3 & /P01 a /P03\newline Cumple; Alta; E1; H5 & El sistema debe desplegar los indicadores operacionales del día con una latencia no superior a 5 minutos desde su recepción en el backend central para terminales con cobertura móvil activa.\newline \emph{SD3, Tabla 3.A.4, fila RNF-\allowbreak{}11.01} \\
\textbf{RNF-\allowbreak{}11.02}\newline Caso, § 17.4, decisión 11; cap. 15, RT-09.01. & Separación OLTP / analítica (Redshift Serverless)\newline \textbf{EDT:} 3.4.11, 3.2.3 & /P01 a /P03\newline Cumple; Alta; E1; H5 & El sistema debe desacoplar el procesamiento analítico y de reportería (OLAP/BI) del motor transaccional operativo (OLTP), garantizando que las consultas analíticas pesadas no degraden los tiempos de respuesta de transacciones críticas (picking $\leq$ 1s, preventa $\leq$ 1,5s).\newline \emph{SD3, Tabla 3.A.4, fila RNF-\allowbreak{}11.02} \\
\textbf{RNF-\allowbreak{}12.01}\newline Cap. 9.9; Cap. 15, RT-16.30. & M11 Canal moderno y Portal de Clientes en producción en el mes 21, antes de enero de 2029\newline \textbf{EDT:} 3.5.1, 3.5.2, 4.3.3 & /P01 a /P05\newline Cumple; Crítica; E2; H10 & El portal de autoatención de clientes (RF-12.15 a RF-12.18) y la recepción EDI (RF-12.01 a RF-12.09) deben estar en producción antes de enero de 2029.\newline \emph{SD3, Tabla 3.A.4, fila RNF-\allowbreak{}12.01} \\
\textbf{RNF-\allowbreak{}12.02}\newline Cap. 15, RT-05.23 (formato exigido por cada cadena). & Perfiles de cadena OpenAS2 4.9 y tabla de equivalencias EDI configurables (edi\_perfil\_cadena)\newline \textbf{EDT:} 3.6.6 & /P01 a /P03\newline Cumple; Alta; E2; H10 & El mecanismo de recepción EDI debe permitir configurar el mapeo de mensajería de una nueva cadena sin desarrollo de software a medida.\newline \emph{SD3, Tabla 3.A.4, fila RNF-\allowbreak{}12.02} \\
\textbf{RNF-\allowbreak{}12.03}\newline Propio LafroX. Fuente: Caso, sección 13.2; RF-12.13. & M11 mediante INT-08/OpenAS2 4.9 envía el acuse comercial digital dentro de 30 min desde la descarga, referido a la misma guía y a su ASN enviado al cargar el camión; distinto del MDN de transporte y del acuse tributario.\newline \textbf{EDT:} 3.5.1, 3.6.5 & /P01 a /P03\newline Cumple; Alta; E2; H10 & El acuse de recibo digital de cada entrega a una cadena se envía dentro de los 30 minutos posteriores a la descarga (RF-12.13), referido a la misma guía y al aviso de despacho enviado al cargar el camión (RF-12.09).\newline \emph{SD3, Tabla 3.A.5b, fila RNF-\allowbreak{}12.03} \\
\textbf{RNF-\allowbreak{}13.01}\newline BTT, RT-03.10, RT-03.11, RT-03.12, RT-03.13 & Perfil wms\_only con PostgreSQL 16.x y RabbitMQ 4.x locales por sitio; 24 h de autonomía\newline \textbf{EDT:} 3.3.6, 6.6.3 & /P01 a /P05\newline Cumple; Crítica; E1 y E2; H5 y H10 & El componente on-premise debe operar de forma autónoma y degradada ante la pérdida total del enlace con la nube, durante un período mínimo de 24 horas continuas, continuando registrando las transacciones operacionales críticas de forma local, con integridad garantizada y sin pérdida de datos.\newline \emph{SD3, Tabla 3.A.5, fila RNF-\allowbreak{}13.01} \\
\textbf{RNF-\allowbreak{}13.02}\newline BTT, RT-03.13 & Declaración de funciones no disponibles sin conexión y procedimiento manual que las suple\newline \textbf{EDT:} 3.3.6, 6.6.3 & /P01 a /P05\newline Cumple; Crítica; E1 y E2; H5 y H10 & El PROPONENTE debe declarar qué funciones NO estarán disponibles en modo desconectado y qué procedimiento manual las suple.\newline \emph{SD3, Tabla 3.A.5, fila RNF-\allowbreak{}13.02} \\
\textbf{RNF-\allowbreak{}13.03}\newline BTT, RT-03.14 & Clúster Talca de tres servidores HPE ProLiant DL345 Gen11 con Proxmox VE 9.2 y Ceph; dos HPE ProLiant DL20 Gen11 activo/espera en Concepción; pares Advantech ARK-1251 activo/espera en cada cross-docking\newline \textbf{EDT:} 6.3.1, 6.3.3, 6.3.4 & /P01 a /P03\newline Cumple; Alta; E1 y E2; H5 y H10 & Los equipos on-premise críticos deben ser redundantes.\newline \emph{SD3, Tabla 3.A.5, fila RNF-\allowbreak{}13.03} \\
\textbf{RNF-\allowbreak{}13.04}\newline BTT, RT-03.14 & Ceph con tres réplicas en Talca; RAID 10 local en Concepción, RAID 1 en cross-docking y RAID 6 en NAS\newline \textbf{EDT:} 6.3.1, 6.3.3, 6.3.4 & /P01 a /P03\newline Cumple; Alta; E1 y E2; H5 y H10 & El almacenamiento local debe tolerar la falla de al menos un disco\newline \emph{SD3, Tabla 3.A.5, fila RNF-\allowbreak{}13.04} \\
\textbf{RNF-\allowbreak{}13.05}\newline BTT, RT-03.14 & Niveles declarados y justificados: Ceph con tres réplicas sin RAID sobre NVMe en Talca, RAID 10 en Concepción, RAID 1 en cross-docking y RAID 6 en NAS\newline \textbf{EDT:} 6.3.1, 6.3.3, 6.3.4 & /P01 a /P03\newline Cumple; Alta; E1 y E2; H5 y H10 & el PROPONENTE debe declarar el nivel RAID escogido y justificarlo frente a las alternativas.\newline \emph{SD3, Tabla 3.A.5, fila RNF-\allowbreak{}13.05} \\
\textbf{RNF-\allowbreak{}13.06}\newline BTT, RT-03.15 & Ventana de parches y líneas base CIS acordada con la Contraparte Técnica antes de cada intervención, con Ansible (ansible-core 2.21), respaldo, verificación y reversión\newline \textbf{EDT:} 6.6.1, 8.2.5 & /P01 a /P03\newline Cumple; Alta; E1 y E2; H5 y H10 & Los sistemas on-premise deben endurecerse conforme a los CIS Benchmarks aplicables, con gestión centralizada de parches y ventana de aplicación acordada con el CLIENTE.\newline \emph{SD3, Tabla 3.A.5, fila RNF-\allowbreak{}13.06} \\
\textbf{RNF-\allowbreak{}13.07}\newline BTT, RT-03.17 & Talca y Concepción: fibra, LTE de otro proveedor y Starlink en espera caliente. Cada cross-docking: Starlink principal y dos operadores LTE. Tres caminos por sitio, con conmutación automática en menos de 30 s y mínimos de subida medidos antes de aceptar.\newline \textbf{EDT:} 6.6.2, 5.3.1, 5.3.2, 5.3.3 & /P01 a /P03\newline Cumple; Alta; E1 y E2; H5 y H10 & El enlace entre el sitio on-premise y la nube debe ser redundante, con caminos físicos y proveedores distintos, y conmutación automática con tiempo de conmutación declarado ($\leq$ 5 min).\newline \emph{SD3, Tabla 3.A.5, fila RNF-\allowbreak{}13.07} \\
\textbf{RNF-\allowbreak{}13.08}\newline BTT, RT-03.20 & Cálculo del ancho de banda por sitio en régimen y peak\newline \textbf{EDT:} 2.1.2 & /P01 a /P03\newline Cumple; Alta; E1 y E2; H5 y H10 & El PROPONENTE debe dimensionar el ancho de banda requerido por sitio, en régimen normal y en peak, y justificarlo con el cálculo de volumen de transacciones y de datos.\newline \emph{SD3, Tabla 3.A.5, fila RNF-\allowbreak{}13.08} \\
\textbf{RNF-\allowbreak{}13.09}\newline BTT, RT-03.23 & Wi-Fi 6E industrial Cisco IW9167I segmentado por VLAN con autenticación EAP-TLS; estudio de cobertura presentado: diseño predictivo por zona de 65 + 7 puntos de acceso, con al menos dos por cámara, verificación en sitio en el mes 2 y medición de aceptación antes de cada marcha blanca\newline \textbf{EDT:} 2.1.2, 6.2.3 & /P01 a /P03\newline Cumple; Alta; E1 y E2; H5 y H10 & El PROPONENTE debe realizar un estudio de cobertura de red inalámbrica en los centros de distribución, que incluya expresamente el interior de las cámaras de refrigerado y de congelado, hoy sin señal, y proponer la solución técnica para resolverlo.\newline \emph{SD3, Tabla 3.A.5, fila RNF-\allowbreak{}13.09} \\
\textbf{RNF-\allowbreak{}14.01}\newline BTT, RT-11.08 & TLS 1.3 en entradas HTTPS y componentes de la solución; TLS 1.0/1.1 prohibidos, HSTS con precarga, certificados con renovación y alerta. El Anexo 4-R admite TLS 1.2 para ERP o terceros incompatibles; esa excepción impide acreditar TLS 1.3 en la totalidad de los enlaces exigidos.\newline \textbf{EDT:} 3.2.1, 3.2.5 & /P01 a /P05\newline Parcial; Crítica; E1 y E2; H5 y H10 & Todo el tráfico en tránsito debe emplear TLS 1.3, con prohibición expresa de TLS 1.0 y 1.1, conjuntos de cifrado modernos, HSTS con precarga y gestión automatizada de certificados con rotación y alerta anticipada de vencimiento.\newline \emph{SD3, Tabla 3.A.5, fila RNF-\allowbreak{}14.01} \\
\textbf{RNF-\allowbreak{}14.02}\newline BTT, RT-11.09 & Cifrado en reposo con AWS KMS administrado por el CLIENTE, rotación y separación de custodia\newline \textbf{EDT:} 3.2.5 & /P01 a /P05\newline Cumple; Crítica; E1 y E2; H5 y H10 & La totalidad de los datos en reposo debe estar cifrada, con claves gestionadas en un servicio de gestión de claves o en un módulo de seguridad de hardware, política de rotación declarada y separación de funciones en la custodia de claves.\newline \emph{SD3, Tabla 3.A.5, fila RNF-\allowbreak{}14.02} \\
\textbf{RNF-\allowbreak{}14.03}\newline BTT, RT-03.04 & VPC de Producción con zona privada para aplicación y datos; subred pública solo con NAT; el canal AS2 llega por CloudFront a un ALB interno; portales en S3 privado tras CloudFront\newline \textbf{EDT:} 3.1.2 & /P01 a /P03\newline Cumple; Alta; E1 y E2; H5 y H10 & En la nube, la red debe segmentarse por capas, con subredes privadas para aplicación y datos, y exposición pública restringida a la capa de borde. Ningún componente de datos debe ser alcanzable desde Internet.\newline \emph{SD3, Tabla 3.A.5, fila RNF-\allowbreak{}14.03} \\
\textbf{RNF-\allowbreak{}14.04}\newline BTT, RT-11.14, RT-11.15 & Security Lake y Security Hub centralizan hallazgos para el SOC 24$\times$7 y correlacionan credencial revocada, consultas comerciales masivas, alteración de temperatura y EDI anómalo. Los eventos son inalterables: 12 meses en línea y 24 adicionales recuperables; auditoría y seguridad se conservan siete años bajo Object Lock.\newline \textbf{EDT:} 8.1.5, 3.2.5 & /P01 a /P03\newline Cumple; Alta; E1 y E2; H5 y H10 & Los eventos de seguridad deben registrarse de forma centralizada e inalterable, y correlacionarse en una plataforma SIEM con casos de uso de detección definidos específicamente para el proceso de negocio del caso (ej. intentos de acceso a trazabilidad de lote, modificaciones de inventario no autorizadas).\newline \emph{SD3, Tabla 3.A.5, fila RNF-\allowbreak{}14.04} \\
\textbf{RNF-\allowbreak{}14.05}\newline BTT, RT-11.16 & CrowdStrike Falcon en los nodos, servidores, máquinas virtuales y estaciones on-premise, y GuardDuty Runtime Monitoring en las tareas ECS Fargate de sa-east-1 y us-east-1, con consola central y respuesta del SOC 24$\times$7. El servidor del ERP es un activo del CLIENTE que no forma parte de la solución; el firewall de Talca solo le permite el tráfico de la ACL y la administración del CLIENTE.\newline \textbf{EDT:} 3.2.5, 5.2.2 & /P01 a /P03\newline Cumple; Alta; E1 y E2; H5 y H10 & Se debe implementar detección y respuesta en puntos finales (EDR) y en cargas de trabajo, tanto en nube como on-premise, con cobertura 24$\times$7.\newline \emph{SD3, Tabla 3.A.5, fila RNF-\allowbreak{}14.05} \\
\textbf{RNF-\allowbreak{}14.06}\newline BTT, RT-11.20 & Tercero independiente antes de cada producción y al menos anualmente en Operación, con informe íntegro y remediación.\newline \textbf{EDT:} 3.8.6, 3.9.5, 8.2.2 & /P01 a /P03\newline Parcial; Alta; E1 y E2; H5 y H10 & Se deben ejecutar pruebas de intrusión (pentesting) por un tercero independiente del ADJUDICATARIO, anualmente y antes de cada paso a producción, con entrega íntegra del informe al CLIENTE y plan de remediación con plazos. LafroX ejecuta estas pruebas de forma semestral, por sobre el mínimo anual (Subdocumento 1, sección «Modelo de Gobierno de Seguridad de la Información»).\newline \emph{SD3, Tabla 3.A.5, fila RNF-\allowbreak{}14.06} \\
\end{tablasimple}

\begin{tablasimple}{>{\hsize=0.82\hsize\linewidth=\hsize}Y >{\hsize=1.04\hsize\linewidth=\hsize}Y >{\hsize=0.65\hsize\linewidth=\hsize}Y >{\hsize=1.49\hsize\linewidth=\hsize}Y}
{\multicolumn{4}{l}{\textbf{Matriz individual: bloque 10 de 26.}} \\
\cab{Requisito y origen} & \cab{Componente y EDT} & \cab{Prueba prevista y estado} & \cab{Criterio y fuente}}
\textbf{RNF-\allowbreak{}14.07}\newline BTT, RT-11.22 & SAST (PHP 8.5Stan/Larastan), SCA (composer audit), DAST y escaneo de imágenes con bloqueo\newline \textbf{EDT:} 3.1.4, 1.5.2 & /P01 a /P03\newline Cumple; Alta; E1 y E2; H5 y H10 & El flujo de integración continua debe incorporar análisis estático de código (SAST), análisis de composición de software (SCA), análisis dinámico (DAST) y escaneo de imágenes de contenedor, con criterios de bloqueo automático del despliegue ante hallazgos críticos.\newline \emph{SD3, Tabla 3.A.5, fila RNF-\allowbreak{}14.07} \\
\textbf{RNF-\allowbreak{}14.08}\newline BTT, RT-11.23 & Cada versión liberada se entrega al CLIENTE junto con su SBOM de 3.11.4 en formato CycloneDX o SPDX, con componentes, versiones y licencias.\newline \textbf{EDT:} 3.11.4 & /P01 a /P03\newline Cumple; Alta; E1 y E2; H5 y H10 & Cada versión liberada debe acompañarse de su inventario de componentes de software en formato CycloneDX o SPDX, entregado al CLIENTE.\newline \emph{SD3, Tabla 3.A.5, fila RNF-\allowbreak{}14.08} \\
\textbf{RNF-\allowbreak{}14.09}\newline BTT, RT-11.24 & Firma de imágenes y procedencia SLSA nivel 3 (GitLab CI + AWS CodeBuild)\newline \textbf{EDT:} 3.1.4 & /P01 a /P03\newline Cumple; Alta; E1 y E2; H5 y H10 & Los artefactos deben firmarse y su procedencia debe verificarse conforme a SLSA nivel 3 o superior.\newline \emph{SD3, Tabla 3.A.5, fila RNF-\allowbreak{}14.09} \\
\textbf{RNF-\allowbreak{}14.10}\newline BTT, RT-11.25 & Datos sintéticos o anonimizados en Desarrollo y QA\newline \textbf{EDT:} 3.1.1 & /P01 a /P03\newline Cumple; Alta; E1 y E2; H5 y H10 & Queda prohibido el uso de datos productivos reales en ambientes no productivos sin anonimización o seudonimización verificable.\newline \emph{SD3, Tabla 3.A.5, fila RNF-\allowbreak{}14.10} \\
\textbf{RNF-\allowbreak{}14.11}\newline Propio LafroX. Fuente: SD2, Anexo 2.2, S-13; Ley N° 21.719. & Posición registrada por vehículo y ruta, no por persona; uso sujeto a acuerdo sindical, acceso por rol autorizado con auditoría y cifrado de campo\newline \textbf{EDT:} 5.4.2, 1.9.3 & /P01 a /P05\newline Cumple; Crítica; E1; H5 & Hasta el acuerdo con el sindicato, la posición se registra por vehículo y ruta, no por persona, y no se usa para controlar la jornada. El acceso a la geolocalización es por rol autorizado, queda auditado (RNF-09.03) y se cifra a nivel de campo (RNF-07.02).\newline \emph{SD3, Tabla 3.A.5b, fila RNF-\allowbreak{}14.11} \\
\textbf{RNF-\allowbreak{}14.12}\newline Propio LafroX. Fuente: Ley N° 21.719; BTT, cap. 11. & Atributos de datos personales de terreno clasificados por categoría; datos sensibles cifrados a nivel de campo y consultas auditadas\newline \textbf{EDT:} 1.9.3, 2.2.4 & /P01 a /P03\newline Cumple; Alta; E1; H5 & El 100 \% de los campos con datos personales capturados en terreno (receptor, conductor, preventista, cliente) tiene categoría asignada. Los de categoría sensible se cifran a nivel de campo (RNF-07.02) y su consulta queda auditada (RNF-09.03).\newline \emph{SD3, Tabla 3.A.5b, fila RNF-\allowbreak{}14.12} \\
\textbf{RNF-\allowbreak{}15.01}\newline BTT, RT-12.07 & Tokens con vida máxima de 30 minutos, inactividad de 15 o 30 minutos, refresco rotatorio de hasta 24 horas o 8 horas para privilegios, revocación y control de sesiones concurrentes.\newline \textbf{EDT:} 3.3.1 & /P01 a /P05\newline Cumple; Crítica; E1 y E2; H5 y H10 & El sistema debe implementar política de sesión con duración máxima, caducidad por inactividad, renovación de la credencial de sesión tras la autenticación, revocación inmediata y control de sesiones concurrentes. Los valores de duración, inactividad y concurrencia se declaran en la arquitectura de la oferta y se verifican antes del cierre de la marcha blanca E1.\newline \emph{SD3, Tabla 3.A.5, fila RNF-\allowbreak{}15.01} \\
\textbf{RNF-\allowbreak{}15.02}\newline BTT, RT-12.08 & Tokens firmados y de vida breve, con credencial de refresco rotatoria; los identificadores de sesión no se transportan en la URL.\newline \textbf{EDT:} 3.3.1 & /P01 a /P05\newline Cumple; Crítica; E1 y E2; H5 y H10 & Las credenciales de sesión deben ser firmadas y de vida breve, con credencial de refresco rotatoria. Queda prohibido transportar identificadores de sesión en la ruta de la dirección web (URL). La vida de cada credencial se declara en la arquitectura de la oferta y se verifica antes del cierre de la marcha blanca E1.\newline \emph{SD3, Tabla 3.A.5, fila RNF-\allowbreak{}15.02} \\
\textbf{RNF-\allowbreak{}15.03}\newline BTT, RT-12.09 & Las altas, cambios de rol, elevaciones, bloqueos y bajas se registran en una bitácora inalterable con retención de siete años.\newline \textbf{EDT:} 3.3.1 & /P01 a /P03\newline Cumple; Alta; E1 y E2; H5 y H10 & El sistema debe registrar auditoría completa del ciclo de vida de la identidad: creación, modificación, elevación, bloqueo y baja de cuentas, con no repudio y retención declarada.\newline \emph{SD3, Tabla 3.A.5, fila RNF-\allowbreak{}15.03} \\
\textbf{RNF-\allowbreak{}16.01}\newline BTT, RT-14.03 & OpenTelemetry (colector ADOT 0.50) y CloudWatch correlacionan solicitudes y transacciones; los percentiles p95 se fijan por operación y se miden sobre la operación percibida y confirmada.\newline \textbf{EDT:} 3.1.5 & /P01 a /P03\newline Cumple; Alta; E1 y E2; H5 y H10 & Los indicadores de nivel de servicio deben medirse sobre la experiencia real de la persona usuaria (percentil 95) y no sobre pruebas sintéticas, sin perjuicio de que estas se empleen como complemento.\newline \emph{SD3, Tabla 3.A.5, fila RNF-\allowbreak{}16.01} \\
\textbf{RNF-\allowbreak{}16.02}\newline BTT, RT-14.05 & Documentación operativa con libro de operación y guía de resolución para cada escenario de falla previsible, mantenidos por el servicio, con automatización progresiva de tareas repetitivas mediante las herramientas previstas.\newline \textbf{EDT:} 3.11.1, 8.1.1 & /P01 a /P03\newline Cumple; Alta; E1 y E2; H5 y H10 & El PROPONENTE debe documentar un libro de operación y una guía de resolución para cada escenario de falla previsible, con automatización progresiva de las tareas repetitivas.\newline \emph{SD3, Tabla 3.A.5, fila RNF-\allowbreak{}16.02} \\
\textbf{RNF-\allowbreak{}16.03}\newline BTT, RT-14.07 & Aplicaciones y colectores excluyen credenciales y datos personales sensibles antes de registrar; acceso limitado y auditado\newline \textbf{EDT:} 3.1.5 & /P01 a /P05\newline Cumple; Crítica; E1 y E2; H5 y H10 & Los registros (logs) de la solución no deben contener datos personales sensibles ni credenciales, y su acceso debe estar controlado y auditado.\newline \emph{SD3, Tabla 3.A.5, fila RNF-\allowbreak{}16.03} \\
\textbf{RNF-\allowbreak{}17.01}\newline BTT, RT-16.10 & Amazon S3 Object Lock conserva auditoría y seguridad por 7 años ($\geq$5 años).\newline \textbf{EDT:} 3.3.5, 3.2.4 & /P01 a /P03\newline Cumple; Alta; E1 y E2; H5 y H10 & El período de retención de la auditoría debe ser el que fije el caso y, en su defecto, no inferior a cinco años.\newline \emph{SD3, Tabla 3.A.5, fila RNF-\allowbreak{}17.01} \\
\textbf{RNF-\allowbreak{}17.02}\newline BTT, RT-16.29 & Exportaciones de gran volumen asíncronas con la sesión disponible; aviso por INT-11 con enlace autenticado al completarse\newline \textbf{EDT:} 3.3.5 & /P01 a /P03\newline Cumple; Alta; E1 y E2; H5 y H10 & Las exportaciones de gran volumen deben procesarse de forma asíncrona, con notificación al completarse y sin bloquear la sesión del usuario.\newline \emph{SD3, Tabla 3.A.5, fila RNF-\allowbreak{}17.02} \\
\textbf{RNF-\allowbreak{}17.03}\newline BTT, RT-16.34 & Portales estáticos en S3 con caché de CloudFront, separados de la API transaccional, con límites de tasa\newline \textbf{EDT:} 3.2.1 & /P01 a /P03\newline Cumple; Alta; E1 y E2; H5 y H10 & El portal público debe resistir picos de tráfico sin degradar los servicios transaccionales internos, mediante aislamiento de recursos y caché.\newline \emph{SD3, Tabla 3.A.5, fila RNF-\allowbreak{}17.03} \\
\textbf{RNF-\allowbreak{}19.01}\newline BTT, RT-09.01 & Cálculo de capacidad: 14,66 TPS peak, usuarios concurrentes, volumen y crecimiento tomados del Caso, Cap. 14\newline \textbf{EDT:} 2.1.2, 8.4.1 & /P01 a /P03\newline Cumple; Alta; E1 y E2; H5 y H10 & El PROPONENTE debe presentar el cálculo de capacidad que sustenta su dimensionamiento, con los supuestos de usuarios concurrentes, transacciones por segundo, volumen de datos y crecimiento anual, tomados de la volumetría del caso.\newline \emph{SD3, Tabla 3.A.5, fila RNF-\allowbreak{}19.01} \\
\textbf{RNF-\allowbreak{}19.02}\newline BTT, RT-02.10, RT-09.04 & ECS Fargate: API con PHP-FPM >70 \%, 2 a 4 tareas y techo 8; integración automática de 2 a 4 tareas por perfil, en ventanas de un minuto: consumidor por edad de mensaje SQS >60 s y trabajos por cola >100 mensajes visibles; costo en la Oferta Económica (Bases Administrativas, art. 50.2, p. 29).\newline \textbf{EDT:} 3.2.1 & /P01 a /P03\newline Cumple; Alta; E1 y E2; H5 y H10 & Las capas de aplicación e integración deben escalar horizontalmente de forma automática, con umbrales, límites superiores y costo asociado declarados en la oferta (arquitectura y Oferta Económica).\newline \emph{SD3, Tabla 3.A.5, fila RNF-\allowbreak{}19.02} \\
\textbf{RNF-\allowbreak{}19.03}\newline BTT, RT-09.05 & Primer cuello de botella identificado con detección y resolución\newline \textbf{EDT:} 2.1.2, 8.4.1 & /P01 a /P03\newline Cumple; Alta; E1 y E2; H5 y H10 & El PROPONENTE debe identificar el componente que primero se convertirá en cuello de botella al crecer la carga y explicar cómo lo detectará y cómo lo resolverá.\newline \emph{SD3, Tabla 3.A.5, fila RNF-\allowbreak{}19.03} \\
\textbf{RNF-\allowbreak{}19.04}\newline BTT, RT-09.06, RT-09.07 & Prueba de carga a 1,5 $\times$ 14,66 = 21,99 TPS en Preproducción y estrés hasta el punto de quiebre\newline \textbf{EDT:} 3.8.4, 3.9.3 & /P01 a /P03\newline Cumple; Alta; E1 y E2; H5 y H10 & Se deben ejecutar pruebas de carga sobre Preproducción con un volumen equivalente a 1,5 veces el peak declarado, y pruebas de estrés hasta identificar el punto de quiebre de la solución.\newline \emph{SD3, Tabla 3.A.5, fila RNF-\allowbreak{}19.04} \\
\textbf{RNF-\allowbreak{}20.01}\newline BTT, RT-10.01 & Disponibilidad $\geq$ 99,9 \% mensual de la transacción crítica de extremo a extremo\newline \textbf{EDT:} 8.1.2 & /P01 a /P05\newline Cumple; Crítica; E1 y E2; H5 y H10 & La solución debe alcanzar una disponibilidad mensual mínima de 99.9\% para los servicios clasificados como críticos, medida sobre la transacción de negocio de extremo a extremo.\newline \emph{SD3, Tabla 3.A.5, fila RNF-\allowbreak{}20.01} \\
\textbf{RNF-\allowbreak{}20.02}\newline BTT, RT-10.02 & Clasificación de servicios en crítico, alto, medio y bajo con su nivel del Art. 78\newline \textbf{EDT:} 2.5.1 & /P01 a /P05\newline Cumple; Crítica; E1 y E2; H5 y H10 & El PROPONENTE debe clasificar cada servicio de la solución en crítico, alto, medio o bajo, justificando la clasificación con el impacto operacional de su indisponibilidad, y aplicar a cada uno el nivel de servicio correspondiente.\newline \emph{SD3, Tabla 3.A.5, fila RNF-\allowbreak{}20.02} \\
\textbf{RNF-\allowbreak{}20.03}\newline BTT, RT-10.03 & Plan de continuidad conforme a ISO 22301 con análisis de impacto, escenarios, procedimientos manuales de respaldo y criterios de activación por proceso aprobados por el CLIENTE; los escenarios de continuidad y recuperación se prueban antes de producción y en Operación.\newline \textbf{EDT:} 2.5.1 & /P01 a /P05\newline Cumple; Crítica; E1 y E2; H5 y H10 & El PROPONENTE debe elaborar un plan de continuidad del negocio conforme a ISO 22301, con análisis de impacto en el negocio, escenarios de contingencia, procedimientos manuales de respaldo y criterios de activación.\newline \emph{SD3, Tabla 3.A.5, fila RNF-\allowbreak{}20.03} \\
\textbf{RNF-\allowbreak{}20.04}\newline BTT, RT-10.07 & Inyección de fallas (instancia, zona, dependencia, latencia, disco) antes de cada paso a producción y semestral\newline \textbf{EDT:} 3.8.4, 3.9.3, 8.2.2 & /P01 a /P03\newline Cumple; Alta; E1 y E2; H5 y H10 & Se deben ejecutar pruebas de resiliencia mediante inyección controlada de fallas (caída de instancia, de zona, de dependencia externa, latencia elevada, saturación de disco) antes de cada paso a producción y al menos una vez por semestre durante la Operación.\newline \emph{SD3, Tabla 3.A.5, fila RNF-\allowbreak{}20.04} \\
\textbf{RNF-\allowbreak{}20.05}\newline BTT, RT-10.08 & Comportamiento por dependencia externa sin respuesta, con error o lenta\newline \textbf{EDT:} 2.5.2 & /P01 a /P03\newline Cumple; Alta; E1 y E2; H5 y H10 & El PROPONENTE debe documentar, por cada dependencia externa, el comportamiento de la solución cuando esa dependencia no responde, responde con error o responde con lentitud.\newline \emph{SD3, Tabla 3.A.5, fila RNF-\allowbreak{}20.05} \\
\textbf{RNF-\allowbreak{}20.06}\newline BTT, RT-07.04, RT-07.07 & Sitio secundario en us-east-1 con RTO de hasta cuatro horas y RPO de hasta quince minutos para los escenarios de diseño, sostenido por la extracción continua por fibra, LTE y Starlink, con dos ensayos regionales por año.\newline \textbf{EDT:} 3.1.3, 2.5.2 & /P01 a /P05\newline Cumple; Crítica; E1 y E2; H5 y H10 & El PROPONENTE debe habilitar un sitio secundario para recuperación ante desastres, con RTO $\leq$ 4 horas y RPO $\leq$ 15 minutos para los servicios críticos, y probar la conmutación al menos 2 veces al año.\newline \emph{SD3, Tabla 3.A.5, fila RNF-\allowbreak{}20.06} \\
\end{tablasimple}

\begin{tablasimple}{>{\hsize=0.82\hsize\linewidth=\hsize}Y >{\hsize=1.04\hsize\linewidth=\hsize}Y >{\hsize=0.65\hsize\linewidth=\hsize}Y >{\hsize=1.49\hsize\linewidth=\hsize}Y}
{\multicolumn{4}{l}{\textbf{Matriz individual: bloque 11 de 26.}} \\
\cab{Requisito y origen} & \cab{Componente y EDT} & \cab{Prueba prevista y estado} & \cab{Criterio y fuente}}
\textbf{RNF-\allowbreak{}20.07}\newline BTT, RT-07.09, RT-07.12 & Esquema 3-2-1-1-0: NAS Synology RS2423RP+ WORM, AWS Backup con Vault Lock, S3 Object Lock y medio físico en custodia\newline \textbf{EDT:} 3.2.4, 6.4.3 & /P01 a /P05\newline Cumple; Crítica; E1 y E2; H5 y H10 & La política de respaldo debe seguir el esquema 3-2-1-1-0: tres copias, en dos medios distintos, una fuera de sitio, una inmutable o fuera de línea y cero errores de verificación de restauración.\newline \emph{SD3, Tabla 3.A.5, fila RNF-\allowbreak{}20.07} \\
\textbf{RNF-\allowbreak{}21.01}\newline BTT, RT-21.01 & NOC 24$\times$7 propio de LafroX en su Centro de Operaciones de Valparaíso (Av. Brasil 2241), fuera de los sitios que monitorea y con continuidad en trabajo remoto mediante Verified Access; un puesto permanente cubierto por cuatro personas en turnos rotativos, cinco desde el 26-04-2028; procedimientos en el libro de operación y las guías de resolución, el escalamiento ante congelamientos y la conmutación regional.\newline \textbf{EDT:} 8.1.1 & /P01 a /P05\newline Cumple; Crítica; E1 y E2; H5 y H10 & El ADJUDICATARIO debe disponer de un centro de operaciones de red (NOC) con cobertura 24$\times$7$\times$365, propio o subcontratado, y declarar su ubicación, dotación por turno y procedimientos.\newline \emph{SD3, Tabla 3.A.5, fila RNF-\allowbreak{}21.01} \\
\textbf{RNF-\allowbreak{}21.02}\newline BTT, RT-21.02 & Guillermo Castillo, Líder de Operación / SRE, designado gerente de servicio dedicado de la Operación con 100 \% de dedicación en los meses 21 a 56; es la contraparte permanente del CLIENTE y representa a LafroX en el Comité de Operación mensual (niveles de servicio, incidentes, problemas y mejora continua)\newline \textbf{EDT:} 8.1.1 & /P01 a /P03\newline Cumple; Alta; E1 y E2; H5 y H10 & El ADJUDICATARIO debe designar un gerente de servicio dedicado como contraparte permanente del CLIENTE durante la fase de Operación.\newline \emph{SD3, Tabla 3.A.5, fila RNF-\allowbreak{}21.02} \\
\textbf{RNF-\allowbreak{}21.03}\newline BTT, RT-21.06 & Mesa con atención mínima de 80 \% antes de 20 segundos, resolución al primer contacto de al menos 70 \% y abandono máximo de 5 \%; medición y reporte mensual con recalibración y refuerzo. Atención 04:00–22:00 de lunes a sábado y cobertura 24$\times$7 en peaks de septiembre y diciembre y ante críticos en la ventana de despacho.\newline \textbf{EDT:} 8.1.1 & /P01 a /P03\newline Cumple; Alta; E1 y E2; H5 y H10 & El centro de atención debe cumplir un tiempo de atención al usuario final de 80\% antes de 20 segundos y una resolución al primer contacto de al menos 70\%. La tasa de abandono no debe superar el 5\%.\newline \emph{SD3, Tabla 3.A.5, fila RNF-\allowbreak{}21.03} \\
\textbf{RNF-\allowbreak{}21.04}\newline BTT, RT-21.07 & Atención 04:00–22:00 de lunes a sábado, que contiene el mínimo de 08:00 a 20:00 en días hábiles e incluye la ventana de despacho de 05:30 a 07:00; incidentes críticos y altos atendidos 24$\times$7$\times$365 (D-10) y 24$\times$7 en los peaks de septiembre y diciembre\newline \textbf{EDT:} 8.1.1 & /P01 a /P03\newline Cumple; Alta; E1 y E2; H5 y H10 & El horario mínimo de atención será de 08:00 a 20:00 en días hábiles, ampliado a 24$\times$7 para los incidentes de severidad crítica y para la ventana operacional crítica (despacho de 05:30 a 07:00).\newline \emph{SD3, Tabla 3.A.5, fila RNF-\allowbreak{}21.04} \\
\textbf{RNF-\allowbreak{}21.05}\newline BTT, RT-21.14 & Dotación de la mesa por Erlang C a partir de los contactos proyectados\newline \textbf{EDT:} 8.1.1 & /P01 a /P03\newline Cumple; Alta; E1 y E2; H5 y H10 & El PROPONENTE debe dimensionar la dotación del centro de atención con fundamento cuantitativo, empleando teoría de colas o el modelo Erlang C, a partir del volumen de contactos proyectado del caso.\newline \emph{SD3, Tabla 3.A.5, fila RNF-\allowbreak{}21.05} \\
\textbf{RNF-\allowbreak{}21.06}\newline BTT, RT-21.15 & Canal único de incidentes y solicitudes, con número de ticket, clasificación por severidad y seguimiento del ciclo completo hasta el cierre conforme del CLIENTE.\newline \textbf{EDT:} 8.1.2 & /P01 a /P03\newline Cumple; Alta; E1 y E2; H5 y H10 & Debe existir un canal único de registro de incidentes y solicitudes, con número de ticket, clasificación por severidad y seguimiento del ciclo de vida completo hasta el cierre conforme.\newline \emph{SD3, Tabla 3.A.5, fila RNF-\allowbreak{}21.06} \\
\textbf{RNF-\allowbreak{}21.07}\newline BTT, RT-21.16 & Traslado de especialistas de nivel 2 (SRE de mantenimiento del hardware, paquete 8.2.3) y de nivel 3 (DES de mantención correctiva, paquete 8.2.1) a Curicó, Chillán, Los Ángeles o el sitio que corresponda, por el medio más rápido disponible y hasta resolver, sin retirar puestos de mesa, NOC ni SOC; la recuperación dentro del Art. 78° sigue siendo remota; el tiempo se imputa al paquete de origen, el exceso lo cubre LafroX y pasajes y viáticos están incluidos en la Oferta Económica\newline \textbf{EDT:} 8.2.3 & /P01 a /P03\newline Cumple; Alta; E1 y E2; H5 y H10 & Los especialistas de niveles 2 y 3 deben trasladarse a los sitios alejados (Curicó, Chillán, Los Ángeles) cuando la resolución lo requiera, por el medio más rápido disponible, sin alterar la atención normal del resto de los sitios. El costo del traslado está incluido en la oferta.\newline \emph{SD3, Tabla 3.A.5, fila RNF-\allowbreak{}21.07} \\
\textbf{RNF-\allowbreak{}21.08}\newline BTT, RT-21.19 & Bolsa anual de mantención evolutiva de 768 HH DES dentro de las 4.608 HH del paquete 8.2.1 (1.536 HH anuales), separada de la mantención correctiva, que no tiene tope ni costo adicional; el saldo no usado se acumula al año siguiente y el excedente se liquida con la tarifa de la Oferta Económica; perfil DES con el mismo pipeline y puertas de calidad; procedimiento de solicitud por la mesa, estimación en cinco días hábiles, aprobación del CLIENTE en el Comité de Operación y liquidación mensual con saldo\newline \textbf{EDT:} 8.2.1 & /P01 a /P03\newline Cumple; Alta; E1 y E2; H5 y H10 & El PROPONENTE debe declarar el tamaño de la bolsa anual de horas de mantención evolutiva, su composición por perfil y el procedimiento de solicitud, estimación, aprobación y liquidación.\newline \emph{SD3, Tabla 3.A.5, fila RNF-\allowbreak{}21.08} \\
\textbf{RNF-\allowbreak{}21.09}\newline BTT, RT-21.22 & Planificación de actualizaciones de componentes base con ventana y reversión\newline \textbf{EDT:} 8.2.5, 8.1.1 & /P01 a /P03\newline Cumple; Alta; E1 y E2; H5 y H10 & Las actualizaciones de versión de los componentes de base (sistemas operativos, bases de datos, frameworks) deben planificarse anualmente, con ventana acordada y plan de reversión.\newline \emph{SD3, Tabla 3.A.5, fila RNF-\allowbreak{}21.09} \\
\textbf{RNF-\allowbreak{}22.01}\newline BTT, RT-22.01 & Capacitación por roles (finales, avanzados, administradores, técnicos, soporte)\newline \textbf{EDT:} 7.1.1 & /P01 a /P03\newline Cumple; Alta; E1 y E2; H5 y H10 & El PROPONENTE debe presentar un plan de capacitación estructurado por perfil: personas usuarias finales, usuarias avanzadas, administradoras, equipo técnico y soporte de niveles 1 y 2.\newline \emph{SD3, Tabla 3.A.5, fila RNF-\allowbreak{}22.01} \\
\textbf{RNF-\allowbreak{}22.02}\newline BTT, RT-22.03 & Manuales por perfil, guías rápidas, preguntas frecuentes, videos tutoriales y base de conocimiento consultable de 3.11.2 y 8.5.3 entregados en español, editables y de propiedad del CLIENTE.\newline \textbf{EDT:} 7.1.5 & /P01 a /P03\newline Cumple; Alta; E1 y E2; H5 y H10 & Todo el material de capacitación debe entregarse en español, en formato editable (Word, PowerPoint, etc.) y de propiedad del CLIENTE: manuales por perfil, guías rápidas, preguntas frecuentes, videos tutoriales y base de conocimiento consultable.\newline \emph{SD3, Tabla 3.A.5, fila RNF-\allowbreak{}22.02} \\
\textbf{RNF-\allowbreak{}22.03}\newline BTT, RT-22.04 & Capacitación por turnos y horarios acordados con el CLIENTE, sin interrumpir despacho de mañana ni preparación nocturna y considerando los peaks de septiembre y diciembre, con acompañamiento y autoformación previstos.\newline \textbf{EDT:} 7.1.1, 7.1.2 & /P01 a /P03\newline Cumple; Alta; E1 y E2; H5 y H10 & La capacitación no podrá afectar la operación del CLIENTE: se programará por turnos y en horarios acordados, considerando la estacionalidad (peaks de septiembre y diciembre) y la ventana operacional (despacho de mañana, preparación nocturna).\newline \emph{SD3, Tabla 3.A.5, fila RNF-\allowbreak{}22.03} \\
\textbf{RNF-\allowbreak{}22.04}\newline BTT, RT-22.05 & Evaluación y certificación como condición de cierre de cada marcha blanca\newline \textbf{EDT:} 7.1.4, 7.3.1 & /P01 a /P03\newline Cumple; Alta; E1 y E2; H5 y H10 & Las personas usuarias administradoras y el equipo técnico del CLIENTE deben ser evaluados y certificados como condición para el cierre de cada marcha blanca.\newline \emph{SD3, Tabla 3.A.5, fila RNF-\allowbreak{}22.04} \\
\textbf{RNF-\allowbreak{}22.05}\newline BTT, RT-22.06 & Jornadas anuales de actualización y capacitación del personal nuevo\newline \textbf{EDT:} 8.5.1 & /P01 a /P03\newline Cumple; Alta; E1 y E2; H5 y H10 & Durante la Operación, el ADJUDICATARIO debe ejecutar al menos dos jornadas anuales de actualización y capacitar al personal nuevo del CLIENTE sin costo adicional.\newline \emph{SD3, Tabla 3.A.5, fila RNF-\allowbreak{}22.05} \\
\textbf{RNF-\allowbreak{}22.06}\newline BTT, RT-22.07 & Acompañamiento y mentoría al equipo técnico tras la Etapa 2\newline \textbf{EDT:} 8.5.2 & /P01 a /P03\newline Cumple; Alta; E1 y E2; H5 y H10 & El ADJUDICATARIO debe proveer acompañamiento y mentoría al equipo técnico del CLIENTE durante al menos seis meses posteriores al paso a producción de la Etapa 2.\newline \emph{SD3, Tabla 3.A.5, fila RNF-\allowbreak{}22.06} \\
\textbf{RNF-\allowbreak{}23.01}\newline BTT, RT-08.10, RT-08.11 & Equipos de terreno y operación (T-11): B-01, B-02, kit B-03 a B-14 y SIM B-15; incluye accesorios y consumibles.\newline \emph{T-12, fila RNF-\allowbreak{}23.01}\newline \textbf{EDT:} 5.1.1, 5.1.4, 2.4.3, 2.4.4 & /P01 a /P03\newline Cumple; Alta; E1 y E2; H5 y H10 & El PROPONENTE debe especificar cada dispositivo de terreno (lectores de código, impresoras térmicas, sensores de temperatura, terminales móviles) con marca, modelo de referencia, cantidad, características mínimas, accesorios, consumibles y costo unitario estimado.\newline \emph{SD3, Tabla 3.A.5, fila RNF-\allowbreak{}23.01} \\
\textbf{RNF-\allowbreak{}23.02}\newline BTT, RT-08.12 & TC58e, MC9400 estándar y Cold Storage Freezer, ZQ620 Plus, termógrafo EYE SENSOR, balanza Dibal, ZT411, PAX A920 Pro y kit de frío B-03 a B-11.\newline \emph{T-12, fila RNF-\allowbreak{}23.02}\newline \textbf{EDT:} 5.1.1, 5.1.4, 2.4.3 & /P01 a /P03\newline Cumple; Alta; E1 y E2; H5 y H10 & Los dispositivos de terreno deben declarar su grado de protección contra polvo y agua (IP) y su resistencia a caídas, coherentes con el entorno de operación (intemperie, humedad, polvo, vibración, temperatura). El grado IP y la altura de caída de cada modelo se declaran en la especificación de hardware de la oferta (RNF-23.01) y se verifican con la ficha del fabricante.\newline \emph{SD3, Tabla 3.A.5, fila RNF-\allowbreak{}23.02} \\
\textbf{RNF-\allowbreak{}23.03}\newline BTT, RT-08.13 & Vida útil declarada de cinco años, reserva del 10 \%, garantía, repuestos y reposición durante los 56 meses del contrato. LafroX mantiene los componentes de cámara y los kits de camión, incluida calibración y revisión de glicol y baterías; las reservas térmicas son 31 por componente unitario, 93 sondas y 93 botellas.\newline \textbf{EDT:} 5.1.1, 5.1.4, 8.2.7 & /P01 a /P03\newline Cumple; Alta; E1 y E2; H5 y H10 & El PROPONENTE debe indicar el ciclo de vida esperado de cada dispositivo de terreno, la disponibilidad de repuestos y el plan de reposición durante los 56 meses del Contrato.\newline \emph{SD3, Tabla 3.A.5, fila RNF-\allowbreak{}23.03} \\
\textbf{RNF-\allowbreak{}23.04}\newline BTT, RT-08.07, RT-08.08, RT-08.09 & Ocho estaciones de trabajo nuevas con dos monitores de 24 pulgadas, ergonomía NCh 2527 y certificación de eficiencia energética, incorporadas antes de su acceso por Verified Access a la gestión central con CrowdStrike Falcon, cifrado de disco, control de extraíbles y actualización automatizada.\newline \textbf{EDT:} 5.1.2; 5.1.3 & /P01 a /P03\newline Cumple; Alta; E1 y E2; H5 y H10 & El PROPONENTE debe especificar las estaciones de trabajo requeridas para la operación y administración de la plataforma, en la cantidad que determine el dimensionamiento, con monitores duales y gestión centralizada con cifrado de disco. Los puestos cumplirán las condiciones ergonómicas exigidas en RT-08.08, que remite a la NCh 2527 (\textit{Bases Técnicas Transversales}, 2026, cap. 8, RT-08.08, p. 19).\newline \emph{SD3, Tabla 3.A.5, fila RNF-\allowbreak{}23.04} \\
\textbf{RT-\allowbreak{}02.01}\newline Bases Técnicas Transversales, RT-02.01 & Arquitectura lógica de ocho capas con los módulos M1–M12 y los servicios transversales del inventario de capacidades (Anexo 4-N)\newline \textbf{EDT:} 2.1.1 & /P01\newline Cumple; RT & La solución se organizará en las ocho capas del numeral 2.1. El PROPONENTE presentará el diagrama de la arquitectura lógica identificando cada capa, sus componentes y las interfaces entre ellas.\newline \emph{T-12, Parte B, fila RT-\allowbreak{}02.01} \\
\textbf{RT-\allowbreak{}02.02}\newline Bases Técnicas Transversales, RT-02.02 & Monolito modular Laravel 13.x / PHP 8.5.x con perfiles desplegables por separado (API, wms\_only, shipper, consumidor, trabajos, planificador, erp-sync) y motor de rutas en contenedor propio\newline \textbf{EDT:} 2.1.1 & /P01\newline Cumple; RT & La arquitectura será modular, con límites de contexto explícitos y acoplamiento débil. Se rechazará toda arquitectura monolítica que no permita desplegar de forma independiente sus componentes críticos.\newline \emph{T-12, Parte B, fila RT-\allowbreak{}02.02} \\
\textbf{RT-\allowbreak{}02.03}\newline Bases Técnicas Transversales, RT-02.03 & Descripción de arquitectura conforme a ISO/IEC/IEEE 42010 con vistas lógica, de procesos, de despliegue, de datos y de seguridad\newline \textbf{EDT:} 2.1.1, 2.1.2 & /P01\newline Cumple; RT & La descripción de la arquitectura se ajustará a ISO/IEC/IEEE 42010, con vistas lógica, de procesos, de despliegue, de datos y de seguridad.\newline \emph{T-12, Parte B, fila RT-\allowbreak{}02.03} \\
\textbf{RT-\allowbreak{}02.04}\newline Bases Técnicas Transversales, RT-02.04 & Registro único ADR-01 a ADR-24, fechado, con decisión, alternativas, criterio, consecuencias y evidencia exigida; entregable actualizado durante la ejecución bajo el Comité de Arquitectura.\newline \textbf{EDT:} 2.1.1, 1.8.4 & /P01\newline Cumple; RT & El PROPONENTE mantendrá un registro de decisiones de arquitectura (ADR) fechado, con la alternativa escogida, las alternativas descartadas y el criterio de decisión. El registro es entregable contractual y se actualizará durante toda la ejecución.\newline \emph{T-12, Parte B, fila RT-\allowbreak{}02.04} \\
\textbf{RT-\allowbreak{}02.05}\newline Bases Técnicas Transversales, RT-02.05 & Capa de negocio sin estado en ECS Fargate; sesión en Keycloak 26.x y estado en Aurora PostgreSQL 16 y ElastiCache con Valkey 8.1\newline \textbf{EDT:} 2.1.1 & /P01\newline Cumple; RT & La capa de servicios de negocio será sin estado. El estado de sesión y el estado de proceso residirán en almacenes externos con alta disponibilidad.\newline \emph{T-12, Parte B, fila RT-\allowbreak{}02.05} \\
\end{tablasimple}

\begin{tablasimple}{>{\hsize=0.82\hsize\linewidth=\hsize}Y >{\hsize=1.04\hsize\linewidth=\hsize}Y >{\hsize=0.65\hsize\linewidth=\hsize}Y >{\hsize=1.49\hsize\linewidth=\hsize}Y}
{\multicolumn{4}{l}{\textbf{Matriz individual: bloque 12 de 26.}} \\
\cab{Requisito y origen} & \cab{Componente y EDT} & \cab{Prueba prevista y estado} & \cab{Criterio y fuente}}
\textbf{RT-\allowbreak{}02.06}\newline Bases Técnicas Transversales, RT-02.06 & UUID declarado por el cliente conservado en reintentos de API y eventos; resultado durable deduplicado por 30 días. Una escritura más antigua se aísla para revisión antes de admitirla de nuevo, según 4.1.16.\newline \textbf{EDT:} 3.3.6 & /P01\newline Cumple; RT & Toda operación de escritura expuesta a reintentos será idempotente, con clave de idempotencia declarada por el cliente y ventana de deduplicación documentada.\newline \emph{T-12, Parte B, fila RT-\allowbreak{}02.06} \\
\textbf{RT-\allowbreak{}02.07}\newline Bases Técnicas Transversales, RT-02.07 & RabbitMQ 4.x con confirmación de publicación y Amazon SQS FIFO con grupos por agregado; consumidor PHP 8.5 con deduplicación\newline \textbf{EDT:} 3.2.2 & /P01\newline Cumple; RT & Los flujos de eventos garantizarán entrega al menos una vez, con deduplicación en el consumidor y orden garantizado dentro de la partición o del agregado cuando el proceso lo exija.\newline \emph{T-12, Parte B, fila RT-\allowbreak{}02.07} \\
\textbf{RT-\allowbreak{}02.08}\newline Bases Técnicas Transversales, RT-02.08 & Tiempos máximos explícitos, reintento con espera exponencial y variación aleatoria, cortacircuitos, colas aisladas y límites de tasa en Amazon API Gateway\newline \textbf{EDT:} 2.1.1 & /P01\newline Cumple; RT & La solución implementará patrones de resiliencia demostrables: reintento con retroceso exponencial y variación aleatoria, cortacircuitos, mamparos de aislamiento, límites de tasa y tiempo de espera explícito en toda llamada remota. No se admiten llamadas remotas sin tiempo de espera.\newline \emph{T-12, Parte B, fila RT-\allowbreak{}02.08} \\
\textbf{RT-\allowbreak{}02.09}\newline Bases Técnicas Transversales, RT-02.09 & Degradación controlada informada a la persona usuaria («modo offline»)\newline \textbf{EDT:} 3.3.6 & /P01\newline Cumple; RT & La solución degradará de forma elegante: ante la indisponibilidad de un componente no crítico deberá continuar operando en modo reducido, informando la degradación a la persona usuaria, y nunca fallar de forma total.\newline \emph{T-12, Parte B, fila RT-\allowbreak{}02.09} \\
\textbf{RT-\allowbreak{}02.10}\newline Bases Técnicas Transversales, RT-02.10 & ECS Fargate: API con PHP-FPM >70 \%, 2 a 4 tareas y techo 8; integración automática de 2 a 4 tareas por perfil, en ventanas de un minuto: consumidor por edad de mensaje SQS >60 s y trabajos por cola >100 mensajes visibles; costo en la Oferta Económica (Bases Administrativas, art. 50.2, p. 29).\newline \textbf{EDT:} 3.2.1 & /P01\newline Cumple; RT & Las capas de aplicación e integración escalarán horizontalmente de forma automática, con umbrales, límites superiores y costo asociado declarados en la oferta.\newline \emph{T-12, Parte B, fila RT-\allowbreak{}02.10} \\
\textbf{RT-\allowbreak{}02.11}\newline Bases Técnicas Transversales, RT-02.11 & Declaración de puntos únicos de falla y riesgos residuales justificados\newline \textbf{EDT:} 2.1.2 & /P01\newline Cumple; RT & El PROPONENTE declarará explícitamente los puntos únicos de falla que subsistan en su arquitectura y justificará por qué son aceptables. Omitir esta declaración cuando existan puntos únicos de falla se evaluará como observación grave.\newline \emph{T-12, Parte B, fila RT-\allowbreak{}02.11} \\
\textbf{RT-\allowbreak{}02.12}\newline Bases Técnicas Transversales, RT-02.12 & Nuevo sitio incorporado parametrizando identificador, bodega, rutas y cartera en los maestros existentes, con perfil local e infraestructura como código; bloques 10.6.0.0/16 y 10.7.0.0/16 reservados\newline \textbf{EDT:} 2.1.1 & /P01\newline Cumple; RT & La solución admitirá su replicación a nuevas unidades, sitios, sucursales o filiales del CLIENTE sin rediseño arquitectónico, mediante parametrización o multi-tenencia. Caso, cap. 15 — Replicación a nuevas unidades: Exigible. La compañía evalúa abrir un centro de distribución en la Región de Los Lagos hacia 2030. La solución debe admitir la incorporación de un nuevo sitio, con su bodega, sus rutas y su cartera, por parametrización.\newline \emph{T-12, Parte B, fila RT-\allowbreak{}02.12} \\
\textbf{RT-\allowbreak{}02.13}\newline Bases Técnicas Transversales, RT-02.13 & Modelo de dominio con entidades, relaciones y eventos canónicos\newline \textbf{EDT:} 2.1.3 & /P01\newline Cumple; RT & El PROPONENTE presentará un modelo de dominio del negocio del caso, con las entidades principales, sus relaciones y los eventos de negocio que las modifican.\newline \emph{T-12, Parte B, fila RT-\allowbreak{}02.13} \\
\textbf{RT-\allowbreak{}02.14}\newline Bases Técnicas Transversales, RT-02.14 & Capa anticorrupción del ERP (A-04, erp-sync), estrangulamiento del WMS de 2013 (ADR-08) y abstracción de proveedores por adaptador\newline \textbf{EDT:} 3.3.2, 3.3.3 & /P01\newline Cumple; RT & Se valorará la aplicación documentada de patrones de arquitectura evolutiva que permitan sustituir un componente sin reescribir la solución: capa anticorrupción frente a sistemas heredados, estrangulamiento progresivo y abstracción de proveedores.\newline \emph{T-12, Parte B, fila RT-\allowbreak{}02.14} \\
\textbf{RT-\allowbreak{}03.01}\newline Bases Técnicas Transversales, RT-03.01 & Amazon Web Services: región primaria sa-east-1 (São Paulo) y secundaria us-east-1\newline \textbf{EDT:} 5.2.1, 2.1.2 & /P01\newline Cumple; RT & El PROPONENTE declarará el proveedor de nube pública, la región primaria y la región secundaria utilizadas. El proveedor deberá contar con presencia de región o zona en Chile o en Sudamérica.\newline \emph{T-12, Parte B, fila RT-\allowbreak{}03.01} \\
\textbf{RT-\allowbreak{}03.02}\newline Bases Técnicas Transversales, RT-03.02 & ECS Fargate, Application Load Balancer, Aurora PostgreSQL 16, ElastiCache con Valkey 8.1 y OpenAS2 4.9 en dos zonas de sa-east-1\newline \textbf{EDT:} 3.1.2 & /P01\newline Cumple; RT & Todos los componentes con requisito de alta disponibilidad se desplegarán en al menos dos zonas de disponibilidad. No se aceptará un diseño en una sola zona.\newline \emph{T-12, Parte B, fila RT-\allowbreak{}03.02} \\
\textbf{RT-\allowbreak{}03.03}\newline Bases Técnicas Transversales, RT-03.03 & Terraform 1.15 y Ansible (ansible-core 2.21) versionados en el repositorio del CLIENTE; creación inicial de la cuenta raíz documentada y sin altas manuales por consola\newline \textbf{EDT:} 3.1.4 & /P01\newline Cumple; RT & La totalidad de la infraestructura se definirá como código, versionada en el repositorio del CLIENTE, revisable y reproducible. No se admite infraestructura creada manualmente por consola, salvo la cuenta raíz inicial, cuya creación deberá documentarse.\newline \emph{T-12, Parte B, fila RT-\allowbreak{}03.03} \\
\textbf{RT-\allowbreak{}03.04}\newline Bases Técnicas Transversales, RT-03.04 & VPC de Producción con zona privada para aplicación y datos; subred pública solo con NAT; el canal AS2 llega por CloudFront a un ALB interno; portales en S3 privado tras CloudFront\newline \textbf{EDT:} 3.1.2 & /P01\newline Cumple; RT & La red se segmentará por capas, con subredes privadas para aplicación y datos, y exposición pública restringida a la capa de borde. Ningún componente de datos será alcanzable desde Internet.\newline \emph{T-12, Parte B, fila RT-\allowbreak{}03.04} \\
\textbf{RT-\allowbreak{}03.05}\newline Bases Técnicas Transversales, RT-03.05 & Servicios administrados de AWS (Aurora, ECS Fargate, SQS, ElastiCache, Secrets Manager) con excepciones justificadas en ADR\newline \textbf{EDT:} 2.1.2 & /P01\newline Cumple; RT & El PROPONENTE privilegiará servicios administrados por sobre servicios autoadministrados cuando ello reduzca el riesgo operacional, y justificará cada excepción.\newline \emph{T-12, Parte B, fila RT-\allowbreak{}03.05} \\
\textbf{RT-\allowbreak{}03.06}\newline Bases Técnicas Transversales, RT-03.06 & Etiquetado obligatorio por ambiente, módulo y centro de costo desde la infraestructura como código; desglose mensual con Cost Explorer y alertas de AWS Budgets\newline \textbf{EDT:} 8.4.2 & /P01\newline Cumple; RT & Se aplicarán prácticas FinOps: etiquetado obligatorio de todos los recursos por ambiente, módulo y centro de costo; presupuestos con alertas de desviación; y reporte mensual de consumo desglosado entregado al CLIENTE.\newline \emph{T-12, Parte B, fila RT-\allowbreak{}03.06} \\
\textbf{RT-\allowbreak{}03.07}\newline Bases Técnicas Transversales, RT-03.07 & Reversibilidad: componentes portables (Laravel 13.x/PHP 8.5.x, Angular 22.x, PostgreSQL 16.x, RabbitMQ 4.x, Keycloak 26.x) y no portables (DynamoDB, Redshift/Glue/QuickSight, API Gateway, Verified Access), con esfuerzo de salida\newline \textbf{EDT:} 1.7.1 & /P01\newline Cumple; RT & El PROPONENTE declarará su estrategia de reversibilidad y de mitigación del bloqueo por proveedor, identificando qué componentes son portables, cuáles no lo son y cuál sería el esfuerzo estimado de una migración.\newline \emph{T-12, Parte B, fila RT-\allowbreak{}03.07} \\
\textbf{RT-\allowbreak{}03.08}\newline Bases Técnicas Transversales, RT-03.08 & Savings Plans para la capacidad base; compromiso y costo reflejados en la estructura de costos de la Oferta Económica (Bases Administrativas, art. 50.2, p. 29).\newline \textbf{EDT:} 5.2.1 & /P01\newline Cumple; RT & La solución empleará instancias reservadas, planes de ahorro o capacidad comprometida cuando el perfil de carga lo justifique, y lo reflejará en la estructura de costos.\newline \emph{T-12, Parte B, fila RT-\allowbreak{}03.08} \\
\textbf{RT-\allowbreak{}03.09}\newline Bases Técnicas Transversales, RT-03.09 & ECS Fargate y autorizador AWS Lambda; comparación de costo frente a instancias permanentes: capacidad bajo demanda para rutas y peak y Savings Plans para la base\newline \textbf{EDT:} 3.2.1 & /P01\newline Cumple; RT & Se valorará el uso de cómputo sin servidor o de contenedores administrados para las cargas de perfil variable, con el análisis comparativo de costo frente a instancias permanentes.\newline \emph{T-12, Parte B, fila RT-\allowbreak{}03.09} \\
\textbf{RT-\allowbreak{}03.10}\newline Bases Técnicas Transversales, RT-03.10 & Perfil wms\_only con PostgreSQL 16.x y RabbitMQ 4.x locales por sitio; 24 h de autonomía\newline \textbf{EDT:} 3.3.6, 6.6.3 & /P01\newline Cumple; RT & El componente on-premise operará de forma autónoma y degradada ante la pérdida total del enlace con la nube, durante un período mínimo de 24 horas continuas o el mayor que fije el caso. Caso, cap. 15 — Operación desconectada del componente on-premise: Centro de distribución: mínimo 24 horas continuas de recepción, preparación y despacho sin enlace. Dispositivos de terreno: un turno completo de 14 horas sin cobertura, con capacidad de registrar la jornada íntegra de un repartidor de ruta rural.\newline \emph{T-12, Parte B, fila RT-\allowbreak{}03.10} \\
\textbf{RT-\allowbreak{}03.11}\newline Bases Técnicas Transversales, RT-03.11 & Registro local cifrado: Room/SQLite en terminales y PostgreSQL 16.x por sitio, con cola RabbitMQ 4 persistente\newline \textbf{EDT:} 3.3.6 & /P01\newline Cumple; RT & Durante la operación desconectada, la solución continuará registrando las transacciones operacionales críticas de forma local, con integridad garantizada y sin pérdida de datos.\newline \emph{T-12, Parte B, fila RT-\allowbreak{}03.11} \\
\textbf{RT-\allowbreak{}03.12}\newline Bases Técnicas Transversales, RT-03.12 & Reconciliación determinista con bitácora (consumidor PHP 8.5 y Amazon SQS FIFO de reconciliación)\newline \textbf{EDT:} 3.3.6 & /P01\newline Cumple; RT & Restablecido el enlace, la sincronización será automática, con reconciliación determinista de conflictos, regla de resolución documentada y bitácora auditable de las decisiones aplicadas.\newline \emph{T-12, Parte B, fila RT-\allowbreak{}03.12} \\
\textbf{RT-\allowbreak{}03.13}\newline Bases Técnicas Transversales, RT-03.13 & Declaración de las funciones no disponibles sin conexión y del procedimiento manual que las suple; Caso, cap. 15: sincronización del dispositivo de reparto en no más de 10 minutos tras un turno sin señal y del CD en no más de 2 horas tras un corte de 24 horas, con resolución determinista de los conflictos de stock\newline \textbf{EDT:} 3.3.6 & /P01\newline Cumple; RT & El PROPONENTE declarará qué funciones NO estarán disponibles en modo desconectado y qué procedimiento manual las suple. La ausencia de esta declaración se evaluará como observación grave. Caso, cap. 15 — Sincronización tras la reconexión: La sincronización de un dispositivo de reparto tras un turno completo sin señal no debe superar 10 minutos. La sincronización del centro de distribución tras un corte de 24 horas no debe superar 2 horas y debe resolver de forma determinista los conflictos de stock.\newline \emph{T-12, Parte B, fila RT-\allowbreak{}03.13} \\
\textbf{RT-\allowbreak{}03.14}\newline Bases Técnicas Transversales, RT-03.14 & Tres nodos HPE ProLiant DL345 Gen11 o Dell PowerEdge R7615 en Talca con Proxmox VE 9.2 y Ceph sin RAID por hardware, con tres réplicas; servidores HPE ProLiant DL20 Gen11 con RAID 10 en Concepción; pares Advantech ARK-1251 o equivalente con RAID 1 en cross-docking; NAS Synology RackStation RS2423RP+ con RAID 6.\newline \textbf{EDT:} 6.3.1, 6.3.3, 6.3.4 & /P01\newline Cumple; RT & Los equipos on-premise críticos serán redundantes. El almacenamiento local tolerará la falla de al menos un disco; el PROPONENTE declarará el nivel RAID escogido y lo justificará frente a las alternativas.\newline \emph{T-12, Parte B, fila RT-\allowbreak{}03.14} \\
\textbf{RT-\allowbreak{}03.15}\newline Bases Técnicas Transversales, RT-03.15 & Ansible (ansible-core 2.21) con líneas base CIS; ventana de aplicación acordada con la Contraparte Técnica, respetando congelamientos, con respaldo, verificación y reversión\newline \textbf{EDT:} 6.6.1, 8.2.5 & /P01\newline Cumple; RT & Los sistemas on-premise se endurecerán conforme a los CIS Benchmarks aplicables, con gestión centralizada de parches y ventana de aplicación acordada con el CLIENTE.\newline \emph{T-12, Parte B, fila RT-\allowbreak{}03.15} \\
\textbf{RT-\allowbreak{}03.16}\newline Bases Técnicas Transversales, RT-03.16 & AWS Distro for OpenTelemetry (colector 0.50) en los sitios hacia Amazon CloudWatch, con alertamiento unificado\newline \textbf{EDT:} 3.1.5 & /P01\newline Cumple; RT & El monitoreo del componente on-premise se integrará a la misma plataforma de observabilidad que la nube, con alertamiento unificado.\newline \emph{T-12, Parte B, fila RT-\allowbreak{}03.16} \\
\end{tablasimple}

\begin{tablasimple}{>{\hsize=0.82\hsize\linewidth=\hsize}Y >{\hsize=1.04\hsize\linewidth=\hsize}Y >{\hsize=0.65\hsize\linewidth=\hsize}Y >{\hsize=1.49\hsize\linewidth=\hsize}Y}
{\multicolumn{4}{l}{\textbf{Matriz individual: bloque 13 de 26.}} \\
\cab{Requisito y origen} & \cab{Componente y EDT} & \cab{Prueba prevista y estado} & \cab{Criterio y fuente}}
\textbf{RT-\allowbreak{}03.17}\newline Bases Técnicas Transversales, RT-03.17 & Talca y Concepción: fibra, LTE de otro proveedor y Starlink en espera caliente. Cada cross-docking: Starlink principal y dos operadores LTE. Tres caminos por sitio, con conmutación automática en menos de 30 s y mínimos de subida medidos antes de aceptar.\newline \textbf{EDT:} 6.6.2, 5.3.1, 5.3.2, 5.3.3 & /P01\newline Cumple; RT & El enlace entre el sitio on-premise y la nube será redundante, con caminos físicos y proveedores distintos, y conmutación automática con tiempo de conmutación declarado.\newline \emph{T-12, Parte B, fila RT-\allowbreak{}03.17} \\
\textbf{RT-\allowbreak{}03.18}\newline Bases Técnicas Transversales, RT-03.18 & Android Enterprise/Zebra DNA cubren inventario, configuración, actualización, bloqueo y borrado remoto de 382 terminales al inicio y 433 en el año 3. Greengrass V2 despliega los gateways. Registradores y sensores se inventarían y mantienen con herramientas de fabricante/TC58e. Sensores y balanzas no tienen sistema operativo administrable: se registran en el inventario y su firmware se actualiza en la mantención; como no guardan datos ni credenciales, su bloqueo es la baja de su registro desde la consola central y no requieren borrado.\newline \textbf{EDT:} 3.2.6, 3.4.5, 8.2.3 & /P01\newline Cumple; RT & Los dispositivos de borde y de terreno se administrarán de forma remota y centralizada: inventario, configuración, actualización de firmware y de aplicación, bloqueo y borrado remoto.\newline \emph{T-12, Parte B, fila RT-\allowbreak{}03.18} \\
\textbf{RT-\allowbreak{}03.19}\newline Bases Técnicas Transversales, RT-03.19 & AWS IoT Greengrass V2 en gateways Moxa UC-8200: lectura, validación y alerta local; lotes con compresión sin pérdida y alertas críticas inmediatas\newline \textbf{EDT:} 6.5.3, 3.2.3 & /P01\newline Cumple; RT & Se valorará el procesamiento en el borde de las cargas que lo admitan —filtrado, agregación previa, inferencia local— reduciendo el volumen transferido y la dependencia del enlace.\newline \emph{T-12, Parte B, fila RT-\allowbreak{}03.19} \\
\textbf{RT-\allowbreak{}03.20}\newline Bases Técnicas Transversales, RT-03.20 & Cálculo del ancho de banda por sitio en régimen y peak\newline \textbf{EDT:} 2.1.2 & /P01\newline Cumple; RT & El PROPONENTE dimensionará el ancho de banda requerido por sitio, en régimen normal y en peak, y lo justificará con el cálculo de volumen de transacciones y de datos.\newline \emph{T-12, Parte B, fila RT-\allowbreak{}03.20} \\
\textbf{RT-\allowbreak{}03.21}\newline Bases Técnicas Transversales, RT-03.21 & AWS Site-to-Site VPN IPsec IKEv2 con BGP sobre Transit Gateway\newline \textbf{EDT:} 6.6.2 & /P01\newline Cumple; RT & La conexión entre la red del CLIENTE y la nube se establecerá mediante enlace privado dedicado o red privada virtual con cifrado, según lo que el volumen y la criticidad justifiquen.\newline \emph{T-12, Parte B, fila RT-\allowbreak{}03.21} \\
\textbf{RT-\allowbreak{}03.22}\newline Bases Técnicas Transversales, RT-03.22 & AWS Verified Access con confianza OIDC de Keycloak 26 y postura de CrowdStrike Falcon (ADR-16)\newline \textbf{EDT:} 3.2.5, 3.3.1 & /P01\newline Cumple; RT & El acceso remoto de las personas trabajadoras del CLIENTE, incluido el trabajo desde el hogar, se resolverá con acceso a la red de confianza cero, con verificación de postura del dispositivo. No se admite exponer servicios internos directamente a Internet.\newline \emph{T-12, Parte B, fila RT-\allowbreak{}03.22} \\
\textbf{RT-\allowbreak{}03.23}\newline Bases Técnicas Transversales, RT-03.23 & Wi-Fi 6E industrial segmentado por VLAN, autenticación EAP-TLS por certificado contra el RADIUS local de cada sitio y cobertura verificada en sitio en el mes 2 y en la aceptación\newline \textbf{EDT:} 2.1.2, 6.2.3 & /P01\newline Cumple; RT & La red inalámbrica de los sitios operacionales, cuando el caso la requiera, contará con segmentación por tipo de dispositivo, autenticación por certificado o credencial de empresa y cobertura verificada mediante estudio de sitio.\newline \emph{T-12, Parte B, fila RT-\allowbreak{}03.23} \\
\textbf{RT-\allowbreak{}03.24}\newline Bases Técnicas Transversales, RT-03.24 & QoS de enlaces principal y de respaldo que prioriza recepción, preparación, despacho, sincronización, identidad, guías y alertas térmicas sobre tráfico administrativo; Caso, cap. 15: estudio de cobertura de los CD presentado, que incluye el interior de las cámaras de refrigerado y congelado, con al menos dos puntos de acceso Wi-Fi 6E industriales por cámara, aptos hasta $-$40 °C, y verificación en sitio en el mes 2 (S-36)\newline \textbf{EDT:} 6.6.2; 2.1.2, 6.2.3 & /P01\newline Cumple; RT & El PROPONENTE declarará la calidad de servicio y la priorización de tráfico aplicada a las transacciones operacionales críticas frente al tráfico administrativo. Caso, cap. 15 — Red inalámbrica de los sitios operacionales: Exigible un estudio de cobertura de los centros de distribución que incluya expresamente el interior de las cámaras de refrigerado y de congelado, hoy sin señal, y la propuesta técnica para resolverlo.\newline \emph{T-12, Parte B, fila RT-\allowbreak{}03.24} \\
\textbf{RT-\allowbreak{}04.01}\newline Bases Técnicas Transversales, RT-04.01 & Cinco ambientes (Desarrollo, QA, Preproducción, Producción y Recuperación) en cuentas AWS separadas, operativos en el hito H3 (mes 6)\newline \textbf{EDT:} 3.1.1, 3.1.2, 3.1.3 & /P01\newline Cumple; RT & Los cinco ambientes del numeral 4.1 estarán habilitados y operativos como condición del hito H3 del Formulario E-25.\newline \emph{T-12, Parte B, fila RT-\allowbreak{}04.01} \\
\textbf{RT-\allowbreak{}04.02}\newline Bases Técnicas Transversales, RT-04.02 & Preproducción equivalente en versiones, configuración, dimensionamiento y topología, con tres diferencias justificadas\newline \textbf{EDT:} 3.1.2 & /P01\newline Cumple; RT & Preproducción será equivalente a producción en topología, versiones de componentes y configuración. Las diferencias que subsistan por costo se declararán expresamente y se justificarán.\newline \emph{T-12, Parte B, fila RT-\allowbreak{}04.02} \\
\textbf{RT-\allowbreak{}04.03}\newline Bases Técnicas Transversales, RT-04.03 & GitLab con ramas protegidas, revisión obligatoria por pares y sin escritura directa en la rama principal\newline \textbf{EDT:} 3.1.4 & /P01\newline Cumple; RT & El código residirá en un sistema de control de versiones con ramas protegidas, revisión obligatoria por pares y prohibición de escritura directa sobre la rama principal.\newline \emph{T-12, Parte B, fila RT-\allowbreak{}04.03} \\
\textbf{RT-\allowbreak{}04.04}\newline Bases Técnicas Transversales, RT-04.04 & Trazabilidad requerimiento–incidencia–cambio–prueba–despliegue en GitLab y el pipeline\newline \textbf{EDT:} 1.2.4, 3.1.4 & /P01\newline Cumple; RT & Existirá trazabilidad completa entre requerimiento, incidencia, cambio de código, prueba ejecutada y despliegue realizado.\newline \emph{T-12, Parte B, fila RT-\allowbreak{}04.04} \\
\textbf{RT-\allowbreak{}04.05}\newline Bases Técnicas Transversales, RT-04.05 & GitLab CI + AWS CodeBuild: PHP 8.5Unit, PHPStan/Larastan, Pint, composer audit, escaneo de secretos e imágenes con bloqueo automático\newline \textbf{EDT:} 3.1.4 & /P01\newline Cumple; RT & El flujo de integración continua ejecutará, como mínimo: compilación, pruebas unitarias, análisis estático de código, análisis de composición de software, escaneo de secretos y escaneo de imágenes de contenedor, con criterios de bloqueo automático del despliegue.\newline \emph{T-12, Parte B, fila RT-\allowbreak{}04.05} \\
\textbf{RT-\allowbreak{}04.06}\newline Bases Técnicas Transversales, RT-04.06 & Despliegue automatizado por ECS Fargate con reversión automática\newline \textbf{EDT:} 3.1.4 & /P01\newline Cumple; RT & Los despliegues serán automatizados y reproducibles, con reversión automatizada y sin intervención manual en el paso a producción.\newline \emph{T-12, Parte B, fila RT-\allowbreak{}04.06} \\
\textbf{RT-\allowbreak{}04.07}\newline Bases Técnicas Transversales, RT-04.07 & Azul-verde con canario en ECS Fargate, demostrado en Preproducción\newline \textbf{EDT:} 3.8.8, 3.9.7 & /P01\newline Cumple; RT & La estrategia de despliegue permitirá liberar sin interrupción del servicio: azul-verde, canario o despliegue progresivo. Se declarará cuál se emplea y se demostrará en Preproducción antes de cada paso a producción.\newline \emph{T-12, Parte B, fila RT-\allowbreak{}04.07} \\
\textbf{RT-\allowbreak{}04.08}\newline Bases Técnicas Transversales, RT-04.08 & Configuración en SSM Parameter Store por ambiente; una sola imagen firmada promovida sin recompilar\newline \textbf{EDT:} 3.1.4 & /P01\newline Cumple; RT & Toda configuración estará externalizada del artefacto y gestionada por ambiente. Un mismo artefacto deberá poder promoverse de QA a Preproducción y a Producción sin recompilación.\newline \emph{T-12, Parte B, fila RT-\allowbreak{}04.08} \\
\textbf{RT-\allowbreak{}04.09}\newline Bases Técnicas Transversales, RT-04.09 & AWS Secrets Manager con rotación automática y auditoría (ADR-15)\newline \textbf{EDT:} 3.2.5 & /P01\newline Cumple; RT & Los secretos residirán en un gestor de secretos con rotación automática y auditoría de acceso. Queda prohibida toda credencial embebida en código, imágenes o archivos de configuración.\newline \emph{T-12, Parte B, fila RT-\allowbreak{}04.09} \\
\textbf{RT-\allowbreak{}04.10}\newline Bases Técnicas Transversales, RT-04.10 & Migraciones Laravel 13 versionadas con estrategia expandir-contraer y reversión declarada\newline \textbf{EDT:} 3.1.4 & /P01\newline Cumple; RT & Las migraciones de esquema de base de datos serán versionadas, reversibles y ejecutadas de forma automatizada, con estrategia de compatibilidad que permita convivir dos versiones de la aplicación durante el despliegue.\newline \emph{T-12, Parte B, fila RT-\allowbreak{}04.10} \\
\textbf{RT-\allowbreak{}04.11}\newline Bases Técnicas Transversales, RT-04.11 & Umbral bloqueante de 70 \% de cobertura de la lógica de negocio en GitLab CI\newline \textbf{EDT:} 1.5.2 & /P01\newline Cumple; RT & La cobertura de pruebas automatizadas del código de lógica de negocio será de al menos 70 \%, con umbral bloqueante en el flujo de integración continua.\newline \emph{T-12, Parte B, fila RT-\allowbreak{}04.11} \\
\textbf{RT-\allowbreak{}04.12}\newline Bases Técnicas Transversales, RT-04.12 & Cadencia quincenal; meta ordinaria $\leq$ 5 días hábiles habilitados desde código a producción, con medición adicional del tiempo calendario y de esperas por congelamiento; fallos $\leq$ 5 \% mensual; restauración crítica $\leq$ 4 h.\newline \textbf{EDT:} 3.1.4 & /P01\newline Cumple; RT & El PROPONENTE declarará su frecuencia de despliegue objetivo, su tiempo desde el compromiso de código hasta producción, su tasa de cambios fallidos y su tiempo de restauración, y los medirá durante la Operación.\newline \emph{T-12, Parte B, fila RT-\allowbreak{}04.12} \\
\textbf{RT-\allowbreak{}04.13}\newline Bases Técnicas Transversales, RT-04.13 & Reducción o apagado de Desarrollo, QA y Preproducción; meta de ahorro $\geq$60 \% de horas de cómputo según S-43; valor monetario en la Oferta Económica (Bases Administrativas, art. 50.2, p. 29).\newline \textbf{EDT:} 3.1.1, 8.4.2 & /P01\newline Cumple; RT & Los ambientes no productivos se apagarán o reducirán fuera del horario de uso, con el ahorro reflejado en la estructura de costos.\newline \emph{T-12, Parte B, fila RT-\allowbreak{}04.13} \\
\textbf{RT-\allowbreak{}04.14}\newline Bases Técnicas Transversales, RT-04.14 & INN-02 crea automáticamente con Terraform 1.15 desde GitLab CI los recursos efímeros por incidencia y los destruye al finalizar, incluso ante fallo, dentro de QA o Preproducción y conforme al horario y excepciones de S-43.\newline \textbf{EDT:} 3.10.2.2 & /P01\newline Cumple; RT & Se valorará la existencia de ambientes efímeros por rama o por incidencia, creados y destruidos automáticamente.\newline \emph{T-12, Parte B, fila RT-\allowbreak{}04.14} \\
\textbf{RT-\allowbreak{}05.01}\newline Bases Técnicas Transversales, RT-05.01 & Modelo de datos por dominios y diccionario 5-A con 121 entidades y 995 atributos: nombre, significado, tipo, dominio, obligatoriedad, referencia, propietario funcional y sensibilidad/política por atributo.\newline \textbf{EDT:} 2.1.3 & /P01\newline Cumple; RT & El PROPONENTE entregará el modelo de datos documentado y un diccionario de datos con el nombre, el tipo, el dominio de valores, la obligatoriedad, el propietario y la sensibilidad de cada atributo.\newline \emph{T-12, Parte B, fila RT-\allowbreak{}05.01} \\
\textbf{RT-\allowbreak{}05.02}\newline Bases Técnicas Transversales, RT-05.02 & PostgreSQL 16.x / PostGIS 3.x por sitio, Aurora PostgreSQL, Amazon DynamoDB y Amazon S3, con posición CAP por dominio\newline \textbf{EDT:} 2.1.3 & /P01\newline Cumple; RT & El PROPONENTE justificará la selección del paradigma y del motor de persistencia: relacional o no relacional, garantías transaccionales, y la posición escogida entre consistencia y disponibilidad conforme al teorema CAP, para cada dominio de datos.\newline \emph{T-12, Parte B, fila RT-\allowbreak{}05.02} \\
\textbf{RT-\allowbreak{}05.03}\newline Bases Técnicas Transversales, RT-05.03 & Contrato de escritura con auditoría y outbox en la misma transacción; tabla gob\_auditoria con actor, dispositivo, instante y valores anterior y posterior\newline \textbf{EDT:} 3.3.5 & /P01\newline Cumple; RT & Toda operación de negocio será trazable: la solución permitirá reconstruir quién, qué, cuándo, desde qué dispositivo y con qué valores anteriores y posteriores, para cualquier registro y en cualquier momento del período de retención.\newline \emph{T-12, Parte B, fila RT-\allowbreak{}05.03} \\
\end{tablasimple}

\begin{tablasimple}{>{\hsize=0.82\hsize\linewidth=\hsize}Y >{\hsize=1.04\hsize\linewidth=\hsize}Y >{\hsize=0.65\hsize\linewidth=\hsize}Y >{\hsize=1.49\hsize\linewidth=\hsize}Y}
{\multicolumn{4}{l}{\textbf{Matriz individual: bloque 14 de 26.}} \\
\cab{Requisito y origen} & \cab{Componente y EDT} & \cab{Prueba prevista y estado} & \cab{Criterio y fuente}}
\textbf{RT-\allowbreak{}05.04}\newline Bases Técnicas Transversales, RT-05.04 & Gestión de calidad de datos con dimensiones ISO/IEC 25012, reglas en captura y tablero en Amazon QuickSight\newline \textbf{EDT:} 3.7.1 & /P01\newline Cumple; RT & La calidad de datos se gestionará conforme a ISO/IEC 25012, con validación en el punto de captura, indicadores de completitud, exactitud y consistencia, y tablero de calidad disponible para el CLIENTE.\newline \emph{T-12, Parte B, fila RT-\allowbreak{}05.04} \\
\textbf{RT-\allowbreak{}05.05}\newline Bases Técnicas Transversales, RT-05.05 & Separación OLTP (Aurora PostgreSQL, PostgreSQL 16.x) y analítica (Amazon S3, AWS Glue, Amazon Redshift Serverless, Amazon QuickSight)\newline \textbf{EDT:} 3.2.3 & /P01\newline Cumple; RT & El almacenamiento transaccional y el analítico estarán separados. Ninguna consulta analítica podrá degradar el desempeño de la operación.\newline \emph{T-12, Parte B, fila RT-\allowbreak{}05.05} \\
\textbf{RT-\allowbreak{}05.06}\newline Bases Técnicas Transversales, RT-05.06 & Exportación íntegra autónoma en formatos abiertos con manifiesto, huella y conteo\newline \textbf{EDT:} 9.1.1, 9.2.1 & /P01\newline Cumple; RT & La solución permitirá exportar la totalidad de la información del CLIENTE en formatos abiertos y documentados, en cualquier momento del Contrato, sin costo adicional y sin intervención del ADJUDICATARIO.\newline \emph{T-12, Parte B, fila RT-\allowbreak{}05.06} \\
\textbf{RT-\allowbreak{}05.07}\newline Bases Técnicas Transversales, RT-05.07 & Política de retención, archivado y eliminación segura por dominio y atributo; AWS Backup y S3 Object Lock\newline \textbf{EDT:} 2.1.3, 9.2.2 & /P01\newline Cumple; RT & El PROPONENTE declarará la política de retención, archivado y eliminación por cada dominio de datos, coherente con la normativa aplicable al caso, e implementará un procedimiento verificable de eliminación segura.\newline \emph{T-12, Parte B, fila RT-\allowbreak{}05.07} \\
\textbf{RT-\allowbreak{}05.08}\newline Bases Técnicas Transversales, RT-05.08 & Clasificación AL/CR/ME por atributo con cifrado a nivel de campo mediante AWS KMS y consulta auditada.\newline \textbf{EDT:} 1.9.2, 1.9.3 & /P01\newline Cumple; RT & Los datos personales se tratarán conforme al Artículo 85° de las Bases Administrativas, con seudonimización o cifrado a nivel de campo para las categorías sensibles que el caso identifique.\newline \emph{T-12, Parte B, fila RT-\allowbreak{}05.08} \\
\textbf{RT-\allowbreak{}05.09}\newline Bases Técnicas Transversales, RT-05.09 & Gestión de datos maestros con dueño, aprobador, vigencia y equivalencias ERP–WMS–EDI\newline \textbf{EDT:} 3.7.2 & /P01\newline Cumple; RT & El PROPONENTE presentará una estrategia de gestión de datos maestros que evite la duplicación de entidades compartidas entre módulos y con sistemas externos.\newline \emph{T-12, Parte B, fila RT-\allowbreak{}05.09} \\
\textbf{RT-\allowbreak{}05.10}\newline Bases Técnicas Transversales, RT-05.10 & Política por dominio: documentos tributarios y respaldo, 6 años; registros sanitarios y trazabilidad, vida útil del producto más 6 meses con mínimo de 5 años; temperatura, 5 años; evidencia de entrega, 6 años; geolocalización de personas, 12 meses.\newline \textbf{EDT:} 2.1.3 & /P01\newline Parcial; RT & Se valorará la implementación de un catálogo de datos con linaje automatizado, que permita rastrear el origen de cada indicador de negocio hasta su fuente. Caso, cap. 15 — Retención de datos históricos y de auditoría: Documentos tributarios y su respaldo: 6 años. Registros sanitarios y de trazabilidad de lote: vida útil del producto más 6 meses, con mínimo de 5 años. Registros de temperatura: 5 años. Evidencia de entrega: 6 años. Datos de geolocalización de personas: 12 meses.\newline \emph{T-12, Parte B, fila RT-\allowbreak{}05.10} \\
\textbf{RT-\allowbreak{}05.11}\newline Bases Técnicas Transversales, RT-05.11 & Plan de migración de 32,11 GB con alcance, origen, reglas, calidad, olas y reversión\newline \textbf{EDT:} 3.7.1, 3.7.2, 3.7.3, 3.7.4, 3.7.5 & /P01\newline Cumple; RT & El PROPONENTE presentará un plan de migración con alcance, origen, volumen, reglas de transformación, criterios de calidad, estrategia de ejecución y plan de reversión.\newline \emph{T-12, Parte B, fila RT-\allowbreak{}05.11} \\
\textbf{RT-\allowbreak{}05.12}\newline Bases Técnicas Transversales, RT-05.12 & Perfilado y saneamiento con informe de defectos y decisión por causa\newline \textbf{EDT:} 3.7.1 & /P01\newline Cumple; RT & La migración incluirá una etapa de perfilado y saneamiento previo, con informe de los defectos detectados en los datos de origen y la decisión adoptada sobre cada uno.\newline \emph{T-12, Parte B, fila RT-\allowbreak{}05.12} \\
\textbf{RT-\allowbreak{}05.13}\newline Bases Técnicas Transversales, RT-05.13 & Dos ensayos completos en Preproducción con medición de tiempo y conciliación\newline \textbf{EDT:} 3.7.5 & /P01\newline Cumple; RT & Se ejecutarán al menos dos ensayos completos de migración sobre Preproducción antes de la migración definitiva, con medición del tiempo total y del resultado de la conciliación.\newline \emph{T-12, Parte B, fila RT-\allowbreak{}05.13} \\
\textbf{RT-\allowbreak{}05.14}\newline Bases Técnicas Transversales, RT-05.14 & Conciliación por recuentos, sumas de control y muestreo dirigido, sin diferencias no explicadas\newline \textbf{EDT:} 3.7.5 & /P01\newline Cumple; RT & La conciliación posterior a la migración será cuantitativa y verificable: recuentos, sumas de control y muestreo dirigido. Toda diferencia deberá quedar explicada.\newline \emph{T-12, Parte B, fila RT-\allowbreak{}05.14} \\
\textbf{RT-\allowbreak{}05.15}\newline Bases Técnicas Transversales, RT-05.15 & Alcance del Caso: maestros de clientes, productos y proveedores completos, ventas y pedidos de tres años, inventario de dos años, trazabilidad sanitaria de cinco años y cartera viva con dos años de historia; lo no migrado queda en un repositorio de consulta durante el período de retención aplicable.\newline \textbf{EDT:} 3.7.3 & /P01\newline Cumple; RT & Los datos históricos que no se migren quedarán accesibles en un repositorio de consulta durante el período de retención que fije el caso. Caso, cap. 15 — Datos históricos a migrar: Maestros de clientes, productos y proveedores: completos. Ventas y pedidos: 3 años. Inventario y movimientos: 2 años. Trazabilidad sanitaria: 5 años. Cuentas por cobrar: saldos vivos más 2 años.\newline \emph{T-12, Parte B, fila RT-\allowbreak{}05.15} \\
\textbf{RT-\allowbreak{}05.16}\newline Bases Técnicas Transversales, RT-05.16 & Contratos OpenAPI 3.1 (swagger-php) y AsyncAPI 2.6 generados desde el código\newline \textbf{EDT:} 3.11.3 & /P01\newline Cumple; RT & Los servicios síncronos se documentarán en OpenAPI 3.1 y los flujos dirigidos por eventos en AsyncAPI 2.6 o superior. La documentación se generará desde el código y se mantendrá actualizada automáticamente.\newline \emph{T-12, Parte B, fila RT-\allowbreak{}05.16} \\
\textbf{RT-\allowbreak{}05.17}\newline Bases Técnicas Transversales, RT-05.17 & Versionado semántico de contratos con preaviso mínimo de seis meses\newline \textbf{EDT:} 3.11.3 & /P01\newline Cumple; RT & Los contratos de interfaz se versionarán semánticamente, con compatibilidad hacia atrás y política de obsolescencia con preaviso mínimo de seis meses.\newline \emph{T-12, Parte B, fila RT-\allowbreak{}05.17} \\
\textbf{RT-\allowbreak{}05.18}\newline Bases Técnicas Transversales, RT-05.18 & Autenticación mutua TLS entre sistemas con credenciales rotadas y sin claves estáticas en la URL; Keycloak 26.x.\newline \textbf{EDT:} 3.3.1 & /P01\newline Cumple; RT & La autenticación entre sistemas empleará OAuth 2.1 con credenciales de cliente o autenticación mutua TLS. Queda prohibida la autenticación por clave estática en la ruta de la dirección web.\newline \emph{T-12, Parte B, fila RT-\allowbreak{}05.18} \\
\textbf{RT-\allowbreak{}05.19}\newline Bases Técnicas Transversales, RT-05.19 & transaction\_id y correlacion\_id en auditoría, eventos y mensajes; trazas OpenTelemetry (colector ADOT 0.50)\newline \textbf{EDT:} 3.3.5 & /P01\newline Cumple; RT & Toda integración registrará la transacción de entrada y de salida, con identificador de correlación común que permita seguir una operación de negocio a través de todos los sistemas involucrados.\newline \emph{T-12, Parte B, fila RT-\allowbreak{}05.19} \\
\textbf{RT-\allowbreak{}05.20}\newline Bases Técnicas Transversales, RT-05.20 & Capa anticorrupción del ERP 2017 (A-04, erp-sync) y adaptadores EDI, DTE, GIS y flota\newline \textbf{EDT:} 3.3.2 & /P01\newline Cumple; RT & Las integraciones con sistemas heredados o de terceros se aislarán mediante una capa anticorrupción, de modo que un cambio en el sistema externo no propague su modelo al núcleo de la solución.\newline \emph{T-12, Parte B, fila RT-\allowbreak{}05.20} \\
\textbf{RT-\allowbreak{}05.21}\newline Bases Técnicas Transversales, RT-05.21 & Catálogo INT-01 a INT-15 con modo, volumen, ventana de la contraparte y comportamiento ante falla\newline \textbf{EDT:} 2.1.1 & /P01\newline Cumple; RT & El PROPONENTE declarará, por cada integración, el modo (síncrono o asíncrono), el volumen esperado, la ventana de disponibilidad del sistema contraparte y el comportamiento de la solución cuando ese sistema no responde.\newline \emph{T-12, Parte B, fila RT-\allowbreak{}05.21} \\
\textbf{RT-\allowbreak{}05.22}\newline Bases Técnicas Transversales, RT-05.22 & Carga y descarga masiva con validación previa, error por registro y procesamiento parcial\newline \textbf{EDT:} 3.7.2 & /P01\newline Cumple; RT & La solución soportará la carga y descarga masiva de información en formatos abiertos, con validación previa, informe de errores por registro y procesamiento parcial.\newline \emph{T-12, Parte B, fila RT-\allowbreak{}05.22} \\
\textbf{RT-\allowbreak{}05.23}\newline Bases Técnicas Transversales, RT-05.23 & GS1 (GTIN, SSCC/GS1-128, GLN), EPCIS 2.0/CBV 2.0, EDI por cadena (OpenAS2 4.9) y formato DTE del SII\newline \textbf{EDT:} 3.6.1, 3.6.5 & /P01\newline Cumple; RT & Se emplearán los estándares sectoriales de intercambio que las Bases Técnicas del caso identifiquen. Caso, cap. 15 — Estándares sectoriales de intercambio: Estándares GS1 para identificación de productos, unidades logísticas y ubicaciones; mensajería electrónica entre socios comerciales en el formato que exija cada cadena; estándar de trazabilidad de eventos para la cadena de custodia; y el formato de la autoridad tributaria para los documentos electrónicos.\newline \emph{T-12, Parte B, fila RT-\allowbreak{}05.23} \\
\textbf{RT-\allowbreak{}05.24}\newline Bases Técnicas Transversales, RT-05.24 & No se oferta.\newline \textbf{EDT:} No aplica & 0; brecha visible\newline No cumple; RT & Se valorará la publicación de un portal de servicios para desarrolladores con documentación navegable, ambiente de pruebas y credenciales de prueba autoservidas.\newline \emph{T-12, Parte B, fila RT-\allowbreak{}05.24} \\
\textbf{RT-\allowbreak{}05.25}\newline Bases Técnicas Transversales, RT-05.25 & Capa analítica con tableros en Amazon QuickSight sobre los indicadores del caso\newline \textbf{EDT:} 3.4.11, 3.2.3 & /P01\newline Cumple; RT & La solución proveerá una capa analítica con tableros operacionales y de gestión, construidos sobre los indicadores que las Bases Técnicas del caso definan.\newline \emph{T-12, Parte B, fila RT-\allowbreak{}05.25} \\
\textbf{RT-\allowbreak{}05.26}\newline Bases Técnicas Transversales, RT-05.26 & Filtros por período, unidad y dimensiones, con profundización hasta pedido, entrega, línea y documento\newline \textbf{EDT:} 3.4.11, 3.2.3 & /P01\newline Cumple; RT & Los tableros permitirán filtrar por período, unidad organizacional y dimensiones propias del caso, y profundizar desde el indicador agregado hasta la transacción de origen.\newline \emph{T-12, Parte B, fila RT-\allowbreak{}05.26} \\
\textbf{RT-\allowbreak{}05.27}\newline Bases Técnicas Transversales, RT-05.27 & Autoría en Amazon QuickSight con modelo semántico documentado\newline \textbf{EDT:} 3.2.3 & /P01\newline Cumple; RT & El CLIENTE podrá construir sus propios informes sin intervención del ADJUDICATARIO, mediante una herramienta de autoservicio con modelo semántico documentado.\newline \emph{T-12, Parte B, fila RT-\allowbreak{}05.27} \\
\textbf{RT-\allowbreak{}05.28}\newline Bases Técnicas Transversales, RT-05.28 & Exportación en formatos abiertos y envío programado de informes\newline \textbf{EDT:} 3.4.11 & /P01\newline Cumple; RT & Todo informe será exportable en formatos abiertos y programable para envío automático por calendario.\newline \emph{T-12, Parte B, fila RT-\allowbreak{}05.28} \\
\end{tablasimple}

\begin{tablasimple}{>{\hsize=0.82\hsize\linewidth=\hsize}Y >{\hsize=1.04\hsize\linewidth=\hsize}Y >{\hsize=0.65\hsize\linewidth=\hsize}Y >{\hsize=1.49\hsize\linewidth=\hsize}Y}
{\multicolumn{4}{l}{\textbf{Matriz individual: bloque 15 de 26.}} \\
\cab{Requisito y origen} & \cab{Componente y EDT} & \cab{Prueba prevista y estado} & \cab{Criterio y fuente}}
\textbf{RT-\allowbreak{}05.29}\newline Bases Técnicas Transversales, RT-05.29 & Latencia $\leq$ 5 min (operación del día), $\leq$ 2 h (cierre comercial) y $\leq$ 4 h (gestión); curva térmica por red móvil cada 1 min e Iridium SBD cada 3 min, indicadores en no más de 5 min; sin ambos caminos se muestra antigüedad de la última curva.\newline \textbf{EDT:} 3.4.4, 3.4.5, 3.4.11, 3.2.3 & /P01\newline Cumple; RT & La latencia máxima entre la ocurrencia de una transacción y su disponibilidad en la capa analítica será la que fije el caso y, en su defecto, no superará las 4 horas. Caso, cap. 15 — Latencia de la capa analítica: Indicadores de la operación del día —entregas cumplidas, faltantes, devoluciones, temperatura— no superior a 5 minutos. Cierre comercial del día no superior a 2 horas tras el retorno del último camión. Indicadores de gestión no superior a 4 horas.\newline \emph{T-12, Parte B, fila RT-\allowbreak{}05.29} \\
\textbf{RT-\allowbreak{}05.30}\newline Bases Técnicas Transversales, RT-05.30 & La solución estima la vida útil remanente mediante un modelo cinético determinista por familia, aplicado a la historia térmica del lote. El SD13 documenta el modelo, sus variables, la métrica I-03B y la recalibración semestral o extraordinaria. INN-03 desarrolla su incorporación distribuida al proceso operacional; no atribuye novedad al modelo cinético ni incorpora modelos aprendidos.\newline \textbf{EDT:} 3.10.3.1, 3.10.3.2, 3.10.3.3, 8.3.2 & /P01\newline Cumple; RT & Se valorará la incorporación de analítica predictiva pertinente al proceso del caso, con el modelo, sus variables, su métrica de desempeño y su plan de reentrenamiento documentados.\newline \emph{T-12, Parte B, fila RT-\allowbreak{}05.30} \\
\textbf{RT-\allowbreak{}06.01}\newline Bases Técnicas Transversales, RT-06.01 & Recinto técnico exclusivo en el CD Talca con acceso propio; gabinetes de borde en Concepción y en los tres cross-docking (tipología del Caso, Cap. 15)\newline \textbf{EDT:} 2.3.1, 6.1.1 & /P01\newline Cumple; RT & El espacio asignado será de uso exclusivo de la solución y estará aislado de otras dependencias del CLIENTE, con acceso independiente. Caso, cap. 15 — Tipología del emplazamiento on-premise: Sala técnica secundaria en el centro de distribución de Talca, dimensionada para sostener recepción, preparación y despacho durante un corte. La sala actual de 25 m² no cumple el Capítulo 6 del documento transversal. Gabinetes de borde en el centro de distribución de Concepción y en las tres plataformas de cross-docking.\newline \emph{T-12, Parte B, fila RT-\allowbreak{}06.01} \\
\textbf{RT-\allowbreak{}06.02}\newline Bases Técnicas Transversales, RT-06.02 & Blindaje perimetral de plancha de acero de 3 mm sobre albañilería reforzada, resistencia al fuego F-120\newline \textbf{EDT:} 6.1.1 & /P01\newline Cumple; RT & Los muros no estructurales del recinto contarán con blindaje perimetral; el PROPONENTE especificará el material y la resistencia.\newline \emph{T-12, Parte B, fila RT-\allowbreak{}06.02} \\
\textbf{RT-\allowbreak{}06.03}\newline Bases Técnicas Transversales, RT-06.03 & Plano en planta acotado y a escala del recinto de Talca, con cinco recintos de superficie declarada ($\approx$ 31 m\textsuperscript{2} útiles) y zonas separadas de generador, baterías, climatización, servidores, comunicaciones, trabajo y respaldo\newline \textbf{EDT:} 2.3.1, 6.1.1 & /P01\newline Cumple; RT & El PROPONENTE entregará el plano de distribución interna del recinto, con la separación de las zonas de generadores, baterías, climatización, servidores, comunicaciones, trabajo y respaldo.\newline \emph{T-12, Parte B, fila RT-\allowbreak{}06.03} \\
\textbf{RT-\allowbreak{}06.04}\newline Bases Técnicas Transversales, RT-06.04 & Piso técnico, canalización, cableado Cat6A y fibra OM4 con certificación de cada enlace, conforme a ANSI/TIA-942-B, ANSI/TIA-569-E, ANSI/TIA-568.2-D, ANSI/TIA-568.3-D y ANSI/TIA-606-C para el etiquetado y la documentación.\newline \textbf{EDT:} 6.1.1, 6.2.1, 6.2.2 & /P01\newline Cumple; RT & El piso técnico, la canalización, el cableado estructurado y el etiquetado se ejecutarán conforme a norma, con documentación de la certificación de cada enlace.\newline \emph{T-12, Parte B, fila RT-\allowbreak{}06.04} \\
\textbf{RT-\allowbreak{}06.05}\newline Bases Técnicas Transversales, RT-06.05 & Racks R01 (servidores, 15U/42U, 36 \%) y R02 (comunicaciones, 13U/42U, 31 \%) independientes, con margen de crecimiento\newline \textbf{EDT:} 2.3.2, 6.3.1, 6.3.2 & /P01\newline Cumple; RT & Los racks de servidores serán independientes de los racks de equipos de comunicación. Se declarará la ocupación proyectada de cada rack y su margen de crecimiento.\newline \emph{T-12, Parte B, fila RT-\allowbreak{}06.05} \\
\textbf{RT-\allowbreak{}06.06}\newline Bases Técnicas Transversales, RT-06.06 & Especificación y coordinación de la obra civil de separación por LafroX; ejecución a cargo del CLIENTE\newline \textbf{EDT:} 6.1.1 & /P01\newline Cumple; RT & La obra civil de separación de las instalaciones es de cargo del CLIENTE; su especificación técnica y su coordinación son de cargo del PROPONENTE.\newline \emph{T-12, Parte B, fila RT-\allowbreak{}06.06} \\
\textbf{RT-\allowbreak{}06.07}\newline Bases Técnicas Transversales, RT-06.07 & UPS modular de 15 kVA, doble conversión en línea, N+1, autonomía $\geq$ 30 min a plena carga\newline \textbf{EDT:} 6.1.2 & /P01\newline Cumple; RT & El suministro eléctrico de los equipos será ininterrumpido, con sistema de alimentación ininterrumpida dimensionado para una autonomía mínima de 30 minutos a plena carga.\newline \emph{T-12, Parte B, fila RT-\allowbreak{}06.07} \\
\textbf{RT-\allowbreak{}06.08}\newline Bases Técnicas Transversales, RT-06.08 & Grupo electrógeno de 25 kVA con estanque para 24 h y contrato de reabastecimiento\newline \textbf{EDT:} 6.1.2 & /P01\newline Cumple; RT & La capacidad de generación autónoma asegurará un rango mínimo de 24 horas continuas de operación, con estanque de combustible dimensionado y contrato de reabastecimiento declarado.\newline \emph{T-12, Parte B, fila RT-\allowbreak{}06.08} \\
\textbf{RT-\allowbreak{}06.09}\newline Bases Técnicas Transversales, RT-06.09 & Instalación eléctrica independiente conforme a la normativa chilena y NCh Elec. 2777\newline \textbf{EDT:} 6.1.2 & /P01\newline Cumple; RT & La instalación eléctrica del recinto será independiente de la del resto del edificio y cumplirá la normativa eléctrica chilena vigente, incluida la NCh Elec. 2777 sobre sistemas de puesta a tierra.\newline \emph{T-12, Parte B, fila RT-\allowbreak{}06.09} \\
\textbf{RT-\allowbreak{}06.10}\newline Bases Técnicas Transversales, RT-06.10 & Revisión y medición semestral con informe al CLIENTE\newline \textbf{EDT:} 8.2.6 & /P01\newline Cumple; RT & Se efectuará revisión y medición semestral de las instalaciones eléctricas del recinto, con informe entregable al CLIENTE.\newline \emph{T-12, Parte B, fila RT-\allowbreak{}06.10} \\
\textbf{RT-\allowbreak{}06.11}\newline Bases Técnicas Transversales, RT-06.11 & Carga TI de diseño 9,2 kW, factor de potencia 0,95, PUE estimado $\approx$ 1,5\newline \textbf{EDT:} 2.3.1, 6.1.1 & /P01\newline Cumple; RT & El PROPONENTE declarará la carga eléctrica proyectada en kW, el factor de potencia y la eficiencia en el uso de la energía (PUE) estimada del recinto.\newline \emph{T-12, Parte B, fila RT-\allowbreak{}06.11} \\
\textbf{RT-\allowbreak{}06.12}\newline Bases Técnicas Transversales, RT-06.12 & UPS N+1 y transferencia automática red-generador.\newline \textbf{EDT:} 6.1.2 & /P01\newline Parcial; RT & Se valorará la redundancia de alimentación en configuración 2N o N+1 con doble acometida y transferencia automática.\newline \emph{T-12, Parte B, fila RT-\allowbreak{}06.12} \\
\textbf{RT-\allowbreak{}06.13}\newline Bases Técnicas Transversales, RT-06.13 & Dos unidades de climatización de precisión de 10 kW en N+1\newline \textbf{EDT:} 6.1.3 & /P01\newline Cumple; RT & El recinto contará con climatización de precisión para operación continua, redundante en configuración N+1, con control de temperatura y de humedad relativa dentro de los rangos que recomienda el fabricante del equipamiento.\newline \emph{T-12, Parte B, fila RT-\allowbreak{}06.13} \\
\textbf{RT-\allowbreak{}06.14}\newline Bases Técnicas Transversales, RT-06.14 & Sensores de temperatura, humedad y agua con alertas a Amazon CloudWatch\newline \textbf{EDT:} 6.1.3, 3.1.5 & /P01\newline Cumple; RT & Se monitorearán en línea la temperatura, la humedad y la presencia de agua, con alertamiento integrado a la plataforma de observabilidad.\newline \emph{T-12, Parte B, fila RT-\allowbreak{}06.14} \\
\textbf{RT-\allowbreak{}06.15}\newline Bases Técnicas Transversales, RT-06.15 & Dos racks con paneles ciegos en unidades libres, pasos de cable sellados y orientación frente-frío / trasera-caliente; menor recirculación y mezcla mejoran la eficiencia de climatización y el PUE estimado.\newline \textbf{EDT:} 2.3.2 & /P01\newline Cumple; RT & El PROPONENTE declarará la estrategia de contención de pasillo frío o caliente y su efecto en la eficiencia energética.\newline \emph{T-12, Parte B, fila RT-\allowbreak{}06.15} \\
\textbf{RT-\allowbreak{}06.16}\newline Bases Técnicas Transversales, RT-06.16 & Detección temprana por aspiración láser tipo AnaLASER\newline \textbf{EDT:} 6.1.4 & /P01\newline Cumple; RT & El recinto contará con detección temprana por aspiración de aire con tecnología láser, tipo AnaLASER o equivalente de prestaciones iguales o superiores.\newline \emph{T-12, Parte B, fila RT-\allowbreak{}06.16} \\
\textbf{RT-\allowbreak{}06.17}\newline Bases Técnicas Transversales, RT-06.17 & Extinción automática con agente limpio tipo FM-200, aprobación UL e instalación NFPA\newline \textbf{EDT:} 6.1.4 & /P01\newline Cumple; RT & La extinción será automática mediante agente limpio tipo FM-200 o equivalente, con aprobación UL e instalación conforme a norma NFPA.\newline \emph{T-12, Parte B, fila RT-\allowbreak{}06.17} \\
\textbf{RT-\allowbreak{}06.18}\newline Bases Técnicas Transversales, RT-06.18 & Extintores portátiles con mantención y certificación vigentes\newline \textbf{EDT:} 6.1.4 & /P01\newline Cumple; RT & Se proveerá un sistema secundario de extintores portátiles habilitados, con mantención y certificación vigentes.\newline \emph{T-12, Parte B, fila RT-\allowbreak{}06.18} \\
\textbf{RT-\allowbreak{}06.19}\newline Bases Técnicas Transversales, RT-06.19 & Detección y extinción integradas al monitoreo, con aviso al NOC y al CLIENTE\newline \textbf{EDT:} 6.1.4 & /P01\newline Cumple; RT & El sistema de detección y extinción se integrará al monitoreo en línea y notificará al NOC y a la contraparte del CLIENTE.\newline \emph{T-12, Parte B, fila RT-\allowbreak{}06.19} \\
\textbf{RT-\allowbreak{}06.20}\newline Bases Técnicas Transversales, RT-06.20 & Biometría facial con AFIS de respaldo\newline \textbf{EDT:} 6.4.1 & /P01\newline Cumple; RT & El ingreso al recinto se controlará mediante seguridad física y control de acceso biométrico basado principalmente en biometría facial, con AFIS como respaldo. Se admite proponer sistemas de mayor seguridad.\newline \emph{T-12, Parte B, fila RT-\allowbreak{}06.20} \\
\textbf{RT-\allowbreak{}06.21}\newline Bases Técnicas Transversales, RT-06.21 & Bitácora electrónica auditable conservada $\geq$ 5 años\newline \textbf{EDT:} 6.4.1 & /P01\newline Cumple; RT & Todo ingreso y egreso quedará registrado en una bitácora auditable, con identificación de la persona, fecha, hora y motivo, conservada por el período de retención declarado.\newline \emph{T-12, Parte B, fila RT-\allowbreak{}06.21} \\
\textbf{RT-\allowbreak{}06.22}\newline Bases Técnicas Transversales, RT-06.22 & Espacio de enrolamiento en el pasillo de control y estación de enrolamiento externa\newline \textbf{EDT:} 6.4.1 & /P01\newline Cumple; RT & Entre el acceso principal y el término del pasillo de la zona de control se dispondrá un espacio para la atención de personas en proceso de enrolamiento. Se evaluará mejor la existencia de una estación de enrolamiento fuera de las instalaciones del recinto técnico.\newline \emph{T-12, Parte B, fila RT-\allowbreak{}06.22} \\
\textbf{RT-\allowbreak{}06.23}\newline Bases Técnicas Transversales, RT-06.23 & Esclusa antipassback de una persona con nueva verificación\newline \textbf{EDT:} 6.4.1 & /P01\newline Cumple; RT & Al término del pasillo se instalará un acceso que impida el paso de más de una persona a la vez, con nueva verificación de identidad previa al ingreso.\newline \emph{T-12, Parte B, fila RT-\allowbreak{}06.23} \\
\end{tablasimple}

\begin{tablasimple}{>{\hsize=0.82\hsize\linewidth=\hsize}Y >{\hsize=1.04\hsize\linewidth=\hsize}Y >{\hsize=0.65\hsize\linewidth=\hsize}Y >{\hsize=1.49\hsize\linewidth=\hsize}Y}
{\multicolumn{4}{l}{\textbf{Matriz individual: bloque 16 de 26.}} \\
\cab{Requisito y origen} & \cab{Componente y EDT} & \cab{Prueba prevista y estado} & \cab{Criterio y fuente}}
\textbf{RT-\allowbreak{}06.24}\newline Bases Técnicas Transversales, RT-06.24 & Seis cámaras IP con 30 días en línea y respaldo recuperable\newline \textbf{EDT:} 6.4.2 & /P01\newline Cumple; RT & El recinto contará con videovigilancia y monitoreo IP, con imágenes en línea y disponibles para visualización de al menos los últimos 30 días. Las grabaciones anteriores se respaldarán en un medio secundario recuperable y auditable.\newline \emph{T-12, Parte B, fila RT-\allowbreak{}06.24} \\
\textbf{RT-\allowbreak{}06.25}\newline Bases Técnicas Transversales, RT-06.25 & Control biométrico y bitácora del recinto técnico; autorización previa, acompañamiento permanente y registro de fabricantes, mantenedores y auditores, con revocación del acceso temporal.\newline \textbf{EDT:} 6.4.1 & /P01\newline Cumple; RT & El PROPONENTE declarará el procedimiento de acceso de terceros —fabricantes, mantenedores, auditores— con acompañamiento obligatorio y registro.\newline \emph{T-12, Parte B, fila RT-\allowbreak{}06.25} \\
\textbf{RT-\allowbreak{}06.26}\newline Bases Técnicas Transversales, RT-06.26 & Medio de respaldo físico rotativo cifrado con custodia externa\newline \textbf{EDT:} 6.4.3 & /P01\newline Cumple; RT & Se habilitará un servicio de custodia de medios de respaldo para el sitio primario, en un medio físico transportable a otro lugar cuando el CLIENTE lo determine. Se admite proponer una solución más segura y eficiente, debidamente justificada.\newline \emph{T-12, Parte B, fila RT-\allowbreak{}06.26} \\
\textbf{RT-\allowbreak{}06.27}\newline Bases Técnicas Transversales, RT-06.27 & Recinto de custodia con iluminación sin UV $\leq$ 300 lux, 40–60 \% HR, ventilación y 18–27 °C\newline \textbf{EDT:} 6.4.3 & /P01\newline Cumple; RT & El recinto de custodia cumplirá exigencias de luminosidad, humedad, ventilación y cualquier otro factor que pueda afectar la calidad y la disponibilidad de los medios.\newline \emph{T-12, Parte B, fila RT-\allowbreak{}06.27} \\
\textbf{RT-\allowbreak{}06.28}\newline Bases Técnicas Transversales, RT-06.28 & Inventario de los dos medios rotativos, registro de cada entrada y salida con cadena de custodia y verificación semanal de legibilidad por restauración\newline \textbf{EDT:} 6.4.3 & /P01\newline Cumple; RT & Se llevará un inventario de medios con rotación, verificación periódica de legibilidad y registro de todo movimiento de entrada y de salida.\newline \emph{T-12, Parte B, fila RT-\allowbreak{}06.28} \\
\textbf{RT-\allowbreak{}06.29}\newline Bases Técnicas Transversales, RT-06.29 & Zona de trabajo en la línea técnica del recinto y ocho estaciones de trabajo nuevas, cada una con telefonía IP por softphone y salida a Internet por el firewall D-01 del sitio.\newline \textbf{EDT:} 2.3.1 & /P01\newline Cumple; RT & El PROPONENTE habilitará el espacio físico necesario para el personal encargado de la operación y administración de la plataforma, con estaciones de trabajo, telefonía, conexión a Internet y todo elemento que permita realizar la labor en condiciones adecuadas.\newline \emph{T-12, Parte B, fila RT-\allowbreak{}06.29} \\
\textbf{RT-\allowbreak{}06.30}\newline Bases Técnicas Transversales, RT-06.30 & Puestos de trabajo en la línea técnica, separados de la sala de equipos\newline \textbf{EDT:} 2.3.1 & /P01\newline Cumple; RT & El espacio de operación estará separado de la sala de equipos y no requerirá el ingreso al recinto técnico para las labores habituales de operación.\newline \emph{T-12, Parte B, fila RT-\allowbreak{}06.30} \\
\textbf{RT-\allowbreak{}06.31}\newline Bases Técnicas Transversales, RT-06.31 & Uso declarado de instalaciones sanitarias, zonas de seguridad y áreas exteriores del edificio\newline \textbf{EDT:} 2.3.1 & /P01\newline Cumple; RT & Las instalaciones sanitarias, las zonas de seguridad ante emergencia y las áreas exteriores existentes en el edificio del CLIENTE podrán utilizarse y no deben implementarse nuevamente. El PROPONENTE declarará de cuáles hará uso.\newline \emph{T-12, Parte B, fila RT-\allowbreak{}06.31} \\
\textbf{RT-\allowbreak{}06.32}\newline Bases Técnicas Transversales, RT-06.32 & Fibra y LTE con ingreso por puntos separados y ductos independientes; Starlink como tercer camino\newline \textbf{EDT:} 6.2.1 & /P01\newline Cumple; RT & El acceso a las redes de comunicaciones estará provisto a través de rutas físicas distintas, con ingreso al edificio por puntos separados.\newline \emph{T-12, Parte B, fila RT-\allowbreak{}06.32} \\
\textbf{RT-\allowbreak{}06.33}\newline Bases Técnicas Transversales, RT-06.33 & Conectividad, seguridad y canalizaciones provistas por LafroX\newline \textbf{EDT:} 6.2.1 & /P01\newline Cumple; RT & El PROPONENTE proveerá toda la conectividad, la seguridad y las canalizaciones, cañerías o ductos requeridos para cumplir los niveles de servicio comprometidos.\newline \emph{T-12, Parte B, fila RT-\allowbreak{}06.33} \\
\textbf{RT-\allowbreak{}06.34}\newline Bases Técnicas Transversales, RT-06.34 & Esquema de respaldo 3-2-1-1-0, con copia inmutable adicional y medio físico cifrado bajo custodia externa, justificado frente a borrado y cifrado malicioso.\newline \textbf{EDT:} 3.2.4 & /P01\newline Cumple; RT & Se privilegiará al PROPONENTE que ofrezca o provea especificaciones nuevas o mejores que las aquí establecidas, debidamente fundamentadas.\newline \emph{T-12, Parte B, fila RT-\allowbreak{}06.34} \\
\textbf{RT-\allowbreak{}07.01}\newline Bases Técnicas Transversales, RT-07.01 & Activo-pasivo en caliente en us-east-1 (ADR-09), justificado frente a activo-activo y restauración en frío por costo, RTO y complejidad\newline \textbf{EDT:} 2.5.2, 3.1.3 & /P01\newline Cumple; RT & El PROPONENTE declarará la modalidad escogida —activo-activo o activo-pasivo— y la justificará frente al costo, al RTO comprometido y a la complejidad operacional que introduce.\newline \emph{T-12, Parte B, fila RT-\allowbreak{}07.01} \\
\textbf{RT-\allowbreak{}07.02}\newline Bases Técnicas Transversales, RT-07.02 & us-east-1 a $\approx$ 7.700 km de sa-east-1; sa-east-1 a $\approx$ 2.700 km de Talca, con amenazas comunes\newline \textbf{EDT:} 3.1.3 & /P01\newline Cumple; RT & El sitio secundario estará emplazado a una distancia suficiente del principal para no verse afectado por el mismo evento de fuerza mayor. El PROPONENTE declarará la distancia y el análisis de amenazas comunes considerado.\newline \emph{T-12, Parte B, fila RT-\allowbreak{}07.02} \\
\textbf{RT-\allowbreak{}07.03}\newline Bases Técnicas Transversales, RT-07.03 & Aurora Global Database, AWS DMS, S3 Replication Time Control y DynamoDB Global Tables, con alarmas de retraso a 60 s, 5 y 15 min\newline \textbf{EDT:} 3.2.4 & /P01\newline Cumple; RT & La replicación de datos será continua, con medición y alertamiento del retraso de replicación.\newline \emph{T-12, Parte B, fila RT-\allowbreak{}07.03} \\
\textbf{RT-\allowbreak{}07.04}\newline Bases Técnicas Transversales, RT-07.04 & Sitio secundario en us-east-1 con RTO de hasta cuatro horas y RPO de hasta quince minutos para los escenarios de diseño, sostenido por la extracción continua por fibra, LTE y Starlink, con dos ensayos regionales por año.\newline \textbf{EDT:} 3.1.3, 2.5.2 & /P01\newline Cumple; RT & El objetivo de tiempo de recuperación (RTO) no superará 4 horas y el objetivo de punto de recuperación (RPO) no superará 15 minutos para los servicios críticos, salvo exigencia superior del caso.\newline \emph{T-12, Parte B, fila RT-\allowbreak{}07.04} \\
\textbf{RT-\allowbreak{}07.05}\newline Bases Técnicas Transversales, RT-07.05 & Secuencia documentada y automatizada con AWS Systems Manager Automation, ejecutable por el CLIENTE\newline \textbf{EDT:} 2.5.2 & /P01\newline Cumple; RT & El procedimiento de conmutación estará documentado, automatizado en la mayor medida posible y ejecutable por el personal del CLIENTE tras la transferencia de conocimiento.\newline \emph{T-12, Parte B, fila RT-\allowbreak{}07.05} \\
\textbf{RT-\allowbreak{}07.06}\newline Bases Técnicas Transversales, RT-07.06 & Procedimiento de retorno con resincronización y conciliación, ensayado cada semestre\newline \textbf{EDT:} 2.5.2, 3.1.3 & /P01\newline Cumple; RT & Existirá un procedimiento de retorno al sitio principal igualmente documentado y probado, con reconciliación de los datos generados durante la contingencia.\newline \emph{T-12, Parte B, fila RT-\allowbreak{}07.06} \\
\textbf{RT-\allowbreak{}07.07}\newline Bases Técnicas Transversales, RT-07.07 & Ensayo semestral de pérdida regional y de la sala de Talca con medición de RTO y RPO\newline \textbf{EDT:} 3.8.5, 3.9.4, 8.1.3 & /P01\newline Cumple; RT & El plan de recuperación ante desastres se probará al menos dos veces al año mediante conmutación real, con informe de resultados, medición del RTO y del RPO efectivamente alcanzados y plan de corrección de las brechas detectadas.\newline \emph{T-12, Parte B, fila RT-\allowbreak{}07.07} \\
\textbf{RT-\allowbreak{}07.08}\newline Bases Técnicas Transversales, RT-07.08 & Detección de indisponibilidad mediante supervisión de salud y autorización del CLIENTE para iniciar la conmutación.\newline \textbf{EDT:} 2.5.2, 3.1.3 & /P01\newline Parcial; RT & Se valorará que la conmutación sea automática ante la detección de indisponibilidad, con criterio de disparo declarado y protección contra conmutación innecesaria.\newline \emph{T-12, Parte B, fila RT-\allowbreak{}07.08} \\
\textbf{RT-\allowbreak{}07.09}\newline Bases Técnicas Transversales, RT-07.09 & Esquema 3-2-1-1-0: NAS Synology RS2423RP+ WORM, AWS Backup con Vault Lock, S3 Object Lock y medio físico en custodia\newline \textbf{EDT:} 3.2.4 & /P01\newline Cumple; RT & La política de respaldo seguirá el esquema 3-2-1-0: tres copias, en dos medios distintos, una fuera de sitio, una inmutable o fuera de línea y cero errores de verificación de restauración.\newline \emph{T-12, Parte B, fila RT-\allowbreak{}07.09} \\
\textbf{RT-\allowbreak{}07.10}\newline Bases Técnicas Transversales, RT-07.10 & Respaldos cifrados con claves AWS KMS administradas por el CLIENTE, independientes de producción\newline \textbf{EDT:} 3.2.4 & /P01\newline Cumple; RT & Los respaldos estarán cifrados en reposo y en tránsito, con clave gestionada de forma independiente de la infraestructura respaldada.\newline \emph{T-12, Parte B, fila RT-\allowbreak{}07.10} \\
\textbf{RT-\allowbreak{}07.11}\newline Bases Técnicas Transversales, RT-07.11 & AWS Backup Vault Lock en modo Compliance y S3 Object Lock\newline \textbf{EDT:} 3.2.4 & /P01\newline Cumple; RT & Las copias inmutables estarán protegidas contra borrado y contra modificación durante su período de retención, incluso frente a credenciales administrativas comprometidas.\newline \emph{T-12, Parte B, fila RT-\allowbreak{}07.11} \\
\textbf{RT-\allowbreak{}07.12}\newline Bases Técnicas Transversales, RT-07.12 & Restauración mensual documentada con medición del tiempo efectivo\newline \textbf{EDT:} 8.1.6 & /P01\newline Cumple; RT & Se ejecutará y documentará una prueba de restauración al menos mensual, sobre una muestra representativa, con medición del tiempo efectivo de restauración.\newline \emph{T-12, Parte B, fila RT-\allowbreak{}07.12} \\
\textbf{RT-\allowbreak{}07.13}\newline Bases Técnicas Transversales, RT-07.13 & Tabla de respaldo y restauración por dominio de datos\newline \textbf{EDT:} 3.2.4 & /P01\newline Cumple; RT & El PROPONENTE declarará, por cada dominio de datos, la frecuencia de respaldo, el período de retención y el tiempo estimado de restauración completa.\newline \emph{T-12, Parte B, fila RT-\allowbreak{}07.13} \\
\textbf{RT-\allowbreak{}07.14}\newline Bases Técnicas Transversales, RT-07.14 & Restauración granular por PITR de 35 días sobre instancia temporal\newline \textbf{EDT:} 3.2.4 & /P01\newline Cumple; RT & Los respaldos permitirán la restauración granular: un registro, una tabla, un módulo o el sistema completo.\newline \emph{T-12, Parte B, fila RT-\allowbreak{}07.14} \\
\end{tablasimple}

\begin{tablasimple}{>{\hsize=0.82\hsize\linewidth=\hsize}Y >{\hsize=1.04\hsize\linewidth=\hsize}Y >{\hsize=0.65\hsize\linewidth=\hsize}Y >{\hsize=1.49\hsize\linewidth=\hsize}Y}
{\multicolumn{4}{l}{\textbf{Matriz individual: bloque 17 de 26.}} \\
\cab{Requisito y origen} & \cab{Componente y EDT} & \cab{Prueba prevista y estado} & \cab{Criterio y fuente}}
\textbf{RT-\allowbreak{}08.01}\newline Bases Técnicas Transversales, RT-08.01 & Servidores HPE ProLiant DL345 Gen11 o Dell PowerEdge R7615 con AMD EPYC 9124 de 16 núcleos, 64 GB DDR5, NVMe y dos puertos de 10 GbE; HPE ProLiant DL20 Gen11 con Intel Xeon E-2468 de 8 núcleos, 32 GB, cuatro SSD en RAID 10 y cuatro puertos de 1 GbE; mini-PC Advantech ARK-1251 con Intel Core Ultra 5-125U, 16 GB DDR5, dos SSD NVMe en RAID 1 y tres puertos Ethernet; consumo por equipo en la tabla de carga eléctrica y en la sección 4-W.8 del Anexo 4-W, y dimensionamiento por máquina virtual en el Anexo 4-W.\newline \textbf{EDT:} 5.1.2 & /P01\newline Cumple; RT & El PROPONENTE especificará el equipamiento de cómputo con su marca, modelo de referencia, procesador, memoria, almacenamiento local, interfaces y consumo, junto con el cálculo de dimensionamiento que lo sustenta.\newline \emph{T-12, Parte B, fila RT-\allowbreak{}08.01} \\
\textbf{RT-\allowbreak{}08.02}\newline Bases Técnicas Transversales, RT-08.02 & Ceph de 3 réplicas sobre NVMe, RAID 10, RAID 1 y RAID 6 según sitio; comprobaciones de integridad con reparación y monitoreo predictivo de salud y desgaste\newline \textbf{EDT:} 6.3.1 & /P01\newline Cumple; RT & El almacenamiento será redundante, con tolerancia declarada a la falla de discos, control de errores y monitoreo predictivo de salud de los medios.\newline \emph{T-12, Parte B, fila RT-\allowbreak{}08.02} \\
\textbf{RT-\allowbreak{}08.03}\newline Bases Técnicas Transversales, RT-08.03 & Conmutadores de núcleo en par; firewalls en par activo/pasivo; balanceadores con alta disponibilidad.\newline \textbf{EDT:} 6.3.2 & /P01\newline Cumple; RT & Los conmutadores de núcleo, los cortafuegos y los balanceadores de carga estarán en configuración de alta disponibilidad, sin punto único de falla.\newline \emph{T-12, Parte B, fila RT-\allowbreak{}08.03} \\
\textbf{RT-\allowbreak{}08.04}\newline Bases Técnicas Transversales, RT-08.04 & Equipamiento de cómputo, almacenamiento y red de la sección 8.1 de las Bases con doble fuente en circuitos distintos: servidores de Talca y Concepción, respaldo local, firewalls, switches de núcleo, acceso y cross-docking, consola KVM, gateway IoT y mini-PC con dos fuentes DIN y módulo de redundancia. Cada gabinete recibe dos circuitos (PDU A/B en Talca; UPS y circuito directo del tablero en borde y piso). Los equipos de una sola fuente (servidor del ERP del CLIENTE, Starlink, ONT y router LTE del operador) se conectan a conmutadores automáticos de transferencia entre ambos circuitos.\newline \textbf{EDT:} 6.3.1 & /P01\newline Cumple; RT & Todo el equipamiento contará con fuentes de poder redundantes y conexión a circuitos eléctricos distintos.\newline \emph{T-12, Parte B, fila RT-\allowbreak{}08.04} \\
\textbf{RT-\allowbreak{}08.05}\newline Bases Técnicas Transversales, RT-08.05 & Márgenes de capacidad sobre la carga del año 3: 155,6 \% WMS Talca, 151,9 \% cada cross-docking y 158,3 \% aplicación en nube, con ampliación por revisión trimestral\newline \textbf{EDT:} 8.4.1 & /P01\newline Cumple; RT & El PROPONENTE declarará el margen de crecimiento del dimensionamiento propuesto, expresado como porcentaje sobre la carga proyectada del caso, y el procedimiento de ampliación.\newline \emph{T-12, Parte B, fila RT-\allowbreak{}08.05} \\
\textbf{RT-\allowbreak{}08.06}\newline Bases Técnicas Transversales, RT-08.06 & Equipamiento nuevo, sin uso previo y con garantía de fábrica vigente desde la recepción conforme.\newline \textbf{EDT:} 5.1.2 & /P01\newline Cumple; RT & El equipamiento será nuevo, sin uso previo, con garantía de fábrica vigente desde la recepción conforme.\newline \emph{T-12, Parte B, fila RT-\allowbreak{}08.06} \\
\textbf{RT-\allowbreak{}08.07}\newline Bases Técnicas Transversales, RT-08.07 & Ocho estaciones nuevas con dos monitores, asignadas a despacho, administración, planificación, calidad y TI conforme al dimensionamiento de sitios.\newline \textbf{EDT:} 5.1.2; 5.1.3 & /P01\newline Cumple; RT & El PROPONENTE especificará las estaciones de trabajo requeridas para la operación y la administración de la plataforma, en la cantidad que determine el dimensionamiento del caso, con monitores duales.\newline \emph{T-12, Parte B, fila RT-\allowbreak{}08.07} \\
\textbf{RT-\allowbreak{}08.08}\newline Bases Técnicas Transversales, RT-08.08 & Estaciones de trabajo con ergonomía NCh 2527 y certificación de eficiencia energética.\newline \textbf{EDT:} 5.1.2; 5.1.3 & /P01\newline Cumple; RT & Los puestos de trabajo cumplirán las condiciones ergonómicas de la NCh 2527 y los equipos contarán con certificación de eficiencia energética.\newline \emph{T-12, Parte B, fila RT-\allowbreak{}08.08} \\
\textbf{RT-\allowbreak{}08.09}\newline Bases Técnicas Transversales, RT-08.09 & Ocho estaciones nuevas y equipos existentes con gestión central, CrowdStrike Falcon, cifrado de disco, control de extraíbles y actualización automatizada antes de Verified Access\newline \textbf{EDT:} 3.2.5, 5.2.2 & /P01\newline Cumple; RT & Las estaciones estarán gestionadas de forma centralizada, con cifrado de disco, control de dispositivos extraíbles, antivirus con detección y respuesta y actualización automatizada.\newline \emph{T-12, Parte B, fila RT-\allowbreak{}08.09} \\
\textbf{RT-\allowbreak{}08.10}\newline Bases Técnicas Transversales, RT-08.10 & T-11 especifica marca, modelo, cantidad, características, accesorios y consumibles de los dispositivos de terreno y de operación; el costo unitario estimado se declara en la Oferta Económica, porque la Oferta Técnica no puede contener valores unitarios (Bases Administrativas, art. 50.2, p. 29). B-06, B-07 y B-10 a B-14 son accesorios y consumibles de B-04 y B-05, especificados por sus características mínimas; la marca y el modelo de referencia corresponden a B-03, B-04, B-05, B-08 y B-09.\newline \textbf{EDT:} 5.1.1, 5.1.4, 2.4.3, 2.4.4 & /P01\newline Cumple; RT & El PROPONENTE especificará cada dispositivo de terreno con marca, modelo de referencia, cantidad, características mínimas, accesorios, consumibles y costo unitario estimado, aun cuando su compra sea de cargo del CLIENTE.\newline \emph{T-12, Parte B, fila RT-\allowbreak{}08.10} \\
\textbf{RT-\allowbreak{}08.11}\newline Bases Técnicas Transversales, RT-08.11 & MC9400 estándar y Cold Storage Freezer, TC58e de reparto y preventa, ZQ620 Plus, PAX A920 Pro, baterías y cargadores, UPS B-13 y protección B-14.\newline \emph{T-12, fila RT-\allowbreak{}08.11}\newline \textbf{EDT:} 5.1.1, 5.1.4, 2.4.3 & /P01\newline Cumple; RT & La especificación considerará las condiciones reales de uso del caso: intemperie, humedad, polvo, vibración, temperatura, uso con guantes, luminosidad y autonomía de batería requerida por turno.\newline \emph{T-12, Parte B, fila RT-\allowbreak{}08.11} \\
\textbf{RT-\allowbreak{}08.12}\newline Bases Técnicas Transversales, RT-08.12 & T-11 declara grado IP y caída por modelo: TC58e de preventa y de reparto IP65/68 y 1,8 m; MC9400 estándar IP65/68 y 2,4 m; MC9400 Cold Storage Freezer IP65/68 y 2,1 m; ZQ620 Plus IP54, con caída mínima de 1,5 m exigida en la compra; termógrafo EYE SENSOR EN12830 IP67 de montaje fijo; balanza Dibal fija con IP65 mínimo exigido en la compra; impresora ZT411 fija de interior; PAX A920 Pro sin grado IP del fabricante, usado bajo el resguardo de la cabina o del local del cliente. El kit es de montaje fijo: B-03, B-08 y B-09 son IP67; B-04 (IP30) y B-05 van en un gabinete IP54 de cabina, y B-06, B-10 y B-11 se compran con IP67 como mínimo.\newline \textbf{EDT:} 5.1.1, 5.1.4, 2.4.3 & /P01\newline Cumple; RT & Los dispositivos declararán su grado de protección contra polvo y agua y su resistencia a caídas, coherentes con el entorno de operación.\newline \emph{T-12, Parte B, fila RT-\allowbreak{}08.12} \\
\textbf{RT-\allowbreak{}08.13}\newline Bases Técnicas Transversales, RT-08.13 & Vida útil esperada de cinco años, reserva del 10 \%, garantía del fabricante y reposición de unidades dadas de baja desde la reserva. Las baterías de los terminales se supervisan con Zebra DNA y se reponen bajo el 80 \% de su capacidad nominal; el MC9400 estándar tiene 81 baterías de repuesto en el año 3 y el PAX A920 Pro, cargador de vehículo. LafroX mantiene los componentes de cámara y los kits de camión, incluida calibración y revisión de glicol y baterías; las reservas térmicas son 31 por componente unitario, 93 sondas y 93 botellas.\newline \textbf{EDT:} 5.1.1, 5.1.4, 8.2.7 & /P01\newline Cumple; RT & El PROPONENTE indicará el ciclo de vida esperado de cada dispositivo, la disponibilidad de repuestos y el plan de reposición durante los 56 meses del Contrato.\newline \emph{T-12, Parte B, fila RT-\allowbreak{}08.13} \\
\textbf{RT-\allowbreak{}08.14}\newline Bases Técnicas Transversales, RT-08.14 & Inventario común con Android Enterprise/Zebra DNA para terminales, Printer Profile Manager para impresoras, PAXSTORE para POS, Greengrass V2 para gateways y herramientas Teltonika/TC58e para registradores. Sensores y balanzas se actualizan en mantenimiento y se bloquean con la baja de su registro desde la consola central.\newline \textbf{EDT:} 3.2.6 & /P01\newline Cumple; RT & Los dispositivos se integrarán a la gestión centralizada de flota exigida en RT-03.18.\newline \emph{T-12, Parte B, fila RT-\allowbreak{}08.14} \\
\textbf{RT-\allowbreak{}08.15}\newline Bases Técnicas Transversales, RT-08.15 & LafroX provee una unidad de cada tipo y variante del T-11 para aceptación antes de la compra masiva; se incorpora al parque o reserva dimensionados.\newline \textbf{EDT:} 5.1.1 & /P01\newline Cumple; RT & El PROPONENTE proveerá una unidad de cada tipo de dispositivo especificado para pruebas de aceptación por parte del CLIENTE, antes de la compra masiva.\newline \emph{T-12, Parte B, fila RT-\allowbreak{}08.15} \\
\textbf{RT-\allowbreak{}08.16}\newline Bases Técnicas Transversales, RT-08.16 & Plan de ciclo de vida con recepción técnica, instalación y pruebas de aceptación, mantenimiento de hardware, actualizaciones anuales y reposición/disposición final.\newline \textbf{EDT:} 5.1.3, 6.5.1, 8.2.7 & /P01\newline Cumple; RT & El PROPONENTE presentará el plan de ciclo de vida del equipamiento: recepción, puesta en servicio, mantención, actualización, retiro y disposición final.\newline \emph{T-12, Parte B, fila RT-\allowbreak{}08.16} \\
\textbf{RT-\allowbreak{}08.17}\newline Bases Técnicas Transversales, RT-08.17 & Procedimiento de sanitización NIST SP 800-88 Rev. 2: inventario y custodia de todo medio retirado, purga o destrucción, verificación y certificado al CLIENTE; incluye servidores, NAS, HDD/SSD, mini-PC, móviles y medio físico rotativo de respaldo.\newline \textbf{EDT:} 8.2.7, 9.2.2 & /P01\newline Cumple; RT & Todo medio de almacenamiento que salga de servicio será borrado de forma segura y verificable, con certificado de destrucción o de sanitización entregado al CLIENTE.\newline \emph{T-12, Parte B, fila RT-\allowbreak{}08.17} \\
\textbf{RT-\allowbreak{}08.18}\newline Bases Técnicas Transversales, RT-08.18 & Gestor autorizado y registrado conforme a Ley N.º 20.920, trazabilidad y certificado de disposición final entregado al CLIENTE; control de sanitización o destrucción de medios.\newline \textbf{EDT:} 8.2.7 & /P01\newline Cumple; RT & La disposición final de equipamiento electrónico se realizará con gestor autorizado, conforme a la normativa de residuos aplicable, con certificado de disposición.\newline \emph{T-12, Parte B, fila RT-\allowbreak{}08.18} \\
\textbf{RT-\allowbreak{}08.19}\newline Bases Técnicas Transversales, RT-08.19 & Extensión de vida útil: equipos reparables en servicio al menos 60 meses desde su recepción mediante mantención, repuestos y reemplazo de baterías, sin renovación general al cierre del contrato; menos residuos electrónicos.\newline \textbf{EDT:} 8.2.7 & /P01\newline Cumple; RT & Se valorará una estrategia de reacondicionamiento o de extensión de vida útil que reduzca el impacto ambiental, cuantificada en la propuesta.\newline \emph{T-12, Parte B, fila RT-\allowbreak{}08.19} \\
\textbf{RT-\allowbreak{}09.01}\newline Bases Técnicas Transversales, RT-09.01 & Dimensionamiento de capacidad sobre la volumetría del caso y sus umbrales: confirmación de línea de picking $\leq$1 s, entrega $\leq$2 s, línea de preventa $\leq$1,5 s y consulta de stock y crédito $\leq$2 s.\newline \textbf{EDT:} 2.1.2 & /P01\newline Cumple; RT & El PROPONENTE presentará el cálculo de capacidad que sustenta su dimensionamiento, con los supuestos de usuarios concurrentes, transacciones por segundo, volumen de datos y crecimiento anual, tomados de la volumetría del caso. Caso, cap. 15 — Transacción operacional crítica: Confirmación de una línea de preparación de pedidos: no superior a 1 segundo. Registro de una entrega en el local del cliente: no superior a 2 segundos. Registro de una línea en la toma de pedido de preventa: no superior a 1,5 segundos. Consulta de stock y crédito en preventa: no superior a 2 segundos.\newline \emph{T-12, Parte B, fila RT-\allowbreak{}09.01} \\
\textbf{RT-\allowbreak{}09.02}\newline Bases Técnicas Transversales, RT-09.02 & Concurrencia y transacciones derivadas de la volumetría de los numerales 14.1 y 14.2 del caso, con su perfil horario.\newline \textbf{EDT:} 3.8.4 & /P01\newline Cumple; RT & La solución soportará la concurrencia y el volumen de transacciones que fije el caso, y mantendrá los umbrales del numeral 9.1 bajo esa carga. Caso, cap. 15 — Concurrencia y volumen de transacciones: El PROPONENTE lo deriva de la volumetría del numeral 14.1, considerando el perfil horario descrito, y lo declara conforme al numeral 14.2.\newline \emph{T-12, Parte B, fila RT-\allowbreak{}09.02} \\
\textbf{RT-\allowbreak{}09.03}\newline Bases Técnicas Transversales, RT-09.03 & Aurora fijada para operar a 3$\times$ y Amazon ECS Fargate con perfiles de escalamiento independientes, para triplicar la volumetría inicial en tres años sin rediseño.\newline \textbf{EDT:} 2.1.2 & /P01\newline Cumple; RT & La solución soportará, sin rediseño, un crecimiento de al menos tres veces la volumetría inicial del caso en un horizonte de tres años.\newline \emph{T-12, Parte B, fila RT-\allowbreak{}09.03} \\
\textbf{RT-\allowbreak{}09.04}\newline Bases Técnicas Transversales, RT-09.04 & Escalamiento horizontal automático por perfil en Amazon ECS Fargate, con métricas de un minuto, nueva tarea en servicio en menos de tres minutos y drenaje del balanceador sin perder transacciones en curso.\newline \textbf{EDT:} 3.2.1 & /P01\newline Cumple; RT & El escalamiento de las capas de aplicación e integración será horizontal y automático, con tiempo de reacción declarado y sin pérdida de transacciones en curso.\newline \emph{T-12, Parte B, fila RT-\allowbreak{}09.04} \\
\textbf{RT-\allowbreak{}09.05}\newline Bases Técnicas Transversales, RT-09.05 & ACL local y erp-sync para emisión documental del ERP, identificados como primer cuello de botella probable, con métricas de detección y medidas de desacople y control de flujo.\newline \textbf{EDT:} 2.1.2 & /P01\newline Cumple; RT & El PROPONENTE identificará el componente que primero se convertirá en cuello de botella al crecer la carga y explicará cómo lo detectará y cómo lo resolverá.\newline \emph{T-12, Parte B, fila RT-\allowbreak{}09.05} \\
\textbf{RT-\allowbreak{}09.06}\newline Bases Técnicas Transversales, RT-09.06 & Pruebas en Preproducción a 1,5 veces el peak declarado y prueba de estrés hasta encontrar el punto de quiebre.\newline \textbf{EDT:} 3.8.4, 3.9.3 & /P01\newline Cumple; RT & Se ejecutarán pruebas de carga sobre Preproducción con un volumen equivalente a 1,5 veces el peak declarado, y pruebas de estrés hasta identificar el punto de quiebre de la solución.\newline \emph{T-12, Parte B, fila RT-\allowbreak{}09.06} \\
\end{tablasimple}

\begin{tablasimple}{>{\hsize=0.82\hsize\linewidth=\hsize}Y >{\hsize=1.04\hsize\linewidth=\hsize}Y >{\hsize=0.65\hsize\linewidth=\hsize}Y >{\hsize=1.49\hsize\linewidth=\hsize}Y}
{\multicolumn{4}{l}{\textbf{Matriz individual: bloque 18 de 26.}} \\
\cab{Requisito y origen} & \cab{Componente y EDT} & \cab{Prueba prevista y estado} & \cab{Criterio y fuente}}
\textbf{RT-\allowbreak{}09.07}\newline Bases Técnicas Transversales, RT-09.07 & Protocolo AL-PERF-01 con curva de respuesta frente a carga, saturación, consumo de recursos y recuperación posterior al peak.\newline \textbf{EDT:} 3.8.4, 3.9.3 & /P01\newline Cumple; RT & El informe de pruebas de carga incluirá la curva de tiempo de respuesta frente a carga, el punto de saturación, el consumo de recursos y el comportamiento durante y después del peak.\newline \emph{T-12, Parte B, fila RT-\allowbreak{}09.07} \\
\textbf{RT-\allowbreak{}09.08}\newline Bases Técnicas Transversales, RT-09.08 & Degradación controlada mediante encolamiento, límites de tasa y mensajes operativos explícitos, con registro de operaciones no confirmadas.\newline \textbf{EDT:} 2.1.1 & /P01\newline Cumple; RT & La solución degradará de forma controlada al superarse la capacidad: encolamiento, limitación de tasa y mensaje explícito a la persona usuaria, nunca error genérico ni pérdida silenciosa de transacciones.\newline \emph{T-12, Parte B, fila RT-\allowbreak{}09.08} \\
\textbf{RT-\allowbreak{}09.09}\newline Bases Técnicas Transversales, RT-09.09 & Gestión de capacidad con proyección trimestral, alertas de agotamiento y propuestas de dimensionamiento que incluyen su efecto en costos de infraestructura, nube, licencias y enlaces respecto de la base contratada, con consumos y condiciones económicas vigentes para decisión del CLIENTE antes de ejecutarlas.\newline \textbf{EDT:} 8.4.1 & /P01\newline Cumple; RT & Se gestionará la capacidad durante la Operación con proyección trimestral de crecimiento, alertas anticipadas de agotamiento y propuesta de ajuste de dimensionamiento y de costo.\newline \emph{T-12, Parte B, fila RT-\allowbreak{}09.09} \\
\textbf{RT-\allowbreak{}09.10}\newline Bases Técnicas Transversales, RT-09.10 & CI ejecuta carga de regresión automatizada en QA en cada versión candidata y al menos semanalmente, y antes de liberar en Preproducción; compara percentil 95, errores y colas con la versión aceptada y bloquea regresiones.\newline \textbf{EDT:} 3.1.4 & /P01\newline Cumple; RT & Se valorará la existencia de pruebas de carga automatizadas ejecutadas de forma periódica en el flujo de integración continua, con detección de regresiones de desempeño.\newline \emph{T-12, Parte B, fila RT-\allowbreak{}09.10} \\
\textbf{RT-\allowbreak{}10.01}\newline Bases Técnicas Transversales, RT-10.01 & Disponibilidad mensual mínima de 99,9 \% medida en la transacción crítica de negocio de extremo a extremo.\newline \textbf{EDT:} 8.1.2 & /P01\newline Cumple; RT & La solución alcanzará una disponibilidad mensual mínima de 99,9 \% para los servicios clasificados como críticos, medida sobre la transacción de negocio de extremo a extremo.\newline \emph{T-12, Parte B, fila RT-\allowbreak{}10.01} \\
\textbf{RT-\allowbreak{}10.02}\newline Bases Técnicas Transversales, RT-10.02 & Clasificación de servicios por criticidad operacional y aplicación del nivel de servicio correspondiente al Artículo 78° de las Bases Administrativas.\newline \textbf{EDT:} 2.5.1 & /P01\newline Cumple; RT & El PROPONENTE clasificará cada servicio de la solución en crítico, alto, medio o bajo, justificando la clasificación con el impacto operacional de su indisponibilidad, y aplicará a cada uno el nivel de servicio correspondiente del Artículo 78° de las Bases Administrativas.\newline \emph{T-12, Parte B, fila RT-\allowbreak{}10.02} \\
\textbf{RT-\allowbreak{}10.03}\newline Bases Técnicas Transversales, RT-10.03 & Plan de continuidad alineado con ISO 22301, análisis de impacto, escenarios de contingencia, procedimientos manuales y criterios de activación; la continuidad TIC se articula con recuperación ante desastres.\newline \textbf{EDT:} 2.5.1 & /P01\newline Cumple; RT & El plan de continuidad del negocio se elaborará conforme a ISO 22301, con análisis de impacto en el negocio, escenarios de contingencia, procedimientos manuales de respaldo y criterios de activación.\newline \emph{T-12, Parte B, fila RT-\allowbreak{}10.03} \\
\textbf{RT-\allowbreak{}10.04}\newline Bases Técnicas Transversales, RT-10.04 & Continuidad TIC conforme a ISO/IEC 27031, articulada con el plan de recuperación ante desastres del Capítulo 7.\newline \textbf{EDT:} 2.5.2 & /P01\newline Cumple; RT & La continuidad TIC se estructurará conforme a ISO/IEC 27031, articulada con el plan de recuperación ante desastres del Capítulo 7.\newline \emph{T-12, Parte B, fila RT-\allowbreak{}10.04} \\
\textbf{RT-\allowbreak{}10.05}\newline Bases Técnicas Transversales, RT-10.05 & Cero indisponibilidad en la ventana crítica 05:30–07:00 de lunes a sábado y en preparación de 22:00–06:00 y recepción de proveedores de 08:00–18:00; cambios fuera de esas ventanas, de los congelamientos del 1 al 25 de septiembre, todo diciembre y los tres primeros días hábiles mensuales, con aviso previo de diez días hábiles.\newline \textbf{EDT:} 1.3.1, 8.2.1 & /P01\newline Cumple; RT & Los mantenimientos programados se ejecutarán fuera de la ventana operacional crítica que defina el caso, con aviso previo mínimo de diez días hábiles. Caso, cap. 15 — Ventana operacional protegida: Ventana crítica de despacho de 05:30 a 07:00 de lunes a sábado: indisponibilidad cero. Preparación de pedidos de 22:00 a 06:00. Recepción de proveedores de 08:00 a 18:00. Congelamiento total del 1 al 25 de septiembre y durante todo diciembre. Congelamiento durante los tres primeros días hábiles de cada mes.\newline \emph{T-12, Parte B, fila RT-\allowbreak{}10.05} \\
\textbf{RT-\allowbreak{}10.06}\newline Bases Técnicas Transversales, RT-10.06 & Despliegue azul-verde con canario, sin interrupción; indisponibilidad programada sólo excepcional y justificada caso a caso.\newline \textbf{EDT:} 3.8.8, 3.9.7 & /P01\newline Cumple; RT & La solución permitirá desplegar cambios sin interrupción del servicio. Las ventanas de indisponibilidad programada serán excepcionales y deberán justificarse caso a caso.\newline \emph{T-12, Parte B, fila RT-\allowbreak{}10.06} \\
\textbf{RT-\allowbreak{}10.07}\newline Bases Técnicas Transversales, RT-10.07 & Inyección controlada de fallas de instancia, zona, dependencia, latencia y disco antes de cada paso a producción y semestral en Operación.\newline \textbf{EDT:} 3.8.4, 3.9.3, 8.2.2 & /P01\newline Cumple; RT & Se ejecutarán pruebas de resiliencia mediante inyección controlada de fallas —caída de instancia, de zona, de dependencia externa, latencia elevada, saturación de disco— antes de cada paso a producción y al menos una vez por semestre durante la Operación.\newline \emph{T-12, Parte B, fila RT-\allowbreak{}10.07} \\
\textbf{RT-\allowbreak{}10.08}\newline Bases Técnicas Transversales, RT-10.08 & Conducta documentada de cada dependencia externa ante falta de respuesta, error o lentitud.\newline \textbf{EDT:} 2.5.2 & /P01\newline Cumple; RT & El PROPONENTE documentará, por cada dependencia externa, el comportamiento de la solución cuando esa dependencia no responde, responde con error o responde con lentitud.\newline \emph{T-12, Parte B, fila RT-\allowbreak{}10.08} \\
\textbf{RT-\allowbreak{}10.09}\newline Bases Técnicas Transversales, RT-10.09 & Presupuesto de error de aproximadamente 43 minutos al mes por servicio crítico, con suspensión de despliegues no correctivos al agotarse.\newline \textbf{EDT:} 8.1.2 & /P01\newline Cumple; RT & Se declarará un presupuesto de error por servicio crítico y su vinculación con el ritmo de despliegue de cambios.\newline \emph{T-12, Parte B, fila RT-\allowbreak{}10.09} \\
\textbf{RT-\allowbreak{}11.01}\newline Bases Técnicas Transversales, RT-11.01 & Modelo Zero Trust NIST SP 800-207 con verificación explícita por solicitud, privilegio mínimo y presunción de compromiso.\newline \textbf{EDT:} 2.2.1 & /P01\newline Cumple; RT & La arquitectura de seguridad se basará en el modelo Zero Trust conforme a NIST SP 800-207: verificación explícita de cada solicitud, privilegio mínimo y presunción de compromiso.\newline \emph{T-12, Parte B, fila RT-\allowbreak{}11.01} \\
\textbf{RT-\allowbreak{}11.02}\newline Bases Técnicas Transversales, RT-11.02 & Modelado de amenazas STRIDE por módulo e integración externa, con actualización ante cambios arquitectónicos relevantes.\newline \textbf{EDT:} 2.2.2 & /P01\newline Cumple; RT & El PROPONENTE entregará un modelado de amenazas documentado por cada componente y por cada integración externa, con metodología declarada (STRIDE u otra), y lo actualizará ante cada cambio arquitectónico relevante.\newline \emph{T-12, Parte B, fila RT-\allowbreak{}11.02} \\
\textbf{RT-\allowbreak{}11.03}\newline Bases Técnicas Transversales, RT-11.03 & Clasificación Restringido, Confidencial, Interno y Público con controles diferenciados documentados en matriz.\newline \textbf{EDT:} 2.2.4 & /P01\newline Cumple; RT & La información del CLIENTE se clasificará por nivel de sensibilidad, con controles diferenciados por nivel documentados en una matriz.\newline \emph{T-12, Parte B, fila RT-\allowbreak{}11.03} \\
\textbf{RT-\allowbreak{}11.04}\newline Bases Técnicas Transversales, RT-11.04 & Escaneo continuo con metas de remediación de 7/15/30 días corridos. Si hay congelamiento se solicita cambio de emergencia; sin autorización del CLIENTE se contiene y remedia en la primera fecha habilitada. El plazo excedido por esa decisión queda registrado como decisión del CLIENTE.\newline \textbf{EDT:} 8.2.2 & /P01\newline Cumple; RT & Existirá un programa de gestión de vulnerabilidades con escaneo continuo y plazos máximos de remediación de 7 días corridos para vulnerabilidades críticas, 15 días para altas y 30 días para medias, contados desde su publicación o detección.\newline \emph{T-12, Parte B, fila RT-\allowbreak{}11.04} \\
\textbf{RT-\allowbreak{}11.05}\newline Bases Técnicas Transversales, RT-11.05 & Matriz de controles ISO/IEC 27001 e ISO/IEC 27002 con implementación concreta y evidencia.\newline \textbf{EDT:} 2.2.4 & /P01\newline Cumple; RT & El PROPONENTE mantendrá una matriz de controles de seguridad trazable a ISO/IEC 27001 e ISO/IEC 27002, indicando el control, su implementación concreta en la solución y la evidencia que lo acredita.\newline \emph{T-12, Parte B, fila RT-\allowbreak{}11.05} \\
\textbf{RT-\allowbreak{}11.06}\newline Bases Técnicas Transversales, RT-11.06 & Controles ISO/IEC 27017 para nube y tratamiento de datos personales conforme a ISO/IEC 27018.\newline \textbf{EDT:} 2.2.1, 1.9.1 & /P01\newline Cumple; RT & Se aplicarán los controles de ISO/IEC 27017 para los servicios en nube y de ISO/IEC 27018 para el tratamiento de datos personales en nube.\newline \emph{T-12, Parte B, fila RT-\allowbreak{}11.06} \\
\textbf{RT-\allowbreak{}11.07}\newline Bases Técnicas Transversales, RT-11.07 & Portales, API e identidad usan CloudFront/WAF y protección Shield Advanced. AS2 publica OpenAS2 4.9 por una distribución propia de CloudFront con AWS WAF (lista de IP de las cadenas, solo POST con cabeceras AS2 y límite de tasa) y Shield Advanced, hacia un ALB interno; firma, cifrado y MDN protegen cada mensaje de extremo a extremo.\newline \textbf{EDT:} 3.2.1 & /P01\newline Cumple; RT & La publicación de servicios se realizará exclusivamente a través de la capa de borde, con red de distribución de contenidos, cortafuegos de aplicaciones web con reglas gestionadas y personalizadas, y protección contra denegación de servicio distribuida en capas 3, 4 y 7.\newline \emph{T-12, Parte B, fila RT-\allowbreak{}11.07} \\
\textbf{RT-\allowbreak{}11.08}\newline Bases Técnicas Transversales, RT-11.08 & TLS 1.3 en entradas HTTPS y componentes de la solución; TLS 1.0/1.1 prohibidos, HSTS con precarga, certificados con renovación y alerta. El Anexo 4-R admite TLS 1.2 para ERP o terceros incompatibles; esa excepción impide acreditar TLS 1.3 en la totalidad de los enlaces exigidos.\newline \textbf{EDT:} 3.2.1 & /P01\newline Parcial; RT & El cifrado en tránsito empleará TLS 1.3, con prohibición expresa de TLS 1.0 y 1.1, conjuntos de cifrado modernos, HSTS con precarga y gestión automatizada de certificados con rotación y alerta anticipada de vencimiento.\newline \emph{T-12, Parte B, fila RT-\allowbreak{}11.08} \\
\textbf{RT-\allowbreak{}11.09}\newline Bases Técnicas Transversales, RT-11.09 & Cifrado de datos en reposo con AWS KMS administrado por el CLIENTE, política de rotación y separación de funciones en custodia.\newline \textbf{EDT:} 3.2.5 & /P01\newline Cumple; RT & La totalidad de los datos en reposo estará cifrada, con claves gestionadas en un servicio de gestión de claves o en un módulo de seguridad de hardware, política de rotación declarada y separación de funciones en la custodia de claves.\newline \emph{T-12, Parte B, fila RT-\allowbreak{}11.09} \\
\textbf{RT-\allowbreak{}11.10}\newline Bases Técnicas Transversales, RT-11.10 & Cifrado a nivel de campo de comportamiento de pago, antecedentes comerciales y geolocalización; PAN tokenizado por la pasarela si se incorporan medios de pago.\newline \textbf{EDT:} 2.2.4 & /P01\newline Cumple; RT & Los datos de categoría sensible que el caso identifique se cifrarán adicionalmente a nivel de campo, de modo que el acceso a la base de datos no revele su contenido. Caso, cap. 15 — Cifrado a nivel de campo: Exigible para el comportamiento de pago y los antecedentes comerciales de los clientes, para los datos de geolocalización de personas y para los datos de medios de pago electrónicos, si se incorporan.\newline \emph{T-12, Parte B, fila RT-\allowbreak{}11.10} \\
\textbf{RT-\allowbreak{}11.11}\newline Bases Técnicas Transversales, RT-11.11 & Amazon API Gateway con autenticación, autorización, cuotas, límites de tasa, validación de esquema e inspección de carga útil.\newline \textbf{EDT:} 3.2.1 & /P01\newline Cumple; RT & La puerta de enlace de servicios aplicará autenticación, autorización, cuotas, límites de tasa, validación de esquema e inspección de carga útil.\newline \emph{T-12, Parte B, fila RT-\allowbreak{}11.11} \\
\textbf{RT-\allowbreak{}11.12}\newline Bases Técnicas Transversales, RT-11.12 & AWS WAF con retos progresivos y alternativa accesible por teclado y tecnologías de asistencia; aceptación del acceso legítimo y ajuste de falsos bloqueos\newline \textbf{EDT:} 3.2.1 & /P01\newline Cumple; RT & Los puntos de entrada públicos contarán con protección contra bots y abuso automatizado, con reto progresivo que no degrade la accesibilidad ni bloquee a personas usuarias legítimas.\newline \emph{T-12, Parte B, fila RT-\allowbreak{}11.12} \\
\end{tablasimple}

\begin{tablasimple}{>{\hsize=0.82\hsize\linewidth=\hsize}Y >{\hsize=1.04\hsize\linewidth=\hsize}Y >{\hsize=0.65\hsize\linewidth=\hsize}Y >{\hsize=1.49\hsize\linewidth=\hsize}Y}
{\multicolumn{4}{l}{\textbf{Matriz individual: bloque 19 de 26.}} \\
\cab{Requisito y origen} & \cab{Componente y EDT} & \cab{Prueba prevista y estado} & \cab{Criterio y fuente}}
\textbf{RT-\allowbreak{}11.13}\newline Bases Técnicas Transversales, RT-11.13 & La sección 4.2.5.2 del Subdocumento 4 declara cada entrada externa con su nombre de dominio bajo el dominio del CLIENTE (\texttt{clientes.}, \texttt{transportistas.}, \texttt{proveedores.}, \texttt{api.}, \texttt{id.}, \texttt{consola.}, \texttt{admin-id.} y \texttt{as2.<dominio>}), su servicio, puerto y control; los túneles IPsec de los sitios usan direcciones IP públicas de AWS, sin nombre de dominio.\newline \textbf{EDT:} 2.1.2 & /P01\newline Cumple; RT & El PROPONENTE declarará la superficie de exposición completa de la solución: cada nombre de dominio, puerto y servicio alcanzable desde fuera de la red del CLIENTE.\newline \emph{T-12, Parte B, fila RT-\allowbreak{}11.13} \\
\textbf{RT-\allowbreak{}11.14}\newline Bases Técnicas Transversales, RT-11.14 & Eventos centralizados e inalterables en Amazon CloudWatch, 12 meses en línea y 24 meses adicionales en archivo recuperable; seguridad y auditoría por siete años con Object Lock.\newline \textbf{EDT:} 8.1.5, 3.2.5 & /P01\newline Cumple; RT & Los eventos de seguridad se registrarán de forma centralizada e inalterable, con retención mínima de doce meses en línea y veinticuatro meses adicionales en archivo recuperable.\newline \emph{T-12, Parte B, fila RT-\allowbreak{}11.14} \\
\textbf{RT-\allowbreak{}11.15}\newline Bases Técnicas Transversales, RT-11.15 & AWS Security Lake y Security Hub con casos de uso de detección ligados a credencial de conductor revocada, consultas comerciales masivas, alteración de temperatura y EDI anómalo.\newline \textbf{EDT:} 8.1.5 & /P01\newline Cumple; RT & Los eventos se correlacionarán en una plataforma SIEM, con casos de uso de detección definidos específicamente para el proceso de negocio del caso, y no sólo genéricos de infraestructura.\newline \emph{T-12, Parte B, fila RT-\allowbreak{}11.15} \\
\textbf{RT-\allowbreak{}11.16}\newline Bases Técnicas Transversales, RT-11.16 & CrowdStrike Falcon en los nodos, servidores, máquinas virtuales y estaciones on-premise, y GuardDuty Runtime Monitoring en las tareas ECS Fargate de sa-east-1 y us-east-1, con consola central y respuesta del SOC 24$\times$7. El servidor del ERP es un activo del CLIENTE que no forma parte de la solución; el firewall de Talca solo le permite el tráfico de la ACL y la administración del CLIENTE.\newline \textbf{EDT:} 3.2.5, 5.2.2 & /P01\newline Cumple; RT & Se implementará detección y respuesta en puntos finales y en cargas de trabajo, tanto en nube como on-premise.\newline \emph{T-12, Parte B, fila RT-\allowbreak{}11.16} \\
\textbf{RT-\allowbreak{}11.17}\newline Bases Técnicas Transversales, RT-11.17 & SOC 24$\times$7 propio de LafroX en su Centro de Operaciones de Valparaíso, desde la primera marcha blanca (mes 13) hasta el mes 56; un puesto permanente cubierto por cuatro analistas en turnos rotativos, cinco desde el 26-04-2028, contratados antes del mes 13 bajo el Encargado de Seguridad; procedimientos de aviso de incidente crítico en dos horas, notificación de brechas en 24 horas y análisis de causa raíz.\newline \textbf{EDT:} 8.1.5 & /P01\newline Cumple; RT & El PROPONENTE dispondrá de un centro de operaciones de seguridad con cobertura 24$\times$7, propio o subcontratado, y declarará su ubicación, dotación y procedimientos.\newline \emph{T-12, Parte B, fila RT-\allowbreak{}11.17} \\
\textbf{RT-\allowbreak{}11.18}\newline Bases Técnicas Transversales, RT-11.18 & Plan de respuesta con clasificación, escalamiento, plazos, responsables y comunicación al CLIENTE en dos horas para incidentes críticos.\newline \textbf{EDT:} 2.2.3 & /P01\newline Cumple; RT & El plan de respuesta a incidentes de seguridad definirá clasificación, cadena de escalamiento, plazos, responsables y protocolo de comunicación al CLIENTE dentro de las dos horas de detectado un incidente de severidad crítica.\newline \emph{T-12, Parte B, fila RT-\allowbreak{}11.18} \\
\textbf{RT-\allowbreak{}11.19}\newline Bases Técnicas Transversales, RT-11.19 & SOC 24$\times$7 de 8.1.5: toda brecha de seguridad o de datos personales se notifica al CLIENTE con informe preliminar dentro de 24 horas de su detección y análisis de causa raíz dentro de los cinco días hábiles siguientes a la detección.\newline \textbf{EDT:} 8.1.5 & /P01\newline Cumple; RT & Toda brecha de seguridad o de datos personales se notificará al CLIENTE dentro de las 24 horas de su detección, con informe preliminar, y el análisis de causa raíz se entregará dentro de los cinco días hábiles siguientes.\newline \emph{T-12, Parte B, fila RT-\allowbreak{}11.19} \\
\textbf{RT-\allowbreak{}11.20}\newline Bases Técnicas Transversales, RT-11.20 & Pruebas de intrusión por tercero independiente de LafroX, al menos anuales durante la Operación y antes de cada paso a producción; entrega al CLIENTE de cada informe íntegro y plan de remediación con responsables y plazos.\newline \textbf{EDT:} 3.8.6, 3.9.5, 8.2.2 & /P01\newline Cumple; RT & Se ejecutarán pruebas de intrusión por un tercero independiente del ADJUDICATARIO, anualmente y antes de cada paso a producción, con entrega íntegra del informe al CLIENTE y plan de remediación con plazos.\newline \emph{T-12, Parte B, fila RT-\allowbreak{}11.20} \\
\textbf{RT-\allowbreak{}11.21}\newline Bases Técnicas Transversales, RT-11.21 & SOC de 8.1.5 realiza al menos un ejercicio de simulación de incidente por año durante Operación, con participación del CLIENTE y registro de escenario, bitácora y acciones de mejora.\newline \textbf{EDT:} 8.1.5 & /P01\newline Cumple; RT & Se realizarán ejercicios de simulación de incidente con participación del CLIENTE al menos una vez al año durante la Operación.\newline \emph{T-12, Parte B, fila RT-\allowbreak{}11.21} \\
\textbf{RT-\allowbreak{}11.22}\newline Bases Técnicas Transversales, RT-11.22 & SAST con PHP 8.5Stan y Larastan, SCA con composer audit, pruebas DAST y escaneo de imágenes, con bloqueo del despliegue ante hallazgos críticos.\newline \textbf{EDT:} 3.1.4 & /P01\newline Cumple; RT & El flujo de integración continua incorporará análisis estático de código, análisis de composición de software, análisis dinámico y escaneo de imágenes de contenedor, con criterios de bloqueo automático del despliegue ante hallazgos críticos.\newline \emph{T-12, Parte B, fila RT-\allowbreak{}11.22} \\
\textbf{RT-\allowbreak{}11.23}\newline Bases Técnicas Transversales, RT-11.23 & Cada versión liberada se entrega al CLIENTE junto con su SBOM de 3.11.4 en formato CycloneDX o SPDX, con componentes, versiones y licencias.\newline \textbf{EDT:} 3.11.4 & /P01\newline Cumple; RT & Cada versión liberada se acompañará de su inventario de componentes de software en formato CycloneDX o SPDX, entregado al CLIENTE.\newline \emph{T-12, Parte B, fila RT-\allowbreak{}11.23} \\
\textbf{RT-\allowbreak{}11.24}\newline Bases Técnicas Transversales, RT-11.24 & Firma de artefactos y verificación de procedencia conforme a SLSA nivel 3 con GitLab CI y AWS CodeBuild.\newline \textbf{EDT:} 3.1.4 & /P01\newline Cumple; RT & Los artefactos se firmarán y su procedencia se verificará conforme a SLSA nivel 3 o superior.\newline \emph{T-12, Parte B, fila RT-\allowbreak{}11.24} \\
\textbf{RT-\allowbreak{}11.25}\newline Bases Técnicas Transversales, RT-11.25 & Datos sintéticos o anonimizados en Desarrollo y QA; se impide el uso no anonimizado de datos productivos reales.\newline \textbf{EDT:} 3.1.1 & /P01\newline Cumple; RT & Queda prohibido el uso de datos productivos reales en ambientes no productivos sin anonimización o seudonimización verificable.\newline \emph{T-12, Parte B, fila RT-\allowbreak{}11.25} \\
\textbf{RT-\allowbreak{}11.26}\newline Bases Técnicas Transversales, RT-11.26 & Proceso de aprobación de dependencias por licencia, mantenimiento activo y vulnerabilidades conocidas, con composer audit y lista de licencias admitidas.\newline \textbf{EDT:} 3.11.4 & /P01\newline Cumple; RT & El PROPONENTE declarará el proceso de aprobación de nuevas dependencias de terceros, incluyendo criterios de licencia, mantención activa y ausencia de vulnerabilidades conocidas.\newline \emph{T-12, Parte B, fila RT-\allowbreak{}11.26} \\
\textbf{RT-\allowbreak{}11.27}\newline Bases Técnicas Transversales, RT-11.27 & Acceso directo de desarrolladores a producción prohibido; acceso excepcional temporal y aprobado mediante ECS Exec o Session Manager con registro de comandos y salida\newline \textbf{EDT:} 3.3.1 & /P01\newline Cumple; RT & Las personas desarrolladoras no tendrán acceso interactivo directo al ambiente de producción. Todo acceso excepcional será temporal, aprobado, registrado y con sesión grabada.\newline \emph{T-12, Parte B, fila RT-\allowbreak{}11.27} \\
\textbf{RT-\allowbreak{}11.28}\newline Bases Técnicas Transversales, RT-11.28 & Seguridad aplica OWASP SAMM al inicio del desarrollo y anualmente; mide madurez por práctica y sigue acciones de mejora con responsable y evidencias.\newline \textbf{EDT:} 2.2.1 & /P01\newline Cumple; RT & Se aplicará el marco OWASP SAMM o equivalente para medir y mejorar la madurez del proceso de desarrollo seguro, con evaluación inicial y reevaluación anual.\newline \emph{T-12, Parte B, fila RT-\allowbreak{}11.28} \\
\textbf{RT-\allowbreak{}12.01}\newline Bases Técnicas Transversales, RT-12.01 & Keycloak 26.x centraliza identidad con OIDC y OAuth 2.1, o SAML 2.0 cuando la integración lo requiera, e integración con el directorio del CLIENTE por LDAP o equivalente en nube\newline \textbf{EDT:} 3.3.1 & /P01\newline Cumple; RT & La gestión de identidad será centralizada, con federación mediante OpenID Connect y OAuth 2.1, o SAML 2.0 cuando la integración con el CLIENTE lo requiera, e integración con el directorio corporativo del CLIENTE por LDAP o su equivalente en la nube.\newline \emph{T-12, Parte B, fila RT-\allowbreak{}12.01} \\
\textbf{RT-\allowbreak{}12.02}\newline Bases Técnicas Transversales, RT-12.02 & Keycloak 26 centraliza la autenticación OIDC y propaga el cierre de sesión a los servicios conectados.\newline \textbf{EDT:} 3.3.1 & /P01\newline Cumple; RT & Existirá inicio de sesión único para todos los módulos de la solución, con cierre de sesión propagado a todos ellos.\newline \emph{T-12, Parte B, fila RT-\allowbreak{}12.02} \\
\textbf{RT-\allowbreak{}12.03}\newline Bases Técnicas Transversales, RT-12.03 & Keycloak 26 exige MFA a administradores y a todo acceso privilegiado, con claves FIDO2 para administradores, y Verified Access exige MFA y postura en consolas externas. Todo acceso desde fuera de la red corporativa exige segundo factor: en los portales, código TOTP o clave de acceso de Keycloak; en terreno, el terminal enrolado por MDM más el PIN personal.\newline \textbf{EDT:} 3.3.1 & /P01\newline Cumple; RT & La autenticación multifactor será obligatoria para personas usuarias administradoras, para todo acceso privilegiado y para todo acceso originado fuera de la red corporativa.\newline \emph{T-12, Parte B, fila RT-\allowbreak{}12.03} \\
\textbf{RT-\allowbreak{}12.04}\newline Bases Técnicas Transversales, RT-12.04 & Factores resistentes a la suplantación FIDO2 o claves de acceso para perfiles administradores.\newline \textbf{EDT:} 3.3.1 & /P01\newline Cumple; RT & Se soportarán factores resistentes a la suplantación de identidad, tipo FIDO2 o claves de acceso, al menos para los perfiles administradores.\newline \emph{T-12, Parte B, fila RT-\allowbreak{}12.04} \\
\textbf{RT-\allowbreak{}12.05}\newline Bases Técnicas Transversales, RT-12.05 & Control RBAC complementado con atributos de instalación, turno, ruta, dispositivo y empresa transportista, con matriz de segregación de funciones.\newline \textbf{EDT:} 3.3.1 & /P01\newline Cumple; RT & El control de acceso será basado en roles, complementado con control basado en atributos donde el proceso lo exija, con matriz de segregación de funciones documentada y verificable.\newline \emph{T-12, Parte B, fila RT-\allowbreak{}12.05} \\
\textbf{RT-\allowbreak{}12.06}\newline Bases Técnicas Transversales, RT-12.06 & IAM Identity Center federado con Keycloak 26 por SAML 2.0 y SCIM, MFA, conjuntos de permisos temporales con aprobación previa y sesiones ECS Exec o Session Manager con registro de comandos y salida.\newline \textbf{EDT:} 3.3.1 & /P01\newline Cumple; RT & Los accesos privilegiados se gestionarán con elevación temporal a demanda, aprobación previa y grabación de sesión para las operaciones de mayor riesgo.\newline \emph{T-12, Parte B, fila RT-\allowbreak{}12.06} \\
\textbf{RT-\allowbreak{}12.07}\newline Bases Técnicas Transversales, RT-12.07 & Tokens con vida máxima de 30 minutos, inactividad de 15 o 30 minutos, refresco rotatorio de hasta 24 horas o 8 horas para privilegios, revocación y control de sesiones concurrentes.\newline \textbf{EDT:} 3.3.1 & /P01\newline Cumple; RT & La política de sesión declarará duración máxima, caducidad por inactividad, renovación de la credencial de sesión tras la autenticación, revocación inmediata y control de sesiones concurrentes.\newline \emph{T-12, Parte B, fila RT-\allowbreak{}12.07} \\
\textbf{RT-\allowbreak{}12.08}\newline Bases Técnicas Transversales, RT-12.08 & Tokens firmados y de vida breve, con credencial de refresco rotatoria; los identificadores de sesión no se transportan en la URL.\newline \textbf{EDT:} 3.3.1 & /P01\newline Cumple; RT & Las credenciales de sesión serán firmadas y de vida breve, con credencial de refresco rotatoria. Queda prohibido transportar identificadores de sesión en la ruta de la dirección web.\newline \emph{T-12, Parte B, fila RT-\allowbreak{}12.08} \\
\textbf{RT-\allowbreak{}12.09}\newline Bases Técnicas Transversales, RT-12.09 & Las altas, cambios de rol, elevaciones, bloqueos y bajas se registran en una bitácora inalterable con retención de siete años.\newline \textbf{EDT:} 3.3.1 & /P01\newline Cumple; RT & Se registrará auditoría completa del ciclo de vida de la identidad: creación, modificación, elevación, bloqueo y baja de cuentas, con no repudio y retención declarada.\newline \emph{T-12, Parte B, fila RT-\allowbreak{}12.09} \\
\end{tablasimple}

\begin{tablasimple}{>{\hsize=0.82\hsize\linewidth=\hsize}Y >{\hsize=1.04\hsize\linewidth=\hsize}Y >{\hsize=0.65\hsize\linewidth=\hsize}Y >{\hsize=1.49\hsize\linewidth=\hsize}Y}
{\multicolumn{4}{l}{\textbf{Matriz individual: bloque 20 de 26.}} \\
\cab{Requisito y origen} & \cab{Componente y EDT} & \cab{Prueba prevista y estado} & \cab{Criterio y fuente}}
\textbf{RT-\allowbreak{}12.10}\newline Bases Técnicas Transversales, RT-12.10 & Aprovisionamiento y desaprovisionamiento automatizados por trabajos Laravel 13 ligados al ciclo laboral; baja efectiva en menos de 24 h con revocación de sesiones\newline \textbf{EDT:} 3.3.1 & /P01\newline Cumple; RT & El aprovisionamiento y el desaprovisionamiento estarán automatizados y ligados al ciclo de vida laboral, con baja efectiva en un plazo no superior a 24 horas desde la desvinculación.\newline \emph{T-12, Parte B, fila RT-\allowbreak{}12.10} \\
\textbf{RT-\allowbreak{}12.11}\newline Bases Técnicas Transversales, RT-12.11 & Credencial de turno firmada, PIN en terminal enrolado y verificador local; pasos guiados para guantes a $-$22 °C, una mano, de pie y a la intemperie; relevo por turno y suplentes enrolados\newline \textbf{EDT:} 3.3.1 & /P01\newline Cumple; RT & El mecanismo de autenticación se adecuará al perfil operacional real descrito en el caso: entornos de terreno, uso con guantes, baja alfabetización digital, dispositivos compartidos por turno y ausencia de correo electrónico personal. Caso, cap. 15 — Autenticación en el perfil operacional: Conductores de empresas transportistas sin correo corporativo y con rotación sin aviso. Dispositivos compartidos entre turnos en bodega. Personal de preparación con 38 \% de rotación anual. Uso con guantes térmicos a $-$22 °C. Uso a una mano durante la descarga. Uso de pie, a la intemperie, en la puerta del local del cliente.\newline \emph{T-12, Parte B, fila RT-\allowbreak{}12.11} \\
\textbf{RT-\allowbreak{}12.12}\newline Bases Técnicas Transversales, RT-12.12 & Registro, verificación y recuperación autoservidos con Keycloak 26 para clientes, cuentas de cadenas, transportistas, conductores y proveedores; OTP para quienes no tienen correo\newline \textbf{EDT:} 3.5.2, 3.3.1 & /P01\newline Cumple; RT & Las personas usuarias externas del CLIENTE —clientes, proveedores, pacientes, productores, según el caso— dispondrán de un mecanismo de registro, verificación de identidad y recuperación de acceso autoservido y seguro. Caso, cap. 15 — Personas usuarias externas: Clientes del canal tradicional y de food service; cadenas de supermercados; empresas transportistas y sus conductores; y proveedores.\newline \emph{T-12, Parte B, fila RT-\allowbreak{}12.12} \\
\textbf{RT-\allowbreak{}12.13}\newline Bases Técnicas Transversales, RT-12.13 & Cuenta de emergencia con doble autorización, custodia fuera de banda, alcance controlado, auditoría y rotación posterior.\newline \textbf{EDT:} 3.3.1 & /P01\newline Cumple; RT & El PROPONENTE declarará el procedimiento de acceso de emergencia (cuenta de último recurso), su custodia, su control y su auditoría.\newline \emph{T-12, Parte B, fila RT-\allowbreak{}12.13} \\
\textbf{RT-\allowbreak{}13.01}\newline Bases Técnicas Transversales, RT-13.01 & WCAG 2.2 AA verificado con herramientas automatizadas y pruebas manuales, con informe de conformidad entregado en los expedientes de aceptación\newline \textbf{EDT:} 3.8.2, 3.9.2 & /P01\newline Cumple; RT & Todas las interfaces destinadas a personas usuarias cumplirán WCAG 2.2 nivel AA, verificado con herramientas automatizadas y con pruebas manuales, con informe de conformidad entregable.\newline \emph{T-12, Parte B, fila RT-\allowbreak{}13.01} \\
\textbf{RT-\allowbreak{}13.02}\newline Bases Técnicas Transversales, RT-13.02 & Portales y consola Angular 22 con Tailwind CSS: puntos sm/md/lg/xl/2xl, menú adaptable, formularios y listados reorganizados en escritorio, tableta y teléfono; Kotlin 2.4 nativo usa diseño por pantalla fija Zebra.\newline \textbf{EDT:} 3.5.2, 2.6.2 & /P01\newline Cumple; RT & El diseño será responsivo y funcionará correctamente en escritorio, tableta y teléfono, con puntos de quiebre coherentes y adaptación inteligente del contenido, no simple reducción.\newline \emph{T-12, Parte B, fila RT-\allowbreak{}13.02} \\
\textbf{RT-\allowbreak{}13.03}\newline Bases Técnicas Transversales, RT-13.03 & Investigación por perfil, prototipos y pruebas con participantes del CLIENTE; hallazgos incorporados antes de construir\newline \textbf{EDT:} 2.6.1, 2.6.2 & /P01\newline Cumple; RT & El PROPONENTE ejecutará investigación con personas usuarias reales del CLIENTE, prototipado y pruebas de usabilidad antes de la construcción definitiva, y documentará los hallazgos y los cambios de diseño que produjeron.\newline \emph{T-12, Parte B, fila RT-\allowbreak{}13.03} \\
\textbf{RT-\allowbreak{}13.04}\newline Bases Técnicas Transversales, RT-13.04 & Indicadores por perfil en la Tabla 3.A.4a: preparador ($\leq$ 1 s, 3 pasos, error $\leq$ 5 \%, aprendizaje $\leq$ 2 h), recepción ($\leq$ 2 s, 3 pasos, $\leq$ 5 \%, $\leq$ 2 h), preventista ($\leq$ 1,5 s, 3 pasos, $\leq$ 5 \%, $\leq$ 4 h) y conductor ($\leq$ 2 s, 3 o 4 pasos, $\leq$ 5 \%, $\leq$ 30 min), medidos en la prueba de usabilidad con personas reales de cada perfil.\newline \textbf{EDT:} 2.6.3 & /P01\newline Cumple; RT & Se comprometerán indicadores de usabilidad medibles: tiempo máximo de la transacción operacional crítica, número máximo de pasos, tasa de error tolerada y tiempo de aprendizaje esperado por perfil.\newline \emph{T-12, Parte B, fila RT-\allowbreak{}13.04} \\
\textbf{RT-\allowbreak{}13.05}\newline Bases Técnicas Transversales, RT-13.05 & Accesos a todas las funciones principales de cada perfil Angular 22/Kotlin 2.4 en un máximo de tres interacciones desde su inicio, comprobados en prototipos y aceptación.\newline \textbf{EDT:} 2.6.2 & /P01\newline Cumple; RT & Ninguna funcionalidad principal requerirá más de tres interacciones desde la pantalla de inicio del perfil correspondiente.\newline \emph{T-12, Parte B, fila RT-\allowbreak{}13.05} \\
\textbf{RT-\allowbreak{}13.06}\newline Bases Técnicas Transversales, RT-13.06 & Resultado visual por acción (confirmada, pendiente, rechazada) y errores comprensibles sin mensajes técnicos\newline \textbf{EDT:} 2.6.2 & /P01\newline Cumple; RT & La solución entregará retroalimentación visual clara ante cada acción, manejará los errores de forma comprensible —indicando qué ocurrió y qué hacer— y evitará mensajes técnicos dirigidos a la persona usuaria final.\newline \emph{T-12, Parte B, fila RT-\allowbreak{}13.06} \\
\textbf{RT-\allowbreak{}13.07}\newline Bases Técnicas Transversales, RT-13.07 & Alto contraste, controles de al menos 44 $\times$ 44 píxeles, iconos con texto y flujos guiados; 15 $\times$ 15 mm en picking\newline \textbf{EDT:} 2.6.2 & /P01\newline Cumple; RT & La solución soportará a personas usuarias con baja alfabetización digital: alto contraste, objetivos táctiles de al menos 44 $\times$ 44 píxeles, iconografía acompañada de texto, y flujos guiados paso a paso.\newline \emph{T-12, Parte B, fila RT-\allowbreak{}13.07} \\
\textbf{RT-\allowbreak{}13.08}\newline Bases Técnicas Transversales, RT-13.08 & Aceptación con flujos reales: MC9400 Cold Storage Freezer a $-$22 °C con guantes; TC58e de preventa y de reparto bajo lluvia, sobre 34 °C, una mano, sol directo y turno sin conexión\newline \textbf{EDT:} 3.8.3 & /P01\newline Cumple; RT & Cuando el caso lo requiera, las interfaces de terreno operarán con guantes, a la intemperie, con luminosidad variable y sin conexión. Caso, cap. 15 — Interfaces de terreno: Exigible operación con guantes térmicos, a $-$22 °C en cámara de congelado, bajo lluvia y sobre 34 °C en verano, a una sola mano, con pantalla legible bajo sol directo, y sin conexión durante un turno completo.\newline \emph{T-12, Parte B, fila RT-\allowbreak{}13.08} \\
\textbf{RT-\allowbreak{}13.09}\newline Bases Técnicas Transversales, RT-13.09 & Sistema de diseño documentado común a Angular 22/Tailwind y Kotlin 2.4: hasta cinco colores principales, a lo más dos familias tipográficas, iconografía, retícula y componentes reutilizables.\newline \textbf{EDT:} 2.6.2 & /P01\newline Cumple; RT & Existirá un sistema de diseño documentado con paleta de colores acotada, jerarquía tipográfica de a lo más dos familias, iconografía coherente, retícula y componentes reutilizables.\newline \emph{T-12, Parte B, fila RT-\allowbreak{}13.09} \\
\textbf{RT-\allowbreak{}13.10}\newline Bases Técnicas Transversales, RT-13.10 & Matriz Chrome, Edge, Firefox y Safari con versión estable mayor vigente e inmediatamente anterior; registro de números exactos, sistemas, pruebas y actualización por versión, compatibles con el Baseline de Angular 22.\newline \textbf{EDT:} 3.8.2 & /P01\newline Cumple; RT & El PROPONENTE declarará la matriz de navegadores y de versiones soportadas y su política de actualización.\newline \emph{T-12, Parte B, fila RT-\allowbreak{}13.10} \\
\textbf{RT-\allowbreak{}13.11}\newline Bases Técnicas Transversales, RT-13.11 & Navegación íntegra por teclado en web y Kotlin 2.4; foco visible y lógico, gestión de diálogos, controles semánticos, atajos documentados y desactivables y aceptación por teclado conforme a WCAG 2.2 AA.\newline \textbf{EDT:} 3.8.2 & /P01\newline Cumple; RT & La solución será navegable íntegramente por teclado, con orden de foco lógico y atajos para las operaciones frecuentes.\newline \emph{T-12, Parte B, fila RT-\allowbreak{}13.11} \\
\textbf{RT-\allowbreak{}13.12}\newline Bases Técnicas Transversales, RT-13.12 & Modo claro/oscuro, tamaño de texto y densidad configurables por persona y conservados en su dispositivo; español para los perfiles del caso, sin exigencia de otros idiomas.\newline \textbf{EDT:} 2.6.2 & /P01\newline Cumple; RT & Se valorará el soporte de modo oscuro, de personalización de la interfaz por persona usuaria y de múltiples idiomas cuando el caso lo justifique.\newline \emph{T-12, Parte B, fila RT-\allowbreak{}13.12} \\
\textbf{RT-\allowbreak{}14.01}\newline Bases Técnicas Transversales, RT-14.01 & OpenTelemetry para PHP 8.5/Laravel 13 y Kotlin 2.4, AWS Distro for OpenTelemetry (colector 0.50) y CloudWatch correlacionados por transaction\_id\newline \textbf{EDT:} 3.1.5 & /P01\newline Cumple; RT & La observabilidad será unificada para nube y on-premise, con métricas, registros y trazas distribuidas correlacionadas por un identificador único de transacción, instrumentadas conforme a OpenTelemetry.\newline \emph{T-12, Parte B, fila RT-\allowbreak{}14.01} \\
\textbf{RT-\allowbreak{}14.02}\newline Bases Técnicas Transversales, RT-14.02 & Tablero del mandante de solo lectura con auditoría de consultas: estado del servicio con métricas de CloudWatch de resolución de un minuto e indicadores de negocio de QuickSight con las latencias de RT-05.29 (operación del día $\leq$ 5 minutos, cierre comercial $\leq$ 2 horas, gestión $\leq$ 4 horas); exportación en CSV y Excel.\newline \textbf{EDT:} 3.1.5 & /P01\newline Cumple; RT & El CLIENTE dispondrá de acceso propio y permanente a los tableros operacionales y de negocio, con datos en tiempo real y capacidad de exportación.\newline \emph{T-12, Parte B, fila RT-\allowbreak{}14.02} \\
\textbf{RT-\allowbreak{}14.03}\newline Bases Técnicas Transversales, RT-14.03 & OpenTelemetry (colector ADOT 0.50) y CloudWatch correlacionan solicitudes y transacciones; los percentiles p95 se fijan por operación y se miden sobre la operación percibida y confirmada.\newline \textbf{EDT:} 3.1.5 & /P01\newline Cumple; RT & Los indicadores de nivel de servicio se medirán sobre la experiencia real de la persona usuaria y no sobre pruebas sintéticas, sin perjuicio de que estas se empleen como complemento.\newline \emph{T-12, Parte B, fila RT-\allowbreak{}14.03} \\
\textbf{RT-\allowbreak{}14.04}\newline Bases Técnicas Transversales, RT-14.04 & Alertas por síntomas de negocio, agrupación y supresión de ruido, escalamiento automático por criticidad y turnos NOC/SOC 24$\times$7\newline \textbf{EDT:} 3.1.5 & /P01\newline Cumple; RT & El alertamiento se basará en síntomas de negocio y no sólo en umbrales de infraestructura, con supresión de ruido, agrupación, escalamiento automático y turnos de disponibilidad declarados.\newline \emph{T-12, Parte B, fila RT-\allowbreak{}14.04} \\
\textbf{RT-\allowbreak{}14.05}\newline Bases Técnicas Transversales, RT-14.05 & Libro de operación y guías por cada escenario de falla previsible, con automatización progresiva de tareas repetitivas\newline \textbf{EDT:} 3.11.1 & /P01\newline Cumple; RT & Existirá un libro de operación y una guía de resolución documentados para cada escenario de falla previsible, con automatización progresiva de las tareas repetitivas.\newline \emph{T-12, Parte B, fila RT-\allowbreak{}14.05} \\
\textbf{RT-\allowbreak{}14.06}\newline Bases Técnicas Transversales, RT-14.06 & Gestión de incidentes y problemas de 8.1.2: análisis obligatorio de causa raíz de cada incidente crítico, entregado al CLIENTE dentro de cinco días hábiles desde el incidente, y seguimiento por SRE de las acciones correctivas con responsable y plazo hasta su cierre verificado.\newline \textbf{EDT:} 8.1.2 & /P01\newline Cumple; RT & Todo incidente crítico dará lugar a un análisis de causa raíz obligatorio, con informe entregable dentro de cinco días hábiles y seguimiento de las acciones correctivas hasta su cierre.\newline \emph{T-12, Parte B, fila RT-\allowbreak{}14.06} \\
\textbf{RT-\allowbreak{}14.07}\newline Bases Técnicas Transversales, RT-14.07 & Exclusión de credenciales y datos personales sensibles en registros; acceso limitado y auditado\newline \textbf{EDT:} 3.1.5 & /P01\newline Cumple; RT & Los registros de la solución no contendrán datos personales sensibles ni credenciales, y su acceso estará controlado y auditado.\newline \emph{T-12, Parte B, fila RT-\allowbreak{}14.07} \\
\textbf{RT-\allowbreak{}14.08}\newline Bases Técnicas Transversales, RT-14.08 & CloudWatch: registros 12 meses en línea y 24 archivados, métricas 13 meses; trazas 30 días en línea en aws/spans mediante Transaction Search, sin archivo adicional; costo por almacenamiento en línea y archivo en la Oferta Económica (Bases Administrativas, art. 50.2, p. 29).\newline \textbf{EDT:} 3.1.5 & /P01\newline Cumple; RT & El PROPONENTE declarará la retención de métricas, registros y trazas, y su costo asociado, distinguiendo el almacenamiento en línea del archivado.\newline \emph{T-12, Parte B, fila RT-\allowbreak{}14.08} \\
\textbf{RT-\allowbreak{}14.09}\newline Bases Técnicas Transversales, RT-14.09 & Trabajo programado con reglas estadísticas deterministas, sin IA, que compara latencia, errores y colas de Amazon CloudWatch con su historia por franja y día; alerta al NOC/SRE antes de alcanzar los límites; aceptación con degradación gradual.\newline \textbf{EDT:} 3.1.5 & /P01\newline Cumple; RT & Se valorará la detección proactiva de anomalías mediante análisis del comportamiento histórico, con alerta antes de que el incidente afecte a la operación.\newline \emph{T-12, Parte B, fila RT-\allowbreak{}14.09} \\
\end{tablasimple}

\begin{tablasimple}{>{\hsize=0.82\hsize\linewidth=\hsize}Y >{\hsize=1.04\hsize\linewidth=\hsize}Y >{\hsize=0.65\hsize\linewidth=\hsize}Y >{\hsize=1.49\hsize\linewidth=\hsize}Y}
{\multicolumn{4}{l}{\textbf{Matriz individual: bloque 21 de 26.}} \\
\cab{Requisito y origen} & \cab{Componente y EDT} & \cab{Prueba prevista y estado} & \cab{Criterio y fuente}}
\textbf{RT-\allowbreak{}15.01}\newline Bases Técnicas Transversales, RT-15.01 & Utilización por perfil/sitio, escalamiento Fargate y revisión trimestral de capacidad para evitar ociosidad\newline \textbf{EDT:} 8.4.1 & /P01\newline Cumple; RT & El PROPONENTE dimensionará la infraestructura ajustada a la demanda real, evitando capacidad ociosa permanente, y declarará el factor de utilización proyectado.\newline \emph{T-12, Parte B, fila RT-\allowbreak{}15.01} \\
\textbf{RT-\allowbreak{}15.02}\newline Bases Técnicas Transversales, RT-15.02 & Ambientes no productivos apagados/reducidos y experiencia logística/cadena de frío.\newline \textbf{EDT:} 3.1.1, 8.4.2 & /P01\newline Parcial; RT & Los ambientes no productivos se apagarán o reducirán fuera del horario de uso. Caso, cap. 15 — Certificaciones sectoriales del adjudicatario: Conocimiento acreditado del Reglamento Sanitario de los Alimentos y de la normativa de documentos tributarios electrónicos. Experiencia comprobable en distribución de consumo masivo o en operación logística con cadena de frío.\newline \emph{T-12, Parte B, fila RT-\allowbreak{}15.02} \\
\textbf{RT-\allowbreak{}15.03}\newline Bases Técnicas Transversales, RT-15.03 & Paquete 8.4.3 con tres informes anuales de huella de carbono de la operación, cada uno con su metodología y comparación interanual (meses 21 a 56).\newline \textbf{EDT:} 8.4.3 & /P01\newline Parcial; RT & El PROPONENTE estimará la huella de carbono anual de la operación de la solución y declarará la metodología empleada.\newline \emph{T-12, Parte B, fila RT-\allowbreak{}15.03} \\
\textbf{RT-\allowbreak{}15.04}\newline Bases Técnicas Transversales, RT-15.04 & PUE estimado de Talca.\newline \textbf{EDT:} 8.4.3 & /P01\newline Parcial; RT & Se declarará la eficiencia en el uso de la energía (PUE) del recinto on-premise y la intensidad de carbono de la región de nube escogida.\newline \emph{T-12, Parte B, fila RT-\allowbreak{}15.04} \\
\textbf{RT-\allowbreak{}15.05}\newline Bases Técnicas Transversales, RT-15.05 & No se oferta.\newline \textbf{EDT:} No aplica & 0; brecha visible\newline No cumple; RT & Se valorará la elección de regiones de nube con menor intensidad de carbono cuando la latencia y la regulación lo permitan, con el análisis comparativo correspondiente.\newline \emph{T-12, Parte B, fila RT-\allowbreak{}15.05} \\
\textbf{RT-\allowbreak{}15.06}\newline Bases Técnicas Transversales, RT-15.06 & Meta de al menos 60 \% menos horas de cómputo no productivo respecto de operación continua; horario S-43 de 60/168 h (64,3 \% potencial), medición mensual de uso de nube y reporte anual con estimación de huella.\newline \textbf{EDT:} 8.4.3 & /P01\newline Cumple; RT & Se valorará la definición de metas de reducción del consumo durante la Operación, con medición y reporte anual.\newline \emph{T-12, Parte B, fila RT-\allowbreak{}15.06} \\
\textbf{RT-\allowbreak{}15.07}\newline Bases Técnicas Transversales, RT-15.07 & ISO 9001:2015, ISO/IEC 27001:2022, CMMI-DEV nivel 3 e ISO/IEC 20000-1:2018 declaradas.\newline \textbf{EDT:} 1.9.4 & /P01\newline Parcial; RT & Las certificaciones se acreditarán con copia del certificado vigente y con el código de verificación del organismo emisor cuando exista.\newline \emph{T-12, Parte B, fila RT-\allowbreak{}15.07} \\
\textbf{RT-\allowbreak{}15.08}\newline Bases Técnicas Transversales, RT-15.08 & Equipo con personas, roles y dedicaciones.\newline \textbf{EDT:} 1.8.1 & /P01\newline Parcial; RT & Las personas certificadas formarán parte del equipo efectivamente asignado al PROYECTO, con dedicación declarada. No se aceptará acreditar personal que no participe.\newline \emph{T-12, Parte B, fila RT-\allowbreak{}15.08} \\
\textbf{RT-\allowbreak{}15.09}\newline Bases Técnicas Transversales, RT-15.09 & Vigencia de certificaciones corporativas y del equipo durante todo el contrato; reposición de persona certificada saliente con perfil igual o superior, certificación vigente y dedicación equivalente, autorización previa escrita, acreditación y traslape mínimo de quince días hábiles sin costo adicional conforme al Art. 76°.\newline \textbf{EDT:} 1.9.4 & /P01\newline Cumple; RT & El ADJUDICATARIO mantendrá vigentes estas certificaciones durante todo el período contractual y las repondrá ante la salida de una persona certificada, conforme al Artículo 76° de las Bases Administrativas.\newline \emph{T-12, Parte B, fila RT-\allowbreak{}15.09} \\
\textbf{RT-\allowbreak{}16.01}\newline Bases Técnicas Transversales, RT-16.01 & Consola Angular 22 para el CLIENTE: personas, roles, permisos y unidades organizacionales con jerarquías sobre Keycloak 26.x, con auditoría\newline \textbf{EDT:} 3.3.1 & /P01\newline Cumple; RT & La solución dispondrá de un módulo de administración que permita al CLIENTE, sin intervención del ADJUDICATARIO, gestionar personas usuarias, roles, permisos, unidades organizacionales y sus jerarquías.\newline \emph{T-12, Parte B, fila RT-\allowbreak{}16.01} \\
\textbf{RT-\allowbreak{}16.02}\newline Bases Técnicas Transversales, RT-16.02 & Consola de administración para parámetros de reglas existentes; mae\_parametro\_version conserva valor, vigencia y actores, y gob\_auditoria registra antes/después, quién y cuándo. Los cambios fuera de las reglas parametrizadas siguen gestión de cambios.\newline \textbf{EDT:} 3.3.5 & /P01\newline Cumple; RT & Las reglas de negocio parametrizables —umbrales, plazos, montos, tolerancias, catálogos, listas de valores, textos de notificación— serán configurables desde la interfaz de administración, con control de versiones y registro de quién cambió qué y cuándo.\newline \emph{T-12, Parte B, fila RT-\allowbreak{}16.02} \\
\textbf{RT-\allowbreak{}16.03}\newline Bases Técnicas Transversales, RT-16.03 & Segundo perfil obligatorio para cambios de impacto operacional; versión y auditoría con proponente, aprobador, valores y justificación\newline \textbf{EDT:} 3.3.5 & /P01\newline Cumple; RT & Todo cambio de parámetro con impacto operacional requerirá aprobación de un segundo perfil y quedará registrado con su justificación.\newline \emph{T-12, Parte B, fila RT-\allowbreak{}16.03} \\
\textbf{RT-\allowbreak{}16.04}\newline Bases Técnicas Transversales, RT-16.04 & Consola Angular 22/Laravel 13 administra valores de reglas implementadas, catálogos, listas, textos y sitios con versión y auditoría; algoritmos, reglas nuevas, estados/transiciones, esquemas, contratos y transportes nuevos requieren desarrollo.\newline \textbf{EDT:} 3.3.5 & /P01\newline Cumple; RT & El PROPONENTE declarará expresamente qué elementos son parametrizables y cuáles requieren desarrollo. Presentar como parametrizable lo que exige desarrollo se evaluará como observación grave.\newline \emph{T-12, Parte B, fila RT-\allowbreak{}16.04} \\
\textbf{RT-\allowbreak{}16.05}\newline Bases Técnicas Transversales, RT-16.05 & QA simula cambios de parámetro con las mismas transacciones y datos controlados, compara valor vigente/propuesto y exige comprobación y aprobación antes de Producción.\newline \textbf{EDT:} 3.1.1 & /P01\newline Cumple; RT & Existirá un ambiente de simulación que permita probar el efecto de un cambio de parámetro antes de aplicarlo a producción.\newline \emph{T-12, Parte B, fila RT-\allowbreak{}16.05} \\
\textbf{RT-\allowbreak{}16.06}\newline Bases Técnicas Transversales, RT-16.06 & gob\_auditoria registra actor/sistema, hora con zona, origen y valores anteriores/posteriores por operación\newline \textbf{EDT:} 3.3.5 & /P01\newline Cumple; RT & Toda operación que cree, modifique o elimine información quedará registrada con identificación de la persona o del sistema que la ejecutó, fecha y hora con zona horaria, origen, valores anteriores y valores posteriores.\newline \emph{T-12, Parte B, fila RT-\allowbreak{}16.06} \\
\textbf{RT-\allowbreak{}16.07}\newline Bases Técnicas Transversales, RT-16.07 & Auditoría inalterable bajo S3 Object Lock, protegida de modificación/eliminación por cualquier perfil\newline \textbf{EDT:} 3.3.5, 3.2.4 & /P01\newline Cumple; RT & El registro de auditoría será inalterable y no podrá ser modificado ni eliminado por ningún perfil, incluido el administrador de la plataforma.\newline \emph{T-12, Parte B, fila RT-\allowbreak{}16.07} \\
\textbf{RT-\allowbreak{}16.08}\newline Bases Técnicas Transversales, RT-16.08 & Consola Angular 22/API Laravel 13 consulta gob\_auditoria por persona, período, entidad y operación; exporta resultado filtrado CSV UTF-8/JSON con permisos, sin acceso a la base; trabajo asíncrono para grandes volúmenes.\newline \textbf{EDT:} 3.3.5 & /P01\newline Cumple; RT & El CLIENTE podrá consultar y exportar la auditoría desde la interfaz, con filtros por persona, período, entidad y tipo de operación, sin requerir acceso a la base de datos.\newline \emph{T-12, Parte B, fila RT-\allowbreak{}16.08} \\
\textbf{RT-\allowbreak{}16.09}\newline Bases Técnicas Transversales, RT-16.09 & Registro de consultas a crédito, pagos y geolocalización, además de modificaciones\newline \textbf{EDT:} 3.3.5 & /P01\newline Cumple; RT & Las consultas a información sensible quedarán registradas, no sólo las modificaciones. Caso, cap. 15 — Registro de consultas a información sensible: Exigible para los antecedentes comerciales y el comportamiento de pago de los clientes, y para los datos de geolocalización de personas, además del registro de modificaciones.\newline \emph{T-12, Parte B, fila RT-\allowbreak{}16.09} \\
\textbf{RT-\allowbreak{}16.10}\newline Bases Técnicas Transversales, RT-16.10 & Auditoría y seguridad retenidas siete años bajo Object Lock\newline \textbf{EDT:} 3.2.4 & /P01\newline Cumple; RT & El período de retención de la auditoría será el que fije el caso y, en su defecto, no inferior a cinco años.\newline \emph{T-12, Parte B, fila RT-\allowbreak{}16.10} \\
\textbf{RT-\allowbreak{}16.11}\newline Bases Técnicas Transversales, RT-16.11 & Flujos con estados, transiciones y responsables; plazos configurables, escalamiento automático por vencimiento y delegación por ausencia en gob\_asignacion\newline \textbf{EDT:} 3.3.1 & /P01\newline Cumple; RT & La solución soportará flujos de trabajo con estados, transiciones, responsables, plazos, escalamiento automático por vencimiento y delegación por ausencia.\newline \emph{T-12, Parte B, fila RT-\allowbreak{}16.11} \\
\textbf{RT-\allowbreak{}16.12}\newline Bases Técnicas Transversales, RT-16.12 & CLIENTE configura en consola responsables por rol/sitio, plazos y niveles de aprobación de flujos existentes; gob\_asignacion y mae\_parametro\_version conservan vigencias y parámetros, con segregación obligatoria.\newline \textbf{EDT:} 3.3.1 & /P01\newline Cumple; RT & Los flujos serán configurables por el CLIENTE sin desarrollo, al menos en lo relativo a responsables, plazos y niveles de aprobación.\newline \emph{T-12, Parte B, fila RT-\allowbreak{}16.12} \\
\textbf{RT-\allowbreak{}16.13}\newline Bases Técnicas Transversales, RT-16.13 & Bandeja unificada por persona con aprobaciones, revisiones y excepciones, priorizadas por criticidad y plazo y con alertas de vencimiento\newline \textbf{EDT:} 3.3.1 & /P01\newline Cumple; RT & Toda solicitud pendiente será visible para su responsable en una bandeja de tareas unificada, con priorización y alerta de vencimiento.\newline \emph{T-12, Parte B, fila RT-\allowbreak{}16.13} \\
\textbf{RT-\allowbreak{}16.14}\newline Bases Técnicas Transversales, RT-16.14 & Motor de reglas del monolito Laravel 13 sobre Symfony ExpressionLanguage 7.4: las condiciones de los puntos de decisión implementados se guardan versionadas en el registro de parámetros, se aprueban por un segundo perfil y se evalúan en ejecución sin recompilar ni desplegar; la auditoría de cada transacción conserva regla, versión y resultado. Caso, cap. 15: POD articulado con la guía de despacho electrónica y el acuse del ERP para el receptor persona natural.\newline \textbf{EDT:} 3.3.5 & /P01\newline Cumple; RT & El motor de reglas permitirá definir condiciones de negocio evaluables sin recompilación, con trazabilidad de qué regla se aplicó a cada transacción. Caso, cap. 15 — Firma electrónica: La prueba de entrega debe articularse con la guía de despacho electrónica y su acuse de recibo conforme a la normativa tributaria vigente. El PROPONENTE deberá investigar y proponer el mecanismo, considerando que el receptor habitual es una persona natural sin firma electrónica avanzada.\newline \emph{T-12, Parte B, fila RT-\allowbreak{}16.14} \\
\textbf{RT-\allowbreak{}16.15}\newline Bases Técnicas Transversales, RT-16.15 & Repositorio S3 versionado y cifrado con metadatos y control de acceso; búsqueda por contenido y metadatos y previsualización sin descarga\newline \textbf{EDT:} 3.3.5 & /P01\newline Cumple; RT & La solución gestionará documentos con versionado, metadatos, control de acceso, búsqueda por contenido y por metadato, y previsualización sin descarga.\newline \emph{T-12, Parte B, fila RT-\allowbreak{}16.15} \\
\textbf{RT-\allowbreak{}16.16}\newline Bases Técnicas Transversales, RT-16.16 & Objetos POD/DTE en Amazon S3 cifrados, con hash y retención por dominio según política del caso\newline \textbf{EDT:} 3.4.8, 3.2.4 & /P01\newline Cumple; RT & Los documentos se almacenarán cifrados, con verificación de integridad y retención conforme a la política del caso.\newline \emph{T-12, Parte B, fila RT-\allowbreak{}16.16} \\
\end{tablasimple}

\begin{tablasimple}{>{\hsize=0.82\hsize\linewidth=\hsize}Y >{\hsize=1.04\hsize\linewidth=\hsize}Y >{\hsize=0.65\hsize\linewidth=\hsize}Y >{\hsize=1.49\hsize\linewidth=\hsize}Y}
{\multicolumn{4}{l}{\textbf{Matriz individual: bloque 22 de 26.}} \\
\cab{Requisito y origen} & \cab{Componente y EDT} & \cab{Prueba prevista y estado} & \cab{Criterio y fuente}}
\textbf{RT-\allowbreak{}16.17}\newline Bases Técnicas Transversales, RT-16.17 & Firma por acto: guía de despacho electrónica con firma electrónica avanzada del CLIENTE aplicada por el ERP ante el SII, con verificación de la firma del XML y de la vigencia del certificado a la fecha de firma antes de vincularla; acuse de receptor electrónico habilitado con su firma verificada del mismo modo; firma electrónica simple del receptor persona natural (Ley N.º 19.799, art. 2, letra f) como evidencia vinculada, con el acuse según el mecanismo que acuerde el CLIENTE (V-18); el caso no exige firma avanzada para otros actos. XML, acuses y POD con hash, marca de tiempo y cadena de verificación en S3 Object Lock por seis años.\newline \textbf{EDT:} 3.3.2 & /P01\newline Cumple; RT & La solución soportará firma electrónica conforme a la Ley N° 19.799, con firma avanzada para los actos que el caso lo requiera, y verificación de validez del certificado al momento de la firma.\newline \emph{T-12, Parte B, fila RT-\allowbreak{}16.17} \\
\textbf{RT-\allowbreak{}16.18}\newline Bases Técnicas Transversales, RT-16.18 & Guía electrónica firmada por el ERP; XML firmado, acuse SII y POD con hash, marca de tiempo del registro como sello de tiempo y cadena de verificación disponible en S3 Object Lock por seis años, para verificar tras vencer el certificado.\newline \textbf{EDT:} 3.3.2 & /P01\newline Cumple; RT & Se generará el sello de tiempo y se conservará la evidencia de firma que permita verificar el documento con posterioridad al vencimiento del certificado.\newline \emph{T-12, Parte B, fila RT-\allowbreak{}16.18} \\
\textbf{RT-\allowbreak{}16.19}\newline Bases Técnicas Transversales, RT-16.19 & CLIENTE administra plantillas versionadas de avisos y comprobantes en Angular 22; Laravel 13 combina datos autorizados de pedido/entrega y descarga HTML abierto, reutilizando not\_plantilla/cuerpo\_ref.\newline \textbf{EDT:} 3.6.4 & /P01\newline Cumple; RT & La solución generará documentos a partir de plantillas administrables por el CLIENTE, con datos de la transacción y salida en formato abierto.\newline \emph{T-12, Parte B, fila RT-\allowbreak{}16.19} \\
\textbf{RT-\allowbreak{}16.20}\newline Bases Técnicas Transversales, RT-16.20 & Trabajador Laravel 13 de notificaciones: correo por Amazon SES, SMS por Amazon SNS, WhatsApp por AWS End User Messaging Social y aviso en portal\newline \textbf{EDT:} 3.6.4 & /P01\newline Cumple; RT & La solución enviará notificaciones por al menos tres canales: correo electrónico, notificación en la aplicación y mensajería instantánea o SMS, según lo que el caso requiera.\newline \emph{T-12, Parte B, fila RT-\allowbreak{}16.20} \\
\textbf{RT-\allowbreak{}16.21}\newline Bases Técnicas Transversales, RT-16.21 & CLIENTE administra plantillas versionadas con variables transaccionales; elige por cliente CORREO, AVISO\_EN\_PORTAL, WHATSAPP y SMS, incluido despacho/llegada. not\_preferencia conserva elección y bajas entre versiones.\newline \textbf{EDT:} 3.6.4 & /P01\newline Cumple; RT & Las plantillas de notificación serán administrables por el CLIENTE, con variables de la transacción, y versionadas. Caso, cap. 15 — Canales de notificación: Correo electrónico, notificación en aplicación, mensajería instantánea al cliente del canal tradicional, y mensaje de texto para avisos de despacho y de llegada. El canal debe elegirse por cliente, porque buena parte del canal tradicional no usa correo.\newline \emph{T-12, Parte B, fila RT-\allowbreak{}16.21} \\
\textbf{RT-\allowbreak{}16.22}\newline Bases Técnicas Transversales, RT-16.22 & Preferencias personales de canal y frecuencia como atributos de identidad en Keycloak 26; avisos obligatorios por el canal que defina el CLIENTE\newline \textbf{EDT:} 3.6.4 & /P01\newline Cumple; RT & Cada persona usuaria podrá configurar sus preferencias de canal y de frecuencia, respetando las notificaciones que el CLIENTE defina como obligatorias.\newline \emph{T-12, Parte B, fila RT-\allowbreak{}16.22} \\
\textbf{RT-\allowbreak{}16.23}\newline Bases Técnicas Transversales, RT-16.23 & INT-11 envía asíncronamente, reintenta y deduplica, y registra entrega y error. La apertura se registra cuando el proveedor o el portal la informa; en SMS, con la primera apertura del enlace corto al aviso en el portal, y sin ella queda NO\_SOPORTADA.\newline \textbf{EDT:} 3.6.4 & /P01\newline Cumple; RT & El envío será asíncrono, con reintento ante falla, control de duplicados y registro de entrega, apertura y error por cada mensaje.\newline \emph{T-12, Parte B, fila RT-\allowbreak{}16.23} \\
\textbf{RT-\allowbreak{}16.24}\newline Bases Técnicas Transversales, RT-16.24 & Correo: Amazon SES; SMS: Amazon SNS; WhatsApp: AWS End User Messaging Social; portal: LafroX, Angular 22/Laravel 13 en ECS Fargate N-04. Volumen conjunto: 2.800/5.200 avisos diarios. Costo unitario por canal y tratamiento del costo variable en la Oferta Económica (Bases Administrativas, art. 50.2, p. 29).\newline \textbf{EDT:} 3.6.4 & /P01\newline Cumple; RT & El PROPONENTE declarará el proveedor de cada canal, su costo unitario, su volumen proyectado y el tratamiento del costo variable en la Oferta Económica.\newline \emph{T-12, Parte B, fila RT-\allowbreak{}16.24} \\
\textbf{RT-\allowbreak{}16.25}\newline Bases Técnicas Transversales, RT-16.25 & INT-11: separación de avisos operacionales y comerciales; remitente y baja expedita en cada mensaje comercial, preferencia registrada y supresión desde la solicitud de envíos, colas y reintentos conforme a Ley N.º 19.496, art. 28 B.\newline \textbf{EDT:} 3.6.4 & /P01\newline Cumple; RT & Las notificaciones respetarán la normativa de comunicaciones comerciales y permitirán la baja cuando corresponda.\newline \emph{T-12, Parte B, fila RT-\allowbreak{}16.25} \\
\textbf{RT-\allowbreak{}16.26}\newline Bases Técnicas Transversales, RT-16.26 & No se oferta.\newline \textbf{EDT:} No aplica & 0; brecha visible\newline No cumple; RT & Se valorará la integración con canales conversacionales que permitan a la persona usuaria responder y ejecutar acciones desde el propio canal.\newline \emph{T-12, Parte B, fila RT-\allowbreak{}16.26} \\
\textbf{RT-\allowbreak{}16.27}\newline Bases Técnicas Transversales, RT-16.27 & Búsqueda global con texto completo en PostgreSQL 16, tolerancia a errores, facetas y control de acceso de quien busca\newline \textbf{EDT:} 3.3.1 & /P01\newline Cumple; RT & La solución dispondrá de búsqueda global con indexación de texto completo, tolerancia a errores de escritura, filtros facetados y respeto del control de acceso de la persona que busca.\newline \emph{T-12, Parte B, fila RT-\allowbreak{}16.27} \\
\textbf{RT-\allowbreak{}16.28}\newline Bases Técnicas Transversales, RT-16.28 & Listados ordenables, filtrables y paginados; exportación en formatos abiertos con los filtros aplicados\newline \textbf{EDT:} 3.3.5 & /P01\newline Cumple; RT & Los listados serán ordenables, filtrables, paginados y exportables en formatos abiertos, con el filtro aplicado reflejado en la exportación.\newline \emph{T-12, Parte B, fila RT-\allowbreak{}16.28} \\
\textbf{RT-\allowbreak{}16.29}\newline Bases Técnicas Transversales, RT-16.29 & Exportación de gran volumen como trabajo asíncrono, sesión disponible para otras operaciones autorizadas y aviso al completarse mediante el servicio de notificaciones existente; descarga bajo permisos del solicitante.\newline \textbf{EDT:} 3.3.5 & /P01\newline Cumple; RT & Las exportaciones de gran volumen se procesarán de forma asíncrona, con notificación al completarse y sin bloquear la sesión.\newline \emph{T-12, Parte B, fila RT-\allowbreak{}16.29} \\
\textbf{RT-\allowbreak{}16.30}\newline Bases Técnicas Transversales, RT-16.30 & Auditoría de exportaciones sensibles, catálogo público sin precios y portales autenticados para clientes, transportistas y proveedores\newline \textbf{EDT:} 3.3.5 & /P01\newline Cumple; RT & Toda exportación de información sensible quedará registrada en la auditoría, con identificación de quién exportó qué y cuándo. Caso, cap. 15 — Portal público: Catálogo público de productos sin precios. Portal autenticado para clientes, con pedido en autoservicio, estado de entrega, documentos y saldo. Portal autenticado para transportistas y para proveedores.\newline \emph{T-12, Parte B, fila RT-\allowbreak{}16.30} \\
\textbf{RT-\allowbreak{}16.31}\newline Bases Técnicas Transversales, RT-16.31 & Catálogo público sin precios en S3/CloudFront; pruebas de accesibilidad y carga de portales\newline \textbf{EDT:} 3.5.2, 3.8.2, 3.8.4, 3.9.2, 3.9.3 & /P01\newline Cumple; RT & La solución dispondrá de un portal público con la información que el caso determine, accesible sin autenticación, con los mismos estándares de accesibilidad y desempeño que el resto de la plataforma.\newline \emph{T-12, Parte B, fila RT-\allowbreak{}16.31} \\
\textbf{RT-\allowbreak{}16.32}\newline Bases Técnicas Transversales, RT-16.32 & Portal de clientes para pedido, estado, documentos y saldo; portales transportista/proveedor\newline \textbf{EDT:} 3.5.2 & /P01\newline Cumple; RT & Las personas usuarias externas dispondrán de autoatención para las consultas de mayor frecuencia, evitando el contacto telefónico para operaciones simples.\newline \emph{T-12, Parte B, fila RT-\allowbreak{}16.32} \\
\textbf{RT-\allowbreak{}16.33}\newline Bases Técnicas Transversales, RT-16.33 & Portal de Clientes: meta S-42 de al menos 15 \% menos consultas asistidas de estado de entrega, saldo y documentos por cliente activo al cierre del primer año de Operación; línea base en marcha blanca y reporte mensual al CLIENTE.\newline \textbf{EDT:} 3.5.2 & /P01\newline Cumple; RT & El PROPONENTE estimará la reducción esperada del volumen de atención asistida por efecto de la autoatención y comprometerá el indicador.\newline \emph{T-12, Parte B, fila RT-\allowbreak{}16.33} \\
\textbf{RT-\allowbreak{}16.34}\newline Bases Técnicas Transversales, RT-16.34 & Portales estáticos S3/CloudFront con caché, aislados de API y límite de tasa\newline \textbf{EDT:} 3.2.1 & /P01\newline Cumple; RT & El portal público resistirá picos de tráfico sin degradar los servicios transaccionales internos, mediante aislamiento de recursos y caché.\newline \emph{T-12, Parte B, fila RT-\allowbreak{}16.34} \\
\textbf{RT-\allowbreak{}17.01}\newline Bases Técnicas Transversales, RT-17.01 & Kotlin 2.4/Android 15 con Room para preventa/reparto/bodega; PWA Angular 22 para autoatención cliente; todos offline\newline \textbf{EDT:} 3.4.1, 3.4.6, 3.4.8, 3.4.3, 3.5.2 & /P01\newline Cumple; RT & La solución proveerá una aplicación móvil para los perfiles operacionales que el caso identifique, con funcionamiento sin conexión y sincronización diferida cuando la operación en terreno lo requiera. Caso, cap. 15 — Aplicación móvil: Exigible en cuatro perfiles, todos con operación desconectada: preventa, reparto, preparación y recepción en bodega, y autoatención del cliente.\newline \emph{T-12, Parte B, fila RT-\allowbreak{}17.01} \\
\textbf{RT-\allowbreak{}17.02}\newline Bases Técnicas Transversales, RT-17.02 & Kotlin 2.4 nativo justificado frente a PWA y Flutter por operación desconectada, periféricos y costo de mantención de un único artefacto\newline \textbf{EDT:} 2.1.1 & /P01\newline Cumple; RT & El PROPONENTE declarará si la aplicación es nativa, híbrida o web progresiva, y justificará la decisión frente al requisito de operación desconectada, al acceso a periféricos y al costo de mantención.\newline \emph{T-12, Parte B, fila RT-\allowbreak{}17.02} \\
\textbf{RT-\allowbreak{}17.03}\newline Bases Técnicas Transversales, RT-17.03 & Versión mayor vigente de Android y las dos anteriores, con matriz por modelo y política de actualización por tandas\newline \textbf{EDT:} 3.2.6 & /P01\newline Cumple; RT & La aplicación soportará las versiones de sistema operativo móvil vigentes y las dos anteriores, con política de actualización declarada.\newline \emph{T-12, Parte B, fila RT-\allowbreak{}17.03} \\
\textbf{RT-\allowbreak{}17.04}\newline Bases Técnicas Transversales, RT-17.04 & Distribución por Android Enterprise y Zebra DNA con paquete firmado y verificación de firma y huella antes de instalar\newline \textbf{EDT:} 3.2.6 & /P01\newline Cumple; RT & La aplicación se distribuirá por las tiendas oficiales o mediante gestión de flota corporativa, con firma de la aplicación y verificación de integridad.\newline \emph{T-12, Parte B, fila RT-\allowbreak{}17.04} \\
\textbf{RT-\allowbreak{}17.05}\newline Bases Técnicas Transversales, RT-17.05 & Cola Room/SQLite cifrada, bloqueo/borrado remoto mediante MDM\newline \textbf{EDT:} 3.2.6 & /P01\newline Cumple; RT & La información almacenada en el dispositivo estará cifrada, con borrado remoto y bloqueo ante pérdida o desvinculación de la persona usuaria.\newline \emph{T-12, Parte B, fila RT-\allowbreak{}17.05} \\
\textbf{RT-\allowbreak{}17.06}\newline Bases Técnicas Transversales, RT-17.06 & Integración Kotlin 2.4 de lectores Zebra, impresoras ZQ620 Plus, POS PAX, balanza Dibal y captura de fotografía/firma. Frío de cámara: B-01/B-02. Frío de camión: gateway Greengrass V2 B-04, ioLogik/PT100 B-05/B-06, simuladores B-07, sensor independiente B-08 y termógrafo B-03 EYE SENSOR EN12830; Wi-Fi local, Bluetooth e Iridium Edge B-09. La posición GPS de flota se mantiene separada de la temperatura.\newline \textbf{EDT:} 6.5.1, 6.5.2, 6.5.3, 6.5.4, 6.5.5, 6.5.6 & /P01\newline Cumple; RT & La aplicación integrará los periféricos que el caso requiera: cámara, lector de códigos, NFC, GPS, impresora de etiquetas, balanza o báscula. Caso, cap. 15 — Periféricos a integrar: Lectores de código de barras unidimensionales y bidimensionales, impresora térmica portátil en cabina, sensores de temperatura en cámaras y en vehículos, telemetría de flota, terminales de pago electrónico, balanzas de recepción, y captura de firma y de fotografía en el dispositivo.\newline \emph{T-12, Parte B, fila RT-\allowbreak{}17.06} \\
\textbf{RT-\allowbreak{}17.07}\newline Bases Técnicas Transversales, RT-17.07 & Consumo de datos de 9,83 MB por turno en la ruta más larga y sincronización de la flota en diez minutos; consumo de batería máximo comprometido de 80 \% por turno o jornada en bodega y preventa, con cargador de vehículo en reparto, sincronización por lotes y lector por gatillo, medido en la prueba de aceptación con la unidad de prueba de cada modelo.\newline \textbf{EDT:} 2.1.2 & /P01\newline Cumple; RT & El consumo de datos móviles y de batería se optimizará y se declarará el consumo estimado por turno de trabajo.\newline \emph{T-12, Parte B, fila RT-\allowbreak{}17.07} \\
\end{tablasimple}

\begin{tablasimple}{>{\hsize=0.82\hsize\linewidth=\hsize}Y >{\hsize=1.04\hsize\linewidth=\hsize}Y >{\hsize=0.65\hsize\linewidth=\hsize}Y >{\hsize=1.49\hsize\linewidth=\hsize}Y}
{\multicolumn{4}{l}{\textbf{Matriz individual: bloque 23 de 26.}} \\
\cab{Requisito y origen} & \cab{Componente y EDT} & \cab{Prueba prevista y estado} & \cab{Criterio y fuente}}
\textbf{RT-\allowbreak{}17.08}\newline Bases Técnicas Transversales, RT-17.08 & Vista ligera del Portal de Clientes para teléfonos de bajo costo o generaciones anteriores: imágenes bajo demanda, paginación y sin animaciones; navegadores/sistemas compatibles y aceptación en esos equipos.\newline \textbf{EDT:} 3.5.2 & /P01\newline Cumple; RT & Se valorará la disponibilidad de una versión de la interfaz para dispositivos de bajo costo o de generaciones anteriores, ampliando la cobertura de personas usuarias.\newline \emph{T-12, Parte B, fila RT-\allowbreak{}17.08} \\
\textbf{RT-\allowbreak{}18.01}\newline Bases Técnicas Transversales, RT-18.01 & D-09 declara que la solución no incorpora componentes de inteligencia artificial; el ruteo es determinista y explicable.\newline \textbf{EDT:} No aplica & /P01\newline Cumple; RT & El PROPONENTE declarará cada componente de inteligencia artificial: propósito, modelo o servicio empleado, proveedor, versión y ubicación de procesamiento de los datos.\newline \emph{T-12, Parte B, fila RT-\allowbreak{}18.01} \\
\textbf{RT-\allowbreak{}18.02}\newline Bases Técnicas Transversales, RT-18.02 & Garantía contractual de no utilizar datos del CLIENTE para entrenar modelos de terceros sin autorización previa, expresa y escrita; obligación extensiva a proveedores y subcontratistas con acceso a esos datos.\newline \textbf{EDT:} No aplica & /P01\newline Cumple; RT & Se garantizará contractualmente que los datos del CLIENTE no serán utilizados para entrenar modelos de terceros, salvo autorización expresa y escrita.\newline \emph{T-12, Parte B, fila RT-\allowbreak{}18.02} \\
\textbf{RT-\allowbreak{}18.03}\newline Bases Técnicas Transversales, RT-18.03 & No hay componente de inteligencia artificial en el alcance; por ello no se requiere fijar límites de uso, validación humana ni supervisión de resultados de IA.\newline \textbf{EDT:} No aplica & /P01\newline Cumple; RT & Se documentarán los límites de uso, los casos en que el resultado requiere validación humana previa y el procedimiento de supervisión.\newline \emph{T-12, Parte B, fila RT-\allowbreak{}18.03} \\
\textbf{RT-\allowbreak{}18.04}\newline Bases Técnicas Transversales, RT-18.04 & La oferta declara que no incorpora componentes de inteligencia artificial ni modelos aprendidos, por lo que no hay resultados de IA cuyos riesgos de sesgo, alucinación, fuga o uso indebido deban mitigarse.\newline \textbf{EDT:} No aplica & /P01\newline Cumple; RT & Los riesgos de sesgo, alucinación, fuga de información y uso indebido se evaluarán y mitigarán conforme al NIST AI Risk Management Framework 1.0 y a la norma ISO/IEC 42001.\newline \emph{T-12, Parte B, fila RT-\allowbreak{}18.04} \\
\textbf{RT-\allowbreak{}18.05}\newline Bases Técnicas Transversales, RT-18.05 & D-09 excluye interacciones con componentes de inteligencia artificial que requieran registrar entradas, salidas y decisiones humanas posteriores.\newline \textbf{EDT:} No aplica & /P01\newline Cumple; RT & Toda interacción relevante con un componente de inteligencia artificial quedará registrada para efectos de auditoría y trazabilidad, incluyendo la entrada, la salida y la decisión humana posterior.\newline \emph{T-12, Parte B, fila RT-\allowbreak{}18.05} \\
\textbf{RT-\allowbreak{}18.06}\newline Bases Técnicas Transversales, RT-18.06 & La solución no automatiza funciones mediante inteligencia artificial; no existe un componente de IA cuya desactivación o respaldo manual se requiera.\newline \textbf{EDT:} No aplica & /P01\newline Cumple; RT & El componente podrá desactivarse sin comprometer la operación del resto de la solución, y existirá un procedimiento manual de respaldo para la función que automatiza.\newline \emph{T-12, Parte B, fila RT-\allowbreak{}18.06} \\
\textbf{RT-\allowbreak{}18.07}\newline Bases Técnicas Transversales, RT-18.07 & La oferta declara que no presenta resultados de inteligencia artificial a las personas usuarias.\newline \textbf{EDT:} No aplica & /P01\newline Cumple; RT & Cuando el resultado se presente a una persona usuaria, se indicará expresamente que fue generado o sugerido de forma automática, y su nivel de confianza cuando el modelo lo provea.\newline \emph{T-12, Parte B, fila RT-\allowbreak{}18.07} \\
\textbf{RT-\allowbreak{}18.08}\newline Bases Técnicas Transversales, RT-18.08 & Modelo cinético determinista de INN-03 con variables de entrada, métrica I-03B, línea base, recalibración semestral y umbral de detección de deriva.\newline \textbf{EDT:} 3.10.3.1, 8.3.2 & /P01\newline Cumple; RT & Los modelos predictivos declararán sus variables de entrada, su métrica de desempeño, su línea base y su plan de reentrenamiento y de detección de deriva.\newline \emph{T-12, Parte B, fila RT-\allowbreak{}18.08} \\
\textbf{RT-\allowbreak{}18.09}\newline Bases Técnicas Transversales, RT-18.09 & D-09 declara que la solución no incorpora componentes de inteligencia artificial cuyos resultados requieran asignación de responsabilidad por IA.\newline \textbf{EDT:} No aplica & /P01\newline Cumple; RT & La responsabilidad por los resultados de los componentes de inteligencia artificial recae íntegramente en el ADJUDICATARIO, conforme al Artículo 86° de las Bases Administrativas.\newline \emph{T-12, Parte B, fila RT-\allowbreak{}18.09} \\
\textbf{RT-\allowbreak{}18.10}\newline Bases Técnicas Transversales, RT-18.10 & No se oferta.\newline \textbf{EDT:} No aplica & 0; brecha visible\newline No cumple; RT & Se valorará el uso de automatización robótica de procesos o de agentes para tareas repetitivas de back office, con el ahorro de horas cuantificado y reflejado en la Oferta Económica.\newline \emph{T-12, Parte B, fila RT-\allowbreak{}18.10} \\
\textbf{RT-\allowbreak{}19.01}\newline Bases Técnicas Transversales, RT-19.01 & Oficina de gestión basada en PMBOK, con tablero Kanban para el flujo diario, justificado por el trabajo incremental y subordinado a las líneas base y a la aceptación formal.\newline \textbf{EDT:} 1.8.1 & /P01\newline Cumple; RT & El ADJUDICATARIO constituirá una oficina de gestión del proyecto con metodología declarada, basada en el PMBOK del Project Management Institute e integrando prácticas ágiles donde el PROPONENTE lo justifique.\newline \emph{T-12, Parte B, fila RT-\allowbreak{}19.01} \\
\textbf{RT-\allowbreak{}19.02}\newline Bases Técnicas Transversales, RT-19.02 & Alex Aravena, Jefe de Proyecto y apoderado, mantiene dedicación exclusiva durante implementación y facultades para comprometer al ADJUDICATARIO en materias de ejecución; cambios de alcance, plazo o costo sujetos a las instancias de aprobación previstas.\newline \textbf{EDT:} 1.8.1 & /P01\newline Cumple; RT & El Jefe de Proyecto tendrá dedicación exclusiva durante toda la fase de implementación y facultades para comprometer al ADJUDICATARIO en materias de ejecución.\newline \emph{T-12, Parte B, fila RT-\allowbreak{}19.02} \\
\textbf{RT-\allowbreak{}19.03}\newline Bases Técnicas Transversales, RT-19.03 & Procedimiento de control de cambios con registro, análisis de impacto y aprobación previa de la autoridad del CLIENTE.\newline \textbf{EDT:} 1.4.3, 1.8.2 & /P01\newline Cumple; RT & Existirá un procedimiento formal de control de cambios conforme al Artículo 72° de las Bases Administrativas, con registro de cambios, análisis de impacto y aprobación previa a la ejecución.\newline \emph{T-12, Parte B, fila RT-\allowbreak{}19.03} \\
\textbf{RT-\allowbreak{}19.04}\newline Bases Técnicas Transversales, RT-19.04 & Gestión de riesgos conforme a ISO 31000, con registro vivo, revisión en cada Comité de Proyecto y planes con responsable, plazo y disparador.\newline \textbf{EDT:} 1.6.1, 1.6.2 & /P01\newline Cumple; RT & La gestión del riesgo seguirá la norma ISO 31000, con registro de riesgos vivo, revisión en cada Comité de Proyecto y planes de mitigación con responsable, plazo y disparador.\newline \emph{T-12, Parte B, fila RT-\allowbreak{}19.04} \\
\textbf{RT-\allowbreak{}19.05}\newline Bases Técnicas Transversales, RT-19.05 & Espacio colaborativo accesible al CLIENTE con documentación, entregables, actas y registros vigentes de riesgos y cambios.\newline \textbf{EDT:} 1.8.8 & /P01\newline Cumple; RT & El ADJUDICATARIO habilitará un espacio colaborativo accesible al CLIENTE, con la documentación del proyecto, los entregables, las actas, el registro de riesgos y el registro de cambios siempre actualizados.\newline \emph{T-12, Parte B, fila RT-\allowbreak{}19.05} \\
\textbf{RT-\allowbreak{}19.06}\newline Bases Técnicas Transversales, RT-19.06 & Informe mensual con estado del cronograma, avance físico y financiero, entregables, desviaciones, riesgos, incidencias y compromisos siguientes con responsables y fechas; sustentado en líneas base, valor ganado y registros de ejecución.\newline \textbf{EDT:} 1.8.6 & /P01\newline Cumple; RT & El ADJUDICATARIO entregará un informe mensual de avance con estado del cronograma, avance físico y financiero, entregables del período, desviaciones, riesgos, incidencias y compromisos del período siguiente.\newline \emph{T-12, Parte B, fila RT-\allowbreak{}19.06} \\
\textbf{RT-\allowbreak{}19.07}\newline Bases Técnicas Transversales, RT-19.07 & Medición de avance mediante valor ganado, con SPI y CPI, sin declarar porcentajes cualitativos.\newline \textbf{EDT:} 1.8.6 & /P01\newline Cumple; RT & El avance se medirá con valor ganado, reportando el índice de desempeño del cronograma y del costo, y no mediante declaración cualitativa de porcentaje de avance.\newline \emph{T-12, Parte B, fila RT-\allowbreak{}19.07} \\
\textbf{RT-\allowbreak{}19.08}\newline Bases Técnicas Transversales, RT-19.08 & Tablero de estado del proyecto publicado para el CLIENTE, accesible en todo momento y sincronizado con el último informe.\newline \textbf{EDT:} 1.8.7 & /P01\newline Cumple; RT & El CLIENTE dispondrá de un tablero de estado del proyecto actualizado, accesible en cualquier momento.\newline \emph{T-12, Parte B, fila RT-\allowbreak{}19.08} \\
\textbf{RT-\allowbreak{}19.09}\newline Bases Técnicas Transversales, RT-19.09 & Actas de comités dentro de dos días hábiles, con acuerdos, responsables y plazos.\newline \textbf{EDT:} 1.8.2, 1.8.3, 1.8.4, 1.8.5 & /P01\newline Cumple; RT & Las actas de todos los comités del Artículo 71° de las Bases Administrativas se levantarán dentro de los dos días hábiles siguientes y registrarán acuerdos, responsables y plazos.\newline \emph{T-12, Parte B, fila RT-\allowbreak{}19.09} \\
\textbf{RT-\allowbreak{}19.10}\newline Bases Técnicas Transversales, RT-19.10 & Desviaciones mayores al 10 \% o SPI/CPI bajo 0,90 se informan con plan de recuperación dentro de cinco días hábiles.\newline \textbf{EDT:} 1.8.6 & /P01\newline Cumple; RT & Toda desviación superior al 10 \% en un hito se comunicará dentro de los cinco días hábiles de detectada, con plan de recuperación.\newline \emph{T-12, Parte B, fila RT-\allowbreak{}19.10} \\
\textbf{RT-\allowbreak{}20.01}\newline Bases Técnicas Transversales, RT-20.01 & Plan de implantación por sitio, ola y perfil, alineado con el cronograma obligatorio del Artículo 17.\newline \textbf{EDT:} 4.1.1 & /P01\newline Cumple; RT & El PROPONENTE presentará un plan de implantación por sitio y por perfil, coherente con el cronograma obligatorio del Artículo 17° de las Bases Administrativas.\newline \emph{T-12, Parte B, fila RT-\allowbreak{}20.01} \\
\textbf{RT-\allowbreak{}20.02}\newline Bases Técnicas Transversales, RT-20.02 & Paso a producción gradual por olas, con criterio de avance y procedimiento técnico y operacional de reversión con tiempo declarado.\newline \textbf{EDT:} 4.1.1, 4.1.2, 4.1.3 & /P01\newline Cumple; RT & La estrategia de paso a producción será gradual y reversible. Se declarará el criterio de avance entre olas y el procedimiento de reversión, con su tiempo de ejecución.\newline \emph{T-12, Parte B, fila RT-\allowbreak{}20.02} \\
\textbf{RT-\allowbreak{}20.03}\newline Bases Técnicas Transversales, RT-20.03 & Convivencia con la operación vigente durante la marcha blanca, conciliación diaria y control para evitar doble digitación.\newline \textbf{EDT:} 4.2.1, 4.3.1 & /P01\newline Cumple; RT & Durante la marcha blanca, la solución convivirá con la operación vigente del CLIENTE, con conciliación diaria y sin doble digitación no declarada.\newline \emph{T-12, Parte B, fila RT-\allowbreak{}20.03} \\
\textbf{RT-\allowbreak{}20.04}\newline Bases Técnicas Transversales, RT-20.04 & Indicadores diarios y umbrales para el avance y cierre de cada marcha blanca.\newline \textbf{EDT:} 4.2.1, 4.3.1 & /P01\newline Cumple; RT & El PROPONENTE definirá los indicadores diarios que se medirán durante la marcha blanca y sus umbrales de avance, conforme al Artículo 17.3 de las Bases Administrativas.\newline \emph{T-12, Parte B, fila RT-\allowbreak{}20.04} \\
\end{tablasimple}

\begin{tablasimple}{>{\hsize=0.82\hsize\linewidth=\hsize}Y >{\hsize=1.04\hsize\linewidth=\hsize}Y >{\hsize=0.65\hsize\linewidth=\hsize}Y >{\hsize=1.49\hsize\linewidth=\hsize}Y}
{\multicolumn{4}{l}{\textbf{Matriz individual: bloque 24 de 26.}} \\
\cab{Requisito y origen} & \cab{Componente y EDT} & \cab{Prueba prevista y estado} & \cab{Criterio y fuente}}
\textbf{RT-\allowbreak{}20.05}\newline Bases Técnicas Transversales, RT-20.05 & Acompañamiento en terreno por ola con dotación declarada y decreciente conforme a la curva de adopción.\newline \textbf{EDT:} 4.2.2 & /P01\newline Cumple; RT & Se dispondrá de acompañamiento en terreno durante las primeras semanas de cada ola, con dotación declarada y decreciente según la curva de adopción.\newline \emph{T-12, Parte B, fila RT-\allowbreak{}20.05} \\
\textbf{RT-\allowbreak{}20.06}\newline Bases Técnicas Transversales, RT-20.06 & Estabilización de cuatro semanas por ola, con dotación declarada y sin costo adicional.\newline \textbf{EDT:} 4.2.2, 4.3.2 & /P01\newline Cumple; RT & Se establecerá un período de estabilización con atención reforzada tras cada paso a producción, con dotación y duración declaradas y sin costo adicional.\newline \emph{T-12, Parte B, fila RT-\allowbreak{}20.06} \\
\textbf{RT-\allowbreak{}20.07}\newline Bases Técnicas Transversales, RT-20.07 & Definición de terminado acordada con el CLIENTE para cada entregable antes de ejecutarlo, con código, pruebas, documentación, seguridad y despliegue verificados mediante evidencias y criterios de aceptación del paquete.\newline \textbf{EDT:} 1.5.1 & /P01\newline Cumple; RT & Existirá una definición de terminado acordada con el CLIENTE, aplicable a cada entregable, que incluya código, pruebas, documentación, seguridad y despliegue.\newline \emph{T-12, Parte B, fila RT-\allowbreak{}20.07} \\
\textbf{RT-\allowbreak{}20.08}\newline Bases Técnicas Transversales, RT-20.08 & Protocolo de aceptación de cada hito conforme al Formulario T-17, con criterios objetivos y verificables y evidencias del hito, firmado por la Contraparte Técnica.\newline \textbf{EDT:} 3.8.7, 3.9.6, 4.2.3, 4.3.3 & /P01\newline Cumple; RT & El protocolo de aceptación de cada hito se formalizará conforme al Formulario T-17, con criterios objetivos y verificables.\newline \emph{T-12, Parte B, fila RT-\allowbreak{}20.08} \\
\textbf{RT-\allowbreak{}21.01}\newline Bases Técnicas Transversales, RT-21.01 & NOC 24$\times$7 propio de LafroX en su Centro de Operaciones de Valparaíso (Av. Brasil 2241), fuera de los sitios que monitorea y con continuidad en trabajo remoto mediante Verified Access; un puesto permanente cubierto por cuatro personas en turnos rotativos, cinco desde el 26-04-2028; procedimientos en el libro de operación y las guías de resolución, el escalamiento ante congelamientos y la conmutación regional.\newline \textbf{EDT:} 8.1.1 & /P01\newline Cumple; RT & El ADJUDICATARIO dispondrá de un centro de operaciones de red con cobertura 24$\times$7$\times$365, propio o subcontratado, y declarará su ubicación, dotación por turno y procedimientos.\newline \emph{T-12, Parte B, fila RT-\allowbreak{}21.01} \\
\textbf{RT-\allowbreak{}21.02}\newline Bases Técnicas Transversales, RT-21.02 & Guillermo Castillo, Líder de Operación / SRE, designado gerente de servicio dedicado de la Operación con 100 \% de dedicación en los meses 21 a 56; es la contraparte permanente del CLIENTE y representa a LafroX en el Comité de Operación mensual (niveles de servicio, incidentes, problemas y mejora continua)\newline \textbf{EDT:} 8.1.1 & /P01\newline Cumple; RT & Se designará un gerente de servicio dedicado como contraparte permanente del CLIENTE durante la fase de Operación.\newline \emph{T-12, Parte B, fila RT-\allowbreak{}21.02} \\
\textbf{RT-\allowbreak{}21.03}\newline Bases Técnicas Transversales, RT-21.03 & Líder de Operación/SRE Guillermo Castillo y Encargado de Seguridad Álvaro Catalán nominados, con dedicación declarada durante la Operación.\newline \textbf{EDT:} 8.1.1 & /P01\newline Parcial; RT & El equipo de operación contará con especialistas por tecnología, nominados y con dedicación declarada.\newline \emph{T-12, Parte B, fila RT-\allowbreak{}21.03} \\
\textbf{RT-\allowbreak{}21.04}\newline Bases Técnicas Transversales, RT-21.04 & Procedimientos de operación documentados y ensayados, transferidos al equipo de TI del CLIENTE junto con la base de conocimiento.\newline \textbf{EDT:} 3.11.1, 7.1.4 & /P01\newline Cumple; RT & Los procedimientos de operación estarán documentados y serán ejecutables por el personal del CLIENTE tras la transferencia de conocimiento.\newline \emph{T-12, Parte B, fila RT-\allowbreak{}21.04} \\
\textbf{RT-\allowbreak{}21.05}\newline Bases Técnicas Transversales, RT-21.05 & Mesa de Servicios propia con escalamiento L1, L2 y L3 y dimensionamiento cuantitativo por Erlang C; sistema de gestión de servicios certificado conforme a ISO/IEC 20000-1:2018.\newline \textbf{EDT:} 8.1.1 & /P01\newline Cumple; RT & El PROPONENTE dispondrá de un centro de atención adecuado para soportar integralmente las consultas, propio o subcontratado con un operador de nivel acreditado.\newline \emph{T-12, Parte B, fila RT-\allowbreak{}21.05} \\
\textbf{RT-\allowbreak{}21.06}\newline Bases Técnicas Transversales, RT-21.06 & Mesa con atención mínima de 80 \% antes de 20 segundos, resolución al primer contacto de al menos 70 \% y abandono máximo de 5 \%; medición y reporte mensual con recalibración y refuerzo. Atención 04:00–22:00 de lunes a sábado y cobertura 24$\times$7 en peaks de septiembre y diciembre y ante críticos en la ventana de despacho.\newline \textbf{EDT:} 8.1.1 & /P01\newline Cumple; RT & El centro de atención cumplirá un tiempo de atención al usuario final de 80 \% antes de 20 segundos y una resolución al primer contacto de al menos 70 \%. La tasa de abandono no superará el 5 \%. Caso, cap. 15 — Horario del centro de atención: De 04:00 a 22:00 de lunes a sábado, con cobertura 24$\times$7 durante los peaks de septiembre y diciembre y ante incidentes de severidad crítica en la ventana de despacho.\newline \emph{T-12, Parte B, fila RT-\allowbreak{}21.06} \\
\textbf{RT-\allowbreak{}21.07}\newline Bases Técnicas Transversales, RT-21.07 & Atención de 04:00 a 22:00 de lunes a sábado, que contiene el mínimo de 08:00 a 20:00 en días hábiles e incluye la ventana de despacho de 05:30 a 07:00; incidentes críticos y altos atendidos 24$\times$7$\times$365 (D-10) y 24$\times$7 en septiembre y diciembre.\newline \textbf{EDT:} 8.1.1 & /P01\newline Cumple; RT & El horario mínimo de atención será de 8:00 a 20:00 en días hábiles, ampliado a 24$\times$7 para los incidentes de severidad crítica y para la ventana operacional que defina el caso.\newline \emph{T-12, Parte B, fila RT-\allowbreak{}21.07} \\
\textbf{RT-\allowbreak{}21.08}\newline Bases Técnicas Transversales, RT-21.08 & Mesa de servicio: orientación funcional simple a compleja y preguntas técnicas frecuentes sobre la aplicación y su entorno; atención N1, escalamiento N2/N3 y base de conocimiento de 8.5.3.\newline \textbf{EDT:} 8.1.1, 8.5.3 & /P01\newline Cumple; RT & La atención cubrirá orientación funcional de distintos grados de complejidad, desde preguntas simples hasta situaciones complejas, y las preguntas técnicas más frecuentes sobre la aplicación y su entorno de uso.\newline \emph{T-12, Parte B, fila RT-\allowbreak{}21.08} \\
\textbf{RT-\allowbreak{}21.09}\newline Bases Técnicas Transversales, RT-21.09 & Registro histórico de todas las actividades de soporte vinculado al ticket, con clasificación, actuaciones, responsables y resultado; informe mensual de niveles de servicio con análisis de tendencia de los principales requerimientos.\newline \textbf{EDT:} 8.1.2 & /P01\newline Cumple; RT & Se mantendrá un registro histórico de las actividades de soporte que permita monitorear y gestionar los principales requerimientos, con análisis de tendencia mensual.\newline \emph{T-12, Parte B, fila RT-\allowbreak{}21.09} \\
\textbf{RT-\allowbreak{}21.10}\newline Bases Técnicas Transversales, RT-21.10 & Centro de soporte revisa tickets y retroalimentación de usuarios para identificar necesidades recurrentes y complejidades operativas; incorpora soluciones y lecciones validadas a la base de conocimiento y revisa trimestralmente uso y utilidad para ajustar procedimientos.\newline \textbf{EDT:} 8.5.3 & /P01\newline Cumple; RT & Existirá un proceso de aprendizaje del centro de soporte que acumule conocimiento sobre las complejidades de la operación y las necesidades de las personas usuarias, reflejado en la base de conocimiento.\newline \emph{T-12, Parte B, fila RT-\allowbreak{}21.10} \\
\textbf{RT-\allowbreak{}21.11}\newline Bases Técnicas Transversales, RT-21.11 & Plataforma de autoformación en línea con cursos por perfil y registro de avance de cada persona.\newline \textbf{EDT:} 7.1.6 & /P01\newline Cumple; RT & Se proveerán servicios de capacitación en línea en modalidad de autoformación, con registro de avance para efectos de monitoreo.\newline \emph{T-12, Parte B, fila RT-\allowbreak{}21.11} \\
\textbf{RT-\allowbreak{}21.12}\newline Bases Técnicas Transversales, RT-21.12 & No se oferta.\newline \textbf{EDT:} No aplica & 0; brecha visible\newline No cumple; RT & Se habilitarán espacios de interacción entre las entidades y personas usuarias del sistema, que permitan compartir experiencias de uso.\newline \emph{T-12, Parte B, fila RT-\allowbreak{}21.12} \\
\textbf{RT-\allowbreak{}21.13}\newline Bases Técnicas Transversales, RT-21.13 & Tablero del CLIENTE de solo lectura con indicadores clave de soporte RT-21.06 abiertos; indicadores específicos según autorización y consultas auditadas.\newline \textbf{EDT:} 3.1.5, 8.1.2 & /P01\newline Cumple; RT & Los indicadores clave del proceso de soporte estarán abiertos al CLIENTE, y los indicadores específicos serán accesibles según nivel de autorización.\newline \emph{T-12, Parte B, fila RT-\allowbreak{}21.13} \\
\textbf{RT-\allowbreak{}21.14}\newline Bases Técnicas Transversales, RT-21.14 & Dotación dimensionada por Erlang C desde una demanda base proyectada de 2.000 contactos mensuales del caso, con límite de capacidad calculado para las franjas de atención.\newline \textbf{EDT:} 8.1.1 & /P01\newline Cumple; RT & El PROPONENTE dimensionará la dotación del centro de atención con fundamento cuantitativo, empleando teoría de colas o el modelo Erlang C, a partir del volumen de contactos proyectado del caso.\newline \emph{T-12, Parte B, fila RT-\allowbreak{}21.14} \\
\textbf{RT-\allowbreak{}21.15}\newline Bases Técnicas Transversales, RT-21.15 & Canal único de mesa para incidentes y solicitudes con número de ticket, clasificación por severidad, escalamiento L1–L3 y seguimiento de actuaciones y estados hasta el cierre conforme.\newline \textbf{EDT:} 8.1.2 & /P01\newline Cumple; RT & Existirá un canal único de registro de incidentes y solicitudes, con número de ticket, clasificación por severidad y seguimiento del ciclo de vida completo hasta el cierre conforme.\newline \emph{T-12, Parte B, fila RT-\allowbreak{}21.15} \\
\textbf{RT-\allowbreak{}21.16}\newline Bases Técnicas Transversales, RT-21.16 & Traslado de especialistas de nivel 2 (SRE de mantenimiento del hardware, paquete 8.2.3) y de nivel 3 (DES de mantención correctiva, paquete 8.2.1) a Curicó, Chillán, Los Ángeles o el sitio que corresponda, por el medio más rápido disponible y hasta resolver, sin retirar puestos de mesa, NOC ni SOC; la recuperación dentro del Art. 78° sigue siendo remota; el tiempo se imputa al paquete de origen, el exceso lo cubre LafroX y pasajes y viáticos están incluidos en la Oferta Económica\newline \textbf{EDT:} 8.2.3 & /P01\newline Cumple; RT & Cuando el caso comprenda sitios alejados, los especialistas de niveles 2 y 3 deberán trasladarse cuando la resolución lo requiera, por el medio más rápido disponible y con disponibilidad de tiempo suficiente, sin alterar la atención normal del resto de los sitios. El costo del traslado está incluido en la oferta. Caso, cap. 15 — Traslado a sitios alejados: Exigible. Seis instalaciones en cuatro regiones, entre Curicó y Los Ángeles, con 340 km entre los extremos.\newline \emph{T-12, Parte B, fila RT-\allowbreak{}21.16} \\
\textbf{RT-\allowbreak{}21.17}\newline Bases Técnicas Transversales, RT-21.17 & Cierre de cada ticket sujeto a confirmación de la persona usuaria registrada en él; sin confirmación permanece pendiente de cierre bajo seguimiento de la mesa.\newline \textbf{EDT:} 8.1.2 & /P01\newline Cumple; RT & El cierre de un ticket requerirá confirmación de la persona usuaria o transcurso del plazo de confirmación automática declarado.\newline \emph{T-12, Parte B, fila RT-\allowbreak{}21.17} \\
\textbf{RT-\allowbreak{}21.18}\newline Bases Técnicas Transversales, RT-21.18 & Informe mensual de niveles de servicio con cumplimiento por severidad y análisis de cada incumplimiento: causa, impacto y acción correctiva con responsable y plazo, revisado en el Comité de Operación.\newline \textbf{EDT:} 8.1.2 & /P01\newline Cumple; RT & El ADJUDICATARIO reportará mensualmente el cumplimiento de los niveles de servicio de atención, con el detalle por severidad y el análisis de los incumplimientos.\newline \emph{T-12, Parte B, fila RT-\allowbreak{}21.18} \\
\textbf{RT-\allowbreak{}21.19}\newline Bases Técnicas Transversales, RT-21.19 & Bolsa anual de mantención evolutiva de 768 HH DES dentro de las 4.608 HH del paquete 8.2.1 (1.536 HH anuales), separada de la mantención correctiva, que no tiene tope ni costo adicional; el saldo no usado se acumula al año siguiente y el excedente se liquida con la tarifa de la Oferta Económica; perfil DES con el mismo pipeline y puertas de calidad; procedimiento de solicitud por la mesa, estimación en cinco días hábiles, aprobación del CLIENTE en el Comité de Operación y liquidación mensual con saldo\newline \textbf{EDT:} 8.2.1 & /P01\newline Cumple; RT & El PROPONENTE declarará el tamaño de la bolsa anual de horas de mantención evolutiva, su composición por perfil y el procedimiento de solicitud, estimación, aprobación y liquidación.\newline \emph{T-12, Parte B, fila RT-\allowbreak{}21.19} \\
\textbf{RT-\allowbreak{}21.20}\newline Bases Técnicas Transversales, RT-21.20 & Mantenimiento evolutivo dentro de la misma cadena de calidad, pruebas y seguridad del desarrollo, con registro de requisito, cambio, prueba, aprobación y despliegue.\newline \textbf{EDT:} 8.2.1 & /P01\newline Cumple; RT & La mantención evolutiva se someterá al mismo estándar de calidad, pruebas y seguridad que el desarrollo original.\newline \emph{T-12, Parte B, fila RT-\allowbreak{}21.20} \\
\textbf{RT-\allowbreak{}21.21}\newline Bases Técnicas Transversales, RT-21.21 & Registro cuantificado e informe trimestral de deuda técnica; 2.304 HH DES de 8.2.4, equivalentes al 1,7881 \% de las 128.853,67 HH de Operación de T-15. Distribución: 64 HH mensuales y 768 HH anuales.\newline \textbf{EDT:} 8.2.4 & /P01\newline Cumple; RT & El PROPONENTE mantendrá un registro de deuda técnica cuantificado y destinará una fracción declarada de la capacidad de la Operación a reducirla.\newline \emph{T-12, Parte B, fila RT-\allowbreak{}21.21} \\
\end{tablasimple}

\begin{tablasimple}{>{\hsize=0.82\hsize\linewidth=\hsize}Y >{\hsize=1.04\hsize\linewidth=\hsize}Y >{\hsize=0.65\hsize\linewidth=\hsize}Y >{\hsize=1.49\hsize\linewidth=\hsize}Y}
{\multicolumn{4}{l}{\textbf{Matriz individual: bloque 25 de 26.}} \\
\cab{Requisito y origen} & \cab{Componente y EDT} & \cab{Prueba prevista y estado} & \cab{Criterio y fuente}}
\textbf{RT-\allowbreak{}21.22}\newline Bases Técnicas Transversales, RT-21.22 & Plan anual de actualización de componentes de base, con ventana aprobada y reversión ensayada.\newline \textbf{EDT:} 8.2.5 & /P01\newline Cumple; RT & Las actualizaciones de versión de los componentes de base se planificarán anualmente, con ventana acordada y plan de reversión.\newline \emph{T-12, Parte B, fila RT-\allowbreak{}21.22} \\
\textbf{RT-\allowbreak{}22.01}\newline Bases Técnicas Transversales, RT-22.01 & Plan de capacitación por perfiles finales, avanzados, administradores, equipo técnico y soporte.\newline \textbf{EDT:} 7.1.1, 7.1.2 & /P01\newline Cumple; RT & El plan de capacitación se estructurará por perfil: personas usuarias finales, usuarias avanzadas, administradoras, equipo técnico y soporte de niveles 1 y 2.\newline \emph{T-12, Parte B, fila RT-\allowbreak{}22.01} \\
\textbf{RT-\allowbreak{}22.02}\newline Bases Técnicas Transversales, RT-22.02 & Capacitación presencial en cada sitio, en línea sincrónica, autoformación y acompañamiento en el puesto durante la marcha blanca.\newline \textbf{EDT:} 7.1.1 & /P01\newline Cumple; RT & Las modalidades incluirán capacitación presencial en cada sitio de operación, sesiones en línea sincrónicas, autoformación y acompañamiento en puesto de trabajo durante la marcha blanca.\newline \emph{T-12, Parte B, fila RT-\allowbreak{}22.02} \\
\textbf{RT-\allowbreak{}22.03}\newline Bases Técnicas Transversales, RT-22.03 & Manuales por perfil, guías rápidas, preguntas frecuentes, videos tutoriales y base de conocimiento consultable de 3.11.2 y 8.5.3 entregados en español, editables y de propiedad del CLIENTE.\newline \textbf{EDT:} 7.1.5 & /P01\newline Cumple; RT & Todo el material de capacitación se entregará en español, en formato editable y de propiedad del CLIENTE: manuales por perfil, guías rápidas, preguntas frecuentes, videos tutoriales y base de conocimiento consultable.\newline \emph{T-12, Parte B, fila RT-\allowbreak{}22.03} \\
\textbf{RT-\allowbreak{}22.04}\newline Bases Técnicas Transversales, RT-22.04 & Calendario de todos los perfiles por turnos y horarios acordados con el CLIENTE, preservando preparación nocturna, despacho 05:30–07:00 y labores en calle de lunes a sábado, considerando peaks de septiembre y diciembre; acuerdo con cada transportista para conductores externos e inducción continua ante rotación anual de preparadores del 38 \%.\newline \textbf{EDT:} 7.1.1, 7.1.2 & /P01\newline Cumple; RT & La capacitación no podrá afectar la operación del CLIENTE: se programará por turnos y en horarios acordados, considerando la estacionalidad y la ventana operacional del caso. Caso, cap. 15 — Restricción de la capacitación: El personal de preparación trabaja en turno nocturno y rota 38 \% al año. Preventistas y conductores están en la calle de lunes a sábado. Los conductores de empresas transportistas no son trabajadores de la compañía y su capacitación requiere acuerdo con cada transportista.\newline \emph{T-12, Parte B, fila RT-\allowbreak{}22.04} \\
\textbf{RT-\allowbreak{}22.05}\newline Bases Técnicas Transversales, RT-22.05 & Evaluación y certificación de administradores y equipo técnico del CLIENTE como condición de cierre de cada marcha blanca.\newline \textbf{EDT:} 7.1.4, 7.3.1 & /P01\newline Cumple; RT & Las personas usuarias administradoras y el equipo técnico del CLIENTE serán evaluados y certificados como condición para el cierre de cada marcha blanca.\newline \emph{T-12, Parte B, fila RT-\allowbreak{}22.05} \\
\textbf{RT-\allowbreak{}22.06}\newline Bases Técnicas Transversales, RT-22.06 & Dos jornadas anuales de actualización y capacitación del personal nuevo del CLIENTE durante la Operación, sin costo adicional.\newline \textbf{EDT:} 8.5.1 & /P01\newline Cumple; RT & Durante la Operación se ejecutarán al menos dos jornadas anuales de actualización y se capacitará al personal nuevo del CLIENTE sin costo adicional.\newline \emph{T-12, Parte B, fila RT-\allowbreak{}22.06} \\
\textbf{RT-\allowbreak{}22.07}\newline Bases Técnicas Transversales, RT-22.07 & Mentoría al equipo técnico del CLIENTE durante seis meses tras la producción de la Etapa 2.\newline \textbf{EDT:} 8.5.2 & /P01\newline Cumple; RT & El ADJUDICATARIO proveerá acompañamiento y mentoría al equipo técnico del CLIENTE durante al menos seis meses posteriores al paso a producción de la Etapa 2.\newline \emph{T-12, Parte B, fila RT-\allowbreak{}22.07} \\
\textbf{RT-\allowbreak{}22.08}\newline Bases Técnicas Transversales, RT-22.08 & 3.11.2 mantiene la base de incidentes y decisiones desde H4 y mide cada mes su uso y utilidad desde la marcha blanca de la Etapa 1; 8.5.3 la actualiza durante los 36 meses de Operación con métricas trimestrales de uso/utilidad.\newline \textbf{EDT:} 3.11.2, 8.5.3 & /P01\newline Cumple; RT & La base de conocimiento será mantenida y actualizada durante todo el Contrato, con métricas de uso y de utilidad percibida.\newline \emph{T-12, Parte B, fila RT-\allowbreak{}22.08} \\
\textbf{RT-\allowbreak{}22.09}\newline Bases Técnicas Transversales, RT-22.09 & Espacio de capacitación con datos ficticios en QA durante todo el contrato, disponible en el horario de S-43\newline \textbf{EDT:} 7.1.6 & /P01\newline Cumple; RT & Se valorará la existencia de un ambiente permanente de entrenamiento, con datos ficticios, disponible para la práctica sin riesgo.\newline \emph{T-12, Parte B, fila RT-\allowbreak{}22.09} \\
\textbf{RT-\allowbreak{}23.01}\newline Bases Técnicas Transversales, RT-23.01 & El sitio publica historia, giro, misión, visión, valores, oficinas y certificaciones declaradas.\newline \textbf{EDT:} No aplica & /P01\newline Parcial; RT & Información institucional: descripción detallada del giro principal, historia y trayectoria de la empresa con al menos tres años de antigüedad comprobable, misión, visión y valores, certificaciones y acreditaciones vigentes, y presencia geográfica y oficinas.\newline \emph{T-12, Parte B, fila RT-\allowbreak{}23.01} \\
\textbf{RT-\allowbreak{}23.02}\newline Bases Técnicas Transversales, RT-23.02 & El sitio publica tres casos con años, métricas e industrias y ofrece contactos bajo solicitud.\newline \textbf{EDT:} No aplica & /P01\newline Parcial; RT & Experiencia y casos de éxito: portafolio de proyectos similares de los últimos cinco años, casos documentados con métricas verificables, testimonios de clientes con autorización de publicación, industrias atendidas con énfasis en el caso asignado, y volumen de transacciones o de personas usuarias gestionadas.\newline \emph{T-12, Parte B, fila RT-\allowbreak{}23.02} \\
\textbf{RT-\allowbreak{}23.03}\newline Bases Técnicas Transversales, RT-23.03 & Organigrama y equipo técnico clave con personas nominadas, roles, dedicaciones y certificaciones individuales de sus líderes, publicados también en el sitio web \url{https://lafrox.github.io/}.\newline \textbf{EDT:} No aplica & /P01\newline Parcial; RT & Equipo profesional: organigrama del equipo directivo, perfiles de socios y directores, currículos resumidos del equipo técnico clave y certificaciones profesionales del personal.\newline \emph{T-12, Parte B, fila RT-\allowbreak{}23.03} \\
\textbf{RT-\allowbreak{}23.04}\newline Bases Técnicas Transversales, RT-23.04 & Se publican servicios, tecnologías, alianzas, infraestructura y metodologías declaradas.\newline \textbf{EDT:} No aplica & /P01\newline Parcial; RT & Capacidades técnicas: servicios y soluciones ofrecidas, conjunto tecnológico dominado, alianzas tecnológicas con proveedores de nube y de plataforma, metodologías de trabajo certificadas, e infraestructura y capacidad instalada.\newline \emph{T-12, Parte B, fila RT-\allowbreak{}23.04} \\
\textbf{RT-\allowbreak{}23.05}\newline Bases Técnicas Transversales, RT-23.05 & El sitio responde por HTTPS con certificado vigente, se adapta al tamaño de pantalla y declara el criterio WCAG 2.2 AA.\newline \textbf{EDT:} No aplica & /P01\newline Parcial; RT & El sitio será accesible conforme a WCAG 2.2 nivel AA, responsivo y con certificado TLS válido y vigente.\newline \emph{T-12, Parte B, fila RT-\allowbreak{}23.05} \\
\textbf{RT-\allowbreak{}23.06}\newline Bases Técnicas Transversales, RT-23.06 & Compromiso de mantener el sitio web \url{https://lafrox.github.io/} activo y disponible desde el registro del participante hasta la adjudicación\newline \textbf{EDT:} No aplica & /P01\newline Cumple; RT & El sitio se mantendrá activo y disponible durante todo el proceso de licitación, desde el registro del participante hasta la adjudicación.\newline \emph{T-12, Parte B, fila RT-\allowbreak{}23.06} \\
\textbf{RT-\allowbreak{}23.07}\newline Bases Técnicas Transversales, RT-23.07 & No se oferta.\newline \textbf{EDT:} No aplica & 0; brecha visible\newline No cumple; RT & Material educativo digital: seminarios en línea, artículos técnicos o centro de recursos y documentación.\newline \emph{T-12, Parte B, fila RT-\allowbreak{}23.07} \\
\textbf{RT-\allowbreak{}23.08}\newline Bases Técnicas Transversales, RT-23.08 & No se oferta.\newline \textbf{EDT:} No aplica & 0; brecha visible\newline No cumple; RT & Demostración en línea, recorrido virtual de la solución o portal de soporte para clientes.\newline \emph{T-12, Parte B, fila RT-\allowbreak{}23.08} \\
\textbf{RT-\allowbreak{}23.09}\newline Bases Técnicas Transversales, RT-23.09 & No se oferta.\newline \textbf{EDT:} No aplica & 0; brecha visible\newline No cumple; RT & Métricas de disponibilidad y desempeño de los servicios del proponente publicadas en tiempo real.\newline \emph{T-12, Parte B, fila RT-\allowbreak{}23.09} \\
\textbf{RT-\allowbreak{}23.10}\newline Bases Técnicas Transversales, RT-23.10 & No se oferta.\newline \textbf{EDT:} No aplica & 0; brecha visible\newline No cumple; RT & Calculadora de retorno de la inversión o herramienta de estimación pertinente a la industria del caso.\newline \emph{T-12, Parte B, fila RT-\allowbreak{}23.10} \\
\textbf{RT-\allowbreak{}24.01}\newline Bases Técnicas Transversales, RT-24.01 & Se entrega con la Propuesta Final.\newline \textbf{EDT:} No aplica & 0; brecha visible\newline No cumple; RT & Todos los integrantes clave del equipo deberán aparecer individualmente presentándose, en toma de medio cuerpo, mirando directamente a la cámara, con vestimenta formal o de negocio informal, y con una duración mínima de aparición de 10 segundos por persona.\newline \emph{T-12, Parte B, fila RT-\allowbreak{}24.01} \\
\textbf{RT-\allowbreak{}24.02}\newline Bases Técnicas Transversales, RT-24.02 & Se entrega con la Propuesta Final.\newline \textbf{EDT:} No aplica & 0; brecha visible\newline No cumple; RT & Cada vez que aparezca un integrante deberá mostrarse en pantalla su nombre completo, su cargo en el proyecto, sus años de experiencia y su especialización relevante.\newline \emph{T-12, Parte B, fila RT-\allowbreak{}24.02} \\
\textbf{RT-\allowbreak{}24.03}\newline Bases Técnicas Transversales, RT-24.03 & Se entrega con la Propuesta Final.\newline \textbf{EDT:} No aplica & 0; brecha visible\newline No cumple; RT & Las certificaciones principales de cada integrante podrán mostrarse en pantalla.\newline \emph{T-12, Parte B, fila RT-\allowbreak{}24.03} \\
\textbf{RT-\allowbreak{}24.04}\newline Bases Técnicas Transversales, RT-24.04 & Se entrega con la Propuesta Final.\newline \textbf{EDT:} No aplica & 0; brecha visible\newline No cumple; RT & El fondo será el mismo para todas las tomas de personas, con iluminación consistente y encuadre similar para todos los participantes.\newline \emph{T-12, Parte B, fila RT-\allowbreak{}24.04} \\
\textbf{RT-\allowbreak{}24.05}\newline Bases Técnicas Transversales, RT-24.05 & Se entrega con la Propuesta Final.\newline \textbf{EDT:} No aplica & 0; brecha visible\newline No cumple; RT & La paleta de colores corporativa, la tipografía y el uso de la marca serán consistentes en todo el video, con logotipo visible en todas las escenas y datos de contacto en el encabezado o pie.\newline \emph{T-12, Parte B, fila RT-\allowbreak{}24.05} \\
\end{tablasimple}

\begin{tablasimple}{>{\hsize=0.82\hsize\linewidth=\hsize}Y >{\hsize=1.04\hsize\linewidth=\hsize}Y >{\hsize=0.65\hsize\linewidth=\hsize}Y >{\hsize=1.49\hsize\linewidth=\hsize}Y}
{\multicolumn{4}{l}{\textbf{Matriz individual: bloque 26 de 26.}} \\
\cab{Requisito y origen} & \cab{Componente y EDT} & \cab{Prueba prevista y estado} & \cab{Criterio y fuente}}
\textbf{RT-\allowbreak{}24.06}\newline Bases Técnicas Transversales, RT-24.06 & Se entrega con la Propuesta Final.\newline \textbf{EDT:} No aplica & 0; brecha visible\newline No cumple; RT & Los niveles de volumen estarán normalizados, con la misma calidad de grabación y sin variaciones bruscas de sonido.\newline \emph{T-12, Parte B, fila RT-\allowbreak{}24.06} \\
\textbf{RT-\allowbreak{}25.01}\newline Bases Técnicas Transversales, RT-25.01 & Prototipos y pruebas de usabilidad por perfil para cada flujo crítico. El prototipo navegable se entrega con la Propuesta Final.\newline \textbf{EDT:} 2.6.2 & /P01\newline Parcial; RT & Navegación intuitiva sin necesidad de manual, con un máximo de tres interacciones para alcanzar cualquier funcionalidad principal.\newline \emph{T-12, Parte B, fila RT-\allowbreak{}25.01} \\
\textbf{RT-\allowbreak{}25.02}\newline Bases Técnicas Transversales, RT-25.02 & Se entrega con la Propuesta Final.\newline \textbf{EDT:} No aplica & 0; brecha visible\newline No cumple; RT & Retroalimentación visual clara ante cada acción, estados de carga y transición fluidos, y manejo elegante de los errores.\newline \emph{T-12, Parte B, fila RT-\allowbreak{}25.02} \\
\textbf{RT-\allowbreak{}25.03}\newline Bases Técnicas Transversales, RT-25.03 & Criterio WCAG 2.2 AA y pruebas de accesibilidad para las interfaces de la solución. El prototipo navegable se entrega con la Propuesta Final.\newline \textbf{EDT:} 2.6.2 & /P01\newline Parcial; RT & Cumplimiento de WCAG 2.2 nivel AA: contraste adecuado, tamaños de fuente legibles, objetivos táctiles de al menos 44 $\times$ 44 píxeles, navegación completa por teclado, y textos alternativos y etiquetas de accesibilidad.\newline \emph{T-12, Parte B, fila RT-\allowbreak{}25.03} \\
\textbf{RT-\allowbreak{}25.04}\newline Bases Técnicas Transversales, RT-25.04 & Se entrega con la Propuesta Final.\newline \textbf{EDT:} No aplica & 0; brecha visible\newline No cumple; RT & Diseño responsivo demostrado en escritorio (1920 $\times$ 1080 y 1366 $\times$ 768), tableta (768 $\times$ 1024 en vertical y horizontal) y teléfono (375 $\times$ 812), con puntos de quiebre coherentes.\newline \emph{T-12, Parte B, fila RT-\allowbreak{}25.04} \\
\textbf{RT-\allowbreak{}25.05}\newline Bases Técnicas Transversales, RT-25.05 & Se entrega con la Propuesta Final.\newline \textbf{EDT:} No aplica & 0; brecha visible\newline No cumple; RT & Sistema de diseño documentado: paleta de a lo más cinco colores principales, a lo más dos familias tipográficas jerarquizadas, iconografía coherente, retícula y espaciado consistentes, y componentes reutilizables.\newline \emph{T-12, Parte B, fila RT-\allowbreak{}25.05} \\
\textbf{RT-\allowbreak{}25.06}\newline Bases Técnicas Transversales, RT-25.06 & Se entrega con la Propuesta Final.\newline \textbf{EDT:} No aplica & 0; brecha visible\newline No cumple; RT & El prototipo se entregará como enlace navegable en línea, sin instalación ni complementos, compatible con los navegadores modernos vigentes.\newline \emph{T-12, Parte B, fila RT-\allowbreak{}25.06} \\
\textbf{RT-\allowbreak{}25.07}\newline Bases Técnicas Transversales, RT-25.07 & Se entrega con la Propuesta Final.\newline \textbf{EDT:} No aplica & 0; brecha visible\newline No cumple; RT & El acceso al prototipo estará garantizado por un mínimo de seis meses desde su entrega.\newline \emph{T-12, Parte B, fila RT-\allowbreak{}25.07} \\
\textbf{RT-\allowbreak{}25.08}\newline Bases Técnicas Transversales, RT-25.08 & Prototipos de flujos críticos con pruebas de usabilidad por perfil. El prototipo navegable se entrega con la Propuesta Final.\newline \textbf{EDT:} 2.6.2 & /P01\newline Parcial; RT & El prototipo permitirá navegación completa entre pantallas, simulación de ingreso de datos, transiciones y animaciones básicas, estados de posado y activo, simulación de carga de archivos, elementos desplegables funcionales y simulación de validaciones.\newline \emph{T-12, Parte B, fila RT-\allowbreak{}25.08} \\
\textbf{RT-\allowbreak{}25.09}\newline Bases Técnicas Transversales, RT-25.09 & Se entrega con la Propuesta Final.\newline \textbf{EDT:} No aplica & 0; brecha visible\newline No cumple; RT & Se entregará la arquitectura de información del prototipo: mapa del sitio completo, diagrama de navegación, taxonomía y nomenclatura, estructura de menús y jerarquía de la información.\newline \emph{T-12, Parte B, fila RT-\allowbreak{}25.09} \\
\textbf{RT-\allowbreak{}25.10}\newline Bases Técnicas Transversales, RT-25.10 & No se oferta.\newline \textbf{EDT:} No aplica & 0; brecha visible\newline No cumple; RT & Se valorará la posibilidad de exportar el prototipo a un documento interactivo para revisión sin conexión.\newline \emph{T-12, Parte B, fila RT-\allowbreak{}25.10} \\
\textbf{RT-\allowbreak{}26.01}\newline Bases Técnicas Transversales, RT-26.01 & Las cinco innovaciones se ubican en capas, módulos e interfaces en SD13. INN-03 usa B-01/B-02 en cámaras y el kit B-03 a B-15 con C-02 en camiones, con M9 y la evaluación local en el camión del SD4.\newline \textbf{EDT:} 3.10 (cuenta) & /P01\newline Cumple; RT & Cada innovación se ubicará explícitamente en la arquitectura: qué capa la contiene, qué componentes la implementan y qué interfaces consume o expone.\newline \emph{T-12, Parte B, fila RT-\allowbreak{}26.01} \\
\textbf{RT-\allowbreak{}26.02}\newline Bases Técnicas Transversales, RT-26.02 & INN-01: 3.10.1.1–3.10.1.4, validación mes 15 y habilitación 16. INN-02: 3.10.2.1–3.10.2.4 y 8.3.1, circuito mes 5 y validación 15. INN-03: 3.10.3.1–3.10.3.4 y 8.3.2, validación 15 y habilitación 16. INN-04: 5.4.3 y 8.3.3–8.3.4, liquidación desde 24. INN-05: 3.10.4.1–3.10.4.4 y 8.3.5–8.3.6, habilitación 21 y bloque condicionado 27.\newline \textbf{EDT:} 3.10 (cuenta), 5.4.3, 8.3 (cuenta) & /P01\newline Cumple; RT & Cada innovación identificará los paquetes de la estructura de descomposición del trabajo que la ejecutan y el mes del cronograma en que se materializa.\newline \emph{T-12, Parte B, fila RT-\allowbreak{}26.02} \\
\textbf{RT-\allowbreak{}26.03}\newline Bases Técnicas Transversales, RT-26.03 & Las innovaciones tecnológicas declaran nivel de madurez con escala TRL de la Comisión Europea y fuentes citadas en formato APA 7.ª edición.\newline \textbf{EDT:} 3.10 (cuenta) & /P01\newline Cumple; RT & Las innovaciones de base tecnológica declararán el nivel de madurez de la tecnología con la escala utilizada y citarán las fuentes en norma APA 7.ª edición.\newline \emph{T-12, Parte B, fila RT-\allowbreak{}26.03} \\
\textbf{RT-\allowbreak{}26.04}\newline Bases Técnicas Transversales, RT-26.04 & Riesgos de adopción de INN-01 a INN-05 con probabilidad, impacto, mitigación y contingencia declarados en cada ficha.\newline \textbf{EDT:} 1.6.1 & /P01\newline Cumple; RT & Cada innovación declarará su riesgo de adopción, su probabilidad, su impacto, la estrategia de mitigación y el plan de contingencia si no rinde lo esperado.\newline \emph{T-12, Parte B, fila RT-\allowbreak{}26.04} \\
\textbf{RT-\allowbreak{}26.05}\newline Bases Técnicas Transversales, RT-26.05 & Cada innovación declara inversión, costo operacional, beneficio esperado e indicadores con línea base, meta y momento de medición.\newline \textbf{EDT:} 3.10.1.4, 3.10.2.4, 3.10.3.4, 3.10.4.4 & /P01\newline Cumple; RT & Cada innovación declarará su indicador de verificación con línea base, meta y momento de medición, y su impacto en inversión, costo operacional y beneficio esperado.\newline \emph{T-12, Parte B, fila RT-\allowbreak{}26.05} \\
\textbf{RT-\allowbreak{}26.06}\newline Bases Técnicas Transversales, RT-26.06 & Ninguna innovación incorpora un modelo aprendido; INN-03 usa un modelo cinético determinista.\newline \textbf{EDT:} No aplica & /P01\newline Cumple; RT & Las innovaciones que incorporen inteligencia artificial cumplirán íntegramente el Capítulo 18 de este documento.\newline \emph{T-12, Parte B, fila RT-\allowbreak{}26.06} \\
\textbf{RT-\allowbreak{}26.07}\newline Bases Técnicas Transversales, RT-26.07 & Modelado de amenazas propio para la sanitización de INN-02 y para el tratamiento de datos de INN-05, con aprobación del Encargado de Seguridad.\newline \textbf{EDT:} 2.2.2 & /P01\newline Cumple; RT & Las innovaciones que modifiquen la arquitectura de seguridad requerirán su propio modelado de amenazas.\newline \emph{T-12, Parte B, fila RT-\allowbreak{}26.07} \\
\textbf{RT-\allowbreak{}26.08}\newline Bases Técnicas Transversales, RT-26.08 & INN-02 mide antes del mes 16 la reducción de horas de diagnóstico mediante I-02A, con línea base de incidentes de prueba en QA de los meses 5 a 10 y meta relativa declarada.\newline \textbf{EDT:} 3.10.2.4 & /P01\newline Cumple; RT & Se valorará que al menos una innovación sea verificable durante la marcha blanca de la Etapa 1, es decir, que su beneficio pueda medirse antes del mes 16.\newline \emph{T-12, Parte B, fila RT-\allowbreak{}26.08} \\
\end{tablasimple}

Las nueve filas absorbidas o alias conservan su disposición en el T-12 y no añaden casos independientes:
\begin{tablasimple}{Y Y Y}
{\cab{Fila} & \cab{Disposición vigente del T-12} & \cab{Casos adicionales}}
RF-\allowbreak{}19.01 & Matriz de clasificación con niveles Restringido, Confidencial, Interno y Público y controles diferenciados por nivel & 0 \\
RF-\allowbreak{}19.02 & Matriz de controles de seguridad con control ISO/IEC 27001/27002, implementación concreta y evidencia trazable & 0 \\
RF-\allowbreak{}19.03 & La sección 4.2.5.2 del Subdocumento 4 declara cada entrada externa con su nombre de dominio bajo el dominio del CLIENTE (\texttt{clientes.}, \texttt{transportistas.}, \texttt{proveedores.}, \texttt{api.}, \texttt{id.}, \texttt{consola.}, \texttt{admin-id.} y \texttt{as2.<dominio>}), su servicio, puerto y control; los túneles IPsec de los sitios usan direcciones IP públicas de AWS, sin nombre de dominio. & 0 \\
RF-\allowbreak{}19.04 & Plan de respuesta con clasificación, escalamiento, plazos y responsables; aviso al CLIENTE dentro de 2 h para incidentes críticos & 0 \\
RF-\allowbreak{}19.05 & SOC 24$\times$7 de 8.1.5: toda brecha de seguridad o de datos personales se notifica al CLIENTE con informe preliminar dentro de 24 horas de su detección y análisis de causa raíz dentro de los cinco días hábiles siguientes a la detección. & 0 \\
RNF-\allowbreak{}06.02 & Absorbido por RT-03.11 (Parte B) y la restricción S-24; sin componente propio: lo satisface el componente de RT-03.11 (bodega 24 h, reparto/campo 14 h) & 0 \\
RNF-\allowbreak{}06.03 & Alias del plazo de sincronización del dispositivo de reparto de RT-03.13 (Caso, cap. 15), asociado al mecanismo RT-03.12: jornada de 14 h sincronizada en no más de 10 min, medida hasta el acuse durable, incluidas evidencias y espera. AL-OFF-01 y Anexo 4-W.6 verifican el plazo; RNF-02.01 verifica separadamente el CD. & 0 \\
RNF-\allowbreak{}11.03 & Absorbido por RF-11.07 (cierre comercial en no más de 2 h desde el retorno del último camión); sin componente propio & 0 \\
RNF-\allowbreak{}11.04 & Absorbido por RT-05.29 (Parte B; latencia máxima de 4 h hacia la capa analítica); sin componente propio & 0 \\
\end{tablasimple}

\begin{tablasimple}{Y Y Y}
{\cab{Contrato} & \cab{Casos previstos} & \cab{Etapa de aceptación / hito}}
INT-01 & /P01 a /P07 & E1; regresión E2; H5 y H10 \\
INT-02 & /P01 a /P07 & E1; regresión E2; H5 y H10 \\
INT-03 & /P01 a /P07 & E1; regresión E2; H5 y H10 \\
INT-04 & /P01 a /P07 & E1; regresión E2; H5 y H10 \\
INT-05 & /P01 a /P07 & E1; regresión E2; H5 y H10 \\
INT-06 & /P01 a /P07 & E1; regresión E2; H5 y H10 \\
INT-07 & /P01 a /P07 & E1; regresión E2; H5 y H10 \\
INT-08 & /P01 a /P07 & E2; H10 \\
INT-09 & /P01 a /P07 & E1; regresión E2; H5 y H10 \\
INT-10 & /P01 a /P07 & E1; regresión E2; H5 y H10 \\
INT-11 & /P01 a /P07 & E1; regresión E2; H5 y H10 \\
INT-12 & /P01 a /P07 & E1; regresión E2; H5 y H10 \\
INT-13 & /P01 a /P07 & E1; regresión E2; H5 y H10 \\
INT-14 & /P01 a /P07 & E1; regresión E2; H5 y H10 \\
INT-15 & /P01 a /P07 & E1; regresión E2; H5 y H10 \\
\end{tablasimple}
